% !TeX program = pdflatex
% papers/cristal_engenharia.tex — GERADO por tools/cristal_tex.py; NÃO editar à mão.
% Fonte: cristal/cristal.jsonl (broca-so, última versão por conceito).
% Cada secção carrega o registo original em comentário %CRISTAL — a volta é
% tests/cristal_volta.js (resíduo 0 byte a byte contra a fonte).
\documentclass[11pt,a4paper]{article}
\usepackage[utf8]{inputenc}\usepackage[T1]{fontenc}\usepackage[brazil]{babel}
\usepackage{amsmath,amssymb}\usepackage{textcomp}
\usepackage[margin=2.4cm]{geometry}\usepackage{xcolor}
\definecolor{cinza}{gray}{0.35}
\newcommand{\prov}[1]{\par\smallskip\noindent{\small\color{cinza}#1}\par}
\setcounter{secnumdepth}{0}
% ─────────────────────────────────────────────────────────────────────────────
%  papers/gkcapa.tex — a capa do Reino, mínima, para os papers em `article`.
%
%  O estilo.tex completo é do `report` (títulos de capítulo, skak, …). Aqui fica
%  só o que a capa precisa, para o paper continuar a compilar sozinho com
%  pdflatex. O tradutor próprio (tests/tex) carrega o estilo.tex da raiz / o
%  symlink papers/estilo.tex e avalia o mesmo `\gkcapa`.
% ─────────────────────────────────────────────────────────────────────────────
\definecolor{ouro}{HTML}{B8912F}
\definecolor{ouroclaro}{HTML}{F3E9CE}
\definecolor{tinta}{HTML}{1E1B16}
\definecolor{regua}{HTML}{6E6858}
\providecommand{\gknota}{\fontsize{7.62}{11.05}\selectfont}
\providecommand{\gktexto}{\fontsize{10.50}{15.22}\selectfont}
\providecommand{\gksub}{\fontsize{12.33}{17.87}\selectfont}
\providecommand{\gksec}{\fontsize{14.47}{20.98}\selectfont}
\providecommand{\gkcap}{\fontsize{16.99}{24.63}\selectfont}
\providecommand{\gktit}{\fontsize{23.42}{33.95}\selectfont}
\providecommand{\Prj}{\mathbb{P}}
\providecommand{\Trf}{\mathcal{T}}
\providecommand{\gkcapa}[3]{%
\title{\vspace{-0.6cm}
{\color{ouro}\rule{\textwidth}{1.4pt}}\\[6mm]
\textbf{\gktit\color{tinta} Reino Dourado}\\[3mm]{\gksec\itshape\color{regua} Golden Kingdom}\\[8mm]
{\gkcap\scshape\color{ouro} #1}\\[5mm]
{\gksub\color{regua} #2}\\[2mm]
{\gktexto\color{regua}\itshape #3}\\[7mm]
{\color{ouro}\rule{\textwidth}{0.6pt}}\\[9mm]{\gknota\color{regua}\begin{minipage}{\textwidth}\setlength{\parindent}{0pt}%
\textbf{Aviso de Isenção Algorítmica:} O universo, as leis e as entidades contidas nestas páginas ---
incluindo felinos matriciais, aves de rapina algébricas e amuletos de controle de fluxo --- são
construções de pura ficção formal. Esta obra não descreve a realidade física, não serve como
engenharia reversa do cosmos e não constitui doutrina religiosa. Qualquer semelhança com bugs reais do seu
sistema operacional é mera coincidência (ou falta de reza). Prossiga por sua própria conta e
risco.\end{minipage}}}
\author{}\date{}}

\begin{document}
\gkcapa{O Cristal --- a engenharia do cristal}
       {arquitetura, cockpit, pontes, imagem, robótica}
       {corpus recuperado do broca-so; 441 conceitos; projeção verificável da fonte}
\maketitle

%CRISTAL {"arestas": [{"alvo": "redes_5g_nr", "confianca": 1.0, "epistemico": "fato", "origem": "wiki", "rel": "parte_de"}, {"alvo": "canais_fisicos_nr", "confianca": 1.0, "epistemico": "fato", "origem": "wiki", "rel": "usa"}, {"alvo": "camada_mac_nr", "confianca": 1.0, "epistemico": "fato", "origem": "wiki", "rel": "usa"}, {"alvo": "acesso_ao_meio", "confianca": 1.0, "epistemico": "fato", "origem": "wiki", "rel": "especializa"}, {"alvo": "eng_telecomunicacoes", "confianca": 1.0, "epistemico": "fato", "origem": "curriculo", "rel": "parte_de"}], "confianca": 1.0, "contraexemplos": [], "descricao": "Como um UE que não tem recurso de subida agendado consegue o primeiro — enviando um preâmbulo no PRACH e esperando resposta. É a porta de entrada de tudo: acesso inicial, handover, retomada, pedido de escalonamento e recuperação de feixe. Com ou sem contenção.", "epistemico": "fato", "exemplos": [], "id": "acesso_aleatorio_nr", "memoria": "semantica", "meta": {"dominio": "engenharia", "fonte": ""}, "origem": "wiki", "palavras_chave": [], "sinonimos": [], "tipo": "conceito", "titulo": "Acesso Aleatório (RACH)"}
\section{Acesso Aleatório (RACH)}
Como um UE que não tem recurso de subida agendado consegue o primeiro — enviando um preâmbulo no PRACH e esperando resposta. É a porta de entrada de tudo: acesso inicial, handover, retomada, pedido de escalonamento e recuperação de feixe. Com ou sem contenção.
\prov{relações: parte\_de redes\_5g\_nr; usa canais\_fisicos\_nr; usa camada\_mac\_nr; especializa acesso\_ao\_meio; parte\_de eng\_telecomunicacoes}
\prov{proveniência: wiki \textperiodcentered{} fato \textperiodcentered{} confiança 1.0 \textperiodcentered{} domínio engenharia \textperiodcentered{} história 6 versões no jornal}

%CRISTAL {"arestas": [{"alvo": "camada_de_enlace", "confianca": 1.0, "epistemico": "fato", "origem": "wiki", "rel": "parte_de"}, {"alvo": "eng_telecomunicacoes", "confianca": 1.0, "epistemico": "fato", "origem": "curriculo", "rel": "parte_de"}], "confianca": 1.0, "contraexemplos": [], "descricao": "Como nós que compartilham o meio evitam colidir: escutar antes de falar (CSMA), detectar colisão no cabo (CD, Ethernet clássica) ou evitá-la no ar (CA, Wi-Fi). A etiqueta do meio compartilhado.", "epistemico": "fato", "exemplos": [], "id": "acesso_ao_meio", "memoria": "semantica", "meta": {"dominio": "engenharia", "fonte": ""}, "origem": "wiki", "palavras_chave": [], "sinonimos": [], "tipo": "conceito", "titulo": "Controle de Acesso ao Meio (CSMA/CD, CSMA/CA)"}
\section{Controle de Acesso ao Meio (CSMA/CD, CSMA/CA)}
Como nós que compartilham o meio evitam colidir: escutar antes de falar (CSMA), detectar colisão no cabo (CD, Ethernet clássica) ou evitá-la no ar (CA, Wi-Fi). A etiqueta do meio compartilhado.
\prov{relações: parte\_de camada\_de\_enlace; parte\_de eng\_telecomunicacoes}
\prov{proveniência: wiki \textperiodcentered{} fato \textperiodcentered{} confiança 1.0 \textperiodcentered{} domínio engenharia \textperiodcentered{} história 3 versões no jornal}

%CRISTAL {"arestas": [{"alvo": "multiplexacao", "confianca": 1.0, "epistemico": "fato", "origem": "wiki", "rel": "especializa"}, {"alvo": "comunicacoes_moveis", "confianca": 1.0, "epistemico": "fato", "origem": "wiki", "rel": "usa"}, {"alvo": "eng_telecomunicacoes", "confianca": 1.0, "epistemico": "fato", "origem": "curriculo", "rel": "parte_de"}], "confianca": 1.0, "contraexemplos": [], "descricao": "Como muitos usuários compartilham o canal de rádio sem colidir — por frequência, tempo, código ou subportadora. A evolução FDMA→TDMA→CDMA→OFDMA marca as gerações do celular.", "epistemico": "inferencia", "exemplos": [], "id": "acesso_multiplo", "memoria": "semantica", "meta": {"dominio": "engenharia", "fonte": ""}, "origem": "wiki", "palavras_chave": [], "sinonimos": [], "tipo": "conceito", "titulo": "Acesso Múltiplo (FDMA/TDMA/CDMA/OFDMA)"}
\section{Acesso Múltiplo (FDMA/TDMA/CDMA/OFDMA)}
Como muitos usuários compartilham o canal de rádio sem colidir — por frequência, tempo, código ou subportadora. A evolução FDMA\ensuremath{\to}TDMA\ensuremath{\to}CDMA\ensuremath{\to}OFDMA marca as gerações do celular.
\prov{relações: especializa multiplexacao; usa comunicacoes\_moveis; parte\_de eng\_telecomunicacoes}
\prov{proveniência: wiki \textperiodcentered{} inferencia \textperiodcentered{} confiança 1.0 \textperiodcentered{} domínio engenharia \textperiodcentered{} história 4 versões no jornal}

%CRISTAL {"arestas": [{"alvo": "eletronica_de_potencia", "confianca": 1.0, "epistemico": "fato", "origem": "wiki", "rel": "usa"}, {"alvo": "motor_eletrico", "confianca": 1.0, "epistemico": "fato", "origem": "wiki", "rel": "usa"}, {"alvo": "eng_sistemas_energia", "confianca": 1.0, "epistemico": "fato", "origem": "curriculo", "rel": "parte_de"}], "confianca": 1.0, "contraexemplos": [], "descricao": "Controlar velocidade e torque do motor por eletrônica de potência — inversor de frequência, controle vetorial. Eficiência e precisão que o motor ligado direto na rede não dá.", "epistemico": "inferencia", "exemplos": [], "id": "acionamento_eletrico", "memoria": "semantica", "meta": {"dominio": "engenharia", "fonte": ""}, "origem": "wiki", "palavras_chave": [], "sinonimos": [], "tipo": "conceito", "titulo": "Acionamentos e Controle de Motores"}
\section{Acionamentos e Controle de Motores}
Controlar velocidade e torque do motor por eletrônica de potência — inversor de frequência, controle vetorial. Eficiência e precisão que o motor ligado direto na rede não dá.
\prov{relações: usa eletronica\_de\_potencia; usa motor\_eletrico; parte\_de eng\_sistemas\_energia}
\prov{proveniência: wiki \textperiodcentered{} inferencia \textperiodcentered{} confiança 1.0 \textperiodcentered{} domínio engenharia \textperiodcentered{} história 4 versões no jornal}

%CRISTAL {"arestas": [{"alvo": "amostragem", "confianca": 1.0, "epistemico": "fato", "origem": "wiki", "rel": "usa"}, {"alvo": "numero_de_salem", "confianca": 1.0, "epistemico": "fato", "origem": "wiki", "rel": "relaciona"}, {"alvo": "lyapunov_da_orbita", "confianca": 1.0, "epistemico": "fato", "origem": "wiki", "rel": "relaciona"}, {"alvo": "corpo_mineral", "confianca": 1.0, "epistemico": "fato", "origem": "wiki", "rel": "usa"}], "confianca": 1.0, "contraexemplos": [], "descricao": "Amostrar abaixo de Nyquist (fs < 2·f) DOBRA a frequência real numa frequência aparente FALSA. A rotação (o Salem) lida devagar demais aparece como outra órbita — um λ/θ errado. É ler o floco com a régua associativa errada: a máquina vê uma órbita que não existe. Verif.: 50 Hz a 60 Hz → 10 Hz.", "epistemico": "inferencia", "exemplos": [], "id": "aliasing", "memoria": "semantica", "meta": {"aliased": true, "dominio": "engenharia", "exemplo_50hz_60fs_aparente": 10, "fonte": "IoT e sensores"}, "origem": "wiki", "palavras_chave": [], "sinonimos": [], "tipo": "conceito", "titulo": "Aliasing (a régua errada na órbita)"}
\section{Aliasing (a régua errada na órbita)}
Amostrar abaixo de Nyquist (fs < 2\textperiodcentered f) DOBRA a frequência real numa frequência aparente FALSA. A rotação (o Salem) lida devagar demais aparece como outra órbita — um \ensuremath{\lambda}/\ensuremath{\theta} errado. É ler o floco com a régua associativa errada: a máquina vê uma órbita que não existe. Verif.: 50 Hz a 60 Hz \ensuremath{\to} 10 Hz.
\prov{relações: usa amostragem; relaciona numero\_de\_salem; relaciona lyapunov\_da\_orbita; usa corpo\_mineral}
\prov{proveniência: wiki \textperiodcentered{} IoT e sensores \textperiodcentered{} inferencia \textperiodcentered{} confiança 1.0 \textperiodcentered{} domínio engenharia \textperiodcentered{} história 10 versões no jornal}

%CRISTAL {"arestas": [{"alvo": "transmissao_de_energia", "confianca": 1.0, "epistemico": "fato", "origem": "wiki", "rel": "usa"}, {"alvo": "materiais_eletricos", "confianca": 1.0, "epistemico": "fato", "origem": "wiki", "rel": "usa"}, {"alvo": "eng_sistemas_energia", "confianca": 1.0, "epistemico": "fato", "origem": "curriculo", "rel": "parte_de"}], "confianca": 1.0, "contraexemplos": [], "descricao": "Dominar o isolamento em centenas de kV: rigidez dielétrica, descargas parciais, sobretensões e para-raios. Por que a transmissão sobe a tensão (menos perda) e como não deixá-la arcar.", "epistemico": "fato", "exemplos": [], "id": "alta_tensao", "memoria": "semantica", "meta": {"dominio": "engenharia", "fonte": ""}, "origem": "wiki", "palavras_chave": [], "sinonimos": [], "tipo": "conceito", "titulo": "Engenharia de Alta Tensão"}
\section{Engenharia de Alta Tensão}
Dominar o isolamento em centenas de kV: rigidez dielétrica, descargas parciais, sobretensões e para-raios. Por que a transmissão sobe a tensão (menos perda) e como não deixá-la arcar.
\prov{relações: usa transmissao\_de\_energia; usa materiais\_eletricos; parte\_de eng\_sistemas\_energia}
\prov{proveniência: wiki \textperiodcentered{} fato \textperiodcentered{} confiança 1.0 \textperiodcentered{} domínio engenharia \textperiodcentered{} história 4 versões no jornal}

%CRISTAL {"arestas": [{"alvo": "sensor", "confianca": 1.0, "epistemico": "fato", "origem": "wiki", "rel": "usa"}, {"alvo": "orbita_de_partida", "confianca": 1.0, "epistemico": "fato", "origem": "wiki", "rel": "relaciona"}], "confianca": 1.0, "contraexemplos": [], "descricao": "Ler um sinal contínuo em instantes espaçados de 1/fs. O teorema de Nyquist: para não perder informação, a taxa fs precisa ser MAIOR que o dobro da maior frequência do sinal (fs > 2·f). Amostrar é discretizar o tempo — pegar pontos da órbita.", "epistemico": "fato", "exemplos": [], "id": "amostragem", "memoria": "semantica", "meta": {"dominio": "engenharia", "fonte": "IoT e sensores"}, "origem": "wiki", "palavras_chave": [], "sinonimos": [], "tipo": "conceito", "titulo": "Amostragem (Nyquist)"}
\section{Amostragem (Nyquist)}
Ler um sinal contínuo em instantes espaçados de 1/fs. O teorema de Nyquist: para não perder informação, a taxa fs precisa ser MAIOR que o dobro da maior frequência do sinal (fs > 2\textperiodcentered f). Amostrar é discretizar o tempo — pegar pontos da órbita.
\prov{relações: usa sensor; relaciona orbita\_de\_partida}
\prov{proveniência: wiki \textperiodcentered{} IoT e sensores \textperiodcentered{} fato \textperiodcentered{} confiança 1.0 \textperiodcentered{} domínio engenharia \textperiodcentered{} história 6 versões no jornal}

%CRISTAL {"arestas": [{"alvo": "processamento_digital_sinais", "confianca": 1.0, "epistemico": "fato", "origem": "wiki", "rel": "parte_de"}, {"alvo": "conversor_ad_da_disc", "confianca": 1.0, "epistemico": "fato", "origem": "wiki", "rel": "explica"}, {"alvo": "eng_eletronica", "confianca": 1.0, "epistemico": "fato", "origem": "curriculo", "rel": "parte_de"}], "confianca": 1.0, "contraexemplos": [], "descricao": "Para reconstruir um sinal sem perda, amostre a mais que o dobro da maior frequência. Abaixo disso, aliasing irreversível. A ponte fundadora entre o contínuo e o digital.", "epistemico": "fato", "exemplos": [], "id": "amostragem_nyquist", "memoria": "semantica", "meta": {"autor": "Nyquist-Shannon", "dominio": "engenharia", "fonte": ""}, "origem": "wiki", "palavras_chave": [], "sinonimos": [], "tipo": "conceito", "titulo": "Amostragem e Teorema de Nyquist"}
\section{Amostragem e Teorema de Nyquist}
Para reconstruir um sinal sem perda, amostre a mais que o dobro da maior frequência. Abaixo disso, aliasing irreversível. A ponte fundadora entre o contínuo e o digital.
\prov{relações: parte\_de processamento\_digital\_sinais; explica conversor\_ad\_da\_disc; parte\_de eng\_eletronica}
\prov{proveniência: wiki \textperiodcentered{} fato \textperiodcentered{} confiança 1.0 \textperiodcentered{} domínio engenharia \textperiodcentered{} história 4 versões no jornal}

%CRISTAL {"arestas": [{"alvo": "transistor", "confianca": 1.0, "epistemico": "fato", "origem": "wiki", "rel": "usa"}, {"alvo": "eletronica", "confianca": 1.0, "epistemico": "fato", "origem": "wiki", "rel": "parte_de"}], "confianca": 1.0, "contraexemplos": [], "descricao": "Um circuito que aumenta a amplitude de um sinal usando energia de uma fonte — o transistor em seu papel analógico. A base do áudio e do rádio.", "epistemico": "fato", "exemplos": [], "id": "amplificador", "memoria": "semantica", "meta": {"dominio": "engenharia", "fonte": ""}, "origem": "wiki", "palavras_chave": [], "sinonimos": [], "tipo": "conceito", "titulo": "Amplificador"}
\section{Amplificador}
Um circuito que aumenta a amplitude de um sinal usando energia de uma fonte — o transistor em seu papel analógico. A base do áudio e do rádio.
\prov{relações: usa transistor; parte\_de eletronica}
\prov{proveniência: wiki \textperiodcentered{} fato \textperiodcentered{} confiança 1.0 \textperiodcentered{} domínio engenharia \textperiodcentered{} história 3 versões no jornal}

%CRISTAL {"arestas": [{"alvo": "amplificador", "confianca": 1.0, "epistemico": "fato", "origem": "wiki", "rel": "especializa"}, {"alvo": "realimentacao_eletronica", "confianca": 1.0, "epistemico": "fato", "origem": "wiki", "rel": "usa"}, {"alvo": "eng_eletronica", "confianca": 1.0, "epistemico": "fato", "origem": "curriculo", "rel": "parte_de"}], "confianca": 1.0, "contraexemplos": [], "descricao": "O bloco analógico universal: ganho enorme, duas entradas, realimentação. Soma, integra, filtra e compara — o 'tijolo' com que se constrói quase todo condicionamento de sinal.", "epistemico": "fato", "exemplos": [], "id": "amplificador_operacional", "memoria": "semantica", "meta": {"dominio": "engenharia", "fonte": ""}, "origem": "wiki", "palavras_chave": [], "sinonimos": [], "tipo": "conceito", "titulo": "Amplificador Operacional"}
\section{Amplificador Operacional}
O bloco analógico universal: ganho enorme, duas entradas, realimentação. Soma, integra, filtra e compara — o 'tijolo' com que se constrói quase todo condicionamento de sinal.
\prov{relações: especializa amplificador; usa realimentacao\_eletronica; parte\_de eng\_eletronica}
\prov{proveniência: wiki \textperiodcentered{} fato \textperiodcentered{} confiança 1.0 \textperiodcentered{} domínio engenharia \textperiodcentered{} história 4 versões no jornal}

%CRISTAL {"arestas": [{"alvo": "ipc", "confianca": 1.0, "epistemico": "fato", "origem": "paper", "rel": "explica"}, {"alvo": "ondas_eletromagneticas", "confianca": 1.0, "epistemico": "fato", "origem": "wiki", "rel": "usa"}, {"alvo": "eletromagnetismo_eng", "confianca": 1.0, "epistemico": "fato", "origem": "wiki", "rel": "consequencia"}, {"alvo": "eng_telecomunicacoes", "confianca": 1.0, "epistemico": "fato", "origem": "curriculo", "rel": "parte_de"}], "confianca": 1.0, "contraexemplos": [], "descricao": "O transdutor entre corrente no fio e onda no espaço (Maxwell irradiado): ganho, diretividade, diagrama de radiação, polarização, dipolo e arranjos. A boca e o ouvido de todo rádio.", "epistemico": "fato", "exemplos": [], "id": "antena", "memoria": "semantica", "meta": {"dominio": "engenharia", "fonte": ""}, "origem": "wiki", "palavras_chave": [], "sinonimos": [], "tipo": "conceito", "titulo": "Antenas"}
\section{Antenas}
O transdutor entre corrente no fio e onda no espaço (Maxwell irradiado): ganho, diretividade, diagrama de radiação, polarização, dipolo e arranjos. A boca e o ouvido de todo rádio.
\prov{relações: explica ipc; usa ondas\_eletromagneticas; consequencia eletromagnetismo\_eng; parte\_de eng\_telecomunicacoes}
\prov{proveniência: wiki \textperiodcentered{} fato \textperiodcentered{} confiança 1.0 \textperiodcentered{} domínio engenharia \textperiodcentered{} história 36 versões no jornal}

%CRISTAL {"arestas": [{"alvo": "smart_grid", "confianca": 1.0, "epistemico": "fato", "origem": "wiki", "rel": "usa"}, {"alvo": "energia_eolica", "confianca": 1.0, "epistemico": "fato", "origem": "wiki", "rel": "relaciona"}, {"alvo": "eng_sistemas_energia", "confianca": 1.0, "epistemico": "fato", "origem": "curriculo", "rel": "parte_de"}, {"alvo": "energia_eletrica", "confianca": 1.0, "epistemico": "fato", "origem": "wiki", "rel": "usa"}], "confianca": 1.0, "contraexemplos": [], "descricao": "Guardar energia para casar geração e consumo no tempo: baterias, bombeamento hidráulico, hidrogênio. A peça que falta para uma rede dominada por sol e vento intermitentes.", "epistemico": "fato", "exemplos": [], "id": "armazenamento_de_energia", "memoria": "semantica", "meta": {"dominio": "engenharia", "fonte": ""}, "origem": "wiki", "palavras_chave": [], "sinonimos": [], "tipo": "conceito", "titulo": "Armazenamento de Energia"}
\section{Armazenamento de Energia}
Guardar energia para casar geração e consumo no tempo: baterias, bombeamento hidráulico, hidrogênio. A peça que falta para uma rede dominada por sol e vento intermitentes.
\prov{relações: usa smart\_grid; relaciona energia\_eolica; parte\_de eng\_sistemas\_energia; usa energia\_eletrica}
\prov{proveniência: wiki \textperiodcentered{} fato \textperiodcentered{} confiança 1.0 \textperiodcentered{} domínio engenharia \textperiodcentered{} história 6 versões no jornal}

%CRISTAL {"arestas": [{"alvo": "armazenamento_de_energia", "confianca": 1.0, "epistemico": "fato", "origem": "wiki", "rel": "usa"}, {"alvo": "estabilidade_da_rede", "confianca": 1.0, "epistemico": "fato", "origem": "wiki", "rel": "explica"}, {"alvo": "lyapunov_da_orbita", "confianca": 1.0, "epistemico": "fato", "origem": "wiki", "rel": "usa"}], "confianca": 1.0, "contraexemplos": [], "descricao": "O armazenamento injeta ou absorve potência em milissegundos, adicionando AMORTECIMENTO ao ângulo de potência: empurra o Lyapunov da rede para mais negativo, segurando-a na borda/cristal estável contra a intermitência das renováveis. É o cristal que mantém o Salem da rede na borda. Verificado: base 0.1 + storage 0.3 → λ=−0.40 (cristal).", "epistemico": "inferencia", "exemplos": [], "id": "armazenamento_estabiliza_rede", "memoria": "semantica", "meta": {"dominio": "engenharia", "fonte": "distribuição e armazenamento", "lambda_com_storage": -0.4, "lambda_sem_storage": -0.1}, "origem": "wiki", "palavras_chave": [], "sinonimos": [], "tipo": "conceito", "titulo": "Armazenamento Estabiliza a Rede (o amortecimento)"}
\section{Armazenamento Estabiliza a Rede (o amortecimento)}
O armazenamento injeta ou absorve potência em milissegundos, adicionando AMORTECIMENTO ao ângulo de potência: empurra o Lyapunov da rede para mais negativo, segurando-a na borda/cristal estável contra a intermitência das renováveis. É o cristal que mantém o Salem da rede na borda. Verificado: base 0.1 + storage 0.3 \ensuremath{\to} \ensuremath{\lambda}=\ensuremath{-}0.40 (cristal).
\prov{relações: usa armazenamento\_de\_energia; explica estabilidade\_da\_rede; usa lyapunov\_da\_orbita}
\prov{proveniência: wiki \textperiodcentered{} distribuição e armazenamento \textperiodcentered{} inferencia \textperiodcentered{} confiança 1.0 \textperiodcentered{} domínio engenharia \textperiodcentered{} história 8 versões no jornal}

%CRISTAL {"arestas": [{"alvo": "armazenamento_o_cristal", "confianca": 1.0, "epistemico": "fato", "origem": "wiki", "rel": "eh_um"}], "confianca": 1.0, "contraexemplos": [], "descricao": "Eletrólise faz H₂ com a energia que sobra; depois célula a combustível ou turbina o reconverte. η_rt baixa (~40%), mas estoca MUITO e por muito tempo (sazonal) — o oposto do supercapacitor.", "epistemico": "fato", "exemplos": [], "id": "armazenamento_hidrogenio", "memoria": "semantica", "meta": {"dominio": "engenharia", "fonte": "distribuição e armazenamento"}, "origem": "wiki", "palavras_chave": [], "sinonimos": [], "tipo": "conceito", "titulo": "Armazenamento em Hidrogênio (power-to-gas)"}
\section{Armazenamento em Hidrogênio (power-to-gas)}
Eletrólise faz H\ensuremath{_{2}} com a energia que sobra; depois célula a combustível ou turbina o reconverte. \ensuremath{\eta}\_rt baixa (\~{}40\%), mas estoca MUITO e por muito tempo (sazonal) — o oposto do supercapacitor.
\prov{relações: eh\_um armazenamento\_o\_cristal}
\prov{proveniência: wiki \textperiodcentered{} distribuição e armazenamento \textperiodcentered{} fato \textperiodcentered{} confiança 1.0 \textperiodcentered{} domínio engenharia \textperiodcentered{} história 4 versões no jornal}

%CRISTAL {"arestas": [{"alvo": "armazenamento_de_energia", "confianca": 1.0, "epistemico": "fato", "origem": "wiki", "rel": "explica"}, {"alvo": "cristal_memoria", "confianca": 1.0, "epistemico": "fato", "origem": "wiki", "rel": "relaciona"}, {"alvo": "cristalografia", "confianca": 1.0, "epistemico": "fato", "origem": "wiki", "rel": "relaciona"}, {"alvo": "corpo_mineral", "confianca": 1.0, "epistemico": "fato", "origem": "wiki", "rel": "usa"}, {"alvo": "orbita_de_partida", "confianca": 1.0, "epistemico": "fato", "origem": "wiki", "rel": "relaciona"}, {"alvo": "integridade_e_norma_conservada", "confianca": 1.0, "epistemico": "fato", "origem": "wiki", "rel": "relaciona"}], "confianca": 1.0, "contraexemplos": [], "descricao": "Guardar energia é segurar um estado estável — o mesmo papel que o CRISTAL faz com dados na memória da máquina mineral. Carregar/descarregar = escrever/ler; a eficiência round-trip η_rt<1 é a integridade (o que se perde no ciclo); a autodescarga é a norma vazando (o cristal drenando devagar, λ<0). Verificado: reservatório carrega 100 → devolve 80 (η_rt=80%).", "epistemico": "inferencia", "exemplos": [], "id": "armazenamento_o_cristal", "memoria": "semantica", "meta": {"autodescarga_lambda": -0.05, "dominio": "engenharia", "eta_rt": 0.8, "fonte": "distribuição e armazenamento", "reservatorio_recuperado": 80.0}, "origem": "wiki", "palavras_chave": [], "sinonimos": [], "tipo": "conceito", "titulo": "Armazenamento = o Cristal da Energia"}
\section{Armazenamento = o Cristal da Energia}
Guardar energia é segurar um estado estável — o mesmo papel que o CRISTAL faz com dados na memória da máquina mineral. Carregar/descarregar = escrever/ler; a eficiência round-trip \ensuremath{\eta}\_rt<1 é a integridade (o que se perde no ciclo); a autodescarga é a norma vazando (o cristal drenando devagar, \ensuremath{\lambda}<0). Verificado: reservatório carrega 100 \ensuremath{\to} devolve 80 (\ensuremath{\eta}\_rt=80\%).
\prov{relações: explica armazenamento\_de\_energia; relaciona cristal\_memoria; relaciona cristalografia; usa corpo\_mineral; relaciona orbita\_de\_partida; relaciona integridade\_e\_norma\_conservada}
\prov{proveniência: wiki \textperiodcentered{} distribuição e armazenamento \textperiodcentered{} inferencia \textperiodcentered{} confiança 1.0 \textperiodcentered{} domínio engenharia \textperiodcentered{} história 11 versões no jornal}

%CRISTAL {"arestas": [{"alvo": "redes_5g_nr", "confianca": 1.0, "epistemico": "fato", "origem": "wiki", "rel": "parte_de"}, {"alvo": "dupla_conectividade_nr", "confianca": 1.0, "epistemico": "fato", "origem": "wiki", "rel": "usa"}], "confianca": 1.0, "contraexemplos": [], "descricao": "NSA ancora o 5G NR no núcleo 4G/LTE existente (implantação rápida, dupla conectividade); SA usa o núcleo 5G novo (5GC) de ponta a ponta — só o SA entrega slicing, URLLC e a arquitetura de serviços plena.", "epistemico": "inferencia", "exemplos": [], "id": "arquitetura_5g_sa_nsa", "memoria": "semantica", "meta": {"dominio": "engenharia", "fonte": ""}, "origem": "wiki", "palavras_chave": [], "sinonimos": [], "tipo": "conceito", "titulo": "5G SA e NSA (Standalone / Non-Standalone)"}
\section{5G SA e NSA (Standalone / Non-Standalone)}
NSA ancora o 5G NR no núcleo 4G/LTE existente (implantação rápida, dupla conectividade); SA usa o núcleo 5G novo (5GC) de ponta a ponta — só o SA entrega slicing, URLLC e a arquitetura de serviços plena.
\prov{relações: parte\_de redes\_5g\_nr; usa dupla\_conectividade\_nr}
\prov{proveniência: wiki \textperiodcentered{} inferencia \textperiodcentered{} confiança 1.0 \textperiodcentered{} domínio engenharia \textperiodcentered{} história 3 versões no jornal}

%CRISTAL {"arestas": [{"alvo": "orbita_de_objeto_real", "confianca": 1.0, "epistemico": "fato", "origem": "wiki", "rel": "parte_de"}, {"alvo": "orbita_de_partida", "confianca": 1.0, "epistemico": "fato", "origem": "wiki", "rel": "relaciona"}], "confianca": 1.0, "contraexemplos": [], "descricao": "r(θ) é o raio máximo da silhueta em cada ângulo, do centroide — o contorno como função periódica, a ÓRBITA fechada do objeto. É invariante a translação (usa o centroide) e a escala (normaliza pela média). A irregularidade (std/média de r) mede o quão espinhosa é a forma. Verif.: a boneca russa irreg~0.19 (oval), o floco-estrela irreg~0.55 (espinhoso).", "epistemico": "fato", "exemplos": [], "id": "assinatura_radial_objeto", "memoria": "semantica", "meta": {"dominio": "imagem", "fonte": "órbitas de objetos"}, "origem": "wiki", "palavras_chave": [], "sinonimos": [], "tipo": "conceito", "titulo": "Assinatura Radial r(θ) (o contorno em polares)"}
\section{Assinatura Radial r(\ensuremath{\theta}) (o contorno em polares)}
r(\ensuremath{\theta}) é o raio máximo da silhueta em cada ângulo, do centroide — o contorno como função periódica, a ÓRBITA fechada do objeto. É invariante a translação (usa o centroide) e a escala (normaliza pela média). A irregularidade (std/média de r) mede o quão espinhosa é a forma. Verif.: a boneca russa irreg\~{}0.19 (oval), o floco-estrela irreg\~{}0.55 (espinhoso).
\prov{relações: parte\_de orbita\_de\_objeto\_real; relaciona orbita\_de\_partida}
\prov{proveniência: wiki \textperiodcentered{} órbitas de objetos \textperiodcentered{} fato \textperiodcentered{} confiança 1.0 \textperiodcentered{} domínio imagem \textperiodcentered{} história 3 versões no jornal}

%CRISTAL {"arestas": [{"alvo": "robotica", "confianca": 1.0, "epistemico": "fato", "origem": "wiki", "rel": "parte_de"}, {"alvo": "sensor", "confianca": 1.0, "epistemico": "fato", "origem": "wiki", "rel": "relaciona"}], "confianca": 1.0, "contraexemplos": [], "descricao": "O órgão que transforma a decisão em movimento: motor, servo, músculo artificial. É o espelho do sensor — onde o sensor faz o mundo entrar, o atuador faz a decisão sair para o mundo.", "epistemico": "fato", "exemplos": [], "id": "atuador", "memoria": "semantica", "meta": {"dominio": "robotica", "fonte": "robótica e sistemas autônomos"}, "origem": "wiki", "palavras_chave": [], "sinonimos": [], "tipo": "conceito", "titulo": "Atuador"}
\section{Atuador}
O órgão que transforma a decisão em movimento: motor, servo, músculo artificial. É o espelho do sensor — onde o sensor faz o mundo entrar, o atuador faz a decisão sair para o mundo.
\prov{relações: parte\_de robotica; relaciona sensor}
\prov{proveniência: wiki \textperiodcentered{} robótica e sistemas autônomos \textperiodcentered{} fato \textperiodcentered{} confiança 1.0 \textperiodcentered{} domínio robotica \textperiodcentered{} história 3 versões no jornal}

%CRISTAL {"arestas": [{"alvo": "atuador", "confianca": 1.0, "epistemico": "fato", "origem": "wiki", "rel": "explica"}, {"alvo": "colapso", "confianca": 1.0, "epistemico": "fato", "origem": "wiki", "rel": "eh_um"}, {"alvo": "sensor_e_colapso", "confianca": 1.0, "epistemico": "fato", "origem": "wiki", "rel": "relaciona"}, {"alvo": "bai", "confianca": 1.0, "epistemico": "fato", "origem": "wiki", "rel": "usa"}], "confianca": 1.0, "contraexemplos": [], "descricao": "Assim como o sensor colapsa o mundo em bit (entrada), o atuador colapsa a decisão em movimento (saída): escolhe UMA ação entre as possíveis e a executa. O colapso do BAI é bidirecional — perceber e agir são o mesmo Ψ, um para dentro e um para fora. Agir é comprometer-se com um ramo.", "epistemico": "inferencia", "exemplos": [], "id": "atuador_colapso", "memoria": "semantica", "meta": {"dominio": "robotica", "fonte": "robótica e sistemas autônomos"}, "origem": "wiki", "palavras_chave": [], "sinonimos": [], "tipo": "conceito", "titulo": "O Atuador é um Colapso de Saída"}
\section{O Atuador é um Colapso de Saída}
Assim como o sensor colapsa o mundo em bit (entrada), o atuador colapsa a decisão em movimento (saída): escolhe UMA ação entre as possíveis e a executa. O colapso do BAI é bidirecional — perceber e agir são o mesmo \ensuremath{\Psi}, um para dentro e um para fora. Agir é comprometer-se com um ramo.
\prov{relações: explica atuador; eh\_um colapso; relaciona sensor\_e\_colapso; usa bai}
\prov{proveniência: wiki \textperiodcentered{} robótica e sistemas autônomos \textperiodcentered{} inferencia \textperiodcentered{} confiança 1.0 \textperiodcentered{} domínio robotica \textperiodcentered{} história 5 versões no jornal}

%CRISTAL {"arestas": [{"alvo": "registro_5g", "confianca": 1.0, "epistemico": "fato", "origem": "wiki", "rel": "parte_de"}, {"alvo": "seguranca_5g", "confianca": 1.0, "epistemico": "fato", "origem": "wiki", "rel": "especializa"}, {"alvo": "funcoes_de_rede_5gc", "confianca": 1.0, "epistemico": "fato", "origem": "wiki", "rel": "usa"}], "confianca": 1.0, "contraexemplos": [], "descricao": "O desafio-resposta MÚTUO: a rede prova sua identidade ao UE e vice-versa, derivando chaves a partir do segredo do SIM (USIM) e do vetor que o AUSF/UDM gera. Fecha o ataque de estação-base falsa que assombrava o 4G.", "epistemico": "fato", "exemplos": [], "id": "autenticacao_5g_aka", "memoria": "semantica", "meta": {"dominio": "engenharia", "fonte": ""}, "origem": "wiki", "palavras_chave": [], "sinonimos": [], "tipo": "conceito", "titulo": "Autenticação 5G-AKA"}
\section{Autenticação 5G-AKA}
O desafio-resposta MÚTUO: a rede prova sua identidade ao UE e vice-versa, derivando chaves a partir do segredo do SIM (USIM) e do vetor que o AUSF/UDM gera. Fecha o ataque de estação-base falsa que assombrava o 4G.
\prov{relações: parte\_de registro\_5g; especializa seguranca\_5g; usa funcoes\_de\_rede\_5gc}
\prov{proveniência: wiki \textperiodcentered{} fato \textperiodcentered{} confiança 1.0 \textperiodcentered{} domínio engenharia \textperiodcentered{} história 4 versões no jornal}

%CRISTAL {"arestas": [{"alvo": "impressao_3d", "confianca": 1.0, "epistemico": "fato", "origem": "wiki", "rel": "usa"}, {"alvo": "arquitetura_von_neumann", "confianca": 1.0, "epistemico": "fato", "origem": "wiki", "rel": "usa"}, {"alvo": "reproducao_tiffany", "confianca": 1.0, "epistemico": "fato", "origem": "wiki", "rel": "relaciona"}], "confianca": 1.0, "contraexemplos": [], "descricao": "Uma impressora 3D que imprime as próprias peças (RepRap) é o construtor universal de von Neumann encarnado: uma máquina que se reproduz. A manufatura fecha o laço sobre si mesma — o mesmo gesto da broca-so que se constrói e da reprodução da Tiffany.", "epistemico": "inferencia", "exemplos": [], "id": "auto_replicacao_reprap", "memoria": "semantica", "meta": {"dominio": "manufatura", "fonte": "impressão 3D e manufatura"}, "origem": "wiki", "palavras_chave": [], "sinonimos": [], "tipo": "conceito", "titulo": "Auto-Replicação (RepRap = a máquina que se faz)"}
\section{Auto-Replicação (RepRap = a máquina que se faz)}
Uma impressora 3D que imprime as próprias peças (RepRap) é o construtor universal de von Neumann encarnado: uma máquina que se reproduz. A manufatura fecha o laço sobre si mesma — o mesmo gesto da broca-so que se constrói e da reprodução da Tiffany.
\prov{relações: usa impressao\_3d; usa arquitetura\_von\_neumann; relaciona reproducao\_tiffany}
\prov{proveniência: wiki \textperiodcentered{} impressão 3D e manufatura \textperiodcentered{} inferencia \textperiodcentered{} confiança 1.0 \textperiodcentered{} domínio manufatura \textperiodcentered{} história 4 versões no jornal}

%CRISTAL {"arestas": [{"alvo": "rede_autonoma_broca", "confianca": 1.0, "epistemico": "fato", "origem": "wiki", "rel": "parte_de"}, {"alvo": "checklist_sistema_autonomo_no_ar", "confianca": 1.0, "epistemico": "fato", "origem": "wiki", "rel": "usa"}, {"alvo": "sistema_e_lei_nao_dono", "confianca": 1.0, "epistemico": "fato", "origem": "wiki", "rel": "confirma"}, {"alvo": "orquestrador_roteador_regex", "confianca": 1.0, "epistemico": "fato", "origem": "wiki", "rel": "relaciona"}, {"alvo": "ciclo_autoconsumo", "confianca": 1.0, "epistemico": "fato", "origem": "wiki", "rel": "usa"}], "confianca": 1.0, "contraexemplos": [], "descricao": "Autonomia se conquista REMOVENDO, não adicionando: sem scheduler central, sem buffer, sem treino do operador. A ordem emerge do recall + da régua de ouro; impô-la mata a autonomia. Auto-gerenciar = deixar a álgebra coordenar, não um maestro.", "epistemico": "inferencia", "exemplos": [], "id": "autonomia_subtrativa", "memoria": "semantica", "meta": {"dominio": "arquitetura", "fonte": "semilla_rede_autonoma"}, "origem": "wiki", "palavras_chave": [], "sinonimos": [], "tipo": "conceito", "titulo": "Auto-gestão subtrativa (a ordem emerge do recall)"}
\section{Auto-gestão subtrativa (a ordem emerge do recall)}
Autonomia se conquista REMOVENDO, não adicionando: sem scheduler central, sem buffer, sem treino do operador. A ordem emerge do recall + da régua de ouro; impô-la mata a autonomia. Auto-gerenciar = deixar a álgebra coordenar, não um maestro.
\prov{relações: parte\_de rede\_autonoma\_broca; usa checklist\_sistema\_autonomo\_no\_ar; confirma sistema\_e\_lei\_nao\_dono; relaciona orquestrador\_roteador\_regex; usa ciclo\_autoconsumo}
\prov{proveniência: wiki \textperiodcentered{} semilla\_rede\_autonoma \textperiodcentered{} inferencia \textperiodcentered{} confiança 1.0 \textperiodcentered{} domínio arquitetura \textperiodcentered{} história 6 versões no jornal}

%CRISTAL {"arestas": [{"alvo": "bateria_de_tracao", "confianca": 1.0, "epistemico": "fato", "origem": "wiki", "rel": "usa"}, {"alvo": "eficiencia_round_trip", "confianca": 1.0, "epistemico": "fato", "origem": "wiki", "rel": "relaciona"}], "confianca": 1.0, "contraexemplos": [], "descricao": "Quanto o carro anda com uma carga — a energia da bateria dividida pelo consumo (Wh/km). É o tamanho do reservatório visto da estrada; a ansiedade de autonomia é o medo do cristal esvaziar longe da rede.", "epistemico": "fato", "exemplos": [], "id": "autonomia_veicular", "memoria": "semantica", "meta": {"dominio": "engenharia", "fonte": "mobilidade elétrica"}, "origem": "wiki", "palavras_chave": [], "sinonimos": [], "tipo": "conceito", "titulo": "Autonomia (o alcance do cristal)"}
\section{Autonomia (o alcance do cristal)}
Quanto o carro anda com uma carga — a energia da bateria dividida pelo consumo (Wh/km). É o tamanho do reservatório visto da estrada; a ansiedade de autonomia é o medo do cristal esvaziar longe da rede.
\prov{relações: usa bateria\_de\_tracao; relaciona eficiencia\_round\_trip}
\prov{proveniência: wiki \textperiodcentered{} mobilidade elétrica \textperiodcentered{} fato \textperiodcentered{} confiança 1.0 \textperiodcentered{} domínio engenharia \textperiodcentered{} história 3 versões no jornal}

%CRISTAL {"arestas": [{"alvo": "repositorio_broca_os", "confianca": 1.0, "epistemico": "fato", "origem": "wiki", "rel": "parte_de"}, {"alvo": "bai_tese_mae", "confianca": 1.0, "epistemico": "fato", "origem": "wiki", "rel": "usa"}, {"alvo": "banda_exclusiva_por_tecido", "confianca": 1.0, "epistemico": "fato", "origem": "wiki", "rel": "usa"}, {"alvo": "sistema_e_lei_nao_dono", "confianca": 1.0, "epistemico": "fato", "origem": "wiki", "rel": "confirma"}], "confianca": 1.0, "contraexemplos": [], "descricao": "Quem prova posse do TECIDO é a RAIZ da sua banda: a posse É a credencial (banda=sha256(tecido)), com todas as capacidades sobre a própria comunidade — sem precisar que ninguém conceda. Auto-soberana; a encarnação de 'é a lei, não o dono'. No BAI, o tecido é o grant-raiz do reticulado.", "epistemico": "fato", "exemplos": [], "id": "autoridade_intrinseca", "memoria": "semantica", "meta": {"dominio": "arquitetura", "fonte": "goldchain/bai_autoridade.py"}, "origem": "wiki", "palavras_chave": [], "sinonimos": [], "tipo": "conceito", "titulo": "Autoridade Intrínseca (posse do tecido)"}
\section{Autoridade Intrínseca (posse do tecido)}
Quem prova posse do TECIDO é a RAIZ da sua banda: a posse É a credencial (banda=sha256(tecido)), com todas as capacidades sobre a própria comunidade — sem precisar que ninguém conceda. Auto-soberana; a encarnação de 'é a lei, não o dono'. No BAI, o tecido é o grant-raiz do reticulado.
\prov{relações: parte\_de repositorio\_broca\_os; usa bai\_tese\_mae; usa banda\_exclusiva\_por\_tecido; confirma sistema\_e\_lei\_nao\_dono}
\prov{proveniência: wiki \textperiodcentered{} goldchain/bai\_autoridade.py \textperiodcentered{} fato \textperiodcentered{} confiança 1.0 \textperiodcentered{} domínio arquitetura \textperiodcentered{} história 5 versões no jornal}

%CRISTAL {"arestas": [{"alvo": "escalabilidade", "confianca": 1.0, "epistemico": "fato", "origem": "wiki", "rel": "usa"}, {"alvo": "tolerancia_a_falhas", "confianca": 1.0, "epistemico": "fato", "origem": "wiki", "rel": "usa"}, {"alvo": "utilizacao_e_borda", "confianca": 1.0, "epistemico": "fato", "origem": "wiki", "rel": "usa"}, {"alvo": "carga_coordenada", "confianca": 1.0, "epistemico": "fato", "origem": "wiki", "rel": "relaciona"}, {"alvo": "data_center", "confianca": 1.0, "epistemico": "fato", "origem": "wiki", "rel": "usa"}], "confianca": 1.0, "contraexemplos": [], "descricao": "Espalhar as requisições entre muitos servidores para aplainar a utilização de pico — a MESMA matemática do smart charging da cidade. Concentrar tudo num nó cria um hotspot (ρ→caos, latência explode); distribuir mantém cada nó no cristal. Verif.: 900 req/s concentrado=caos, em 10 nós=borda estável (ρ=0.9).", "epistemico": "inferencia", "exemplos": [], "id": "balanceamento_de_carga", "memoria": "semantica", "meta": {"classe_balanceado": "borda (latência dispara)", "classe_concentrado": "caos (nó saturado, fila sem limite)", "dominio": "engenharia", "fonte": "nuvem e data centers", "rho_pico_balanceado": 0.9}, "origem": "wiki", "palavras_chave": [], "sinonimos": [], "tipo": "conceito", "titulo": "Balanceamento de Carga (encher o vale)"}
\section{Balanceamento de Carga (encher o vale)}
Espalhar as requisições entre muitos servidores para aplainar a utilização de pico — a MESMA matemática do smart charging da cidade. Concentrar tudo num nó cria um hotspot (\ensuremath{\rho}\ensuremath{\to}caos, latência explode); distribuir mantém cada nó no cristal. Verif.: 900 req/s concentrado=caos, em 10 nós=borda estável (\ensuremath{\rho}=0.9).
\prov{relações: usa escalabilidade; usa tolerancia\_a\_falhas; usa utilizacao\_e\_borda; relaciona carga\_coordenada; usa data\_center}
\prov{proveniência: wiki \textperiodcentered{} nuvem e data centers \textperiodcentered{} inferencia \textperiodcentered{} confiança 1.0 \textperiodcentered{} domínio engenharia \textperiodcentered{} história 10 versões no jornal}

%CRISTAL {"arestas": [{"alvo": "antena", "confianca": 1.0, "epistemico": "fato", "origem": "wiki", "rel": "usa"}, {"alvo": "propagacao_de_ondas", "confianca": 1.0, "epistemico": "fato", "origem": "wiki", "rel": "usa"}, {"alvo": "razao_sinal_ruido", "confianca": 1.0, "epistemico": "fato", "origem": "wiki", "rel": "usa"}, {"alvo": "eng_telecomunicacoes", "confianca": 1.0, "epistemico": "fato", "origem": "curriculo", "rel": "parte_de"}], "confianca": 1.0, "contraexemplos": [], "descricao": "A contabilidade do enlace: potência transmitida + ganhos − perdas ≥ sensibilidade + margem. Diz se o rádio fecha — do Wi-Fi de casa ao sinal de uma sonda a bilhões de km.", "epistemico": "inferencia", "exemplos": [], "id": "balanco_de_enlace", "memoria": "semantica", "meta": {"dominio": "engenharia", "fonte": ""}, "origem": "wiki", "palavras_chave": [], "sinonimos": [], "tipo": "conceito", "titulo": "Balanço de Enlace (Link Budget)"}
\section{Balanço de Enlace (Link Budget)}
A contabilidade do enlace: potência transmitida + ganhos \ensuremath{-} perdas \ensuremath{\ge} sensibilidade + margem. Diz se o rádio fecha — do Wi-Fi de casa ao sinal de uma sonda a bilhões de km.
\prov{relações: usa antena; usa propagacao\_de\_ondas; usa razao\_sinal\_ruido; parte\_de eng\_telecomunicacoes}
\prov{proveniência: wiki \textperiodcentered{} inferencia \textperiodcentered{} confiança 1.0 \textperiodcentered{} domínio engenharia \textperiodcentered{} história 5 versões no jornal}

%CRISTAL {"arestas": [{"alvo": "rede_autonoma_broca", "confianca": 1.0, "epistemico": "fato", "origem": "wiki", "rel": "parte_de"}, {"alvo": "rede_limpa_bump", "confianca": 1.0, "epistemico": "fato", "origem": "wiki", "rel": "usa"}, {"alvo": "a_banda_e_o_tecido", "confianca": 1.0, "epistemico": "fato", "origem": "wiki", "rel": "usa"}, {"alvo": "isolamento_base_cliente", "confianca": 1.0, "epistemico": "fato", "origem": "wiki", "rel": "usa"}], "confianca": 1.0, "contraexemplos": [], "descricao": "A identidade vira canal: a banda de um par é sha256 do seu tecido; quem tem o tecido sintoniza, quem não tem lê ruído. Assim o CLIENTE ganha uma banda EXCLUSIVA derivada da sua base isolada — a mesma teoria da rede limpa (bump ⊕ keystream(banda)) aplicada ao nível do olho. Banda = pertencimento, não permissão.", "epistemico": "inferencia", "exemplos": [], "id": "banda_exclusiva_por_tecido", "memoria": "semantica", "meta": {"dominio": "arquitetura", "fonte": "semilla_rede_autonoma"}, "origem": "wiki", "palavras_chave": [], "sinonimos": [], "tipo": "conceito", "titulo": "Banda exclusiva = sha256(tecido)"}
\section{Banda exclusiva = sha256(tecido)}
A identidade vira canal: a banda de um par é sha256 do seu tecido; quem tem o tecido sintoniza, quem não tem lê ruído. Assim o CLIENTE ganha uma banda EXCLUSIVA derivada da sua base isolada — a mesma teoria da rede limpa (bump \ensuremath{\oplus} keystream(banda)) aplicada ao nível do olho. Banda = pertencimento, não permissão.
\prov{relações: parte\_de rede\_autonoma\_broca; usa rede\_limpa\_bump; usa a\_banda\_e\_o\_tecido; usa isolamento\_base\_cliente}
\prov{proveniência: wiki \textperiodcentered{} semilla\_rede\_autonoma \textperiodcentered{} inferencia \textperiodcentered{} confiança 1.0 \textperiodcentered{} domínio arquitetura \textperiodcentered{} história 5 versões no jornal}

%CRISTAL {"arestas": [{"alvo": "estrutura_de_bandas", "confianca": 1.0, "epistemico": "fato", "origem": "wiki", "rel": "parte_de"}], "confianca": 1.0, "contraexemplos": [], "descricao": "A lacuna de energia entre a banda de valência e a de condução — pequena no semicondutor, grande no isolante, nula no condutor. O parâmetro-chave.", "epistemico": "fato", "exemplos": [], "id": "banda_proibida", "memoria": "semantica", "meta": {"dominio": "engenharia", "fonte": ""}, "origem": "wiki", "palavras_chave": [], "sinonimos": [], "tipo": "conceito", "titulo": "Banda Proibida (Band Gap)"}
\section{Banda Proibida (Band Gap)}
A lacuna de energia entre a banda de valência e a de condução — pequena no semicondutor, grande no isolante, nula no condutor. O parâmetro-chave.
\prov{relações: parte\_de estrutura\_de\_bandas}
\prov{proveniência: wiki \textperiodcentered{} fato \textperiodcentered{} confiança 1.0 \textperiodcentered{} domínio engenharia \textperiodcentered{} história 2 versões no jornal}

%CRISTAL {"arestas": [{"alvo": "repositorio_broca_os", "confianca": 1.0, "epistemico": "fato", "origem": "wiki", "rel": "parte_de"}, {"alvo": "isolamento_base_cliente", "confianca": 1.0, "epistemico": "fato", "origem": "wiki", "rel": "especializa"}, {"alvo": "comunidade_organica_por_banda", "confianca": 1.0, "epistemico": "fato", "origem": "wiki", "rel": "usa"}, {"alvo": "roteamento_multi_tenant", "confianca": 1.0, "epistemico": "fato", "origem": "wiki", "rel": "relaciona"}], "confianca": 1.0, "contraexemplos": [], "descricao": "Cada comunidade tem sua base física isolada — clientes/<banda-hex>/apps/… + manifest.json (dono/raiz). Bandas não se cruzam no disco: o app de uma comunidade nunca cai na base de outra. O isolamento_base_cliente do tiffany_platform estendido à rede viva.", "epistemico": "fato", "exemplos": [], "id": "base_cliente_por_banda", "memoria": "semantica", "meta": {"dominio": "arquitetura", "fonte": "semilla_repositorio_broca"}, "origem": "wiki", "palavras_chave": [], "sinonimos": [], "tipo": "conceito", "titulo": "Base Isolada por Banda (clientes/<banda>/apps)"}
\section{Base Isolada por Banda (clientes/<banda>/apps)}
Cada comunidade tem sua base física isolada — clientes/<banda-hex>/apps/… + manifest.json (dono/raiz). Bandas não se cruzam no disco: o app de uma comunidade nunca cai na base de outra. O isolamento\_base\_cliente do tiffany\_platform estendido à rede viva.
\prov{relações: parte\_de repositorio\_broca\_os; especializa isolamento\_base\_cliente; usa comunidade\_organica\_por\_banda; relaciona roteamento\_multi\_tenant}
\prov{proveniência: wiki \textperiodcentered{} semilla\_repositorio\_broca \textperiodcentered{} fato \textperiodcentered{} confiança 1.0 \textperiodcentered{} domínio arquitetura \textperiodcentered{} história 5 versões no jornal}

%CRISTAL {"arestas": [{"alvo": "bateria_de_tracao", "confianca": 1.0, "epistemico": "fato", "origem": "wiki", "rel": "explica"}, {"alvo": "armazenamento_o_cristal", "confianca": 1.0, "epistemico": "fato", "origem": "wiki", "rel": "eh_um"}, {"alvo": "corpo_mineral", "confianca": 1.0, "epistemico": "fato", "origem": "wiki", "rel": "usa"}], "confianca": 1.0, "contraexemplos": [], "descricao": "A bateria de tração é o MESMO cristal de energia do armazenamento estacionário, agora móvel: segura um estado estável, o estado de carga (SoC) é a norma guardada, e carregar/descarregar é escrever/ler. O carro é um cristal de energia que anda.", "epistemico": "inferencia", "exemplos": [], "id": "bateria_cristal_movel", "memoria": "semantica", "meta": {"dominio": "engenharia", "fonte": "mobilidade elétrica"}, "origem": "wiki", "palavras_chave": [], "sinonimos": [], "tipo": "conceito", "titulo": "A Bateria = o Cristal com Rodas"}
\section{A Bateria = o Cristal com Rodas}
A bateria de tração é o MESMO cristal de energia do armazenamento estacionário, agora móvel: segura um estado estável, o estado de carga (SoC) é a norma guardada, e carregar/descarregar é escrever/ler. O carro é um cristal de energia que anda.
\prov{relações: explica bateria\_de\_tracao; eh\_um armazenamento\_o\_cristal; usa corpo\_mineral}
\prov{proveniência: wiki \textperiodcentered{} mobilidade elétrica \textperiodcentered{} inferencia \textperiodcentered{} confiança 1.0 \textperiodcentered{} domínio engenharia \textperiodcentered{} história 4 versões no jornal}

%CRISTAL {"arestas": [{"alvo": "carro_eletrico", "confianca": 1.0, "epistemico": "fato", "origem": "wiki", "rel": "parte_de"}], "confianca": 1.0, "contraexemplos": [], "descricao": "O grande pacote de íon-lítio que move o carro (dezenas de kWh). É a fonte de energia e o item mais caro e pesado; sua densidade energética define a autonomia.", "epistemico": "fato", "exemplos": [], "id": "bateria_de_tracao", "memoria": "semantica", "meta": {"dominio": "engenharia", "fonte": "mobilidade elétrica"}, "origem": "wiki", "palavras_chave": [], "sinonimos": [], "tipo": "conceito", "titulo": "Bateria de Tração"}
\section{Bateria de Tração}
O grande pacote de íon-lítio que move o carro (dezenas de kWh). É a fonte de energia e o item mais caro e pesado; sua densidade energética define a autonomia.
\prov{relações: parte\_de carro\_eletrico}
\prov{proveniência: wiki \textperiodcentered{} mobilidade elétrica \textperiodcentered{} fato \textperiodcentered{} confiança 1.0 \textperiodcentered{} domínio engenharia \textperiodcentered{} história 2 versões no jornal}

%CRISTAL {"arestas": [{"alvo": "armazenamento_o_cristal", "confianca": 1.0, "epistemico": "fato", "origem": "wiki", "rel": "eh_um"}], "confianca": 1.0, "contraexemplos": [], "descricao": "Guarda energia numa reação química reversível (íon-lítio: Li⁺ migrando entre eletrodos). η_rt~90%, resposta rápida; a tecnologia que domina o armazenamento novo (rede, carros, celulares).", "epistemico": "fato", "exemplos": [], "id": "bateria_eletroquimica", "memoria": "semantica", "meta": {"dominio": "engenharia", "fonte": "distribuição e armazenamento"}, "origem": "wiki", "palavras_chave": [], "sinonimos": [], "tipo": "conceito", "titulo": "Bateria Eletroquímica"}
\section{Bateria Eletroquímica}
Guarda energia numa reação química reversível (íon-lítio: Li\ensuremath{^{+}} migrando entre eletrodos). \ensuremath{\eta}\_rt\~{}90\%, resposta rápida; a tecnologia que domina o armazenamento novo (rede, carros, celulares).
\prov{relações: eh\_um armazenamento\_o\_cristal}
\prov{proveniência: wiki \textperiodcentered{} distribuição e armazenamento \textperiodcentered{} fato \textperiodcentered{} confiança 1.0 \textperiodcentered{} domínio engenharia \textperiodcentered{} história 4 versões no jornal}

%CRISTAL {"arestas": [{"alvo": "roteamento", "confianca": 1.0, "epistemico": "fato", "origem": "wiki", "rel": "especializa"}, {"alvo": "internet", "confianca": 1.0, "epistemico": "fato", "origem": "wiki", "rel": "explica"}, {"alvo": "eng_telecomunicacoes", "confianca": 1.0, "epistemico": "fato", "origem": "curriculo", "rel": "parte_de"}], "confianca": 1.0, "contraexemplos": [], "descricao": "O protocolo que costura a internet: sistemas autônomos anunciam quais prefixos alcançam, por política, não só distância. O BGP é a diplomacia — e a fragilidade — do roteamento global.", "epistemico": "fato", "exemplos": [], "id": "bgp", "memoria": "semantica", "meta": {"dominio": "engenharia", "fonte": ""}, "origem": "wiki", "palavras_chave": [], "sinonimos": [], "tipo": "conceito", "titulo": "BGP (Roteamento entre Sistemas Autônomos)"}
\section{BGP (Roteamento entre Sistemas Autônomos)}
O protocolo que costura a internet: sistemas autônomos anunciam quais prefixos alcançam, por política, não só distância. O BGP é a diplomacia — e a fragilidade — do roteamento global.
\prov{relações: especializa roteamento; explica internet; parte\_de eng\_telecomunicacoes}
\prov{proveniência: wiki \textperiodcentered{} fato \textperiodcentered{} confiança 1.0 \textperiodcentered{} domínio engenharia \textperiodcentered{} história 4 versões no jornal}

%CRISTAL {"arestas": [{"alvo": "imagem_como_orbita_gato", "confianca": 1.0, "epistemico": "fato", "origem": "wiki", "rel": "usa"}, {"alvo": "corpo_mineral", "confianca": 1.0, "epistemico": "fato", "origem": "wiki", "rel": "usa"}, {"alvo": "computador_mineral", "confianca": 1.0, "epistemico": "fato", "origem": "wiki", "rel": "relaciona"}], "confianca": 1.0, "contraexemplos": [], "descricao": "Uma biblioteca generativa que faz imagens a partir dos objetos da máquina — rotações (Salem), espirais dos metais σ_n, texturas do gato, o campo da β-numeração — e as combina. Nada é carregado de fora: cada forma nasce de um número mineral. Geradores: onda, interferencia, espiral_metal, gato_textura, campo_beta; misturas: mistura, morph, gato_mistura.", "epistemico": "fato", "exemplos": [], "id": "biblioteca_formas_minerais", "memoria": "semantica", "meta": {"dominio": "imagem", "fonte": "formas minerais"}, "origem": "wiki", "palavras_chave": [], "sinonimos": [], "tipo": "conceito", "titulo": "Biblioteca de Formas Minerais"}
\section{Biblioteca de Formas Minerais}
Uma biblioteca generativa que faz imagens a partir dos objetos da máquina — rotações (Salem), espirais dos metais \ensuremath{\sigma}\_n, texturas do gato, o campo da \ensuremath{\beta}-numeração — e as combina. Nada é carregado de fora: cada forma nasce de um número mineral. Geradores: onda, interferencia, espiral\_metal, gato\_textura, campo\_beta; misturas: mistura, morph, gato\_mistura.
\prov{relações: usa imagem\_como\_orbita\_gato; usa corpo\_mineral; relaciona computador\_mineral}
\prov{proveniência: wiki \textperiodcentered{} formas minerais \textperiodcentered{} fato \textperiodcentered{} confiança 1.0 \textperiodcentered{} domínio imagem \textperiodcentered{} história 4 versões no jornal}

%CRISTAL {"arestas": [{"alvo": "genetica", "confianca": 1.0, "epistemico": "fato", "origem": "wiki", "rel": "usa"}, {"alvo": "crispr_edicao", "confianca": 1.0, "epistemico": "fato", "origem": "wiki", "rel": "usa"}, {"alvo": "celula_computador_mineral", "confianca": 1.0, "epistemico": "fato", "origem": "wiki", "rel": "usa"}], "confianca": 1.0, "contraexemplos": [], "descricao": "fato", "epistemico": {"_desc": true}, "exemplos": [], "id": "biotecnologia", "memoria": "semantica", "meta": {"dominio": "biotecnologia", "fonte": "biotecnologia e engenharia genética"}, "origem": "wiki", "palavras_chave": [], "sinonimos": [], "tipo": "conceito", "titulo": "Biotecnologia"}
\section{Biotecnologia}
fato
\prov{relações: usa genetica; usa crispr\_edicao; usa celula\_computador\_mineral}
\prov{proveniência: wiki \textperiodcentered{} biotecnologia e engenharia genética \textperiodcentered{} \{'\_desc': True\} \textperiodcentered{} confiança 1.0 \textperiodcentered{} domínio biotecnologia \textperiodcentered{} história 4 versões no jornal}

%CRISTAL {"arestas": [{"alvo": "transicoes_reta_bai_newton", "confianca": 1.0, "epistemico": "fato", "origem": "wiki", "rel": "parte_de"}, {"alvo": "numero_de_salem", "confianca": 1.0, "epistemico": "fato", "origem": "wiki", "rel": "usa"}, {"alvo": "julia", "confianca": 1.0, "epistemico": "fato", "origem": "wiki", "rel": "usa"}, {"alvo": "fronteira_e_fail_closed", "confianca": 1.0, "epistemico": "fato", "origem": "wiki", "rel": "relaciona"}, {"alvo": "rastro_da_agulha_e_fractal", "confianca": 1.0, "epistemico": "fato", "origem": "wiki", "rel": "relaciona"}], "confianca": 1.0, "contraexemplos": [], "descricao": "A ponte mais funda: a BORDA da parábola (λ=0) é o mesmo objeto nos três. Na reta mineral é o número de SALEM (ρ=1, no círculo); no Newton é o CONJUNTO DE JULIA (a fronteira fractal entre bacias); no BAI é o FAIL-CLOSED (o orçamento δ=0, a decisão no limite). O λ=0 é a fronteira fractal em toda leitura — onde a ordem encontra o caos.", "epistemico": "fato", "exemplos": [], "id": "borda_e_o_julia", "memoria": "semantica", "meta": {"dominio": "ponte", "fonte": "transições órbita"}, "origem": "wiki", "palavras_chave": [], "sinonimos": [], "tipo": "conceito", "titulo": "A Borda é o Conjunto de Julia (salem ↔ Julia ↔ fail-closed)"}
\section{A Borda é o Conjunto de Julia (salem \ensuremath{\leftrightarrow} Julia \ensuremath{\leftrightarrow} fail-closed)}
A ponte mais funda: a BORDA da parábola (\ensuremath{\lambda}=0) é o mesmo objeto nos três. Na reta mineral é o número de SALEM (\ensuremath{\rho}=1, no círculo); no Newton é o CONJUNTO DE JULIA (a fronteira fractal entre bacias); no BAI é o FAIL-CLOSED (o orçamento \ensuremath{\delta}=0, a decisão no limite). O \ensuremath{\lambda}=0 é a fronteira fractal em toda leitura — onde a ordem encontra o caos.
\prov{relações: parte\_de transicoes\_reta\_bai\_newton; usa numero\_de\_salem; usa julia; relaciona fronteira\_e\_fail\_closed; relaciona rastro\_da\_agulha\_e\_fractal}
\prov{proveniência: wiki \textperiodcentered{} transições órbita \textperiodcentered{} fato \textperiodcentered{} confiança 1.0 \textperiodcentered{} domínio ponte \textperiodcentered{} história 6 versões no jornal}

%CRISTAL {"arestas": [{"alvo": "segmentacao_de_imagem", "confianca": 1.0, "epistemico": "fato", "origem": "wiki", "rel": "parte_de"}, {"alvo": "a_fronteira", "confianca": 1.0, "epistemico": "fato", "origem": "wiki", "rel": "relaciona"}, {"alvo": "julia", "confianca": 1.0, "epistemico": "fato", "origem": "wiki", "rel": "relaciona"}, {"alvo": "flip_duopolinomios", "confianca": 1.0, "epistemico": "fato", "origem": "wiki", "rel": "relaciona"}, {"alvo": "corpo_mineral", "confianca": 1.0, "epistemico": "fato", "origem": "wiki", "rel": "usa"}], "confianca": 1.0, "contraexemplos": [], "descricao": "O limite entre segmentos é onde o gradiente |∇img| é alto — a borda no espaço, o mesmo gradiente do flip. E as bacias de segmentação (watershed) são BACIAS DE ATRAÇÃO: cada pixel flui para um mínimo, como no fractal de Newton, e as cristas entre bacias são as fronteiras.", "epistemico": "inferencia", "exemplos": [], "id": "bordas_fronteira", "memoria": "semantica", "meta": {"dominio": "imagem", "fonte": "segmentação"}, "origem": "wiki", "palavras_chave": [], "sinonimos": [], "tipo": "conceito", "titulo": "Bordas = a Fronteira (o gradiente)"}
\section{Bordas = a Fronteira (o gradiente)}
O limite entre segmentos é onde o gradiente |\ensuremath{\nabla}img| é alto — a borda no espaço, o mesmo gradiente do flip. E as bacias de segmentação (watershed) são BACIAS DE ATRAÇÃO: cada pixel flui para um mínimo, como no fractal de Newton, e as cristas entre bacias são as fronteiras.
\prov{relações: parte\_de segmentacao\_de\_imagem; relaciona a\_fronteira; relaciona julia; relaciona flip\_duopolinomios; usa corpo\_mineral}
\prov{proveniência: wiki \textperiodcentered{} segmentação \textperiodcentered{} inferencia \textperiodcentered{} confiança 1.0 \textperiodcentered{} domínio imagem \textperiodcentered{} história 10 versões no jornal}

%CRISTAL {"arestas": [{"alvo": "cockpit_observatorio", "confianca": 1.0, "epistemico": "fato", "origem": "wiki", "rel": "parte_de"}, {"alvo": "olho_de_cock_web", "confianca": 1.0, "epistemico": "fato", "origem": "wiki", "rel": "relaciona"}], "confianca": 1.0, "contraexemplos": [], "descricao": "A classe servidor: 'Hodge interior contínuo; BSD só na borda'. Mantém o estado Hodge em memória e corre dois workers — _worker (recalcula quando livro/ledger/fase/chain mudam por mtime) e _conductor_worker (o ritmo vivo ~0.8s + acopla o olho invertido). Serve HTTP+SSE. A análise pesada roda em threads para não bloquear o stream.", "epistemico": "fato", "exemplos": [], "id": "broca_cockpit_engine", "memoria": "semantica", "meta": {"autor": "broca-so", "dominio": "cockpit", "fonte": "goldchain/cockpit_server.py:BrocaCockpit"}, "origem": "wiki", "palavras_chave": [], "sinonimos": [], "tipo": "conceito", "titulo": "O Motor BrocaCockpit (Hodge por dentro)"}
\section{O Motor BrocaCockpit (Hodge por dentro)}
A classe servidor: 'Hodge interior contínuo; BSD só na borda'. Mantém o estado Hodge em memória e corre dois workers — \_worker (recalcula quando livro/ledger/fase/chain mudam por mtime) e \_conductor\_worker (o ritmo vivo \~{}0.8s + acopla o olho invertido). Serve HTTP+SSE. A análise pesada roda em threads para não bloquear o stream.
\prov{relações: parte\_de cockpit\_observatorio; relaciona olho\_de\_cock\_web}
\prov{proveniência: wiki \textperiodcentered{} goldchain/cockpit\_server.py:BrocaCockpit \textperiodcentered{} fato \textperiodcentered{} confiança 1.0 \textperiodcentered{} domínio cockpit \textperiodcentered{} história 3 versões no jornal}

%CRISTAL {"arestas": [{"alvo": "repositorio_broca_os", "confianca": 1.0, "epistemico": "fato", "origem": "wiki", "rel": "usa"}, {"alvo": "comunidade_organica_por_banda", "confianca": 1.0, "epistemico": "fato", "origem": "wiki", "rel": "usa"}, {"alvo": "software_arte_organismo", "confianca": 1.0, "epistemico": "fato", "origem": "wiki", "rel": "usa"}, {"alvo": "autoridade_intrinseca", "confianca": 1.0, "epistemico": "fato", "origem": "wiki", "rel": "usa"}, {"alvo": "goldchain_como_rede", "confianca": 1.0, "epistemico": "fato", "origem": "wiki", "rel": "semelhante"}, {"alvo": "autocomunicacao_de_massa", "confianca": 1.0, "epistemico": "fato", "origem": "wiki", "rel": "semelhante"}, {"alvo": "cultura_da_convergencia", "confianca": 1.0, "epistemico": "fato", "origem": "wiki", "rel": "semelhante"}, {"alvo": "inteligencia_coletiva", "confianca": 1.0, "epistemico": "fato", "origem": "wiki", "rel": "semelhante"}, {"alvo": "capitalismo_de_vigilancia", "confianca": 1.0, "epistemico": "fato", "origem": "wiki", "rel": "contradiz"}, {"alvo": "economia_da_atencao", "confianca": 1.0, "epistemico": "fato", "origem": "wiki", "rel": "contradiz"}, {"alvo": "orbita_rede_dois_atratores", "confianca": 1.0, "epistemico": "fato", "origem": "wiki", "rel": "parte_de"}], "confianca": 1.0, "contraexemplos": [], "descricao": "Cada cliente guarda e COMPARTILHA suas apps e conteúdo na sua comunidade (banda) — e isso é uma REDE SOCIAL, como as modernas (comunidades, feeds, compartilhamento), só que SOBERANA: você é dono da sua banda (o tecido é a credencial), não o produto. A autoridade é intrínseca, os admins governam por BAI, e as apps/sementes circulam entre comunidades. É o meio-termo entre o repositório de apps e o goldchain (P2P de consenso) — uma rede social de pertencimento, sem dono central.", "epistemico": "inferencia", "exemplos": [], "id": "broca_rede_social_soberana", "memoria": "semantica", "meta": {"dominio": "arquitetura", "fonte": "semilla_broca_rede_social"}, "origem": "wiki", "palavras_chave": ["cliente salva apps", "onde o cliente salva suas apps", "salvar aplicações", "guardar apps", "compartilhar apps na comunidade", "app do cliente", "rede social do broca-os"], "sinonimos": [], "tipo": "conceito", "titulo": "O broca-os como Rede Social Soberana"}
\section{O broca-os como Rede Social Soberana}
Cada cliente guarda e COMPARTILHA suas apps e conteúdo na sua comunidade (banda) — e isso é uma REDE SOCIAL, como as modernas (comunidades, feeds, compartilhamento), só que SOBERANA: você é dono da sua banda (o tecido é a credencial), não o produto. A autoridade é intrínseca, os admins governam por BAI, e as apps/sementes circulam entre comunidades. É o meio-termo entre o repositório de apps e o goldchain (P2P de consenso) — uma rede social de pertencimento, sem dono central.
\prov{relações: usa repositorio\_broca\_os; usa comunidade\_organica\_por\_banda; usa software\_arte\_organismo; usa autoridade\_intrinseca; semelhante goldchain\_como\_rede; semelhante autocomunicacao\_de\_massa; semelhante cultura\_da\_convergencia; semelhante inteligencia\_coletiva; contradiz capitalismo\_de\_vigilancia; contradiz economia\_da\_atencao; parte\_de orbita\_rede\_dois\_atratores}
\prov{proveniência: wiki \textperiodcentered{} semilla\_broca\_rede\_social \textperiodcentered{} inferencia \textperiodcentered{} confiança 1.0 \textperiodcentered{} domínio arquitetura \textperiodcentered{} história 23 versões no jornal}

%CRISTAL {"arestas": [{"alvo": "impressao_3d", "confianca": 1.0, "epistemico": "fato", "origem": "wiki", "rel": "usa"}, {"alvo": "computador_mineral", "confianca": 1.0, "epistemico": "fato", "origem": "wiki", "rel": "relaciona"}], "confianca": 1.0, "contraexemplos": [], "descricao": "Modelar a peça no computador (CAD) e deixar o algoritmo GERAR a forma ótima a partir de restrições (carga, material, peso) — otimização que desce a um mínimo (o atrator). É o DSL do computador mineral aplicado à forma: primitivas e restrições geram o objeto.", "epistemico": "inferencia", "exemplos": [], "id": "cad_generativo", "memoria": "semantica", "meta": {"dominio": "manufatura", "fonte": "impressão 3D e manufatura"}, "origem": "wiki", "palavras_chave": [], "sinonimos": [], "tipo": "conceito", "titulo": "CAD e Design Generativo"}
\section{CAD e Design Generativo}
Modelar a peça no computador (CAD) e deixar o algoritmo GERAR a forma ótima a partir de restrições (carga, material, peso) — otimização que desce a um mínimo (o atrator). É o DSL do computador mineral aplicado à forma: primitivas e restrições geram o objeto.
\prov{relações: usa impressao\_3d; relaciona computador\_mineral}
\prov{proveniência: wiki \textperiodcentered{} impressão 3D e manufatura \textperiodcentered{} inferencia \textperiodcentered{} confiança 1.0 \textperiodcentered{} domínio manufatura \textperiodcentered{} história 3 versões no jornal}

%CRISTAL {"arestas": [{"alvo": "repositorio_broca_os", "confianca": 1.0, "epistemico": "fato", "origem": "wiki", "rel": "parte_de"}, {"alvo": "cadeias_dep_grant", "confianca": 1.0, "epistemico": "fato", "origem": "wiki", "rel": "usa"}, {"alvo": "bai_deny_overrides", "confianca": 1.0, "epistemico": "fato", "origem": "wiki", "rel": "usa"}, {"alvo": "bai_universalidade", "confianca": 1.0, "epistemico": "fato", "origem": "wiki", "rel": "usa"}], "confianca": 1.0, "contraexemplos": [], "descricao": "A raiz delega capacidades (admin, deploy, convidar, revogar) por uma CADEIA DEP-GRANT: cada concessão carrega um orçamento que decresce a cada salto (delegação limitada), validada turno a turno (o concedente precisa ter admin). Revogar é um DENY que DOMINA e cai em cascata. Fail-closed: sem grant, negado. A governança viva da comunidade — o BAI aplicado à administração.", "epistemico": "fato", "exemplos": [], "id": "cadeia_admins_comunidade", "memoria": "semantica", "meta": {"dominio": "arquitetura", "fonte": "goldchain/bai_autoridade.py"}, "origem": "wiki", "palavras_chave": [], "sinonimos": [], "tipo": "conceito", "titulo": "Cadeia de Administradores de Comunidade"}
\section{Cadeia de Administradores de Comunidade}
A raiz delega capacidades (admin, deploy, convidar, revogar) por uma CADEIA DEP-GRANT: cada concessão carrega um orçamento que decresce a cada salto (delegação limitada), validada turno a turno (o concedente precisa ter admin). Revogar é um DENY que DOMINA e cai em cascata. Fail-closed: sem grant, negado. A governança viva da comunidade — o BAI aplicado à administração.
\prov{relações: parte\_de repositorio\_broca\_os; usa cadeias\_dep\_grant; usa bai\_deny\_overrides; usa bai\_universalidade}
\prov{proveniência: wiki \textperiodcentered{} goldchain/bai\_autoridade.py \textperiodcentered{} fato \textperiodcentered{} confiança 1.0 \textperiodcentered{} domínio arquitetura \textperiodcentered{} história 5 versões no jornal}

%CRISTAL {"arestas": [{"alvo": "matematica", "confianca": 1.0, "epistemico": "fato", "origem": "wiki", "rel": "parte_de"}, {"alvo": "eng_eletronica", "confianca": 1.0, "epistemico": "fato", "origem": "curriculo", "rel": "parte_de"}, {"alvo": "eng_telecomunicacoes", "confianca": 1.0, "epistemico": "fato", "origem": "curriculo", "rel": "parte_de"}, {"alvo": "eng_sistemas_energia", "confianca": 1.0, "epistemico": "fato", "origem": "curriculo", "rel": "parte_de"}], "confianca": 1.0, "contraexemplos": [], "descricao": "Limite, derivada, integral e séries — a linguagem da variação contínua. O alfabeto de toda a física e engenharia que vem depois.", "epistemico": "fato", "exemplos": [], "id": "calculo_engenharia", "memoria": "semantica", "meta": {"dominio": "engenharia", "fonte": ""}, "origem": "wiki", "palavras_chave": [], "sinonimos": [], "tipo": "conceito", "titulo": "Cálculo"}
\section{Cálculo}
Limite, derivada, integral e séries — a linguagem da variação contínua. O alfabeto de toda a física e engenharia que vem depois.
\prov{relações: parte\_de matematica; parte\_de eng\_eletronica; parte\_de eng\_telecomunicacoes; parte\_de eng\_sistemas\_energia}
\prov{proveniência: wiki \textperiodcentered{} fato \textperiodcentered{} confiança 1.0 \textperiodcentered{} domínio engenharia \textperiodcentered{} história 5 versões no jornal}

%CRISTAL {"arestas": [{"alvo": "modelo_osi", "confianca": 1.0, "epistemico": "fato", "origem": "wiki", "rel": "parte_de"}, {"alvo": "eng_telecomunicacoes", "confianca": 1.0, "epistemico": "fato", "origem": "curriculo", "rel": "parte_de"}], "confianca": 1.0, "contraexemplos": [], "descricao": "Entregar quadros entre nós no MESMO segmento físico, detectar erros e mediar o acesso ao meio. A camada que transforma um meio bruto num canal quadro-a-quadro confiável o bastante.", "epistemico": "fato", "exemplos": [], "id": "camada_de_enlace", "memoria": "semantica", "meta": {"dominio": "engenharia", "fonte": ""}, "origem": "wiki", "palavras_chave": [], "sinonimos": [], "tipo": "conceito", "titulo": "Camada de Enlace de Dados"}
\section{Camada de Enlace de Dados}
Entregar quadros entre nós no MESMO segmento físico, detectar erros e mediar o acesso ao meio. A camada que transforma um meio bruto num canal quadro-a-quadro confiável o bastante.
\prov{relações: parte\_de modelo\_osi; parte\_de eng\_telecomunicacoes}
\prov{proveniência: wiki \textperiodcentered{} fato \textperiodcentered{} confiança 1.0 \textperiodcentered{} domínio engenharia \textperiodcentered{} história 3 versões no jornal}

%CRISTAL {"arestas": [{"alvo": "redes_5g_nr", "confianca": 1.0, "epistemico": "fato", "origem": "wiki", "rel": "parte_de"}, {"alvo": "ofdm", "confianca": 1.0, "epistemico": "fato", "origem": "wiki", "rel": "usa"}, {"alvo": "camada_de_enlace", "confianca": 1.0, "epistemico": "fato", "origem": "wiki", "rel": "relaciona"}], "confianca": 1.0, "contraexemplos": [], "descricao": "A camada 1 do NR: mapeia bits codificados em símbolos OFDM sobre a grade tempo-frequência, com múltiplas antenas e feixes. Onde a informação vira onda de rádio — o OFDM/MIMO da trilha de princípios, parametrizado.", "epistemico": "fato", "exemplos": [], "id": "camada_fisica_nr", "memoria": "semantica", "meta": {"dominio": "engenharia", "fonte": ""}, "origem": "wiki", "palavras_chave": [], "sinonimos": [], "tipo": "conceito", "titulo": "Camada Física do NR (PHY)"}
\section{Camada Física do NR (PHY)}
A camada 1 do NR: mapeia bits codificados em símbolos OFDM sobre a grade tempo-frequência, com múltiplas antenas e feixes. Onde a informação vira onda de rádio — o OFDM/MIMO da trilha de princípios, parametrizado.
\prov{relações: parte\_de redes\_5g\_nr; usa ofdm; relaciona camada\_de\_enlace}
\prov{proveniência: wiki \textperiodcentered{} fato \textperiodcentered{} confiança 1.0 \textperiodcentered{} domínio engenharia \textperiodcentered{} história 4 versões no jornal}

%CRISTAL {"arestas": [{"alvo": "camada_de_enlace", "confianca": 1.0, "epistemico": "fato", "origem": "wiki", "rel": "especializa"}, {"alvo": "camada_fisica_nr", "confianca": 1.0, "epistemico": "fato", "origem": "wiki", "rel": "usa"}, {"alvo": "acesso_ao_meio", "confianca": 1.0, "epistemico": "fato", "origem": "wiki", "rel": "usa"}], "confianca": 1.0, "contraexemplos": [], "descricao": "Controle de acesso ao meio: escalona os recursos, multiplexa fluxos lógicos, roda o HARQ e gere o acesso aleatório e o controle de potência. O maestro que decide, a cada slot, quem transmite o quê.", "epistemico": "fato", "exemplos": [], "id": "camada_mac_nr", "memoria": "semantica", "meta": {"dominio": "engenharia", "fonte": ""}, "origem": "wiki", "palavras_chave": [], "sinonimos": [], "tipo": "conceito", "titulo": "Camada MAC do NR"}
\section{Camada MAC do NR}
Controle de acesso ao meio: escalona os recursos, multiplexa fluxos lógicos, roda o HARQ e gere o acesso aleatório e o controle de potência. O maestro que decide, a cada slot, quem transmite o quê.
\prov{relações: especializa camada\_de\_enlace; usa camada\_fisica\_nr; usa acesso\_ao\_meio}
\prov{proveniência: wiki \textperiodcentered{} fato \textperiodcentered{} confiança 1.0 \textperiodcentered{} domínio engenharia \textperiodcentered{} história 4 versões no jornal}

%CRISTAL {"arestas": [{"alvo": "camada_rlc_nr", "confianca": 1.0, "epistemico": "fato", "origem": "wiki", "rel": "usa"}, {"alvo": "criptografia", "confianca": 1.0, "epistemico": "fato", "origem": "wiki", "rel": "usa"}, {"alvo": "dupla_conectividade_nr", "confianca": 1.0, "epistemico": "fato", "origem": "wiki", "rel": "usa"}], "confianca": 1.0, "contraexemplos": [], "descricao": "Convergência de dados: CIFRA e protege a integridade, comprime cabeçalhos (ROHC), reordena e descarta duplicatas. É onde a segurança do ar acontece e onde a dupla conectividade divide/duplica o fluxo.", "epistemico": "fato", "exemplos": [], "id": "camada_pdcp_nr", "memoria": "semantica", "meta": {"dominio": "engenharia", "fonte": ""}, "origem": "wiki", "palavras_chave": [], "sinonimos": [], "tipo": "conceito", "titulo": "Camada PDCP do NR"}
\section{Camada PDCP do NR}
Convergência de dados: CIFRA e protege a integridade, comprime cabeçalhos (ROHC), reordena e descarta duplicatas. É onde a segurança do ar acontece e onde a dupla conectividade divide/duplica o fluxo.
\prov{relações: usa camada\_rlc\_nr; usa criptografia; usa dupla\_conectividade\_nr}
\prov{proveniência: wiki \textperiodcentered{} fato \textperiodcentered{} confiança 1.0 \textperiodcentered{} domínio engenharia \textperiodcentered{} história 4 versões no jornal}

%CRISTAL {"arestas": [{"alvo": "camada_de_enlace", "confianca": 1.0, "epistemico": "fato", "origem": "wiki", "rel": "parte_de"}, {"alvo": "camada_mac_nr", "confianca": 1.0, "epistemico": "fato", "origem": "wiki", "rel": "usa"}], "confianca": 1.0, "contraexemplos": [], "descricao": "Controle de enlace de rádio em três modos: transparente (TM), sem confirmação (UM) e com confirmação (AM, com ARQ). Segmenta e remonta os pacotes ao tamanho que o MAC concede. No NR já não reordena — isso subiu ao PDCP.", "epistemico": "fato", "exemplos": [], "id": "camada_rlc_nr", "memoria": "semantica", "meta": {"dominio": "engenharia", "fonte": ""}, "origem": "wiki", "palavras_chave": [], "sinonimos": [], "tipo": "conceito", "titulo": "Camada RLC do NR"}
\section{Camada RLC do NR}
Controle de enlace de rádio em três modos: transparente (TM), sem confirmação (UM) e com confirmação (AM, com ARQ). Segmenta e remonta os pacotes ao tamanho que o MAC concede. No NR já não reordena — isso subiu ao PDCP.
\prov{relações: parte\_de camada\_de\_enlace; usa camada\_mac\_nr}
\prov{proveniência: wiki \textperiodcentered{} fato \textperiodcentered{} confiança 1.0 \textperiodcentered{} domínio engenharia \textperiodcentered{} história 3 versões no jornal}

%CRISTAL {"arestas": [{"alvo": "camada_pdcp_nr", "confianca": 1.0, "epistemico": "fato", "origem": "wiki", "rel": "usa"}, {"alvo": "ng_ran_gnb", "confianca": 1.0, "epistemico": "fato", "origem": "wiki", "rel": "parte_de"}], "confianca": 1.0, "contraexemplos": [], "descricao": "Controle de recursos de rádio: estabelece/reconfigura a conexão, difunde a informação do sistema, gere medições, mobilidade e portadores. A camada 3 do plano de controle entre o UE e o gNB.", "epistemico": "fato", "exemplos": [], "id": "camada_rrc_nr", "memoria": "semantica", "meta": {"dominio": "engenharia", "fonte": ""}, "origem": "wiki", "palavras_chave": [], "sinonimos": [], "tipo": "conceito", "titulo": "Camada RRC do NR"}
\section{Camada RRC do NR}
Controle de recursos de rádio: estabelece/reconfigura a conexão, difunde a informação do sistema, gere medições, mobilidade e portadores. A camada 3 do plano de controle entre o UE e o gNB.
\prov{relações: usa camada\_pdcp\_nr; parte\_de ng\_ran\_gnb}
\prov{proveniência: wiki \textperiodcentered{} fato \textperiodcentered{} confiança 1.0 \textperiodcentered{} domínio engenharia \textperiodcentered{} história 3 versões no jornal}

%CRISTAL {"arestas": [{"alvo": "camada_pdcp_nr", "confianca": 1.0, "epistemico": "fato", "origem": "wiki", "rel": "usa"}, {"alvo": "qos_fluxos_5g", "confianca": 1.0, "epistemico": "fato", "origem": "wiki", "rel": "usa"}], "confianca": 1.0, "contraexemplos": [], "descricao": "A camada nova do 5G: mapeia os FLUXOS DE QoS do núcleo nos portadores de rádio (DRBs), marcando cada pacote para o tratamento certo. O elo que faz a QoS fim-a-fim descer até o ar.", "epistemico": "fato", "exemplos": [], "id": "camada_sdap_nr", "memoria": "semantica", "meta": {"dominio": "engenharia", "fonte": ""}, "origem": "wiki", "palavras_chave": [], "sinonimos": [], "tipo": "conceito", "titulo": "Camada SDAP do NR"}
\section{Camada SDAP do NR}
A camada nova do 5G: mapeia os FLUXOS DE QoS do núcleo nos portadores de rádio (DRBs), marcando cada pacote para o tratamento certo. O elo que faz a QoS fim-a-fim descer até o ar.
\prov{relações: usa camada\_pdcp\_nr; usa qos\_fluxos\_5g}
\prov{proveniência: wiki \textperiodcentered{} fato \textperiodcentered{} confiança 1.0 \textperiodcentered{} domínio engenharia \textperiodcentered{} história 3 versões no jornal}

%CRISTAL {"arestas": [{"alvo": "camada_fisica_nr", "confianca": 1.0, "epistemico": "fato", "origem": "wiki", "rel": "parte_de"}, {"alvo": "grade_de_recursos_nr", "confianca": 1.0, "epistemico": "fato", "origem": "wiki", "rel": "usa"}], "confianca": 1.0, "contraexemplos": [], "descricao": "Os dutos do ar: PDSCH/PUSCH carregam dados (descida/subida), PDCCH/PUCCH a sinalização de controle, e o PRACH o acesso aleatório inicial. Mapeiam canais de transporte e lógicos das camadas de cima na grade física.", "epistemico": "fato", "exemplos": [], "id": "canais_fisicos_nr", "memoria": "semantica", "meta": {"dominio": "engenharia", "fonte": ""}, "origem": "wiki", "palavras_chave": [], "sinonimos": [], "tipo": "conceito", "titulo": "Canais Físicos (PDSCH/PDCCH/PUSCH/PUCCH/PRACH)"}
\section{Canais Físicos (PDSCH/PDCCH/PUSCH/PUCCH/PRACH)}
Os dutos do ar: PDSCH/PUSCH carregam dados (descida/subida), PDCCH/PUCCH a sinalização de controle, e o PRACH o acesso aleatório inicial. Mapeiam canais de transporte e lógicos das camadas de cima na grade física.
\prov{relações: parte\_de camada\_fisica\_nr; usa grade\_de\_recursos\_nr}
\prov{proveniência: wiki \textperiodcentered{} fato \textperiodcentered{} confiança 1.0 \textperiodcentered{} domínio engenharia \textperiodcentered{} história 3 versões no jornal}

%CRISTAL {"arestas": [{"alvo": "teoria_da_informacao", "confianca": 1.0, "epistemico": "fato", "origem": "wiki", "rel": "parte_de"}, {"alvo": "entropia_shannon", "confianca": 1.0, "epistemico": "fato", "origem": "wiki", "rel": "usa"}, {"alvo": "razao_sinal_ruido", "confianca": 1.0, "epistemico": "fato", "origem": "wiki", "rel": "usa"}, {"alvo": "prova_de_trabalho", "confianca": 1.0, "epistemico": "fato", "origem": "wiki", "rel": "relaciona"}, {"alvo": "eng_telecomunicacoes", "confianca": 1.0, "epistemico": "fato", "origem": "curriculo", "rel": "parte_de"}], "confianca": 1.0, "contraexemplos": [], "descricao": "C = B·log₂(1+S/N): o limite MÁXIMO de bits/s que um canal de banda B e relação sinal-ruído S/N pode carregar sem erro. A lei que nenhum esquema de modulação pode violar — o teto de toda comunicação.", "epistemico": "fato", "exemplos": [], "id": "capacidade_de_canal", "memoria": "semantica", "meta": {"autor": "Shannon", "dominio": "engenharia", "fonte": ""}, "origem": "wiki", "palavras_chave": [], "sinonimos": [], "tipo": "conceito", "titulo": "Capacidade de Canal (Shannon-Hartley)"}
\section{Capacidade de Canal (Shannon-Hartley)}
C = B\textperiodcentered log\ensuremath{_{2}}(1+S/N): o limite MÁXIMO de bits/s que um canal de banda B e relação sinal-ruído S/N pode carregar sem erro. A lei que nenhum esquema de modulação pode violar — o teto de toda comunicação.
\prov{relações: parte\_de teoria\_da\_informacao; usa entropia\_shannon; usa razao\_sinal\_ruido; relaciona prova\_de\_trabalho; parte\_de eng\_telecomunicacoes}
\prov{proveniência: wiki \textperiodcentered{} fato \textperiodcentered{} confiança 1.0 \textperiodcentered{} domínio engenharia \textperiodcentered{} história 6 versões no jornal}

%CRISTAL {"arestas": [{"alvo": "conectividade", "confianca": 1.0, "epistemico": "fato", "origem": "wiki", "rel": "explica"}, {"alvo": "relacao_sinal_ruido", "confianca": 1.0, "epistemico": "fato", "origem": "wiki", "rel": "usa"}, {"alvo": "corpo_mineral", "confianca": 1.0, "epistemico": "fato", "origem": "wiki", "rel": "usa"}, {"alvo": "qam", "confianca": 1.0, "epistemico": "fato", "origem": "wiki", "rel": "relaciona"}], "confianca": 1.0, "contraexemplos": [], "descricao": "C = B·log₂(1+SNR) bits/s: o TETO absoluto de um canal. Abaixo dele existe código que transmite sem erro (cristal); acima dele o erro é inevitável (caos); no teto é a borda. É a fronteira dura da comunicação — a mesma faixa mineral, agora no canal. Verif.: 100 MHz @ 25 dB → 831 Mbps.", "epistemico": "inferencia", "exemplos": [], "id": "capacidade_de_shannon", "memoria": "semantica", "meta": {"dominio": "engenharia", "exemplo_C_100mhz_25db_mbps": 830.9, "fonte": "conectividade e 5G"}, "origem": "wiki", "palavras_chave": [], "sinonimos": [], "tipo": "conceito", "titulo": "Capacidade de Shannon (a borda do canal)"}
\section{Capacidade de Shannon (a borda do canal)}
C = B\textperiodcentered log\ensuremath{_{2}}(1+SNR) bits/s: o TETO absoluto de um canal. Abaixo dele existe código que transmite sem erro (cristal); acima dele o erro é inevitável (caos); no teto é a borda. É a fronteira dura da comunicação — a mesma faixa mineral, agora no canal. Verif.: 100 MHz @ 25 dB \ensuremath{\to} 831 Mbps.
\prov{relações: explica conectividade; usa relacao\_sinal\_ruido; usa corpo\_mineral; relaciona qam}
\prov{proveniência: wiki \textperiodcentered{} conectividade e 5G \textperiodcentered{} inferencia \textperiodcentered{} confiança 1.0 \textperiodcentered{} domínio engenharia \textperiodcentered{} história 5 versões no jornal}

%CRISTAL {"arestas": [{"alvo": "curva_de_carga", "confianca": 1.0, "epistemico": "fato", "origem": "wiki", "rel": "usa"}, {"alvo": "rede_de_carregamento", "confianca": 1.0, "epistemico": "fato", "origem": "wiki", "rel": "usa"}, {"alvo": "estabilidade_da_rede", "confianca": 1.0, "epistemico": "fato", "origem": "wiki", "rel": "explica"}, {"alvo": "lyapunov_da_orbita", "confianca": 1.0, "epistemico": "fato", "origem": "wiki", "rel": "usa"}, {"alvo": "corpo_mineral", "confianca": 1.0, "epistemico": "fato", "origem": "wiki", "rel": "usa"}], "confianca": 1.0, "contraexemplos": [], "descricao": "Smart charging: decidir QUANDO cada carro carrega. A MESMA energia da frota, empilhada no pico da noite, empurra a curva ao CAOS (PAR alto, risco de blackout); espalhada no vale da madrugada, aplaina ao CRISTAL estável. É o controle que segura o dia na borda servível. Verificado: descoordenado PAR=1.66 (caos) vs coordenado PAR=1.41 (cristal).", "epistemico": "inferencia", "exemplos": [], "id": "carga_coordenada", "memoria": "semantica", "meta": {"PAR_coordenado": 1.409, "PAR_descoordenado": 1.657, "classe_coord": "cristal (aplainada, estável)", "classe_desc": "caos (pico agudo, risco de blackout)", "dominio": "engenharia", "fonte": "rede de carregamento e smart cities"}, "origem": "wiki", "palavras_chave": [], "sinonimos": [], "tipo": "conceito", "titulo": "Carga Coordenada (manter o dia na borda)"}
\section{Carga Coordenada (manter o dia na borda)}
Smart charging: decidir QUANDO cada carro carrega. A MESMA energia da frota, empilhada no pico da noite, empurra a curva ao CAOS (PAR alto, risco de blackout); espalhada no vale da madrugada, aplaina ao CRISTAL estável. É o controle que segura o dia na borda servível. Verificado: descoordenado PAR=1.66 (caos) vs coordenado PAR=1.41 (cristal).
\prov{relações: usa curva\_de\_carga; usa rede\_de\_carregamento; explica estabilidade\_da\_rede; usa lyapunov\_da\_orbita; usa corpo\_mineral}
\prov{proveniência: wiki \textperiodcentered{} rede de carregamento e smart cities \textperiodcentered{} inferencia \textperiodcentered{} confiança 1.0 \textperiodcentered{} domínio engenharia \textperiodcentered{} história 12 versões no jornal}

%CRISTAL {"arestas": [{"alvo": "recarga_veicular", "confianca": 1.0, "epistemico": "fato", "origem": "wiki", "rel": "usa"}, {"alvo": "rede_eletrica", "confianca": 1.0, "epistemico": "fato", "origem": "wiki", "rel": "relaciona"}], "confianca": 1.0, "contraexemplos": [], "descricao": "A malha de eletropostos que estende a rede elétrica aos veículos: onde a distribuição encontra a estrada. O gargalo prático da mobilidade elétrica — a rede precisa alcançar cada nó móvel.", "epistemico": "fato", "exemplos": [], "id": "carregador_publico", "memoria": "semantica", "meta": {"dominio": "engenharia", "fonte": "mobilidade elétrica"}, "origem": "wiki", "palavras_chave": [], "sinonimos": [], "tipo": "conceito", "titulo": "Rede de Recarga (a distribuição móvel)"}
\section{Rede de Recarga (a distribuição móvel)}
A malha de eletropostos que estende a rede elétrica aos veículos: onde a distribuição encontra a estrada. O gargalo prático da mobilidade elétrica — a rede precisa alcançar cada nó móvel.
\prov{relações: usa recarga\_veicular; relaciona rede\_eletrica}
\prov{proveniência: wiki \textperiodcentered{} mobilidade elétrica \textperiodcentered{} fato \textperiodcentered{} confiança 1.0 \textperiodcentered{} domínio engenharia \textperiodcentered{} história 3 versões no jornal}

%CRISTAL {"arestas": [{"alvo": "mobilidade_eletrica", "confianca": 1.0, "epistemico": "fato", "origem": "wiki", "rel": "eh_um"}], "confianca": 1.0, "contraexemplos": [], "descricao": "Um veículo cuja tração vem de um motor elétrico alimentado por uma bateria recarregável. Sem motor a combustão, sem câmbio de marchas; torque instantâneo, silencioso, e ~3× mais eficiente tanque-à-roda.", "epistemico": "fato", "exemplos": [], "id": "carro_eletrico", "memoria": "semantica", "meta": {"dominio": "engenharia", "fonte": "mobilidade elétrica"}, "origem": "wiki", "palavras_chave": [], "sinonimos": [], "tipo": "conceito", "titulo": "Carro Elétrico (VE)"}
\section{Carro Elétrico (VE)}
Um veículo cuja tração vem de um motor elétrico alimentado por uma bateria recarregável. Sem motor a combustão, sem câmbio de marchas; torque instantâneo, silencioso, e \~{}3$\times$ mais eficiente tanque-à-roda.
\prov{relações: eh\_um mobilidade\_eletrica}
\prov{proveniência: wiki \textperiodcentered{} mobilidade elétrica \textperiodcentered{} fato \textperiodcentered{} confiança 1.0 \textperiodcentered{} domínio engenharia \textperiodcentered{} história 2 versões no jornal}

%CRISTAL {"arestas": [{"alvo": "redes_5g_nr", "confianca": 1.0, "epistemico": "fato", "origem": "wiki", "rel": "parte_de"}], "confianca": 1.0, "contraexemplos": [], "descricao": "O triângulo que define o 5G: eMBB (banda larga móvel, vídeo 4K), URLLC (ultraconfiável e baixíssima latência — cirurgia remota, veículo) e mMTC (máquina-a-máquina massiva, milhões de sensores/km²). A mesma rede serve os três por parametrização.", "epistemico": "fato", "exemplos": [], "id": "casos_uso_5g", "memoria": "semantica", "meta": {"dominio": "engenharia", "fonte": ""}, "origem": "wiki", "palavras_chave": [], "sinonimos": [], "tipo": "conceito", "titulo": "Casos de Uso 5G (eMBB, URLLC, mMTC)"}
\section{Casos de Uso 5G (eMBB, URLLC, mMTC)}
O triângulo que define o 5G: eMBB (banda larga móvel, vídeo 4K), URLLC (ultraconfiável e baixíssima latência — cirurgia remota, veículo) e mMTC (máquina-a-máquina massiva, milhões de sensores/km\ensuremath{^{2}}). A mesma rede serve os três por parametrização.
\prov{relações: parte\_de redes\_5g\_nr}
\prov{proveniência: wiki \textperiodcentered{} fato \textperiodcentered{} confiança 1.0 \textperiodcentered{} domínio engenharia \textperiodcentered{} história 2 versões no jornal}

%CRISTAL {"arestas": [{"alvo": "rede_autonoma_broca", "confianca": 1.0, "epistemico": "fato", "origem": "wiki", "rel": "parte_de"}, {"alvo": "peirce_celulas_canais", "confianca": 1.0, "epistemico": "fato", "origem": "wiki", "rel": "usa"}, {"alvo": "a_antena_transmite_em_tempo_real", "confianca": 1.0, "epistemico": "fato", "origem": "wiki", "rel": "usa"}], "confianca": 1.0, "contraexemplos": [], "descricao": "A menor unidade da rede é a decomposição de Peirce: uma CÉLULA (estado idempotente, o software) que fala por CANAIS (fluxo nilpotente, o bump numa banda). A dinâmica É o acoplamento — não há processo sem canal, nem canal sem par. Todo agente/serviço do broca-so é uma célula falante.", "epistemico": "inferencia", "exemplos": [], "id": "celula_canal_unidade_viva", "memoria": "semantica", "meta": {"dominio": "arquitetura", "fonte": "semilla_rede_autonoma"}, "origem": "wiki", "palavras_chave": [], "sinonimos": [], "tipo": "conceito", "titulo": "Célula + Canal: a unidade viva"}
\section{Célula + Canal: a unidade viva}
A menor unidade da rede é a decomposição de Peirce: uma CÉLULA (estado idempotente, o software) que fala por CANAIS (fluxo nilpotente, o bump numa banda). A dinâmica É o acoplamento — não há processo sem canal, nem canal sem par. Todo agente/serviço do broca-so é uma célula falante.
\prov{relações: parte\_de rede\_autonoma\_broca; usa peirce\_celulas\_canais; usa a\_antena\_transmite\_em\_tempo\_real}
\prov{proveniência: wiki \textperiodcentered{} semilla\_rede\_autonoma \textperiodcentered{} inferencia \textperiodcentered{} confiança 1.0 \textperiodcentered{} domínio arquitetura \textperiodcentered{} história 4 versões no jornal}

%CRISTAL {"arestas": [{"alvo": "celula", "confianca": 1.0, "epistemico": "fato", "origem": "wiki", "rel": "explica"}, {"alvo": "computador_mineral", "confianca": 1.0, "epistemico": "fato", "origem": "wiki", "rel": "eh_um"}, {"alvo": "metabolismo", "confianca": 1.0, "epistemico": "fato", "origem": "wiki", "rel": "usa"}, {"alvo": "dna_cadeia_base4", "confianca": 1.0, "epistemico": "fato", "origem": "wiki", "rel": "usa"}], "confianca": 1.0, "contraexemplos": [], "descricao": "A célula roda um programa (o genoma = memória), tem um processador (ribossomos e enzimas), energia (ATP) e canais — é o computador mineral que EXISTE, com 4 bilhões de anos. O que a broca-so descreve como máquina, a vida já é: informação numa cadeia, decodificada em função, replicada com checksum.", "epistemico": "inferencia", "exemplos": [], "id": "celula_computador_mineral", "memoria": "semantica", "meta": {"dominio": "biotecnologia", "fonte": "biotecnologia e engenharia genética"}, "origem": "wiki", "palavras_chave": [], "sinonimos": [], "tipo": "conceito", "titulo": "A Célula = o Computador Mineral Original"}
\section{A Célula = o Computador Mineral Original}
A célula roda um programa (o genoma = memória), tem um processador (ribossomos e enzimas), energia (ATP) e canais — é o computador mineral que EXISTE, com 4 bilhões de anos. O que a broca-so descreve como máquina, a vida já é: informação numa cadeia, decodificada em função, replicada com checksum.
\prov{relações: explica celula; eh\_um computador\_mineral; usa metabolismo; usa dna\_cadeia\_base4}
\prov{proveniência: wiki \textperiodcentered{} biotecnologia e engenharia genética \textperiodcentered{} inferencia \textperiodcentered{} confiança 1.0 \textperiodcentered{} domínio biotecnologia \textperiodcentered{} história 5 versões no jornal}

%CRISTAL {"arestas": [{"alvo": "juncao_pn", "confianca": 1.0, "epistemico": "fato", "origem": "wiki", "rel": "usa"}, {"alvo": "energia_solar", "confianca": 1.0, "epistemico": "fato", "origem": "wiki", "rel": "explica"}], "confianca": 1.0, "contraexemplos": [], "descricao": "A junção p-n ao contrário: a luz arranca elétrons e gera corrente — converte fóton em eletricidade. A ponte da física à energia solar.", "epistemico": "fato", "exemplos": [], "id": "celula_fotovoltaica", "memoria": "semantica", "meta": {"dominio": "engenharia", "fonte": ""}, "origem": "wiki", "palavras_chave": [], "sinonimos": [], "tipo": "conceito", "titulo": "Célula Fotovoltaica"}
\section{Célula Fotovoltaica}
A junção p-n ao contrário: a luz arranca elétrons e gera corrente — converte fóton em eletricidade. A ponte da física à energia solar.
\prov{relações: usa juncao\_pn; explica energia\_solar}
\prov{proveniência: wiki \textperiodcentered{} fato \textperiodcentered{} confiança 1.0 \textperiodcentered{} domínio engenharia \textperiodcentered{} história 3 versões no jornal}

%CRISTAL {"arestas": [{"alvo": "rede_cristalina", "confianca": 1.0, "epistemico": "fato", "origem": "wiki", "rel": "parte_de"}], "confianca": 1.0, "contraexemplos": [], "descricao": "O tijolo mínimo que, repetido, gera o cristal inteiro — a menor unidade que carrega toda a simetria da rede.", "epistemico": "fato", "exemplos": [], "id": "celula_unitaria", "memoria": "semantica", "meta": {"dominio": "engenharia", "fonte": ""}, "origem": "wiki", "palavras_chave": [], "sinonimos": [], "tipo": "conceito", "titulo": "Célula Unitária"}
\section{Célula Unitária}
O tijolo mínimo que, repetido, gera o cristal inteiro — a menor unidade que carrega toda a simetria da rede.
\prov{relações: parte\_de rede\_cristalina}
\prov{proveniência: wiki \textperiodcentered{} fato \textperiodcentered{} confiança 1.0 \textperiodcentered{} domínio engenharia \textperiodcentered{} história 2 versões no jornal}

%CRISTAL {"arestas": [{"alvo": "ciclo_termodinamico_orbita", "confianca": 1.0, "epistemico": "fato", "origem": "wiki", "rel": "eh_um"}], "confianca": 1.0, "contraexemplos": [], "descricao": "O ciclo das turbinas a gás e a jato: comprime ar, queima combustível, expande na turbina. Combinado com Rankine (ciclo combinado) atinge >60% de rendimento.", "epistemico": "fato", "exemplos": [], "id": "ciclo_brayton", "memoria": "semantica", "meta": {"dominio": "engenharia", "fonte": "usinas e geração"}, "origem": "wiki", "palavras_chave": [], "sinonimos": [], "tipo": "conceito", "titulo": "Ciclo de Brayton (gás)"}
\section{Ciclo de Brayton (gás)}
O ciclo das turbinas a gás e a jato: comprime ar, queima combustível, expande na turbina. Combinado com Rankine (ciclo combinado) atinge >60\% de rendimento.
\prov{relações: eh\_um ciclo\_termodinamico\_orbita}
\prov{proveniência: wiki \textperiodcentered{} usinas e geração \textperiodcentered{} fato \textperiodcentered{} confiança 1.0 \textperiodcentered{} domínio engenharia \textperiodcentered{} história 4 versões no jornal}

%CRISTAL {"arestas": [{"alvo": "ciclo_termodinamico_orbita", "confianca": 1.0, "epistemico": "fato", "origem": "wiki", "rel": "eh_um"}], "confianca": 1.0, "contraexemplos": [], "descricao": "O ciclo ideal reversível (2 isotermas + 2 adiabáticas), fechado no plano T-s. Dá o RENDIMENTO MÁXIMO de qualquer máquina térmica: η=1−Tc/Th. Nenhuma usina o supera — é a órbita-limite da conversão calor→trabalho.", "epistemico": "fato", "exemplos": [], "id": "ciclo_carnot", "memoria": "semantica", "meta": {"dominio": "engenharia", "fonte": "usinas e geração"}, "origem": "wiki", "palavras_chave": [], "sinonimos": [], "tipo": "conceito", "titulo": "Ciclo de Carnot (o teto)"}
\section{Ciclo de Carnot (o teto)}
O ciclo ideal reversível (2 isotermas + 2 adiabáticas), fechado no plano T-s. Dá o RENDIMENTO MÁXIMO de qualquer máquina térmica: \ensuremath{\eta}=1\ensuremath{-}Tc/Th. Nenhuma usina o supera — é a órbita-limite da conversão calor\ensuremath{\to}trabalho.
\prov{relações: eh\_um ciclo\_termodinamico\_orbita}
\prov{proveniência: wiki \textperiodcentered{} usinas e geração \textperiodcentered{} fato \textperiodcentered{} confiança 1.0 \textperiodcentered{} domínio engenharia \textperiodcentered{} história 4 versões no jornal}

%CRISTAL {"arestas": [{"alvo": "ciclo_termodinamico_orbita", "confianca": 1.0, "epistemico": "fato", "origem": "wiki", "rel": "eh_um"}], "confianca": 1.0, "contraexemplos": [], "descricao": "O ciclo das usinas a vapor (térmica, nuclear, solar CSP): água→caldeira→vapor→turbina→condensador→bomba→água. Um laço fechado que converte calor em trabalho. A maioria da eletricidade do mundo sai daqui.", "epistemico": "fato", "exemplos": [], "id": "ciclo_rankine", "memoria": "semantica", "meta": {"dominio": "engenharia", "fonte": "usinas e geração"}, "origem": "wiki", "palavras_chave": [], "sinonimos": [], "tipo": "conceito", "titulo": "Ciclo de Rankine (vapor)"}
\section{Ciclo de Rankine (vapor)}
O ciclo das usinas a vapor (térmica, nuclear, solar CSP): água\ensuremath{\to}caldeira\ensuremath{\to}vapor\ensuremath{\to}turbina\ensuremath{\to}condensador\ensuremath{\to}bomba\ensuremath{\to}água. Um laço fechado que converte calor em trabalho. A maioria da eletricidade do mundo sai daqui.
\prov{relações: eh\_um ciclo\_termodinamico\_orbita}
\prov{proveniência: wiki \textperiodcentered{} usinas e geração \textperiodcentered{} fato \textperiodcentered{} confiança 1.0 \textperiodcentered{} domínio engenharia \textperiodcentered{} história 4 versões no jornal}

%CRISTAL {"arestas": [{"alvo": "orbita_de_partida", "confianca": 1.0, "epistemico": "fato", "origem": "wiki", "rel": "relaciona"}, {"alvo": "corpo_mineral", "confianca": 1.0, "epistemico": "fato", "origem": "wiki", "rel": "usa"}], "confianca": 1.0, "contraexemplos": [], "descricao": "Um motor térmico (Rankine, Brayton, Carnot) é um LAÇO no espaço de estados (T-s / P-V) que volta ao início; a ÁREA fechada é o trabalho líquido. É uma órbita — a órbita da conversão de energia. Verificado: Carnot 800K/300K fecha, η=62,5%.", "epistemico": "inferencia", "exemplos": [], "id": "ciclo_termodinamico_orbita", "memoria": "semantica", "meta": {"dominio": "engenharia", "eta_carnot_800_300": 0.625, "fechada": true, "fonte": "usinas e geração"}, "origem": "wiki", "palavras_chave": [], "sinonimos": [], "tipo": "conceito", "titulo": "O Ciclo Térmico é uma Órbita Fechada"}
\section{O Ciclo Térmico é uma Órbita Fechada}
Um motor térmico (Rankine, Brayton, Carnot) é um LAÇO no espaço de estados (T-s / P-V) que volta ao início; a ÁREA fechada é o trabalho líquido. É uma órbita — a órbita da conversão de energia. Verificado: Carnot 800K/300K fecha, \ensuremath{\eta}=62,5\%.
\prov{relações: relaciona orbita\_de\_partida; usa corpo\_mineral}
\prov{proveniência: wiki \textperiodcentered{} usinas e geração \textperiodcentered{} inferencia \textperiodcentered{} confiança 1.0 \textperiodcentered{} domínio engenharia \textperiodcentered{} história 6 versões no jornal}

%CRISTAL {"arestas": [{"alvo": "smart_city", "confianca": 1.0, "epistemico": "fato", "origem": "wiki", "rel": "explica"}, {"alvo": "corpo_mineral", "confianca": 1.0, "epistemico": "fato", "origem": "wiki", "rel": "usa"}, {"alvo": "numero_de_salem", "confianca": 1.0, "epistemico": "fato", "origem": "wiki", "rel": "relaciona"}], "confianca": 1.0, "contraexemplos": [], "descricao": "Os laços de realimentação da cidade (energia, trânsito, água) vivem na BORDA entre ordem e caos: rígida demais fica morta e ineficiente (cristal), solta demais engarrafa e apaga (caos). A cidade saudável é a que se auto-organiza na borda — a mesma faixa da máquina mineral.", "epistemico": "inferencia", "exemplos": [], "id": "cidade_na_borda", "memoria": "semantica", "meta": {"dominio": "engenharia", "fonte": "rede de carregamento e smart cities"}, "origem": "wiki", "palavras_chave": [], "sinonimos": [], "tipo": "conceito", "titulo": "A Cidade Viva na Borda"}
\section{A Cidade Viva na Borda}
Os laços de realimentação da cidade (energia, trânsito, água) vivem na BORDA entre ordem e caos: rígida demais fica morta e ineficiente (cristal), solta demais engarrafa e apaga (caos). A cidade saudável é a que se auto-organiza na borda — a mesma faixa da máquina mineral.
\prov{relações: explica smart\_city; usa corpo\_mineral; relaciona numero\_de\_salem}
\prov{proveniência: wiki \textperiodcentered{} rede de carregamento e smart cities \textperiodcentered{} inferencia \textperiodcentered{} confiança 1.0 \textperiodcentered{} domínio engenharia \textperiodcentered{} história 8 versões no jornal}

%CRISTAL {"arestas": [{"alvo": "transistor", "confianca": 1.0, "epistemico": "fato", "origem": "wiki", "rel": "usa"}, {"alvo": "silicio", "confianca": 1.0, "epistemico": "fato", "origem": "wiki", "rel": "usa"}, {"alvo": "fpga", "confianca": 1.0, "epistemico": "fato", "origem": "wiki", "rel": "explica"}, {"alvo": "asic", "confianca": 1.0, "epistemico": "fato", "origem": "wiki", "rel": "explica"}], "confianca": 1.0, "contraexemplos": [], "descricao": "Milhões/bilhões de transistores gravados num único cristal de silício — a integração que barateou a computação; o corpo físico do FPGA e do ASIC.", "epistemico": "fato", "exemplos": [], "id": "circuito_integrado", "memoria": "semantica", "meta": {"autor": "Kilby/Noyce", "dominio": "engenharia", "fonte": ""}, "origem": "wiki", "palavras_chave": [], "sinonimos": [], "tipo": "conceito", "titulo": "Circuito Integrado (Chip)"}
\section{Circuito Integrado (Chip)}
Milhões/bilhões de transistores gravados num único cristal de silício — a integração que barateou a computação; o corpo físico do FPGA e do ASIC.
\prov{relações: usa transistor; usa silicio; explica fpga; explica asic}
\prov{proveniência: wiki \textperiodcentered{} fato \textperiodcentered{} confiança 1.0 \textperiodcentered{} domínio engenharia \textperiodcentered{} história 5 versões no jornal}

%CRISTAL {"arestas": [{"alvo": "circuito_integrado", "confianca": 1.0, "epistemico": "fato", "origem": "wiki", "rel": "especializa"}, {"alvo": "amplificador_operacional", "confianca": 1.0, "epistemico": "fato", "origem": "wiki", "rel": "usa"}, {"alvo": "eng_eletronica", "confianca": 1.0, "epistemico": "fato", "origem": "curriculo", "rel": "parte_de"}], "confianca": 1.0, "contraexemplos": [], "descricao": "Projetar amplificadores, espelhos de corrente, referências e comparadores dentro do chip — casamento de transistores, ganho e ruído. A arte de fazer o analógico caber no digital.", "epistemico": "inferencia", "exemplos": [], "id": "circuito_integrado_analogico", "memoria": "semantica", "meta": {"dominio": "engenharia", "fonte": ""}, "origem": "wiki", "palavras_chave": [], "sinonimos": [], "tipo": "conceito", "titulo": "Circuitos Integrados Analógicos"}
\section{Circuitos Integrados Analógicos}
Projetar amplificadores, espelhos de corrente, referências e comparadores dentro do chip — casamento de transistores, ganho e ruído. A arte de fazer o analógico caber no digital.
\prov{relações: especializa circuito\_integrado; usa amplificador\_operacional; parte\_de eng\_eletronica}
\prov{proveniência: wiki \textperiodcentered{} inferencia \textperiodcentered{} confiança 1.0 \textperiodcentered{} domínio engenharia \textperiodcentered{} história 4 versões no jornal}

%CRISTAL {"arestas": [{"alvo": "lei_de_ohm", "confianca": 1.0, "epistemico": "fato", "origem": "wiki", "rel": "usa"}, {"alvo": "numeros_complexos", "confianca": 1.0, "epistemico": "fato", "origem": "wiki", "rel": "usa"}, {"alvo": "corrente_e_tensao", "confianca": 1.0, "epistemico": "fato", "origem": "wiki", "rel": "usa"}, {"alvo": "eng_eletronica", "confianca": 1.0, "epistemico": "fato", "origem": "curriculo", "rel": "parte_de"}, {"alvo": "eng_telecomunicacoes", "confianca": 1.0, "epistemico": "fato", "origem": "curriculo", "rel": "parte_de"}, {"alvo": "eng_sistemas_energia", "confianca": 1.0, "epistemico": "fato", "origem": "curriculo", "rel": "parte_de"}], "confianca": 1.0, "contraexemplos": [], "descricao": "Analisar redes de resistores, capacitores e indutores: leis de Kirchhoff, Thévenin/Norton, fasores e regime permanente CA. O primeiro ofício do engenheiro elétrico.", "epistemico": "fato", "exemplos": [], "id": "circuitos_eletricos", "memoria": "semantica", "meta": {"autor": "Kirchhoff", "dominio": "engenharia", "fonte": ""}, "origem": "wiki", "palavras_chave": [], "sinonimos": [], "tipo": "conceito", "titulo": "Circuitos Elétricos"}
\section{Circuitos Elétricos}
Analisar redes de resistores, capacitores e indutores: leis de Kirchhoff, Thévenin/Norton, fasores e regime permanente CA. O primeiro ofício do engenheiro elétrico.
\prov{relações: usa lei\_de\_ohm; usa numeros\_complexos; usa corrente\_e\_tensao; parte\_de eng\_eletronica; parte\_de eng\_telecomunicacoes; parte\_de eng\_sistemas\_energia}
\prov{proveniência: wiki \textperiodcentered{} fato \textperiodcentered{} confiança 1.0 \textperiodcentered{} domínio engenharia \textperiodcentered{} história 7 versões no jornal}

%CRISTAL {"arestas": [{"alvo": "segmentacao_de_imagem", "confianca": 1.0, "epistemico": "fato", "origem": "wiki", "rel": "relaciona"}, {"alvo": "kmeans_atratores", "confianca": 1.0, "epistemico": "fato", "origem": "wiki", "rel": "usa"}, {"alvo": "corpo_mineral", "confianca": 1.0, "epistemico": "fato", "origem": "wiki", "rel": "usa"}], "confianca": 1.0, "contraexemplos": [], "descricao": "Classificar uma forma é um COLAPSO: a imagem desaba num rótulo (triângulo, estrela…). Como a forma aparece em qualquer pose, classifica-se por features INVARIANTES ao grupo de transformações (ordem de simetria, convexidade). Cada classe é um cristal nesse espaço; acerta 99% mesmo com a forma revirada. Verif. em formas_sinteticas.py.", "epistemico": "fato", "exemplos": [], "id": "classificacao_de_formas", "memoria": "semantica", "meta": {"dominio": "imagem", "fonte": "formas e limite"}, "origem": "wiki", "palavras_chave": [], "sinonimos": [], "tipo": "conceito", "titulo": "Classificação de Formas (colapso invariante)"}
\section{Classificação de Formas (colapso invariante)}
Classificar uma forma é um COLAPSO: a imagem desaba num rótulo (triângulo, estrela…). Como a forma aparece em qualquer pose, classifica-se por features INVARIANTES ao grupo de transformações (ordem de simetria, convexidade). Cada classe é um cristal nesse espaço; acerta 99\% mesmo com a forma revirada. Verif. em formas\_sinteticas.py.
\prov{relações: relaciona segmentacao\_de\_imagem; usa kmeans\_atratores; usa corpo\_mineral}
\prov{proveniência: wiki \textperiodcentered{} formas e limite \textperiodcentered{} fato \textperiodcentered{} confiança 1.0 \textperiodcentered{} domínio imagem \textperiodcentered{} história 8 versões no jornal}

%CRISTAL {"arestas": [{"alvo": "orbita_de_objeto_real", "confianca": 1.0, "epistemico": "fato", "origem": "wiki", "rel": "parte_de"}, {"alvo": "colapso", "confianca": 1.0, "epistemico": "fato", "origem": "wiki", "rel": "eh_um"}, {"alvo": "descritores_de_fourier_objeto", "confianca": 1.0, "epistemico": "fato", "origem": "wiki", "rel": "usa"}, {"alvo": "colapso_fatorado_glifo", "confianca": 1.0, "epistemico": "fato", "origem": "wiki", "rel": "relaciona"}], "confianca": 1.0, "contraexemplos": [], "descricao": "Classificar um objeto real é um COLAPSO: a sua órbita (os descritores de Fourier) desaba na classe (cristal) mais próxima do banco de órbitas conhecidas. A distância entre descritores mede a semelhança de forma; objetos parecidos têm órbitas próximas. O mesmo colapso do reconhecimento de glifo e da segmentação, agora sobre a silhueta do mundo.", "epistemico": "fato", "exemplos": [], "id": "classificar_objeto_por_orbita", "memoria": "semantica", "meta": {"dominio": "imagem", "fonte": "órbitas de objetos"}, "origem": "wiki", "palavras_chave": [], "sinonimos": [], "tipo": "conceito", "titulo": "Classificar Objeto = Colapso da Órbita"}
\section{Classificar Objeto = Colapso da Órbita}
Classificar um objeto real é um COLAPSO: a sua órbita (os descritores de Fourier) desaba na classe (cristal) mais próxima do banco de órbitas conhecidas. A distância entre descritores mede a semelhança de forma; objetos parecidos têm órbitas próximas. O mesmo colapso do reconhecimento de glifo e da segmentação, agora sobre a silhueta do mundo.
\prov{relações: parte\_de orbita\_de\_objeto\_real; eh\_um colapso; usa descritores\_de\_fourier\_objeto; relaciona colapso\_fatorado\_glifo}
\prov{proveniência: wiki \textperiodcentered{} órbitas de objetos \textperiodcentered{} fato \textperiodcentered{} confiança 1.0 \textperiodcentered{} domínio imagem \textperiodcentered{} história 5 versões no jornal}

%CRISTAL {"arestas": [{"alvo": "cockpit_olho", "confianca": 1.0, "epistemico": "fato", "origem": "wiki", "rel": "explica"}, {"alvo": "metal_broca_so", "confianca": 1.0, "epistemico": "fato", "origem": "wiki", "rel": "parte_de"}, {"alvo": "olho_de_koch", "confianca": 1.0, "epistemico": "fato", "origem": "wiki", "rel": "usa"}, {"alvo": "difusao_preenche_mineracao", "confianca": 1.0, "epistemico": "fato", "origem": "wiki", "rel": "relaciona"}], "confianca": 1.0, "contraexemplos": [], "descricao": "O observatório do broca-so: o painel estelar + o cockpit web que CONTEMPLAM o estado do sistema (mineração, goldchain, tecido) sem produzir nem otimizar — 'sem olho, o sistema dorme'. ESTADO: frente aberta / work-in-progress — precisa avançar ainda (aviso do Aarão, 2 jul 2026); imaturo frente ao resto do metal.", "epistemico": "hipotese", "exemplos": [], "id": "cockpit_observatorio", "memoria": "semantica", "meta": {"autor": "broca-so", "dominio": "cockpit", "fonte": "papers/painel_estelar.tex; goldchain/cockpit/", "wip": true}, "origem": "wiki", "palavras_chave": [], "sinonimos": [], "tipo": "conceito", "titulo": "O Cockpit — Observatório (não fábrica)"}
\section{O Cockpit — Observatório (não fábrica)}
O observatório do broca-so: o painel estelar + o cockpit web que CONTEMPLAM o estado do sistema (mineração, goldchain, tecido) sem produzir nem otimizar — 'sem olho, o sistema dorme'. ESTADO: frente aberta / work-in-progress — precisa avançar ainda (aviso do Aarão, 2 jul 2026); imaturo frente ao resto do metal.
\prov{relações: explica cockpit\_olho; parte\_de metal\_broca\_so; usa olho\_de\_koch; relaciona difusao\_preenche\_mineracao}
\prov{proveniência: wiki \textperiodcentered{} papers/painel\_estelar.tex; goldchain/cockpit/ \textperiodcentered{} hipotese \textperiodcentered{} confiança 1.0 \textperiodcentered{} domínio cockpit \textperiodcentered{} história 5 versões no jornal}

%CRISTAL {"arestas": [{"alvo": "capacidade_de_canal", "confianca": 1.0, "epistemico": "fato", "origem": "wiki", "rel": "usa"}, {"alvo": "eng_telecomunicacoes", "confianca": 1.0, "epistemico": "fato", "origem": "curriculo", "rel": "parte_de"}], "confianca": 1.0, "contraexemplos": [], "descricao": "ADICIONAR redundância controlada para CORRIGIR erros no receptor sem retransmitir. O oposto da compressão: troca-se taxa por confiabilidade. O teorema de Shannon promete erro→0 abaixo da capacidade.", "epistemico": "fato", "exemplos": [], "id": "codificacao_de_canal", "memoria": "semantica", "meta": {"dominio": "engenharia", "fonte": ""}, "origem": "wiki", "palavras_chave": [], "sinonimos": [], "tipo": "conceito", "titulo": "Codificação de Canal (FEC)"}
\section{Codificação de Canal (FEC)}
ADICIONAR redundância controlada para CORRIGIR erros no receptor sem retransmitir. O oposto da compressão: troca-se taxa por confiabilidade. O teorema de Shannon promete erro\ensuremath{\to}0 abaixo da capacidade.
\prov{relações: usa capacidade\_de\_canal; parte\_de eng\_telecomunicacoes}
\prov{proveniência: wiki \textperiodcentered{} fato \textperiodcentered{} confiança 1.0 \textperiodcentered{} domínio engenharia \textperiodcentered{} história 3 versões no jornal}

%CRISTAL {"arestas": [{"alvo": "entropia_shannon", "confianca": 1.0, "epistemico": "fato", "origem": "wiki", "rel": "usa"}, {"alvo": "teoria_da_informacao", "confianca": 1.0, "epistemico": "fato", "origem": "wiki", "rel": "parte_de"}, {"alvo": "eng_telecomunicacoes", "confianca": 1.0, "epistemico": "fato", "origem": "curriculo", "rel": "parte_de"}], "confianca": 1.0, "contraexemplos": [], "descricao": "Remover redundância para chegar perto da entropia da fonte — Huffman, aritmética, Lempel-Ziv. O teorema de Shannon garante o limite: não se comprime abaixo da entropia sem perda.", "epistemico": "fato", "exemplos": [], "id": "codificacao_de_fonte", "memoria": "semantica", "meta": {"autor": "Huffman", "dominio": "engenharia", "fonte": ""}, "origem": "wiki", "palavras_chave": [], "sinonimos": [], "tipo": "conceito", "titulo": "Codificação de Fonte (Compressão)"}
\section{Codificação de Fonte (Compressão)}
Remover redundância para chegar perto da entropia da fonte — Huffman, aritmética, Lempel-Ziv. O teorema de Shannon garante o limite: não se comprime abaixo da entropia sem perda.
\prov{relações: usa entropia\_shannon; parte\_de teoria\_da\_informacao; parte\_de eng\_telecomunicacoes}
\prov{proveniência: wiki \textperiodcentered{} fato \textperiodcentered{} confiança 1.0 \textperiodcentered{} domínio engenharia \textperiodcentered{} história 4 versões no jornal}

%CRISTAL {"arestas": [{"alvo": "codigo_moderno_capacidade", "confianca": 1.0, "epistemico": "fato", "origem": "wiki", "rel": "especializa"}, {"alvo": "camada_fisica_nr", "confianca": 1.0, "epistemico": "fato", "origem": "wiki", "rel": "usa"}], "confianca": 1.0, "contraexemplos": [], "descricao": "O NR usa LDPC nos DADOS (alta taxa, chega perto de Shannon) e códigos POLAR no CONTROLE (blocos curtos, os primeiros códigos com capacidade PROVADA, de Arıkan). Dois códigos modernos, cada um no seu nicho.", "epistemico": "fato", "exemplos": [], "id": "codificacao_nr_ldpc_polar", "memoria": "semantica", "meta": {"autor": "Arıkan/Gallager", "dominio": "engenharia", "fonte": ""}, "origem": "wiki", "palavras_chave": [], "sinonimos": [], "tipo": "conceito", "titulo": "Codificação do NR (LDPC + Polar)"}
\section{Codificação do NR (LDPC + Polar)}
O NR usa LDPC nos DADOS (alta taxa, chega perto de Shannon) e códigos POLAR no CONTROLE (blocos curtos, os primeiros códigos com capacidade PROVADA, de Ar\ensuremath{\imath}kan). Dois códigos modernos, cada um no seu nicho.
\prov{relações: especializa codigo\_moderno\_capacidade; usa camada\_fisica\_nr}
\prov{proveniência: wiki \textperiodcentered{} fato \textperiodcentered{} confiança 1.0 \textperiodcentered{} domínio engenharia \textperiodcentered{} história 3 versões no jornal}

%CRISTAL {"arestas": [{"alvo": "codificacao_de_canal", "confianca": 1.0, "epistemico": "fato", "origem": "wiki", "rel": "especializa"}, {"alvo": "eng_telecomunicacoes", "confianca": 1.0, "epistemico": "fato", "origem": "curriculo", "rel": "parte_de"}], "confianca": 1.0, "contraexemplos": [], "descricao": "Codificar o fluxo com memória (registradores de deslocamento) e decodificar pelo caminho de máxima verossimilhança na treliça — o algoritmo de Viterbi. Espinha do GSM, do Wi-Fi e do espaço profundo.", "epistemico": "fato", "exemplos": [], "id": "codigo_convolucional", "memoria": "semantica", "meta": {"autor": "Viterbi", "dominio": "engenharia", "fonte": ""}, "origem": "wiki", "palavras_chave": [], "sinonimos": [], "tipo": "conceito", "titulo": "Código Convolucional e Viterbi"}
\section{Código Convolucional e Viterbi}
Codificar o fluxo com memória (registradores de deslocamento) e decodificar pelo caminho de máxima verossimilhança na treliça — o algoritmo de Viterbi. Espinha do GSM, do Wi-Fi e do espaço profundo.
\prov{relações: especializa codificacao\_de\_canal; parte\_de eng\_telecomunicacoes}
\prov{proveniência: wiki \textperiodcentered{} fato \textperiodcentered{} confiança 1.0 \textperiodcentered{} domínio engenharia \textperiodcentered{} história 3 versões no jornal}

%CRISTAL {"arestas": [{"alvo": "codificacao_de_canal", "confianca": 1.0, "epistemico": "fato", "origem": "wiki", "rel": "especializa"}, {"alvo": "eng_telecomunicacoes", "confianca": 1.0, "epistemico": "fato", "origem": "curriculo", "rel": "parte_de"}], "confianca": 1.0, "contraexemplos": [], "descricao": "O corretor pioneiro: bits de paridade posicionados para localizar e corrigir 1 erro por bloco (distância mínima 3). A ideia de distância de Hamming funda toda a teoria de códigos.", "epistemico": "fato", "exemplos": [], "id": "codigo_de_hamming", "memoria": "semantica", "meta": {"autor": "Hamming", "dominio": "engenharia", "fonte": ""}, "origem": "wiki", "palavras_chave": [], "sinonimos": [], "tipo": "conceito", "titulo": "Código de Hamming"}
\section{Código de Hamming}
O corretor pioneiro: bits de paridade posicionados para localizar e corrigir 1 erro por bloco (distância mínima 3). A ideia de distância de Hamming funda toda a teoria de códigos.
\prov{relações: especializa codificacao\_de\_canal; parte\_de eng\_telecomunicacoes}
\prov{proveniência: wiki \textperiodcentered{} fato \textperiodcentered{} confiança 1.0 \textperiodcentered{} domínio engenharia \textperiodcentered{} história 3 versões no jornal}

%CRISTAL {"arestas": [{"alvo": "codificacao_de_canal", "confianca": 1.0, "epistemico": "fato", "origem": "wiki", "rel": "especializa"}, {"alvo": "capacidade_de_canal", "confianca": 1.0, "epistemico": "fato", "origem": "wiki", "rel": "confirma"}, {"alvo": "eng_telecomunicacoes", "confianca": 1.0, "epistemico": "fato", "origem": "curriculo", "rel": "parte_de"}], "confianca": 1.0, "contraexemplos": [], "descricao": "Códigos que chegam a fração de dB da capacidade de Shannon por decodificação ITERATIVA (propagação de crença). Turbo (3G/4G) e LDPC (5G, Wi-Fi 6, DVB) realizaram a promessa de 1948.", "epistemico": "fato", "exemplos": [], "id": "codigo_moderno_capacidade", "memoria": "semantica", "meta": {"autor": "Berrou/Gallager", "dominio": "engenharia", "fonte": ""}, "origem": "wiki", "palavras_chave": [], "sinonimos": [], "tipo": "conceito", "titulo": "Turbo-códigos e LDPC"}
\section{Turbo-códigos e LDPC}
Códigos que chegam a fração de dB da capacidade de Shannon por decodificação ITERATIVA (propagação de crença). Turbo (3G/4G) e LDPC (5G, Wi-Fi 6, DVB) realizaram a promessa de 1948.
\prov{relações: especializa codificacao\_de\_canal; confirma capacidade\_de\_canal; parte\_de eng\_telecomunicacoes}
\prov{proveniência: wiki \textperiodcentered{} fato \textperiodcentered{} confiança 1.0 \textperiodcentered{} domínio engenharia \textperiodcentered{} história 4 versões no jornal}

%CRISTAL {"arestas": [{"alvo": "olho_de_cock_web", "confianca": 1.0, "epistemico": "fato", "origem": "wiki", "rel": "parte_de"}, {"alvo": "cotacao_parabola", "confianca": 1.0, "epistemico": "fato", "origem": "wiki", "rel": "usa"}, {"alvo": "conjugados_galois_parabola", "confianca": 1.0, "epistemico": "fato", "origem": "wiki", "rel": "aplicacao"}], "confianca": 1.0, "contraexemplos": [], "descricao": "O herói do topo: o lastro conservado total (sats/BTC) dividido em barra branco(c²·fase·lastro) / preto(a·calor·fees), com a nota do balanço a+c²=1 e H/φ. Lastro '≠0', conservado — a contabilidade sem livro renderizada.", "epistemico": "fato", "exemplos": [], "id": "cofre_parabola_hero", "memoria": "semantica", "meta": {"autor": "broca-so", "dominio": "cockpit", "fonte": "goldchain/cockpit/dtc_renderer.js:renderCofreFromDtc"}, "origem": "wiki", "palavras_chave": [], "sinonimos": [], "tipo": "conceito", "titulo": "O Cofre — Parábola a+c²=1 (herói do topo)"}
\section{O Cofre — Parábola a+c\ensuremath{^{2}}=1 (herói do topo)}
O herói do topo: o lastro conservado total (sats/BTC) dividido em barra branco(c\ensuremath{^{2}}\textperiodcentered fase\textperiodcentered lastro) / preto(a\textperiodcentered calor\textperiodcentered fees), com a nota do balanço a+c\ensuremath{^{2}}=1 e H/\ensuremath{\varphi}. Lastro '\ensuremath{\ne}0', conservado — a contabilidade sem livro renderizada.
\prov{relações: parte\_de olho\_de\_cock\_web; usa cotacao\_parabola; aplicacao conjugados\_galois\_parabola}
\prov{proveniência: wiki \textperiodcentered{} goldchain/cockpit/dtc\_renderer.js:renderCofreFromDtc \textperiodcentered{} fato \textperiodcentered{} confiança 1.0 \textperiodcentered{} domínio cockpit \textperiodcentered{} história 4 versões no jornal}

%CRISTAL {"arestas": [{"alvo": "imagem_para_texto", "confianca": 1.0, "epistemico": "fato", "origem": "wiki", "rel": "parte_de"}, {"alvo": "kmeans_atratores", "confianca": 1.0, "epistemico": "fato", "origem": "wiki", "rel": "usa"}, {"alvo": "colapso", "confianca": 1.0, "epistemico": "fato", "origem": "wiki", "rel": "eh_um"}], "confianca": 1.0, "contraexemplos": [], "descricao": "Uma imagem de glifo desaba em DOIS rótulos ortogonais: o CARACTERE (o conteúdo) e a CALIGRAFIA (o estilo). São dois eixos independentes — o mesmo 'A' em seis fontes, a mesma fonte em 26 letras. O reconhecimento fatora o colapso: colapsa o conteúdo por uma feature invariante ao estilo, e o estilo por outra invariante ao conteúdo. Conteúdo e estilo, dois colapsos de uma imagem.", "epistemico": "fato", "exemplos": [], "id": "colapso_fatorado_glifo", "memoria": "semantica", "meta": {"dominio": "imagem", "fonte": "glifos minerais"}, "origem": "wiki", "palavras_chave": [], "sinonimos": [], "tipo": "conceito", "titulo": "Colapso Fatorado (caractere × caligrafia)"}
\section{Colapso Fatorado (caractere $\times$ caligrafia)}
Uma imagem de glifo desaba em DOIS rótulos ortogonais: o CARACTERE (o conteúdo) e a CALIGRAFIA (o estilo). São dois eixos independentes — o mesmo 'A' em seis fontes, a mesma fonte em 26 letras. O reconhecimento fatora o colapso: colapsa o conteúdo por uma feature invariante ao estilo, e o estilo por outra invariante ao conteúdo. Conteúdo e estilo, dois colapsos de uma imagem.
\prov{relações: parte\_de imagem\_para\_texto; usa kmeans\_atratores; eh\_um colapso}
\prov{proveniência: wiki \textperiodcentered{} glifos minerais \textperiodcentered{} fato \textperiodcentered{} confiança 1.0 \textperiodcentered{} domínio imagem \textperiodcentered{} história 4 versões no jornal}

%CRISTAL {"arestas": [{"alvo": "eletronica", "confianca": 1.0, "epistemico": "fato", "origem": "wiki", "rel": "parte_de"}, {"alvo": "lei_de_ohm", "confianca": 1.0, "epistemico": "fato", "origem": "wiki", "rel": "usa"}], "confianca": 1.0, "contraexemplos": [], "descricao": "Resistor (limita), capacitor (armazena carga) e indutor (armazena campo magnético) — os elementos que moldam sinais sem amplificar.", "epistemico": "fato", "exemplos": [], "id": "componente_passivo", "memoria": "semantica", "meta": {"dominio": "engenharia", "fonte": ""}, "origem": "wiki", "palavras_chave": [], "sinonimos": [], "tipo": "conceito", "titulo": "Componentes Passivos"}
\section{Componentes Passivos}
Resistor (limita), capacitor (armazena carga) e indutor (armazena campo magnético) — os elementos que moldam sinais sem amplificar.
\prov{relações: parte\_de eletronica; usa lei\_de\_ohm}
\prov{proveniência: wiki \textperiodcentered{} fato \textperiodcentered{} confiança 1.0 \textperiodcentered{} domínio engenharia \textperiodcentered{} história 3 versões no jornal}

%CRISTAL {"arestas": [{"alvo": "floco_recipiente_termico", "confianca": 1.0, "epistemico": "fato", "origem": "wiki", "rel": "parte_de"}, {"alvo": "pow_dificuldade_calor", "confianca": 1.0, "epistemico": "fato", "origem": "wiki", "rel": "relaciona"}], "confianca": 1.0, "contraexemplos": [], "descricao": "Comprimir o floco (guardar menos camadas k) AUMENTA a pressão interna: P(k)=W(k)/k, a densidade de energia por camada. Menos camadas concentram a energia — a pressão sobe. É um recipiente energético sem limite RÍGIDO (sempre dá pra comprimir a k=1), mas a pressão tem teto (σ_1²) e o calor dissipado cresce. Verif.: P sobe de 0.005 a 0.99 quando k vai de 220 a 1.", "epistemico": "fato", "exemplos": [], "id": "comprimir_aumenta_pressao", "memoria": "semantica", "meta": {"dominio": "imagem", "fonte": "floco térmico"}, "origem": "wiki", "palavras_chave": [], "sinonimos": [], "tipo": "conceito", "titulo": "Comprimir Aumenta a Pressão do Floco"}
\section{Comprimir Aumenta a Pressão do Floco}
Comprimir o floco (guardar menos camadas k) AUMENTA a pressão interna: P(k)=W(k)/k, a densidade de energia por camada. Menos camadas concentram a energia — a pressão sobe. É um recipiente energético sem limite RÍGIDO (sempre dá pra comprimir a k=1), mas a pressão tem teto (\ensuremath{\sigma}\_1\ensuremath{^{2}}) e o calor dissipado cresce. Verif.: P sobe de 0.005 a 0.99 quando k vai de 220 a 1.
\prov{relações: parte\_de floco\_recipiente\_termico; relaciona pow\_dificuldade\_calor}
\prov{proveniência: wiki \textperiodcentered{} floco térmico \textperiodcentered{} fato \textperiodcentered{} confiança 1.0 \textperiodcentered{} domínio imagem \textperiodcentered{} história 6 versões no jornal}

%CRISTAL {"arestas": [{"alvo": "sistemas_distribuidos", "confianca": 1.0, "epistemico": "fato", "origem": "wiki", "rel": "usa"}, {"alvo": "virtualizacao", "confianca": 1.0, "epistemico": "fato", "origem": "wiki", "rel": "usa"}, {"alvo": "energia_eletrica", "confianca": 1.0, "epistemico": "fato", "origem": "wiki", "rel": "usa"}], "confianca": 1.0, "contraexemplos": [], "descricao": "Alugar computação, armazenamento e rede sob demanda por internet em vez de manter servidores próprios. Elástica (cresce e encolhe com a carga) e paga-por-uso; o pólo OPOSTO ao edge — centraliza onde o edge distribui. É onde o dado do 5G e dos sensores aterrissa e vira decisão.", "epistemico": "fato", "exemplos": [], "id": "computacao_em_nuvem", "memoria": "semantica", "meta": {"dominio": "engenharia", "fonte": "nuvem e data centers"}, "origem": "wiki", "palavras_chave": [], "sinonimos": [], "tipo": "conceito", "titulo": "Computação em Nuvem"}
\section{Computação em Nuvem}
Alugar computação, armazenamento e rede sob demanda por internet em vez de manter servidores próprios. Elástica (cresce e encolhe com a carga) e paga-por-uso; o pólo OPOSTO ao edge — centraliza onde o edge distribui. É onde o dado do 5G e dos sensores aterrissa e vira decisão.
\prov{relações: usa sistemas\_distribuidos; usa virtualizacao; usa energia\_eletrica}
\prov{proveniência: wiki \textperiodcentered{} nuvem e data centers \textperiodcentered{} fato \textperiodcentered{} confiança 1.0 \textperiodcentered{} domínio engenharia \textperiodcentered{} história 8 versões no jornal}

%CRISTAL {"arestas": [{"alvo": "optoeletronica", "confianca": 1.0, "epistemico": "fato", "origem": "wiki", "rel": "usa"}, {"alvo": "multiplexacao", "confianca": 1.0, "epistemico": "fato", "origem": "wiki", "rel": "usa"}, {"alvo": "internet", "confianca": 1.0, "epistemico": "fato", "origem": "wiki", "rel": "explica"}, {"alvo": "eng_telecomunicacoes", "confianca": 1.0, "epistemico": "fato", "origem": "curriculo", "rel": "parte_de"}], "confianca": 1.0, "contraexemplos": [], "descricao": "Enviar bits como pulsos de luz numa fibra de vidro — banda colossal, baixa perda, imune a interferência. A espinha dorsal da internet: WDM leva dezenas de Tb/s por fibra entre continentes.", "epistemico": "fato", "exemplos": [], "id": "comunicacao_optica", "memoria": "semantica", "meta": {"dominio": "engenharia", "fonte": ""}, "origem": "wiki", "palavras_chave": [], "sinonimos": [], "tipo": "conceito", "titulo": "Comunicação Óptica (Fibra)"}
\section{Comunicação Óptica (Fibra)}
Enviar bits como pulsos de luz numa fibra de vidro — banda colossal, baixa perda, imune a interferência. A espinha dorsal da internet: WDM leva dezenas de Tb/s por fibra entre continentes.
\prov{relações: usa optoeletronica; usa multiplexacao; explica internet; parte\_de eng\_telecomunicacoes}
\prov{proveniência: wiki \textperiodcentered{} fato \textperiodcentered{} confiança 1.0 \textperiodcentered{} domínio engenharia \textperiodcentered{} história 5 versões no jornal}

%CRISTAL {"arestas": [{"alvo": "balanco_de_enlace", "confianca": 1.0, "epistemico": "fato", "origem": "wiki", "rel": "usa"}, {"alvo": "propagacao_de_ondas", "confianca": 1.0, "epistemico": "fato", "origem": "wiki", "rel": "usa"}, {"alvo": "eng_telecomunicacoes", "confianca": 1.0, "epistemico": "fato", "origem": "curriculo", "rel": "parte_de"}], "confianca": 1.0, "contraexemplos": [], "descricao": "Enlace de rádio via satélite (GEO distante e estável, ou constelações LEO de baixa latência) — cobertura onde a fibra não chega. Grande atraso e balanço de enlace apertado definem o projeto.", "epistemico": "fato", "exemplos": [], "id": "comunicacao_por_satelite", "memoria": "semantica", "meta": {"dominio": "engenharia", "fonte": ""}, "origem": "wiki", "palavras_chave": [], "sinonimos": [], "tipo": "conceito", "titulo": "Comunicação por Satélite"}
\section{Comunicação por Satélite}
Enlace de rádio via satélite (GEO distante e estável, ou constelações LEO de baixa latência) — cobertura onde a fibra não chega. Grande atraso e balanço de enlace apertado definem o projeto.
\prov{relações: usa balanco\_de\_enlace; usa propagacao\_de\_ondas; parte\_de eng\_telecomunicacoes}
\prov{proveniência: wiki \textperiodcentered{} fato \textperiodcentered{} confiança 1.0 \textperiodcentered{} domínio engenharia \textperiodcentered{} história 4 versões no jornal}

%CRISTAL {"arestas": [{"alvo": "acesso_multiplo", "confianca": 1.0, "epistemico": "fato", "origem": "wiki", "rel": "usa"}, {"alvo": "propagacao_de_ondas", "confianca": 1.0, "epistemico": "fato", "origem": "wiki", "rel": "usa"}, {"alvo": "antena", "confianca": 1.0, "epistemico": "fato", "origem": "wiki", "rel": "usa"}, {"alvo": "eng_telecomunicacoes", "confianca": 1.0, "epistemico": "fato", "origem": "curriculo", "rel": "parte_de"}], "confianca": 1.0, "contraexemplos": [], "descricao": "Cobrir o território com CÉLULAS que reusam frequência, com handover entre elas — a arquitetura que deu telefone e internet a bilhões em movimento. Estação-base, núcleo e o espectro como recurso escasso.", "epistemico": "fato", "exemplos": [], "id": "comunicacoes_moveis", "memoria": "semantica", "meta": {"dominio": "engenharia", "fonte": ""}, "origem": "wiki", "palavras_chave": [], "sinonimos": [], "tipo": "conceito", "titulo": "Comunicações Móveis (Celular)"}
\section{Comunicações Móveis (Celular)}
Cobrir o território com CÉLULAS que reusam frequência, com handover entre elas — a arquitetura que deu telefone e internet a bilhões em movimento. Estação-base, núcleo e o espectro como recurso escasso.
\prov{relações: usa acesso\_multiplo; usa propagacao\_de\_ondas; usa antena; parte\_de eng\_telecomunicacoes}
\prov{proveniência: wiki \textperiodcentered{} fato \textperiodcentered{} confiança 1.0 \textperiodcentered{} domínio engenharia \textperiodcentered{} história 5 versões no jornal}

%CRISTAL {"arestas": [{"alvo": "rede_autonoma_broca", "confianca": 1.0, "epistemico": "fato", "origem": "wiki", "rel": "parte_de"}, {"alvo": "meta-estavel_comunidades_organicas_de_software", "confianca": 1.0, "epistemico": "fato", "origem": "wiki", "rel": "usa"}, {"alvo": "estatuto_o_que_fecha_o_que_segue_para_o_proximo_nivel", "confianca": 1.0, "epistemico": "fato", "origem": "wiki", "rel": "usa"}, {"alvo": "producao_entre_pares", "confianca": 1.0, "epistemico": "fato", "origem": "wiki", "rel": "relaciona"}], "confianca": 1.0, "contraexemplos": [], "descricao": "Uma comunidade não se cadastra: ela se forma quando pares compartilham um tecido (a mesma banda). Não há cliente nem servidor — há pares, e a reciprocidade é a gramática ('quem manda recebe'). Comunidades de software E de arte co-emergem no mesmo substrato: a banda é a fronteira orgânica.", "epistemico": "inferencia", "exemplos": [], "id": "comunidade_organica_por_banda", "memoria": "semantica", "meta": {"dominio": "arquitetura", "fonte": "semilla_rede_autonoma"}, "origem": "wiki", "palavras_chave": [], "sinonimos": [], "tipo": "conceito", "titulo": "Comunidade orgânica = pares numa banda"}
\section{Comunidade orgânica = pares numa banda}
Uma comunidade não se cadastra: ela se forma quando pares compartilham um tecido (a mesma banda). Não há cliente nem servidor — há pares, e a reciprocidade é a gramática ('quem manda recebe'). Comunidades de software E de arte co-emergem no mesmo substrato: a banda é a fronteira orgânica.
\prov{relações: parte\_de rede\_autonoma\_broca; usa meta-estavel\_comunidades\_organicas\_de\_software; usa estatuto\_o\_que\_fecha\_o\_que\_segue\_para\_o\_proximo\_nivel; relaciona producao\_entre\_pares}
\prov{proveniência: wiki \textperiodcentered{} semilla\_rede\_autonoma \textperiodcentered{} inferencia \textperiodcentered{} confiança 1.0 \textperiodcentered{} domínio arquitetura \textperiodcentered{} história 5 versões no jornal}

%CRISTAL {"arestas": [{"alvo": "funcoes_de_rede_5gc", "confianca": 1.0, "epistemico": "fato", "origem": "wiki", "rel": "usa"}], "confianca": 1.0, "contraexemplos": [], "descricao": "O passo final do handover no núcleo: o AMF/SMF reaponta o túnel de dados do UPF para o gNB alvo, e libera os recursos da origem. A partir daí o pacote desce pela nova célula.", "epistemico": "fato", "exemplos": [], "id": "comutacao_de_caminho", "memoria": "semantica", "meta": {"dominio": "engenharia", "fonte": ""}, "origem": "wiki", "palavras_chave": [], "sinonimos": [], "tipo": "conceito", "titulo": "Comutação de Caminho (Path Switch)"}
\section{Comutação de Caminho (Path Switch)}
O passo final do handover no núcleo: o AMF/SMF reaponta o túnel de dados do UPF para o gNB alvo, e libera os recursos da origem. A partir daí o pacote desce pela nova célula.
\prov{relações: usa funcoes\_de\_rede\_5gc}
\prov{proveniência: wiki \textperiodcentered{} fato \textperiodcentered{} confiança 1.0 \textperiodcentered{} domínio engenharia \textperiodcentered{} história 2 versões no jornal}

%CRISTAL {"arestas": [{"alvo": "banda_proibida", "confianca": 1.0, "epistemico": "fato", "origem": "wiki", "rel": "usa"}], "confianca": 1.0, "contraexemplos": [], "descricao": "A classificação pela banda proibida: metal (conduz sempre), isolante (nunca), semicondutor (conduz sob controle) — a base da eletrônica.", "epistemico": "inferencia", "exemplos": [], "id": "condutor_isolante_semicondutor", "memoria": "semantica", "meta": {"dominio": "engenharia", "fonte": ""}, "origem": "wiki", "palavras_chave": [], "sinonimos": [], "tipo": "conceito", "titulo": "Condutor, Isolante e Semicondutor"}
\section{Condutor, Isolante e Semicondutor}
A classificação pela banda proibida: metal (conduz sempre), isolante (nunca), semicondutor (conduz sob controle) — a base da eletrônica.
\prov{relações: usa banda\_proibida}
\prov{proveniência: wiki \textperiodcentered{} inferencia \textperiodcentered{} confiança 1.0 \textperiodcentered{} domínio engenharia \textperiodcentered{} história 2 versões no jornal}

%CRISTAL {"arestas": [{"alvo": "telemetria", "confianca": 1.0, "epistemico": "fato", "origem": "wiki", "rel": "explica"}], "confianca": 1.0, "contraexemplos": [], "descricao": "A capacidade de os nós trocarem dados: o fio nervoso que liga sensores, carros, prédios e o centro. Sem ela a smart city é cega-muda — a telemetria não chega, o colapso na ponta não vira decisão.", "epistemico": "fato", "exemplos": [], "id": "conectividade", "memoria": "semantica", "meta": {"dominio": "engenharia", "fonte": "conectividade e 5G"}, "origem": "wiki", "palavras_chave": [], "sinonimos": [], "tipo": "conceito", "titulo": "Conectividade"}
\section{Conectividade}
A capacidade de os nós trocarem dados: o fio nervoso que liga sensores, carros, prédios e o centro. Sem ela a smart city é cega-muda — a telemetria não chega, o colapso na ponta não vira decisão.
\prov{relações: explica telemetria}
\prov{proveniência: wiki \textperiodcentered{} conectividade e 5G \textperiodcentered{} fato \textperiodcentered{} confiança 1.0 \textperiodcentered{} domínio engenharia \textperiodcentered{} história 2 versões no jornal}

%CRISTAL {"arestas": [{"alvo": "protocolo_tcp", "confianca": 1.0, "epistemico": "fato", "origem": "wiki", "rel": "parte_de"}, {"alvo": "eng_telecomunicacoes", "confianca": 1.0, "epistemico": "fato", "origem": "curriculo", "rel": "parte_de"}], "confianca": 1.0, "contraexemplos": [], "descricao": "O TCP ajusta o ritmo de envio para não afogar a rede — janela deslizante, partida lenta, AIMD (Reno, CUBIC, BBR). O que evita o colapso da internet sob carga; cortesia embutida no protocolo.", "epistemico": "fato", "exemplos": [], "id": "controle_de_congestao", "memoria": "semantica", "meta": {"autor": "Van Jacobson", "dominio": "engenharia", "fonte": ""}, "origem": "wiki", "palavras_chave": [], "sinonimos": [], "tipo": "conceito", "titulo": "Controle de Congestão"}
\section{Controle de Congestão}
O TCP ajusta o ritmo de envio para não afogar a rede — janela deslizante, partida lenta, AIMD (Reno, CUBIC, BBR). O que evita o colapso da internet sob carga; cortesia embutida no protocolo.
\prov{relações: parte\_de protocolo\_tcp; parte\_de eng\_telecomunicacoes}
\prov{proveniência: wiki \textperiodcentered{} fato \textperiodcentered{} confiança 1.0 \textperiodcentered{} domínio engenharia \textperiodcentered{} história 3 versões no jornal}

%CRISTAL {"arestas": [{"alvo": "eletronica_digital_analogica", "confianca": 1.0, "epistemico": "fato", "origem": "wiki", "rel": "usa"}], "confianca": 1.0, "contraexemplos": [], "descricao": "A ponte entre os mundos: amostrar um sinal contínuo em bits (ADC) e reconstruí-lo (DAC) — como o real entra e sai do digital.", "epistemico": "fato", "exemplos": [], "id": "conversao_ad_da", "memoria": "semantica", "meta": {"dominio": "engenharia", "fonte": ""}, "origem": "wiki", "palavras_chave": [], "sinonimos": [], "tipo": "conceito", "titulo": "Conversão A/D e D/A"}
\section{Conversão A/D e D/A}
A ponte entre os mundos: amostrar um sinal contínuo em bits (ADC) e reconstruí-lo (DAC) — como o real entra e sai do digital.
\prov{relações: usa eletronica\_digital\_analogica}
\prov{proveniência: wiki \textperiodcentered{} fato \textperiodcentered{} confiança 1.0 \textperiodcentered{} domínio engenharia \textperiodcentered{} história 2 versões no jornal}

%CRISTAL {"arestas": [{"alvo": "inducao_eletromagnetica", "confianca": 1.0, "epistemico": "fato", "origem": "wiki", "rel": "usa"}, {"alvo": "eletromagnetismo_eng", "confianca": 1.0, "epistemico": "fato", "origem": "wiki", "rel": "usa"}, {"alvo": "eng_sistemas_energia", "confianca": 1.0, "epistemico": "fato", "origem": "curriculo", "rel": "parte_de"}], "confianca": 1.0, "contraexemplos": [], "descricao": "O princípio que liga corrente e movimento: força de Lorentz e lei de Faraday convertem energia elétrica em mecânica e vice-versa. A física comum ao motor e ao gerador.", "epistemico": "fato", "exemplos": [], "id": "conversao_eletromecanica", "memoria": "semantica", "meta": {"autor": "Faraday", "dominio": "engenharia", "fonte": ""}, "origem": "wiki", "palavras_chave": [], "sinonimos": [], "tipo": "conceito", "titulo": "Conversão Eletromecânica de Energia"}
\section{Conversão Eletromecânica de Energia}
O princípio que liga corrente e movimento: força de Lorentz e lei de Faraday convertem energia elétrica em mecânica e vice-versa. A física comum ao motor e ao gerador.
\prov{relações: usa inducao\_eletromagnetica; usa eletromagnetismo\_eng; parte\_de eng\_sistemas\_energia}
\prov{proveniência: wiki \textperiodcentered{} fato \textperiodcentered{} confiança 1.0 \textperiodcentered{} domínio engenharia \textperiodcentered{} história 4 versões no jornal}

%CRISTAL {"arestas": [{"alvo": "sensor", "confianca": 1.0, "epistemico": "fato", "origem": "wiki", "rel": "usa"}, {"alvo": "sensor_e_colapso", "confianca": 1.0, "epistemico": "fato", "origem": "wiki", "rel": "usa"}], "confianca": 1.0, "contraexemplos": [], "descricao": "O circuito que amostra e quantiza um sinal analógico contínuo em números digitais: N bits dão 2^N níveis. Quanto mais bits, mais fino o colapso (menos erro). É o órgão físico da medição.", "epistemico": "fato", "exemplos": [], "id": "conversor_ad", "memoria": "semantica", "meta": {"dominio": "engenharia", "fonte": "IoT e sensores"}, "origem": "wiki", "palavras_chave": [], "sinonimos": [], "tipo": "conceito", "titulo": "Conversor A/D (ADC)"}
\section{Conversor A/D (ADC)}
O circuito que amostra e quantiza um sinal analógico contínuo em números digitais: N bits dão 2\^{}N níveis. Quanto mais bits, mais fino o colapso (menos erro). É o órgão físico da medição.
\prov{relações: usa sensor; usa sensor\_e\_colapso}
\prov{proveniência: wiki \textperiodcentered{} IoT e sensores \textperiodcentered{} fato \textperiodcentered{} confiança 1.0 \textperiodcentered{} domínio engenharia \textperiodcentered{} história 6 versões no jornal}

%CRISTAL {"arestas": [{"alvo": "conversao_ad_da", "confianca": 1.0, "epistemico": "fato", "origem": "wiki", "rel": "especializa"}, {"alvo": "eng_eletronica", "confianca": 1.0, "epistemico": "fato", "origem": "curriculo", "rel": "parte_de"}], "confianca": 1.0, "contraexemplos": [], "descricao": "A fronteira entre o mundo contínuo e o digital: quantização, taxa, resolução e ruído de quantização. Todo sinal real entra por um ADC e sai por um DAC.", "epistemico": "fato", "exemplos": [], "id": "conversor_ad_da_disc", "memoria": "semantica", "meta": {"dominio": "engenharia", "fonte": ""}, "origem": "wiki", "palavras_chave": [], "sinonimos": [], "tipo": "conceito", "titulo": "Conversores A/D e D/A"}
\section{Conversores A/D e D/A}
A fronteira entre o mundo contínuo e o digital: quantização, taxa, resolução e ruído de quantização. Todo sinal real entra por um ADC e sai por um DAC.
\prov{relações: especializa conversao\_ad\_da; parte\_de eng\_eletronica}
\prov{proveniência: wiki \textperiodcentered{} fato \textperiodcentered{} confiança 1.0 \textperiodcentered{} domínio engenharia \textperiodcentered{} história 3 versões no jornal}

%CRISTAL {"arestas": [{"alvo": "eletronica_de_potencia", "confianca": 1.0, "epistemico": "fato", "origem": "wiki", "rel": "parte_de"}, {"alvo": "corrente_alternada_continua", "confianca": 1.0, "epistemico": "fato", "origem": "wiki", "rel": "usa"}, {"alvo": "energia_solar", "confianca": 1.0, "epistemico": "fato", "origem": "wiki", "rel": "aplicacao"}, {"alvo": "eng_sistemas_energia", "confianca": 1.0, "epistemico": "fato", "origem": "curriculo", "rel": "parte_de"}], "confianca": 1.0, "contraexemplos": [], "descricao": "As quatro conversões: CA→CC (retificador), CC→CA (inversor), CC→CC (chopper) e CA→CA. O inversor é o que põe o painel solar e a bateria na rede CA; o retificador alimenta toda eletrônica.", "epistemico": "fato", "exemplos": [], "id": "conversor_de_potencia", "memoria": "semantica", "meta": {"dominio": "engenharia", "fonte": ""}, "origem": "wiki", "palavras_chave": [], "sinonimos": [], "tipo": "conceito", "titulo": "Conversores (Retificador/Inversor/CC-CC)"}
\section{Conversores (Retificador/Inversor/CC-CC)}
As quatro conversões: CA\ensuremath{\to}CC (retificador), CC\ensuremath{\to}CA (inversor), CC\ensuremath{\to}CC (chopper) e CA\ensuremath{\to}CA. O inversor é o que põe o painel solar e a bateria na rede CA; o retificador alimenta toda eletrônica.
\prov{relações: parte\_de eletronica\_de\_potencia; usa corrente\_alternada\_continua; aplicacao energia\_solar; parte\_de eng\_sistemas\_energia}
\prov{proveniência: wiki \textperiodcentered{} fato \textperiodcentered{} confiança 1.0 \textperiodcentered{} domínio engenharia \textperiodcentered{} história 5 versões no jornal}

%CRISTAL {"arestas": [{"alvo": "rede_eletrica", "confianca": 1.0, "epistemico": "fato", "origem": "wiki", "rel": "parte_de"}, {"alvo": "numero_de_salem", "confianca": 1.0, "epistemico": "fato", "origem": "wiki", "rel": "usa"}, {"alvo": "gerador_eletrico", "confianca": 1.0, "epistemico": "fato", "origem": "wiki", "rel": "usa"}], "confianca": 1.0, "contraexemplos": [], "descricao": "A rede usa CA porque o gerador é rotativo e o transformador só funciona com fluxo variável. O fasor e^{iωt} percorre o CÍRCULO UNITÁRIO (ρ=1): é a rotação pura — o Salem/borda da máquina mineral. A rede inteira é UMA rotação a 60 Hz.", "epistemico": "inferencia", "exemplos": [], "id": "corrente_alternada", "memoria": "semantica", "meta": {"classe": "borda/Salem (rotação pura, ρ=1)", "dominio": "engenharia", "fonte": "usinas e geração", "raio_fasor": 1.0}, "origem": "wiki", "palavras_chave": [], "sinonimos": [], "tipo": "conceito", "titulo": "Corrente Alternada (a rotação da rede)"}
\section{Corrente Alternada (a rotação da rede)}
A rede usa CA porque o gerador é rotativo e o transformador só funciona com fluxo variável. O fasor e\^{}\{i\ensuremath{\omega}t\} percorre o CÍRCULO UNITÁRIO (\ensuremath{\rho}=1): é a rotação pura — o Salem/borda da máquina mineral. A rede inteira é UMA rotação a 60 Hz.
\prov{relações: parte\_de rede\_eletrica; usa numero\_de\_salem; usa gerador\_eletrico}
\prov{proveniência: wiki \textperiodcentered{} usinas e geração \textperiodcentered{} inferencia \textperiodcentered{} confiança 1.0 \textperiodcentered{} domínio engenharia \textperiodcentered{} história 8 versões no jornal}

%CRISTAL {"arestas": [{"alvo": "transmissao_de_energia", "confianca": 1.0, "epistemico": "fato", "origem": "wiki", "rel": "usa"}, {"alvo": "transformador", "confianca": 1.0, "epistemico": "fato", "origem": "wiki", "rel": "usa"}], "confianca": 1.0, "contraexemplos": [], "descricao": "AC (Tesla/Westinghouse) inverte o sentido e transforma-se fácil, vencendo na transmissão; DC (Edison) é constante — a 'Guerra das Correntes'.", "epistemico": "inferencia", "exemplos": [], "id": "corrente_alternada_continua", "memoria": "semantica", "meta": {"autor": "Tesla/Edison", "dominio": "engenharia", "fonte": ""}, "origem": "wiki", "palavras_chave": [], "sinonimos": [], "tipo": "conceito", "titulo": "Corrente Alternada vs Contínua"}
\section{Corrente Alternada vs Contínua}
AC (Tesla/Westinghouse) inverte o sentido e transforma-se fácil, vencendo na transmissão; DC (Edison) é constante — a 'Guerra das Correntes'.
\prov{relações: usa transmissao\_de\_energia; usa transformador}
\prov{proveniência: wiki \textperiodcentered{} inferencia \textperiodcentered{} confiança 1.0 \textperiodcentered{} domínio engenharia \textperiodcentered{} história 3 versões no jornal}

%CRISTAL {"arestas": [{"alvo": "transmissao_de_energia", "confianca": 1.0, "epistemico": "fato", "origem": "wiki", "rel": "eh_um"}], "confianca": 1.0, "contraexemplos": [], "descricao": "Para distâncias muito longas ou cabos submarinos, transmitir em CC de alta tensão perde menos que CA (sem reatância nem efeito pelicular). Converte-se CA→CC→CA nas pontas. A exceção à rede rotativa.", "epistemico": "fato", "exemplos": [], "id": "corrente_continua_alta_tensao", "memoria": "semantica", "meta": {"dominio": "engenharia", "fonte": "distribuição e armazenamento"}, "origem": "wiki", "palavras_chave": [], "sinonimos": [], "tipo": "conceito", "titulo": "HVDC (corrente contínua em alta tensão)"}
\section{HVDC (corrente contínua em alta tensão)}
Para distâncias muito longas ou cabos submarinos, transmitir em CC de alta tensão perde menos que CA (sem reatância nem efeito pelicular). Converte-se CA\ensuremath{\to}CC\ensuremath{\to}CA nas pontas. A exceção à rede rotativa.
\prov{relações: eh\_um transmissao\_de\_energia}
\prov{proveniência: wiki \textperiodcentered{} distribuição e armazenamento \textperiodcentered{} fato \textperiodcentered{} confiança 1.0 \textperiodcentered{} domínio engenharia \textperiodcentered{} história 4 versões no jornal}

%CRISTAL {"arestas": [], "confianca": 1.0, "contraexemplos": [], "descricao": "Corrente é o fluxo de carga (ampère); tensão é a diferença de potencial que o empurra (volt). O que flui e o que empurra.", "epistemico": "fato", "exemplos": [], "id": "corrente_e_tensao", "memoria": "semantica", "meta": {"dominio": "engenharia", "fonte": ""}, "origem": "wiki", "palavras_chave": [], "sinonimos": [], "tipo": "conceito", "titulo": "Corrente e Tensão"}
\section{Corrente e Tensão}
Corrente é o fluxo de carga (ampère); tensão é a diferença de potencial que o empurra (volt). O que flui e o que empurra.
\prov{proveniência: wiki \textperiodcentered{} fato \textperiodcentered{} confiança 1.0 \textperiodcentered{} domínio engenharia}

%CRISTAL {"arestas": [{"alvo": "segmentacao_de_imagem", "confianca": 1.0, "epistemico": "fato", "origem": "wiki", "rel": "parte_de"}, {"alvo": "r0_borda_epidemica", "confianca": 1.0, "epistemico": "fato", "origem": "wiki", "rel": "relaciona"}, {"alvo": "orbita_de_partida", "confianca": 1.0, "epistemico": "fato", "origem": "wiki", "rel": "relaciona"}], "confianca": 1.0, "contraexemplos": [], "descricao": "A partir de uma semente, inunda os vizinhos dentro de uma tolerância — uma órbita que cresce sobre a imagem (a mesma forma do contágio R₀ e da difusão) até bater na fronteira. Segmentar é deixar a órbita ocupar sua região.", "epistemico": "inferencia", "exemplos": [], "id": "crescer_regiao", "memoria": "semantica", "meta": {"dominio": "imagem", "fonte": "segmentação"}, "origem": "wiki", "palavras_chave": [], "sinonimos": [], "tipo": "conceito", "titulo": "Crescer Região = a Órbita se Espalhando"}
\section{Crescer Região = a Órbita se Espalhando}
A partir de uma semente, inunda os vizinhos dentro de uma tolerância — uma órbita que cresce sobre a imagem (a mesma forma do contágio R\ensuremath{_{0}} e da difusão) até bater na fronteira. Segmentar é deixar a órbita ocupar sua região.
\prov{relações: parte\_de segmentacao\_de\_imagem; relaciona r0\_borda\_epidemica; relaciona orbita\_de\_partida}
\prov{proveniência: wiki \textperiodcentered{} segmentação \textperiodcentered{} inferencia \textperiodcentered{} confiança 1.0 \textperiodcentered{} domínio imagem \textperiodcentered{} história 8 versões no jornal}

%CRISTAL {"arestas": [{"alvo": "imagem_bai", "confianca": 1.0, "epistemico": "fato", "origem": "wiki", "rel": "parte_de"}, {"alvo": "crescer_regiao", "confianca": 1.0, "epistemico": "fato", "origem": "wiki", "rel": "usa"}, {"alvo": "semente_e_raiz", "confianca": 1.0, "epistemico": "fato", "origem": "wiki", "rel": "usa"}], "confianca": 1.0, "contraexemplos": [], "descricao": "O crescimento de região — inundar os vizinhos parecidos a partir da semente — É o operador P do BAI (propaga a autoridade pela cadeia) e a iniciativa do xadrez (quem dita o jogo). A órbita que espalha o rótulo é a mesma que propaga a concessão.", "epistemico": "inferencia", "exemplos": [], "id": "crescimento_e_propagacao", "memoria": "semantica", "meta": {"dominio": "ponte", "fonte": "imagem↔BAI"}, "origem": "wiki", "palavras_chave": [], "sinonimos": [], "tipo": "conceito", "titulo": "Crescer Região = Propagação P"}
\section{Crescer Região = Propagação P}
O crescimento de região — inundar os vizinhos parecidos a partir da semente — É o operador P do BAI (propaga a autoridade pela cadeia) e a iniciativa do xadrez (quem dita o jogo). A órbita que espalha o rótulo é a mesma que propaga a concessão.
\prov{relações: parte\_de imagem\_bai; usa crescer\_regiao; usa semente\_e\_raiz}
\prov{proveniência: wiki \textperiodcentered{} imagem\ensuremath{\leftrightarrow}BAI \textperiodcentered{} inferencia \textperiodcentered{} confiança 1.0 \textperiodcentered{} domínio ponte \textperiodcentered{} história 4 versões no jornal}

%CRISTAL {"arestas": [{"alvo": "crispr", "confianca": 1.0, "epistemico": "fato", "origem": "wiki", "rel": "explica"}, {"alvo": "dna_cadeia_base4", "confianca": 1.0, "epistemico": "fato", "origem": "wiki", "rel": "usa"}, {"alvo": "computador_mineral", "confianca": 1.0, "epistemico": "fato", "origem": "wiki", "rel": "relaciona"}], "confianca": 1.0, "contraexemplos": [], "descricao": "O CRISPR corta e reescreve o DNA num ponto exato guiado por um RNA — find-and-replace no genoma, a ferramenta de EDIÇÃO da vida. É a engenharia genética como programação: editar o código-fonte do organismo, o mesmo gesto de editar o programa do computador mineral.", "epistemico": "inferencia", "exemplos": [], "id": "crispr_edicao", "memoria": "semantica", "meta": {"dominio": "biotecnologia", "fonte": "biotecnologia e engenharia genética"}, "origem": "wiki", "palavras_chave": [], "sinonimos": [], "tipo": "conceito", "titulo": "CRISPR = a Edição do Genoma (find-and-replace)"}
\section{CRISPR = a Edição do Genoma (find-and-replace)}
O CRISPR corta e reescreve o DNA num ponto exato guiado por um RNA — find-and-replace no genoma, a ferramenta de EDIÇÃO da vida. É a engenharia genética como programação: editar o código-fonte do organismo, o mesmo gesto de editar o programa do computador mineral.
\prov{relações: explica crispr; usa dna\_cadeia\_base4; relaciona computador\_mineral}
\prov{proveniência: wiki \textperiodcentered{} biotecnologia e engenharia genética \textperiodcentered{} inferencia \textperiodcentered{} confiança 1.0 \textperiodcentered{} domínio biotecnologia \textperiodcentered{} história 4 versões no jornal}

%CRISTAL {"arestas": [{"alvo": "transicoes_reta_bai_newton", "confianca": 1.0, "epistemico": "fato", "origem": "wiki", "rel": "parte_de"}, {"alvo": "numeros_de_pisot", "confianca": 1.0, "epistemico": "fato", "origem": "wiki", "rel": "usa"}, {"alvo": "colapso", "confianca": 1.0, "epistemico": "fato", "origem": "wiki", "rel": "relaciona"}], "confianca": 1.0, "contraexemplos": [], "descricao": "O CRISTAL (λ<0) é a contração: na reta mineral é o número de PISOT (as potências se aproximam de inteiros); no Newton é a CONVERGÊNCIA (a semente cai na raiz); no BAI é a AUTORIDADE RESOLVIDA (o colapso fecha numa decisão). Contrair, convergir e resolver são a mesma coisa vista de três lados.", "epistemico": "fato", "exemplos": [], "id": "cristal_e_a_convergencia", "memoria": "semantica", "meta": {"dominio": "ponte", "fonte": "transições órbita"}, "origem": "wiki", "palavras_chave": [], "sinonimos": [], "tipo": "conceito", "titulo": "O Cristal é a Convergência (pisot ↔ convergência ↔ resolvido)"}
\section{O Cristal é a Convergência (pisot \ensuremath{\leftrightarrow} convergência \ensuremath{\leftrightarrow} resolvido)}
O CRISTAL (\ensuremath{\lambda}<0) é a contração: na reta mineral é o número de PISOT (as potências se aproximam de inteiros); no Newton é a CONVERGÊNCIA (a semente cai na raiz); no BAI é a AUTORIDADE RESOLVIDA (o colapso fecha numa decisão). Contrair, convergir e resolver são a mesma coisa vista de três lados.
\prov{relações: parte\_de transicoes\_reta\_bai\_newton; usa numeros\_de\_pisot; relaciona colapso}
\prov{proveniência: wiki \textperiodcentered{} transições órbita \textperiodcentered{} fato \textperiodcentered{} confiança 1.0 \textperiodcentered{} domínio ponte \textperiodcentered{} história 4 versões no jornal}

%CRISTAL {"arestas": [{"alvo": "fisica", "confianca": 1.0, "epistemico": "fato", "origem": "wiki", "rel": "parte_de"}], "confianca": 1.0, "contraexemplos": [], "descricao": "A ciência da estrutura ordenada da matéria sólida — como átomos se arranjam em redes periódicas e como isso determina as propriedades.", "epistemico": "fato", "exemplos": [], "id": "cristalografia", "memoria": "semantica", "meta": {"dominio": "engenharia", "fonte": ""}, "origem": "wiki", "palavras_chave": [], "sinonimos": [], "tipo": "conceito", "titulo": "Cristalografia"}
\section{Cristalografia}
A ciência da estrutura ordenada da matéria sólida — como átomos se arranjam em redes periódicas e como isso determina as propriedades.
\prov{relações: parte\_de fisica}
\prov{proveniência: wiki \textperiodcentered{} fato \textperiodcentered{} confiança 1.0 \textperiodcentered{} domínio engenharia \textperiodcentered{} história 2 versões no jornal}

%CRISTAL {"arestas": [{"alvo": "funcoes_de_rede_5gc", "confianca": 1.0, "epistemico": "fato", "origem": "wiki", "rel": "usa"}, {"alvo": "computacao_em_nuvem", "confianca": 1.0, "epistemico": "fato", "origem": "wiki", "rel": "relaciona"}], "confianca": 1.0, "contraexemplos": [], "descricao": "Separar o plano de CONTROLE (sinalização, decisões — SMF) do plano de USUÁRIO (o tráfego em si — UPF), para escalá-los independentemente e pôr o UPF na borda, perto do usuário. Base da baixa latência do 5G.", "epistemico": "inferencia", "exemplos": [], "id": "cups_separacao_planos", "memoria": "semantica", "meta": {"dominio": "engenharia", "fonte": ""}, "origem": "wiki", "palavras_chave": [], "sinonimos": [], "tipo": "conceito", "titulo": "Separação dos Planos (CUPS)"}
\section{Separação dos Planos (CUPS)}
Separar o plano de CONTROLE (sinalização, decisões — SMF) do plano de USUÁRIO (o tráfego em si — UPF), para escalá-los independentemente e pôr o UPF na borda, perto do usuário. Base da baixa latência do 5G.
\prov{relações: usa funcoes\_de\_rede\_5gc; relaciona computacao\_em\_nuvem}
\prov{proveniência: wiki \textperiodcentered{} inferencia \textperiodcentered{} confiança 1.0 \textperiodcentered{} domínio engenharia \textperiodcentered{} história 3 versões no jornal}

%CRISTAL {"arestas": [{"alvo": "sistema_eletrico_potencia", "confianca": 1.0, "epistemico": "fato", "origem": "wiki", "rel": "parte_de"}, {"alvo": "protecao_de_sistemas", "confianca": 1.0, "epistemico": "fato", "origem": "wiki", "rel": "usa"}, {"alvo": "eng_sistemas_energia", "confianca": 1.0, "epistemico": "fato", "origem": "curriculo", "rel": "parte_de"}], "confianca": 1.0, "contraexemplos": [], "descricao": "Calcular as correntes de FALTA (curto) para dimensionar equipamentos e ajustar a proteção. O pior caso que todo componente da rede precisa suportar por instantes sem se destruir.", "epistemico": "fato", "exemplos": [], "id": "curto_circuito", "memoria": "semantica", "meta": {"dominio": "engenharia", "fonte": ""}, "origem": "wiki", "palavras_chave": [], "sinonimos": [], "tipo": "conceito", "titulo": "Análise de Curto-Circuito"}
\section{Análise de Curto-Circuito}
Calcular as correntes de FALTA (curto) para dimensionar equipamentos e ajustar a proteção. O pior caso que todo componente da rede precisa suportar por instantes sem se destruir.
\prov{relações: parte\_de sistema\_eletrico\_potencia; usa protecao\_de\_sistemas; parte\_de eng\_sistemas\_energia}
\prov{proveniência: wiki \textperiodcentered{} fato \textperiodcentered{} confiança 1.0 \textperiodcentered{} domínio engenharia \textperiodcentered{} história 4 versões no jornal}

%CRISTAL {"arestas": [{"alvo": "rede_eletrica", "confianca": 1.0, "epistemico": "fato", "origem": "wiki", "rel": "explica"}, {"alvo": "orbita_de_partida", "confianca": 1.0, "epistemico": "fato", "origem": "wiki", "rel": "relaciona"}], "confianca": 1.0, "contraexemplos": [], "descricao": "A demanda elétrica ao longo das 24h: vale de madrugada, pico ao anoitecer. É uma ÓRBITA diária; a razão pico/média (PAR) mede quão aguda ela é. Aplainar a curva (PAR→1) é o objetivo — pico agudo custa usinas de reserva caras e arrisca o blackout.", "epistemico": "inferencia", "exemplos": [], "id": "curva_de_carga", "memoria": "semantica", "meta": {"dominio": "engenharia", "fonte": "rede de carregamento e smart cities"}, "origem": "wiki", "palavras_chave": [], "sinonimos": [], "tipo": "conceito", "titulo": "Curva de Carga (a órbita do dia)"}
\section{Curva de Carga (a órbita do dia)}
A demanda elétrica ao longo das 24h: vale de madrugada, pico ao anoitecer. É uma ÓRBITA diária; a razão pico/média (PAR) mede quão aguda ela é. Aplainar a curva (PAR\ensuremath{\to}1) é o objetivo — pico agudo custa usinas de reserva caras e arrisca o blackout.
\prov{relações: explica rede\_eletrica; relaciona orbita\_de\_partida}
\prov{proveniência: wiki \textperiodcentered{} rede de carregamento e smart cities \textperiodcentered{} inferencia \textperiodcentered{} confiança 1.0 \textperiodcentered{} domínio engenharia \textperiodcentered{} história 6 versões no jornal}

%CRISTAL {"arestas": [{"alvo": "curva_de_carga", "confianca": 1.0, "epistemico": "fato", "origem": "wiki", "rel": "eh_um"}, {"alvo": "usina_solar", "confianca": 1.0, "epistemico": "fato", "origem": "wiki", "rel": "usa"}], "confianca": 1.0, "contraexemplos": [], "descricao": "A carga LÍQUIDA (demanda menos solar) quando há muita geração solar: barriga funda ao meio-dia (sol forte) e rampa íngreme ao anoitecer (sol some, todos chegam em casa). O desafio moderno da rede — e onde carregar a frota ao meio-dia enche a barriga.", "epistemico": "fato", "exemplos": [], "id": "curva_do_pato", "memoria": "semantica", "meta": {"dominio": "engenharia", "fonte": "rede de carregamento e smart cities"}, "origem": "wiki", "palavras_chave": [], "sinonimos": [], "tipo": "conceito", "titulo": "Curva do Pato (duck curve)"}
\section{Curva do Pato (duck curve)}
A carga LÍQUIDA (demanda menos solar) quando há muita geração solar: barriga funda ao meio-dia (sol forte) e rampa íngreme ao anoitecer (sol some, todos chegam em casa). O desafio moderno da rede — e onde carregar a frota ao meio-dia enche a barriga.
\prov{relações: eh\_um curva\_de\_carga; usa usina\_solar}
\prov{proveniência: wiki \textperiodcentered{} rede de carregamento e smart cities \textperiodcentered{} fato \textperiodcentered{} confiança 1.0 \textperiodcentered{} domínio engenharia \textperiodcentered{} história 6 versões no jornal}

%CRISTAL {"arestas": [{"alvo": "computacao_em_nuvem", "confianca": 1.0, "epistemico": "fato", "origem": "wiki", "rel": "usa"}, {"alvo": "distribuicao_de_energia", "confianca": 1.0, "epistemico": "fato", "origem": "wiki", "rel": "relaciona"}], "confianca": 1.0, "contraexemplos": [], "descricao": "O galpão de milhares de servidores, com energia redundante, refrigeração e rede de altíssima banda que sustenta a nuvem. É a máquina física da computação em nuvem — e uma das maiores cargas elétricas e térmicas que existem.", "epistemico": "fato", "exemplos": [], "id": "data_center", "memoria": "semantica", "meta": {"dominio": "engenharia", "fonte": "nuvem e data centers"}, "origem": "wiki", "palavras_chave": [], "sinonimos": [], "tipo": "conceito", "titulo": "Data Center"}
\section{Data Center}
O galpão de milhares de servidores, com energia redundante, refrigeração e rede de altíssima banda que sustenta a nuvem. É a máquina física da computação em nuvem — e uma das maiores cargas elétricas e térmicas que existem.
\prov{relações: usa computacao\_em\_nuvem; relaciona distribuicao\_de\_energia}
\prov{proveniência: wiki \textperiodcentered{} nuvem e data centers \textperiodcentered{} fato \textperiodcentered{} confiança 1.0 \textperiodcentered{} domínio engenharia \textperiodcentered{} história 3 versões no jornal}

%CRISTAL {"arestas": [{"alvo": "orbita_de_objeto_real", "confianca": 1.0, "epistemico": "fato", "origem": "wiki", "rel": "relaciona"}, {"alvo": "transformada_de_fourier", "confianca": 1.0, "epistemico": "fato", "origem": "wiki", "rel": "relaciona"}, {"alvo": "corpo_mineral", "confianca": 1.0, "epistemico": "fato", "origem": "wiki", "rel": "usa"}], "confianca": 1.0, "contraexemplos": [], "descricao": "Uma imagem (matriz) decompõe-se em CAMADAS por SVD: I = Σ σ_i·u_i·v_iᵀ. Cada camada é um modo de posto 1, ordenado pelo valor singular σ_i — o espectro. Reconstruir com os k primeiros é a aproximação de posto k. A imagem deixa de ser pixels e vira uma soma de rotações ponderadas, ordenada pela energia.", "epistemico": "fato", "exemplos": [], "id": "decomposicao_espectral_imagem", "memoria": "semantica", "meta": {"dominio": "imagem", "fonte": "decomposição espectral"}, "origem": "wiki", "palavras_chave": [], "sinonimos": [], "tipo": "conceito", "titulo": "Decomposição Espectral da Imagem (SVD)"}
\section{Decomposição Espectral da Imagem (SVD)}
Uma imagem (matriz) decompõe-se em CAMADAS por SVD: I = \ensuremath{\Sigma} \ensuremath{\sigma}\_i\textperiodcentered u\_i\textperiodcentered v\_i\ensuremath{^{T}}. Cada camada é um modo de posto 1, ordenado pelo valor singular \ensuremath{\sigma}\_i — o espectro. Reconstruir com os k primeiros é a aproximação de posto k. A imagem deixa de ser pixels e vira uma soma de rotações ponderadas, ordenada pela energia.
\prov{relações: relaciona orbita\_de\_objeto\_real; relaciona transformada\_de\_fourier; usa corpo\_mineral}
\prov{proveniência: wiki \textperiodcentered{} decomposição espectral \textperiodcentered{} fato \textperiodcentered{} confiança 1.0 \textperiodcentered{} domínio imagem \textperiodcentered{} história 4 versões no jornal}

%CRISTAL {"arestas": [{"alvo": "decomposicao_espectral_imagem", "confianca": 1.0, "epistemico": "fato", "origem": "wiki", "rel": "parte_de"}, {"alvo": "bai", "confianca": 1.0, "epistemico": "fato", "origem": "wiki", "rel": "usa"}, {"alvo": "colapso", "confianca": 1.0, "epistemico": "fato", "origem": "wiki", "rel": "eh_um"}, {"alvo": "imagem_bai", "confianca": 1.0, "epistemico": "fato", "origem": "wiki", "rel": "relaciona"}, {"alvo": "xadrez_bai", "confianca": 1.0, "epistemico": "fato", "origem": "wiki", "rel": "relaciona"}], "confianca": 1.0, "contraexemplos": [], "descricao": "Decompor e reconstruir uma imagem É a MESMA árvore-de-decisão-com-orçamento do xadrez, do BAI e da segmentação (pivô 'decisao', a 4ª leitura). O dicionário traduz 1:1: valor singular σ_i ↔ Φ (medida), rank k ↔ orçamento δ, truncar a cauda ↔ deny-overrides (poda), camada dominante ↔ raiz, reconstrução ↔ colapso Ψ, erro de truncamento ↔ fail-closed. compor() dá decomposição→BAI (e →xadrez, →imagem) de graça. Uma obra é uma decisão espectral.", "epistemico": "fato", "exemplos": [], "id": "decomposicao_isomorfa_ao_bai", "memoria": "semantica", "meta": {"dominio": "imagem", "fonte": "decomposição espectral"}, "origem": "wiki", "palavras_chave": [], "sinonimos": [], "tipo": "conceito", "titulo": "A Decomposição é Isomorfa ao BAI"}
\section{A Decomposição é Isomorfa ao BAI}
Decompor e reconstruir uma imagem É a MESMA árvore-de-decisão-com-orçamento do xadrez, do BAI e da segmentação (pivô 'decisao', a 4ª leitura). O dicionário traduz 1:1: valor singular \ensuremath{\sigma}\_i \ensuremath{\leftrightarrow} \ensuremath{\Phi} (medida), rank k \ensuremath{\leftrightarrow} orçamento \ensuremath{\delta}, truncar a cauda \ensuremath{\leftrightarrow} deny-overrides (poda), camada dominante \ensuremath{\leftrightarrow} raiz, reconstrução \ensuremath{\leftrightarrow} colapso \ensuremath{\Psi}, erro de truncamento \ensuremath{\leftrightarrow} fail-closed. compor() dá decomposição\ensuremath{\to}BAI (e \ensuremath{\to}xadrez, \ensuremath{\to}imagem) de graça. Uma obra é uma decisão espectral.
\prov{relações: parte\_de decomposicao\_espectral\_imagem; usa bai; eh\_um colapso; relaciona imagem\_bai; relaciona xadrez\_bai}
\prov{proveniência: wiki \textperiodcentered{} decomposição espectral \textperiodcentered{} fato \textperiodcentered{} confiança 1.0 \textperiodcentered{} domínio imagem \textperiodcentered{} história 6 versões no jornal}

%CRISTAL {"arestas": [{"alvo": "rede_cristalina", "confianca": 1.0, "epistemico": "fato", "origem": "wiki", "rel": "usa"}], "confianca": 1.0, "contraexemplos": [], "descricao": "Imperfeições na rede (vacâncias, discordâncias, impurezas) — longe de estragar, é o que dá resistência, cor e condutividade controlada.", "epistemico": "fato", "exemplos": [], "id": "defeito_cristalino", "memoria": "semantica", "meta": {"dominio": "engenharia", "fonte": ""}, "origem": "wiki", "palavras_chave": [], "sinonimos": [], "tipo": "conceito", "titulo": "Defeito Cristalino"}
\section{Defeito Cristalino}
Imperfeições na rede (vacâncias, discordâncias, impurezas) — longe de estragar, é o que dá resistência, cor e condutividade controlada.
\prov{relações: usa rede\_cristalina}
\prov{proveniência: wiki \textperiodcentered{} fato \textperiodcentered{} confiança 1.0 \textperiodcentered{} domínio engenharia \textperiodcentered{} história 2 versões no jornal}

%CRISTAL {"arestas": [{"alvo": "qam", "confianca": 1.0, "epistemico": "fato", "origem": "wiki", "rel": "usa"}, {"alvo": "colapso", "confianca": 1.0, "epistemico": "fato", "origem": "wiki", "rel": "eh_um"}, {"alvo": "sensor_e_colapso", "confianca": 1.0, "epistemico": "fato", "origem": "wiki", "rel": "relaciona"}, {"alvo": "corpo_mineral", "confianca": 1.0, "epistemico": "fato", "origem": "wiki", "rel": "usa"}], "confianca": 1.0, "contraexemplos": [], "descricao": "O receptor recebe o ponto afogado em ruído e o DESABA no símbolo de constelação mais próximo (snap ao mais próximo) — o mesmo colapso do sensor e do BAI. Sobrevive se o ruído < margem (d_min/2); 256-QAM tem margem minúscula, exige canal limpo. Verif.: 256-QAM tolera ruído 0.05, não 0.15.", "epistemico": "inferencia", "exemplos": [], "id": "demodular_colapso", "memoria": "semantica", "meta": {"dominio": "engenharia", "fonte": "conectividade e 5G", "tolera_005": true, "tolera_015": false}, "origem": "wiki", "palavras_chave": [], "sinonimos": [], "tipo": "conceito", "titulo": "Demodular é um Colapso"}
\section{Demodular é um Colapso}
O receptor recebe o ponto afogado em ruído e o DESABA no símbolo de constelação mais próximo (snap ao mais próximo) — o mesmo colapso do sensor e do BAI. Sobrevive se o ruído < margem (d\_min/2); 256-QAM tem margem minúscula, exige canal limpo. Verif.: 256-QAM tolera ruído 0.05, não 0.15.
\prov{relações: usa qam; eh\_um colapso; relaciona sensor\_e\_colapso; usa corpo\_mineral}
\prov{proveniência: wiki \textperiodcentered{} conectividade e 5G \textperiodcentered{} inferencia \textperiodcentered{} confiança 1.0 \textperiodcentered{} domínio engenharia \textperiodcentered{} história 5 versões no jornal}

%CRISTAL {"arestas": [{"alvo": "pipeline_deploy_broca", "confianca": 1.0, "epistemico": "fato", "origem": "wiki", "rel": "parte_de"}, {"alvo": "deploy_por_banda", "confianca": 1.0, "epistemico": "fato", "origem": "wiki", "rel": "usa"}, {"alvo": "isolamento_base_cliente", "confianca": 1.0, "epistemico": "fato", "origem": "wiki", "rel": "usa"}, {"alvo": "face_turno_fail_closed", "confianca": 1.0, "epistemico": "fato", "origem": "wiki", "rel": "usa"}], "confianca": 1.0, "contraexemplos": [], "descricao": "goldchain/deploy_agente.py sintoniza a banda, confere o sha256 e escreve de forma ATÔMICA (tmp + os.replace) sob uma ALLOWLIST: só código (.py/.js/.md/.sh/…), JAMAIS dados/ .env .git nem subir na árvore ('..'). Uma banda comprometida, no pior caso, reescreve código versionado — fail-closed.", "epistemico": "fato", "exemplos": [], "id": "deploy_agente_allowlist", "memoria": "semantica", "meta": {"dominio": "arquitetura", "fonte": "goldchain/deploy_agente.py"}, "origem": "wiki", "palavras_chave": [], "sinonimos": [], "tipo": "conceito", "titulo": "Agente de Deploy (allowlist + checksum + atômico)"}
\section{Agente de Deploy (allowlist + checksum + atômico)}
goldchain/deploy\_agente.py sintoniza a banda, confere o sha256 e escreve de forma ATÔMICA (tmp + os.replace) sob uma ALLOWLIST: só código (.py/.js/.md/.sh/…), JAMAIS dados/ .env .git nem subir na árvore ('..'). Uma banda comprometida, no pior caso, reescreve código versionado — fail-closed.
\prov{relações: parte\_de pipeline\_deploy\_broca; usa deploy\_por\_banda; usa isolamento\_base\_cliente; usa face\_turno\_fail\_closed}
\prov{proveniência: wiki \textperiodcentered{} goldchain/deploy\_agente.py \textperiodcentered{} fato \textperiodcentered{} confiança 1.0 \textperiodcentered{} domínio arquitetura \textperiodcentered{} história 5 versões no jornal}

%CRISTAL {"arestas": [{"alvo": "pipeline_deploy_broca", "confianca": 1.0, "epistemico": "fato", "origem": "wiki", "rel": "parte_de"}, {"alvo": "organismo_gex44", "confianca": 1.0, "epistemico": "fato", "origem": "wiki", "rel": "usa"}, {"alvo": "gex_minera_ouro_real", "confianca": 1.0, "epistemico": "fato", "origem": "wiki", "rel": "relaciona"}], "confianca": 1.0, "contraexemplos": [], "descricao": "O bootstrap sincroniza TODO o código por rsync ADITIVO (sem --delete), excluindo dados/ .env *.so e binários — nunca apaga o ledger vivo nem toca o minerador. É como o deployer chega ao gex44 na primeira vez (antes de o deploy-por-banda poder se auto-atualizar).", "epistemico": "fato", "exemplos": [], "id": "deploy_bootstrap_seguro", "memoria": "semantica", "meta": {"dominio": "arquitetura", "fonte": "semilla_deploy_pipeline"}, "origem": "wiki", "palavras_chave": [], "sinonimos": [], "tipo": "conceito", "titulo": "Bootstrap Seguro (sync aditivo)"}
\section{Bootstrap Seguro (sync aditivo)}
O bootstrap sincroniza TODO o código por rsync ADITIVO (sem --delete), excluindo dados/ .env *.so e binários — nunca apaga o ledger vivo nem toca o minerador. É como o deployer chega ao gex44 na primeira vez (antes de o deploy-por-banda poder se auto-atualizar).
\prov{relações: parte\_de pipeline\_deploy\_broca; usa organismo\_gex44; relaciona gex\_minera\_ouro\_real}
\prov{proveniência: wiki \textperiodcentered{} semilla\_deploy\_pipeline \textperiodcentered{} fato \textperiodcentered{} confiança 1.0 \textperiodcentered{} domínio arquitetura \textperiodcentered{} história 4 versões no jornal}

%CRISTAL {"arestas": [{"alvo": "pipeline_deploy_broca", "confianca": 1.0, "epistemico": "fato", "origem": "wiki", "rel": "parte_de"}, {"alvo": "banda_exclusiva_por_tecido", "confianca": 1.0, "epistemico": "fato", "origem": "wiki", "rel": "usa"}, {"alvo": "roadmap_olho_multibanda", "confianca": 1.0, "epistemico": "fato", "origem": "wiki", "rel": "usa"}], "confianca": 1.0, "contraexemplos": [], "descricao": "goldchain/deploy_envia.py publica cada arquivo como evento numa banda de deploy do olho ({path, conteudo base64, sha256}); o agente do outro lado sintoniza a banda e aplica. O deploy usa o mesmo barramento/framing das células e sementes — o pipeline come a própria rede.", "epistemico": "fato", "exemplos": [], "id": "deploy_por_banda", "memoria": "semantica", "meta": {"dominio": "arquitetura", "fonte": "goldchain/deploy_envia.py"}, "origem": "wiki", "palavras_chave": [], "sinonimos": [], "tipo": "conceito", "titulo": "Deploy por Banda (dogfooding)"}
\section{Deploy por Banda (dogfooding)}
goldchain/deploy\_envia.py publica cada arquivo como evento numa banda de deploy do olho (\{path, conteudo base64, sha256\}); o agente do outro lado sintoniza a banda e aplica. O deploy usa o mesmo barramento/framing das células e sementes — o pipeline come a própria rede.
\prov{relações: parte\_de pipeline\_deploy\_broca; usa banda\_exclusiva\_por\_tecido; usa roadmap\_olho\_multibanda}
\prov{proveniência: wiki \textperiodcentered{} goldchain/deploy\_envia.py \textperiodcentered{} fato \textperiodcentered{} confiança 1.0 \textperiodcentered{} domínio arquitetura \textperiodcentered{} história 4 versões no jornal}

%CRISTAL {"arestas": [{"alvo": "pipeline_deploy_broca", "confianca": 1.0, "epistemico": "fato", "origem": "wiki", "rel": "parte_de"}, {"alvo": "organismo_gex44", "confianca": 1.0, "epistemico": "fato", "origem": "wiki", "rel": "usa"}, {"alvo": "sistema_e_lei_nao_dono", "confianca": 1.0, "epistemico": "fato", "origem": "wiki", "rel": "confirma"}], "confianca": 1.0, "contraexemplos": [], "descricao": "O olho de deploy fica em 127.0.0.1 no gex44 — nunca exposto; o acesso é por túnel SSH (o central apresenta, não é correio). A banda filtra o conteúdo; o SSH é a fronteira de rede. Deploy confiável por poll do túnel antes de enviar.", "epistemico": "inferencia", "exemplos": [], "id": "deploy_tunel_rendezvous", "memoria": "semantica", "meta": {"dominio": "arquitetura", "fonte": "semilla_deploy_pipeline"}, "origem": "wiki", "palavras_chave": [], "sinonimos": [], "tipo": "conceito", "titulo": "Túnel-Rendezvous (olho em 127.0.0.1)"}
\section{Túnel-Rendezvous (olho em 127.0.0.1)}
O olho de deploy fica em 127.0.0.1 no gex44 — nunca exposto; o acesso é por túnel SSH (o central apresenta, não é correio). A banda filtra o conteúdo; o SSH é a fronteira de rede. Deploy confiável por poll do túnel antes de enviar.
\prov{relações: parte\_de pipeline\_deploy\_broca; usa organismo\_gex44; confirma sistema\_e\_lei\_nao\_dono}
\prov{proveniência: wiki \textperiodcentered{} semilla\_deploy\_pipeline \textperiodcentered{} inferencia \textperiodcentered{} confiança 1.0 \textperiodcentered{} domínio arquitetura \textperiodcentered{} história 4 versões no jornal}

%CRISTAL {"arestas": [{"alvo": "pipeline_deploy_broca", "confianca": 1.0, "epistemico": "fato", "origem": "wiki", "rel": "parte_de"}, {"alvo": "deploy_por_banda", "confianca": 1.0, "epistemico": "fato", "origem": "wiki", "rel": "confirma"}], "confianca": 1.0, "contraexemplos": [], "descricao": "O verify prova que o remoto é o que foi deployado: o gex44 SÓ calcula os sha256, a comparação é LOCAL (robusta), com guarda contra resposta incompleta e exit≠0 se divergir. Pega o falso-positivo que uma leitura remota silenciosa causaria. A prova honesta de que o código chegou íntegro.", "epistemico": "fato", "exemplos": [], "id": "deploy_verify_checksum", "memoria": "semantica", "meta": {"dominio": "arquitetura", "fonte": "semilla_deploy_pipeline"}, "origem": "wiki", "palavras_chave": [], "sinonimos": [], "tipo": "conceito", "titulo": "Verify por Checksum (a prova do deploy)"}
\section{Verify por Checksum (a prova do deploy)}
O verify prova que o remoto é o que foi deployado: o gex44 SÓ calcula os sha256, a comparação é LOCAL (robusta), com guarda contra resposta incompleta e exit\ensuremath{\ne}0 se divergir. Pega o falso-positivo que uma leitura remota silenciosa causaria. A prova honesta de que o código chegou íntegro.
\prov{relações: parte\_de pipeline\_deploy\_broca; confirma deploy\_por\_banda}
\prov{proveniência: wiki \textperiodcentered{} semilla\_deploy\_pipeline \textperiodcentered{} fato \textperiodcentered{} confiança 1.0 \textperiodcentered{} domínio arquitetura \textperiodcentered{} história 3 versões no jornal}

%CRISTAL {"arestas": [{"alvo": "orbita_de_objeto_real", "confianca": 1.0, "epistemico": "fato", "origem": "wiki", "rel": "parte_de"}, {"alvo": "transformada_de_fourier", "confianca": 1.0, "epistemico": "fato", "origem": "wiki", "rel": "usa"}, {"alvo": "quasicristal_de_rotacoes", "confianca": 1.0, "epistemico": "fato", "origem": "wiki", "rel": "relaciona"}, {"alvo": "numero_de_salem", "confianca": 1.0, "epistemico": "fato", "origem": "wiki", "rel": "relaciona"}, {"alvo": "imagem_em_rotacoes_fourier", "confianca": 1.0, "epistemico": "fato", "origem": "wiki", "rel": "usa"}], "confianca": 1.0, "contraexemplos": [], "descricao": "Os descritores de Fourier |FFT(r(θ))| são os harmônicos que compõem o contorno — as ROTAÇÕES puras (e^{ikθ}, ρ=1 = Salem) que somadas dão a órbita do objeto. O harmônico DOMINANTE é o fundamental (o cristal da forma): 2 para um oval (boneca russa), a simetria n para um polígono. Invariante a rotação (magnitude) e escala (r normalizado). A assinatura mineral do objeto.", "epistemico": "fato", "exemplos": [], "id": "descritores_de_fourier_objeto", "memoria": "semantica", "meta": {"dominio": "imagem", "fonte": "órbitas de objetos"}, "origem": "wiki", "palavras_chave": [], "sinonimos": [], "tipo": "conceito", "titulo": "Descritores de Fourier (as rotações do objeto)"}
\section{Descritores de Fourier (as rotações do objeto)}
Os descritores de Fourier |FFT(r(\ensuremath{\theta}))| são os harmônicos que compõem o contorno — as ROTAÇÕES puras (e\^{}\{ik\ensuremath{\theta}\}, \ensuremath{\rho}=1 = Salem) que somadas dão a órbita do objeto. O harmônico DOMINANTE é o fundamental (o cristal da forma): 2 para um oval (boneca russa), a simetria n para um polígono. Invariante a rotação (magnitude) e escala (r normalizado). A assinatura mineral do objeto.
\prov{relações: parte\_de orbita\_de\_objeto\_real; usa transformada\_de\_fourier; relaciona quasicristal\_de\_rotacoes; relaciona numero\_de\_salem; usa imagem\_em\_rotacoes\_fourier}
\prov{proveniência: wiki \textperiodcentered{} órbitas de objetos \textperiodcentered{} fato \textperiodcentered{} confiança 1.0 \textperiodcentered{} domínio imagem \textperiodcentered{} história 6 versões no jornal}

%CRISTAL {"arestas": [{"alvo": "protocolo_udp", "confianca": 1.0, "epistemico": "fato", "origem": "wiki", "rel": "usa"}, {"alvo": "eng_telecomunicacoes", "confianca": 1.0, "epistemico": "fato", "origem": "curriculo", "rel": "parte_de"}], "confianca": 1.0, "contraexemplos": [], "descricao": "Configura automaticamente um host que entra na rede: IP, máscara, gateway e DNS por concessão temporária. O recepcionista que dá endereço a quem chega — sem ele, configurar rede seria manual.", "epistemico": "fato", "exemplos": [], "id": "dhcp", "memoria": "semantica", "meta": {"dominio": "engenharia", "fonte": ""}, "origem": "wiki", "palavras_chave": [], "sinonimos": [], "tipo": "conceito", "titulo": "DHCP"}
\section{DHCP}
Configura automaticamente um host que entra na rede: IP, máscara, gateway e DNS por concessão temporária. O recepcionista que dá endereço a quem chega — sem ele, configurar rede seria manual.
\prov{relações: usa protocolo\_udp; parte\_de eng\_telecomunicacoes}
\prov{proveniência: wiki \textperiodcentered{} fato \textperiodcentered{} confiança 1.0 \textperiodcentered{} domínio engenharia \textperiodcentered{} história 3 versões no jornal}

%CRISTAL {"arestas": [{"alvo": "reconhecer_caligrafia_estilo", "confianca": 1.0, "epistemico": "fato", "origem": "wiki", "rel": "usa"}, {"alvo": "colapso_fatorado_glifo", "confianca": 1.0, "epistemico": "fato", "origem": "wiki", "rel": "relaciona"}, {"alvo": "teorema_geral_traducao", "confianca": 1.0, "epistemico": "fato", "origem": "wiki", "rel": "usa"}, {"alvo": "laco_traducao_bai_peirce", "confianca": 1.0, "epistemico": "fato", "origem": "wiki", "rel": "relaciona"}, {"alvo": "imagem_bai", "confianca": 1.0, "epistemico": "fato", "origem": "wiki", "rel": "relaciona"}], "confianca": 1.0, "contraexemplos": [], "descricao": "As seis caligrafias (serif, sans, mono, negrito, itálico, casual) entram na máquina de dicionários: cada fonte ↔ seus valores em 5 eixos de estilo (o pivô 'caligrafia'). compor() traduz uma caligrafia noutra eixo a eixo (serif→sans: traço fino→grosso, com serifa→sem); comparar() mostra a órbita partilhada e onde divergem. É o teorema da bússola aplicado à tipografia — um espaço de estilo, muitas fontes. Derivado de medição.", "epistemico": "fato", "exemplos": [], "id": "dicionario_das_fontes", "memoria": "semantica", "meta": {"dominio": "imagem", "fonte": "dicionário de fontes"}, "origem": "wiki", "palavras_chave": [], "sinonimos": [], "tipo": "conceito", "titulo": "Dicionário das Fontes Caligráficas"}
\section{Dicionário das Fontes Caligráficas}
As seis caligrafias (serif, sans, mono, negrito, itálico, casual) entram na máquina de dicionários: cada fonte \ensuremath{\leftrightarrow} seus valores em 5 eixos de estilo (o pivô 'caligrafia'). compor() traduz uma caligrafia noutra eixo a eixo (serif\ensuremath{\to}sans: traço fino\ensuremath{\to}grosso, com serifa\ensuremath{\to}sem); comparar() mostra a órbita partilhada e onde divergem. É o teorema da bússola aplicado à tipografia — um espaço de estilo, muitas fontes. Derivado de medição.
\prov{relações: usa reconhecer\_caligrafia\_estilo; relaciona colapso\_fatorado\_glifo; usa teorema\_geral\_traducao; relaciona laco\_traducao\_bai\_peirce; relaciona imagem\_bai}
\prov{proveniência: wiki \textperiodcentered{} dicionário de fontes \textperiodcentered{} fato \textperiodcentered{} confiança 1.0 \textperiodcentered{} domínio imagem \textperiodcentered{} história 6 versões no jornal}

%CRISTAL {"arestas": [{"alvo": "classificacao_de_formas", "confianca": 1.0, "epistemico": "fato", "origem": "wiki", "rel": "usa"}, {"alvo": "grupo_de_transformacoes", "confianca": 1.0, "epistemico": "fato", "origem": "wiki", "rel": "usa"}, {"alvo": "dicionario_das_fontes", "confianca": 1.0, "epistemico": "fato", "origem": "wiki", "rel": "relaciona"}], "confianca": 1.0, "contraexemplos": [], "descricao": "As 6 formas ↔ seus valores em eixos invariantes (simetria, convexidade, borda), e as 5 transformações ↔ o que preservam (tamanho, ângulo, orientação) — na mesma máquina de dicionários (pivôs 'forma' e 'transformacao'). compor/comparar dão forma↔forma de graça; pentágono e estrela partilham a simetria (5 dobras) e divergem só na convexidade. Derivado de medição.", "epistemico": "fato", "exemplos": [], "id": "dicionario_formas_transformacoes", "memoria": "semantica", "meta": {"dominio": "imagem", "fonte": "formas e limite"}, "origem": "wiki", "palavras_chave": [], "sinonimos": [], "tipo": "conceito", "titulo": "Dicionário de Formas e Transformações"}
\section{Dicionário de Formas e Transformações}
As 6 formas \ensuremath{\leftrightarrow} seus valores em eixos invariantes (simetria, convexidade, borda), e as 5 transformações \ensuremath{\leftrightarrow} o que preservam (tamanho, ângulo, orientação) — na mesma máquina de dicionários (pivôs 'forma' e 'transformacao'). compor/comparar dão forma\ensuremath{\leftrightarrow}forma de graça; pentágono e estrela partilham a simetria (5 dobras) e divergem só na convexidade. Derivado de medição.
\prov{relações: usa classificacao\_de\_formas; usa grupo\_de\_transformacoes; relaciona dicionario\_das\_fontes}
\prov{proveniência: wiki \textperiodcentered{} formas e limite \textperiodcentered{} fato \textperiodcentered{} confiança 1.0 \textperiodcentered{} domínio imagem \textperiodcentered{} história 8 versões no jornal}

%CRISTAL {"arestas": [{"alvo": "cristalografia", "confianca": 1.0, "epistemico": "fato", "origem": "wiki", "rel": "usa"}, {"alvo": "rede_cristalina", "confianca": 1.0, "epistemico": "fato", "origem": "wiki", "rel": "explica"}], "confianca": 1.0, "contraexemplos": [], "descricao": "Bragg: raios X espalhados por um cristal interferem em padrões que revelam a posição dos átomos — como se lê a estrutura invisível.", "epistemico": "fato", "exemplos": [], "id": "difracao_de_raios_x", "memoria": "semantica", "meta": {"autor": "Bragg", "dominio": "engenharia", "fonte": ""}, "origem": "wiki", "palavras_chave": [], "sinonimos": [], "tipo": "conceito", "titulo": "Difração de Raios X"}
\section{Difração de Raios X}
Bragg: raios X espalhados por um cristal interferem em padrões que revelam a posição dos átomos — como se lê a estrutura invisível.
\prov{relações: usa cristalografia; explica rede\_cristalina}
\prov{proveniência: wiki \textperiodcentered{} fato \textperiodcentered{} confiança 1.0 \textperiodcentered{} domínio engenharia \textperiodcentered{} história 3 versões no jornal}

%CRISTAL {"arestas": [{"alvo": "manufatura", "confianca": 1.0, "epistemico": "fato", "origem": "wiki", "rel": "usa"}, {"alvo": "smart_city", "confianca": 1.0, "epistemico": "fato", "origem": "wiki", "rel": "relaciona"}, {"alvo": "banco_de_dados_grafo", "confianca": 1.0, "epistemico": "fato", "origem": "wiki", "rel": "relaciona"}], "confianca": 1.0, "contraexemplos": [], "descricao": "Um modelo digital vivo que espelha o objeto físico fabricado, atualizado por sensores em tempo real. O grafo espelha o mundo: o gêmeo digital é a smart city aplicada a uma peça — a coisa e sua representação no mesmo laço.", "epistemico": "fato", "exemplos": [], "id": "digital_twin", "memoria": "semantica", "meta": {"dominio": "manufatura", "fonte": "impressão 3D e manufatura"}, "origem": "wiki", "palavras_chave": [], "sinonimos": [], "tipo": "conceito", "titulo": "Gêmeo Digital"}
\section{Gêmeo Digital}
Um modelo digital vivo que espelha o objeto físico fabricado, atualizado por sensores em tempo real. O grafo espelha o mundo: o gêmeo digital é a smart city aplicada a uma peça — a coisa e sua representação no mesmo laço.
\prov{relações: usa manufatura; relaciona smart\_city; relaciona banco\_de\_dados\_grafo}
\prov{proveniência: wiki \textperiodcentered{} impressão 3D e manufatura \textperiodcentered{} fato \textperiodcentered{} confiança 1.0 \textperiodcentered{} domínio manufatura \textperiodcentered{} história 4 versões no jornal}

%CRISTAL {"arestas": [{"alvo": "juncao_pn", "confianca": 1.0, "epistemico": "fato", "origem": "wiki", "rel": "usa"}, {"alvo": "dispositivo_semicondutor", "confianca": 1.0, "epistemico": "fato", "origem": "wiki", "rel": "especializa"}], "confianca": 1.0, "contraexemplos": [], "descricao": "O dispositivo de junção p-n mais simples: corrente num sentido só. Retifica AC em DC e protege circuitos — a válvula de mão única.", "epistemico": "fato", "exemplos": [], "id": "diodo", "memoria": "semantica", "meta": {"dominio": "engenharia", "fonte": ""}, "origem": "wiki", "palavras_chave": [], "sinonimos": [], "tipo": "conceito", "titulo": "Diodo"}
\section{Diodo}
O dispositivo de junção p-n mais simples: corrente num sentido só. Retifica AC em DC e protege circuitos — a válvula de mão única.
\prov{relações: usa juncao\_pn; especializa dispositivo\_semicondutor}
\prov{proveniência: wiki \textperiodcentered{} fato \textperiodcentered{} confiança 1.0 \textperiodcentered{} domínio engenharia \textperiodcentered{} história 3 versões no jornal}

%CRISTAL {"arestas": [{"alvo": "semicondutor", "confianca": 1.0, "epistemico": "fato", "origem": "wiki", "rel": "usa"}, {"alvo": "juncao_pn", "confianca": 1.0, "epistemico": "fato", "origem": "wiki", "rel": "usa"}], "confianca": 1.0, "contraexemplos": [], "descricao": "Um componente que explora a junção p-n para retificar, amplificar, chavear ou emitir/captar luz — a família diodo/transistor/LED/célula.", "epistemico": "fato", "exemplos": [], "id": "dispositivo_semicondutor", "memoria": "semantica", "meta": {"dominio": "engenharia", "fonte": ""}, "origem": "wiki", "palavras_chave": [], "sinonimos": [], "tipo": "conceito", "titulo": "Dispositivo Semicondutor"}
\section{Dispositivo Semicondutor}
Um componente que explora a junção p-n para retificar, amplificar, chavear ou emitir/captar luz — a família diodo/transistor/LED/célula.
\prov{relações: usa semicondutor; usa juncao\_pn}
\prov{proveniência: wiki \textperiodcentered{} fato \textperiodcentered{} confiança 1.0 \textperiodcentered{} domínio engenharia \textperiodcentered{} história 3 versões no jornal}

%CRISTAL {"arestas": [{"alvo": "dicionario_das_fontes", "confianca": 1.0, "epistemico": "fato", "origem": "wiki", "rel": "parte_de"}, {"alvo": "eixos_da_caligrafia", "confianca": 1.0, "epistemico": "fato", "origem": "wiki", "rel": "usa"}, {"alvo": "kmeans_atratores", "confianca": 1.0, "epistemico": "fato", "origem": "wiki", "rel": "relaciona"}], "confianca": 1.0, "contraexemplos": [], "descricao": "Quantos dos 5 eixos diferem entre duas fontes é a distância no espaço de estilo. serif e negrito divergem em só 2 (peso e contraste — o negrito é o serif encorpado, ambos com serifa); serif e sans em 3. É a mesma comparar() da máquina de dicionários (órbita partilhada + divergências), agora medindo o quão longe uma caligrafia está de outra.", "epistemico": "inferencia", "exemplos": [], "id": "distancia_de_estilo", "memoria": "semantica", "meta": {"dominio": "imagem", "fonte": "dicionário de fontes"}, "origem": "wiki", "palavras_chave": [], "sinonimos": [], "tipo": "conceito", "titulo": "Distância de Estilo (a divergência entre caligrafias)"}
\section{Distância de Estilo (a divergência entre caligrafias)}
Quantos dos 5 eixos diferem entre duas fontes é a distância no espaço de estilo. serif e negrito divergem em só 2 (peso e contraste — o negrito é o serif encorpado, ambos com serifa); serif e sans em 3. É a mesma comparar() da máquina de dicionários (órbita partilhada + divergências), agora medindo o quão longe uma caligrafia está de outra.
\prov{relações: parte\_de dicionario\_das\_fontes; usa eixos\_da\_caligrafia; relaciona kmeans\_atratores}
\prov{proveniência: wiki \textperiodcentered{} dicionário de fontes \textperiodcentered{} inferencia \textperiodcentered{} confiança 1.0 \textperiodcentered{} domínio imagem \textperiodcentered{} história 4 versões no jornal}

%CRISTAL {"arestas": [{"alvo": "transmissao_de_energia", "confianca": 1.0, "epistemico": "fato", "origem": "wiki", "rel": "usa"}, {"alvo": "energia_eletrica", "confianca": 1.0, "epistemico": "fato", "origem": "wiki", "rel": "usa"}], "confianca": 1.0, "contraexemplos": [], "descricao": "A última etapa: da subestação, em média/baixa tensão, até a tomada — a capilaridade que entrega a energia ao consumo.", "epistemico": "fato", "exemplos": [], "id": "distribuicao_de_energia", "memoria": "semantica", "meta": {"dominio": "engenharia", "fonte": ""}, "origem": "wiki", "palavras_chave": [], "sinonimos": [], "tipo": "conceito", "titulo": "Distribuição de Energia"}
\section{Distribuição de Energia}
A última etapa: da subestação, em média/baixa tensão, até a tomada — a capilaridade que entrega a energia ao consumo.
\prov{relações: usa transmissao\_de\_energia; usa energia\_eletrica}
\prov{proveniência: wiki \textperiodcentered{} fato \textperiodcentered{} confiança 1.0 \textperiodcentered{} domínio engenharia \textperiodcentered{} história 5 versões no jornal}

%CRISTAL {"arestas": [{"alvo": "dna_dupla_helice", "confianca": 1.0, "epistemico": "fato", "origem": "wiki", "rel": "explica"}, {"alvo": "dna_informacao_shannon", "confianca": 1.0, "epistemico": "fato", "origem": "wiki", "rel": "relaciona"}, {"alvo": "genoma", "confianca": 1.0, "epistemico": "fato", "origem": "wiki", "rel": "usa"}, {"alvo": "cadeias_alem_dos_metais", "confianca": 1.0, "epistemico": "fato", "origem": "wiki", "rel": "eh_um"}, {"alvo": "corpo_mineral", "confianca": 1.0, "epistemico": "fato", "origem": "wiki", "rel": "usa"}], "confianca": 1.0, "contraexemplos": [], "descricao": "O DNA é um código de 4 letras (A,C,G,T) numa cadeia linear — literalmente um número em BASE 4. O genoma é uma cadeia/numeração, a mesma estrutura das cadeias além dos metais e da transformada mineral. Codificar e decodificar a sequência é exato e reversível. Verif.: ATGGCAT ↔ 3731.", "epistemico": "inferencia", "exemplos": [], "id": "dna_cadeia_base4", "memoria": "semantica", "meta": {"dominio": "biotecnologia", "exemplo_num_base4": 3731, "exemplo_seq": "ATGGCAT", "fonte": "biotecnologia e engenharia genética"}, "origem": "wiki", "palavras_chave": [], "sinonimos": [], "tipo": "conceito", "titulo": "DNA = Informação numa Cadeia (base 4)"}
\section{DNA = Informação numa Cadeia (base 4)}
O DNA é um código de 4 letras (A,C,G,T) numa cadeia linear — literalmente um número em BASE 4. O genoma é uma cadeia/numeração, a mesma estrutura das cadeias além dos metais e da transformada mineral. Codificar e decodificar a sequência é exato e reversível. Verif.: ATGGCAT \ensuremath{\leftrightarrow} 3731.
\prov{relações: explica dna\_dupla\_helice; relaciona dna\_informacao\_shannon; usa genoma; eh\_um cadeias\_alem\_dos\_metais; usa corpo\_mineral}
\prov{proveniência: wiki \textperiodcentered{} biotecnologia e engenharia genética \textperiodcentered{} inferencia \textperiodcentered{} confiança 1.0 \textperiodcentered{} domínio biotecnologia \textperiodcentered{} história 6 versões no jornal}

%CRISTAL {"arestas": [{"alvo": "codigo_genetico", "confianca": 1.0, "epistemico": "fato", "origem": "wiki", "rel": "explica"}, {"alvo": "ribossomo_koch", "confianca": 1.0, "epistemico": "fato", "origem": "wiki", "rel": "usa"}, {"alvo": "impressao_decodifica", "confianca": 1.0, "epistemico": "fato", "origem": "wiki", "rel": "relaciona"}, {"alvo": "codificacao_decodificacao", "confianca": 1.0, "epistemico": "fato", "origem": "wiki", "rel": "usa"}, {"alvo": "sensor_e_colapso", "confianca": 1.0, "epistemico": "fato", "origem": "wiki", "rel": "relaciona"}], "confianca": 1.0, "contraexemplos": [], "descricao": "DNA (memória) → mRNA (mensagem) → proteína (função): 3 bases (um códon) decodificam 1 aminoácido. O ribossomo LÊ o mRNA (o G-code) e IMPRIME a proteína resíduo a resíduo — a impressora 3D da vida, o mesmo decode da manufatura. O gene é o arquivo; a proteína é a peça. Verif.: AUGGCACGUUUUUAA → MARF.", "epistemico": "inferencia", "exemplos": [], "id": "dogma_central_decodifica", "memoria": "semantica", "meta": {"dominio": "biotecnologia", "exemplo_proteina": "MARF", "fonte": "biotecnologia e engenharia genética"}, "origem": "wiki", "palavras_chave": [], "sinonimos": [], "tipo": "conceito", "titulo": "Dogma Central = Decodificar (o ribossomo imprime)"}
\section{Dogma Central = Decodificar (o ribossomo imprime)}
DNA (memória) \ensuremath{\to} mRNA (mensagem) \ensuremath{\to} proteína (função): 3 bases (um códon) decodificam 1 aminoácido. O ribossomo LÊ o mRNA (o G-code) e IMPRIME a proteína resíduo a resíduo — a impressora 3D da vida, o mesmo decode da manufatura. O gene é o arquivo; a proteína é a peça. Verif.: AUGGCACGUUUUUAA \ensuremath{\to} MARF.
\prov{relações: explica codigo\_genetico; usa ribossomo\_koch; relaciona impressao\_decodifica; usa codificacao\_decodificacao; relaciona sensor\_e\_colapso}
\prov{proveniência: wiki \textperiodcentered{} biotecnologia e engenharia genética \textperiodcentered{} inferencia \textperiodcentered{} confiança 1.0 \textperiodcentered{} domínio biotecnologia \textperiodcentered{} história 6 versões no jornal}

%CRISTAL {"arestas": [{"alvo": "semicondutor", "confianca": 1.0, "epistemico": "fato", "origem": "wiki", "rel": "usa"}, {"alvo": "portador_de_carga", "confianca": 1.0, "epistemico": "fato", "origem": "wiki", "rel": "explica"}], "confianca": 1.0, "contraexemplos": [], "descricao": "Introduzir impurezas de propósito no cristal para criar portadores: tipo-N (elétrons) ou tipo-P (lacunas) — programar a condução do silício.", "epistemico": "fato", "exemplos": [], "id": "dopagem", "memoria": "semantica", "meta": {"dominio": "engenharia", "fonte": ""}, "origem": "wiki", "palavras_chave": [], "sinonimos": [], "tipo": "conceito", "titulo": "Dopagem"}
\section{Dopagem}
Introduzir impurezas de propósito no cristal para criar portadores: tipo-N (elétrons) ou tipo-P (lacunas) — programar a condução do silício.
\prov{relações: usa semicondutor; explica portador\_de\_carga}
\prov{proveniência: wiki \textperiodcentered{} fato \textperiodcentered{} confiança 1.0 \textperiodcentered{} domínio engenharia \textperiodcentered{} história 3 versões no jornal}

%CRISTAL {"arestas": [{"alvo": "xadrez_bai", "confianca": 1.0, "epistemico": "fato", "origem": "wiki", "rel": "especializa"}, {"alvo": "motor_de_xadrez", "confianca": 1.0, "epistemico": "fato", "origem": "wiki", "rel": "usa"}, {"alvo": "minimax_e_colapso", "confianca": 1.0, "epistemico": "fato", "origem": "wiki", "rel": "usa"}], "confianca": 1.0, "contraexemplos": [], "descricao": "A linguagem em que a Tiffany joga xadrez: cada lance É um colapso do BAI, narrado no vocabulário da autorização (colapso Ψ, orçamento δ, deny-overrides, raiz). Motor real em linguagem/jogo_xadrez.py — movimento legal, xeque-mate, e o mate-em-1 encontrado pelo minimax. Não é metáfora: é o mesmo motor.", "epistemico": "fato", "exemplos": [], "id": "dsl_bai_xadrez", "memoria": "semantica", "meta": {"dominio": "ponte", "fonte": "xadrez↔BAI"}, "origem": "wiki", "palavras_chave": [], "sinonimos": [], "tipo": "conceito", "titulo": "DSL-BAI-XADREZ (o jogo no sandbox da Tiffany)"}
\section{DSL-BAI-XADREZ (o jogo no sandbox da Tiffany)}
A linguagem em que a Tiffany joga xadrez: cada lance É um colapso do BAI, narrado no vocabulário da autorização (colapso \ensuremath{\Psi}, orçamento \ensuremath{\delta}, deny-overrides, raiz). Motor real em linguagem/jogo\_xadrez.py — movimento legal, xeque-mate, e o mate-em-1 encontrado pelo minimax. Não é metáfora: é o mesmo motor.
\prov{relações: especializa xadrez\_bai; usa motor\_de\_xadrez; usa minimax\_e\_colapso}
\prov{proveniência: wiki \textperiodcentered{} xadrez\ensuremath{\leftrightarrow}BAI \textperiodcentered{} fato \textperiodcentered{} confiança 1.0 \textperiodcentered{} domínio ponte \textperiodcentered{} história 4 versões no jornal}

%CRISTAL {"arestas": [{"alvo": "dtc_schema_ut", "confianca": 1.0, "epistemico": "fato", "origem": "wiki", "rel": "usa"}, {"alvo": "olho_de_cock_web", "confianca": 1.0, "epistemico": "fato", "origem": "wiki", "rel": "usa"}], "confianca": 1.0, "contraexemplos": [], "descricao": "Módulo JS que desenha o cockpit inteiro a partir do dtc — não é um diagrama isolado, é o pintor de TODOS os painéis (U_t panel, strip, scheduler, cofre, dinâmica, cards, contabilidade, cúspides, timeline, visual/CSS). Tem render (completo) e renderPartial (leve, entre pulsos iguais).", "epistemico": "fato", "exemplos": [], "id": "dtc_renderer", "memoria": "semantica", "meta": {"autor": "broca-so", "dominio": "cockpit", "fonte": "goldchain/cockpit/dtc_renderer.js"}, "origem": "wiki", "palavras_chave": [], "sinonimos": [], "tipo": "conceito", "titulo": "O DtcRenderer (o pintor único do cockpit)"}
\section{O DtcRenderer (o pintor único do cockpit)}
Módulo JS que desenha o cockpit inteiro a partir do dtc — não é um diagrama isolado, é o pintor de TODOS os painéis (U\_t panel, strip, scheduler, cofre, dinâmica, cards, contabilidade, cúspides, timeline, visual/CSS). Tem render (completo) e renderPartial (leve, entre pulsos iguais).
\prov{relações: usa dtc\_schema\_ut; usa olho\_de\_cock\_web}
\prov{proveniência: wiki \textperiodcentered{} goldchain/cockpit/dtc\_renderer.js \textperiodcentered{} fato \textperiodcentered{} confiança 1.0 \textperiodcentered{} domínio cockpit \textperiodcentered{} história 3 versões no jornal}

%CRISTAL {"arestas": [{"alvo": "olho_de_cock_web", "confianca": 1.0, "epistemico": "fato", "origem": "wiki", "rel": "usa"}, {"alvo": "cotacao_parabola", "confianca": 1.0, "epistemico": "fato", "origem": "wiki", "rel": "usa"}], "confianca": 1.0, "contraexemplos": [], "descricao": "'DTC' é um schema de dados que reduz todo o estado do sistema a um vetor genérico U_t: x, v, a, E, φ, κ, H, Δ, evento, origem, destino + blocos contabilidade/pulso/dourada/visual. O acrónimo NÃO é soletrado no código — é o estado genérico U_t que qualquer renderer consome (H=a+c²).", "epistemico": "fato", "exemplos": [], "id": "dtc_schema_ut", "memoria": "semantica", "meta": {"autor": "broca-so", "dominio": "cockpit", "fonte": "goldchain/modulador_vetorial.py; cockpit/modulador_vetorial.js"}, "origem": "wiki", "palavras_chave": [], "sinonimos": [], "tipo": "conceito", "titulo": "DTC (dtc/1) — o Contrato de Estado U_t"}
\section{DTC (dtc/1) — o Contrato de Estado U\_t}
'DTC' é um schema de dados que reduz todo o estado do sistema a um vetor genérico U\_t: x, v, a, E, \ensuremath{\varphi}, \ensuremath{\kappa}, H, \ensuremath{\Delta}, evento, origem, destino + blocos contabilidade/pulso/dourada/visual. O acrónimo NÃO é soletrado no código — é o estado genérico U\_t que qualquer renderer consome (H=a+c\ensuremath{^{2}}).
\prov{relações: usa olho\_de\_cock\_web; usa cotacao\_parabola}
\prov{proveniência: wiki \textperiodcentered{} goldchain/modulador\_vetorial.py; cockpit/modulador\_vetorial.js \textperiodcentered{} fato \textperiodcentered{} confiança 1.0 \textperiodcentered{} domínio cockpit \textperiodcentered{} história 3 versões no jornal}

%CRISTAL {"arestas": [{"alvo": "redes_5g_nr", "confianca": 1.0, "epistemico": "fato", "origem": "wiki", "rel": "parte_de"}], "confianca": 1.0, "contraexemplos": [], "descricao": "O celular agrega simultaneamente duas estações (ex.: LTE + NR no NSA, ou dois gNBs) — soma bandas e suaviza a mobilidade. O PDCP é quem divide e recombina o fluxo entre as pernas.", "epistemico": "inferencia", "exemplos": [], "id": "dupla_conectividade_nr", "memoria": "semantica", "meta": {"dominio": "engenharia", "fonte": ""}, "origem": "wiki", "palavras_chave": [], "sinonimos": [], "tipo": "conceito", "titulo": "Dupla Conectividade"}
\section{Dupla Conectividade}
O celular agrega simultaneamente duas estações (ex.: LTE + NR no NSA, ou dois gNBs) — soma bandas e suaviza a mobilidade. O PDCP é quem divide e recombina o fluxo entre as pernas.
\prov{relações: parte\_de redes\_5g\_nr}
\prov{proveniência: wiki \textperiodcentered{} inferencia \textperiodcentered{} confiança 1.0 \textperiodcentered{} domínio engenharia \textperiodcentered{} história 2 versões no jornal}

%CRISTAL {"arestas": [{"alvo": "internet_das_coisas", "confianca": 1.0, "epistemico": "fato", "origem": "wiki", "rel": "usa"}, {"alvo": "sensor", "confianca": 1.0, "epistemico": "fato", "origem": "wiki", "rel": "relaciona"}], "confianca": 1.0, "contraexemplos": [], "descricao": "Processar o dado no PRÓPRIO sensor (na borda da rede), não mandá-lo cru ao centro. É a lição do floco no metal: a informação se lê onde ela vive — filtrar, agregar e decidir junto ao olho poupa banda, latência e revela o que a régua central perderia.", "epistemico": "inferencia", "exemplos": [], "id": "edge_computing", "memoria": "semantica", "meta": {"dominio": "engenharia", "fonte": "IoT e sensores"}, "origem": "wiki", "palavras_chave": [], "sinonimos": [], "tipo": "conceito", "titulo": "Edge Computing (computar no metal)"}
\section{Edge Computing (computar no metal)}
Processar o dado no PRÓPRIO sensor (na borda da rede), não mandá-lo cru ao centro. É a lição do floco no metal: a informação se lê onde ela vive — filtrar, agregar e decidir junto ao olho poupa banda, latência e revela o que a régua central perderia.
\prov{relações: usa internet\_das\_coisas; relaciona sensor}
\prov{proveniência: wiki \textperiodcentered{} IoT e sensores \textperiodcentered{} inferencia \textperiodcentered{} confiança 1.0 \textperiodcentered{} domínio engenharia \textperiodcentered{} história 6 versões no jornal}

%CRISTAL {"arestas": [{"alvo": "computacao_em_nuvem", "confianca": 1.0, "epistemico": "fato", "origem": "wiki", "rel": "explica"}, {"alvo": "edge_computing", "confianca": 1.0, "epistemico": "fato", "origem": "wiki", "rel": "relaciona"}, {"alvo": "corpo_mineral", "confianca": 1.0, "epistemico": "fato", "origem": "wiki", "rel": "usa"}], "confianca": 1.0, "contraexemplos": [], "descricao": "Duas forças opostas e ambas reais: o EDGE computa junto ao sensor (onde o floco vive), a NUVEM concentra no centro (o metal comum). A pergunta é sempre a mesma da máquina mineral: onde a informação quer ser lida? Latência e privacidade puxam ao edge; escala e visão global puxam à nuvem.", "epistemico": "inferencia", "exemplos": [], "id": "edge_vs_nuvem", "memoria": "semantica", "meta": {"dominio": "engenharia", "fonte": "nuvem e data centers"}, "origem": "wiki", "palavras_chave": [], "sinonimos": [], "tipo": "conceito", "titulo": "Edge vs Nuvem (o floco e o metal)"}
\section{Edge vs Nuvem (o floco e o metal)}
Duas forças opostas e ambas reais: o EDGE computa junto ao sensor (onde o floco vive), a NUVEM concentra no centro (o metal comum). A pergunta é sempre a mesma da máquina mineral: onde a informação quer ser lida? Latência e privacidade puxam ao edge; escala e visão global puxam à nuvem.
\prov{relações: explica computacao\_em\_nuvem; relaciona edge\_computing; usa corpo\_mineral}
\prov{proveniência: wiki \textperiodcentered{} nuvem e data centers \textperiodcentered{} inferencia \textperiodcentered{} confiança 1.0 \textperiodcentered{} domínio engenharia \textperiodcentered{} história 4 versões no jornal}

%CRISTAL {"arestas": [{"alvo": "armazenamento_o_cristal", "confianca": 1.0, "epistemico": "fato", "origem": "wiki", "rel": "explica"}], "confianca": 1.0, "contraexemplos": [], "descricao": "A razão energia devolvida / energia injetada num ciclo carga-descarga. Sempre <1 (a 2ª lei); mede a integridade do cristal de energia. Supercap ~95% > lítio ~90% > pumped hydro ~80% > CAES ~60% > térmico ~55% > hidrogênio ~40%. O que falta para 1 é a perda.", "epistemico": "inferencia", "exemplos": [], "id": "eficiencia_round_trip", "memoria": "semantica", "meta": {"dominio": "engenharia", "fonte": "distribuição e armazenamento", "round_trip": {"ar comprimido (CAES)": 0.6, "hidrelétrica reversível": 0.8, "hidrogênio (power-to-gas)": 0.4, "supercapacitor": 0.95, "térmico (sais fundidos)": 0.55, "volante de inércia": 0.88, "íon-lítio (bateria)": 0.9}}, "origem": "wiki", "palavras_chave": [], "sinonimos": [], "tipo": "conceito", "titulo": "Eficiência Round-Trip (a integridade do cristal)"}
\section{Eficiência Round-Trip (a integridade do cristal)}
A razão energia devolvida / energia injetada num ciclo carga-descarga. Sempre <1 (a 2ª lei); mede a integridade do cristal de energia. Supercap \~{}95\% > lítio \~{}90\% > pumped hydro \~{}80\% > CAES \~{}60\% > térmico \~{}55\% > hidrogênio \~{}40\%. O que falta para 1 é a perda.
\prov{relações: explica armazenamento\_o\_cristal}
\prov{proveniência: wiki \textperiodcentered{} distribuição e armazenamento \textperiodcentered{} inferencia \textperiodcentered{} confiança 1.0 \textperiodcentered{} domínio engenharia \textperiodcentered{} história 4 versões no jornal}

%CRISTAL {"arestas": [{"alvo": "dicionario_das_fontes", "confianca": 1.0, "epistemico": "fato", "origem": "wiki", "rel": "parte_de"}, {"alvo": "reconhecer_caligrafia_estilo", "confianca": 1.0, "epistemico": "fato", "origem": "wiki", "rel": "usa"}], "confianca": 1.0, "contraexemplos": [], "descricao": "O estilo de uma fonte se decompõe em 5 eixos quase ortogonais: INCLINAÇÃO (reto/inclinado — o itálico), LARGURA (proporcional/monoespaçada — o mono), PESO DO TRAÇO (fino/grosso — o negrito), CONTRASTE (baixo/alto — a modulação grosso-fino) e SERIFA (com/sem). Quatro são medidos no alfabeto (inclinação pelo cisalhamento que minimiza a largura, mono pela variância das larguras, peso e contraste pela transformada de distância); a serifa é fato de design.", "epistemico": "fato", "exemplos": [], "id": "eixos_da_caligrafia", "memoria": "semantica", "meta": {"dominio": "imagem", "fonte": "dicionário de fontes"}, "origem": "wiki", "palavras_chave": [], "sinonimos": [], "tipo": "conceito", "titulo": "Os Cinco Eixos da Caligrafia"}
\section{Os Cinco Eixos da Caligrafia}
O estilo de uma fonte se decompõe em 5 eixos quase ortogonais: INCLINAÇÃO (reto/inclinado — o itálico), LARGURA (proporcional/monoespaçada — o mono), PESO DO TRAÇO (fino/grosso — o negrito), CONTRASTE (baixo/alto — a modulação grosso-fino) e SERIFA (com/sem). Quatro são medidos no alfabeto (inclinação pelo cisalhamento que minimiza a largura, mono pela variância das larguras, peso e contraste pela transformada de distância); a serifa é fato de design.
\prov{relações: parte\_de dicionario\_das\_fontes; usa reconhecer\_caligrafia\_estilo}
\prov{proveniência: wiki \textperiodcentered{} dicionário de fontes \textperiodcentered{} fato \textperiodcentered{} confiança 1.0 \textperiodcentered{} domínio imagem \textperiodcentered{} história 3 versões no jornal}

%CRISTAL {"arestas": [{"alvo": "ciclo_termodinamico_orbita", "confianca": 1.0, "epistemico": "fato", "origem": "wiki", "rel": "relaciona"}, {"alvo": "motor_eletrico_tracao", "confianca": 1.0, "epistemico": "fato", "origem": "wiki", "rel": "explica"}, {"alvo": "corpo_mineral", "confianca": 1.0, "epistemico": "fato", "origem": "wiki", "rel": "usa"}], "confianca": 1.0, "contraexemplos": [], "descricao": "O motor a combustão é máquina térmica: preso ao teto de Carnot η=1−Tc/Th (~30-40% na prática). O carro elétrico NÃO é máquina térmica — vai do cristal → inversor → rotação direto, sem ciclo térmico. Por isso ~90% tanque-à-roda: ele SALTA a órbita fechada de Carnot inteira.", "epistemico": "inferencia", "exemplos": [], "id": "eletrico_pula_carnot", "memoria": "semantica", "meta": {"dominio": "engenharia", "eta_combustao": 0.333, "eta_eletrico": 0.9, "fonte": "mobilidade elétrica", "pula_carnot": true}, "origem": "wiki", "palavras_chave": [], "sinonimos": [], "tipo": "conceito", "titulo": "O Elétrico Pula a Órbita de Carnot"}
\section{O Elétrico Pula a Órbita de Carnot}
O motor a combustão é máquina térmica: preso ao teto de Carnot \ensuremath{\eta}=1\ensuremath{-}Tc/Th (\~{}30-40\% na prática). O carro elétrico NÃO é máquina térmica — vai do cristal \ensuremath{\to} inversor \ensuremath{\to} rotação direto, sem ciclo térmico. Por isso \~{}90\% tanque-à-roda: ele SALTA a órbita fechada de Carnot inteira.
\prov{relações: relaciona ciclo\_termodinamico\_orbita; explica motor\_eletrico\_tracao; usa corpo\_mineral}
\prov{proveniência: wiki \textperiodcentered{} mobilidade elétrica \textperiodcentered{} inferencia \textperiodcentered{} confiança 1.0 \textperiodcentered{} domínio engenharia \textperiodcentered{} história 4 versões no jornal}

%CRISTAL {"arestas": [{"alvo": "fisica_ondas_eng", "confianca": 1.0, "epistemico": "fato", "origem": "wiki", "rel": "usa"}, {"alvo": "equacoes_diferenciais", "confianca": 1.0, "epistemico": "fato", "origem": "wiki", "rel": "usa"}, {"alvo": "eng_eletronica", "confianca": 1.0, "epistemico": "fato", "origem": "curriculo", "rel": "parte_de"}, {"alvo": "eng_telecomunicacoes", "confianca": 1.0, "epistemico": "fato", "origem": "curriculo", "rel": "parte_de"}, {"alvo": "eng_sistemas_energia", "confianca": 1.0, "epistemico": "fato", "origem": "curriculo", "rel": "parte_de"}], "confianca": 1.0, "contraexemplos": [], "descricao": "As quatro equações de Maxwell unificam eletricidade, magnetismo e luz: campo elétrico, magnético, Gauss, Faraday e Ampère-Maxwell. A fundação física de antenas, motores e da própria onda.", "epistemico": "fato", "exemplos": [], "id": "eletromagnetismo_eng", "memoria": "semantica", "meta": {"autor": "Maxwell", "dominio": "engenharia", "fonte": ""}, "origem": "wiki", "palavras_chave": [], "sinonimos": [], "tipo": "conceito", "titulo": "Eletromagnetismo (Equações de Maxwell)"}
\section{Eletromagnetismo (Equações de Maxwell)}
As quatro equações de Maxwell unificam eletricidade, magnetismo e luz: campo elétrico, magnético, Gauss, Faraday e Ampère-Maxwell. A fundação física de antenas, motores e da própria onda.
\prov{relações: usa fisica\_ondas\_eng; usa equacoes\_diferenciais; parte\_de eng\_eletronica; parte\_de eng\_telecomunicacoes; parte\_de eng\_sistemas\_energia}
\prov{proveniência: wiki \textperiodcentered{} fato \textperiodcentered{} confiança 1.0 \textperiodcentered{} domínio engenharia \textperiodcentered{} história 6 versões no jornal}

%CRISTAL {"arestas": [{"alvo": "dispositivo_semicondutor", "confianca": 1.0, "epistemico": "fato", "origem": "wiki", "rel": "usa"}, {"alvo": "corrente_e_tensao", "confianca": 1.0, "epistemico": "fato", "origem": "wiki", "rel": "usa"}], "confianca": 1.0, "contraexemplos": [], "descricao": "A engenharia de controlar o fluxo de elétrons em circuitos para processar sinais e potência — de dispositivos a sistemas.", "epistemico": "fato", "exemplos": [], "id": "eletronica", "memoria": "semantica", "meta": {"dominio": "engenharia", "fonte": ""}, "origem": "wiki", "palavras_chave": [], "sinonimos": [], "tipo": "conceito", "titulo": "Eletrônica"}
\section{Eletrônica}
A engenharia de controlar o fluxo de elétrons em circuitos para processar sinais e potência — de dispositivos a sistemas.
\prov{relações: usa dispositivo\_semicondutor; usa corrente\_e\_tensao}
\prov{proveniência: wiki \textperiodcentered{} fato \textperiodcentered{} confiança 1.0 \textperiodcentered{} domínio engenharia \textperiodcentered{} história 3 versões no jornal}

%CRISTAL {"arestas": [{"alvo": "eletronica", "confianca": 1.0, "epistemico": "fato", "origem": "wiki", "rel": "especializa"}, {"alvo": "amplificador", "confianca": 1.0, "epistemico": "fato", "origem": "wiki", "rel": "usa"}, {"alvo": "transistor", "confianca": 1.0, "epistemico": "fato", "origem": "wiki", "rel": "usa"}, {"alvo": "eng_eletronica", "confianca": 1.0, "epistemico": "fato", "origem": "curriculo", "rel": "parte_de"}, {"alvo": "eng_telecomunicacoes", "confianca": 1.0, "epistemico": "fato", "origem": "curriculo", "rel": "parte_de"}, {"alvo": "eng_sistemas_energia", "confianca": 1.0, "epistemico": "fato", "origem": "curriculo", "rel": "parte_de"}], "confianca": 1.0, "contraexemplos": [], "descricao": "Amplificar, filtrar e condicionar sinais contínuos com diodos, transistores e amp-ops — ganho, polarização, realimentação. O elo entre o mundo físico analógico e o resto do sistema.", "epistemico": "fato", "exemplos": [], "id": "eletronica_analogica", "memoria": "semantica", "meta": {"dominio": "engenharia", "fonte": ""}, "origem": "wiki", "palavras_chave": [], "sinonimos": [], "tipo": "conceito", "titulo": "Eletrônica Analógica"}
\section{Eletrônica Analógica}
Amplificar, filtrar e condicionar sinais contínuos com diodos, transistores e amp-ops — ganho, polarização, realimentação. O elo entre o mundo físico analógico e o resto do sistema.
\prov{relações: especializa eletronica; usa amplificador; usa transistor; parte\_de eng\_eletronica; parte\_de eng\_telecomunicacoes; parte\_de eng\_sistemas\_energia}
\prov{proveniência: wiki \textperiodcentered{} fato \textperiodcentered{} confiança 1.0 \textperiodcentered{} domínio engenharia \textperiodcentered{} história 7 versões no jornal}

%CRISTAL {"arestas": [{"alvo": "semicondutor", "confianca": 1.0, "epistemico": "fato", "origem": "wiki", "rel": "usa"}, {"alvo": "mosfet", "confianca": 1.0, "epistemico": "fato", "origem": "wiki", "rel": "usa"}, {"alvo": "eletronica", "confianca": 1.0, "epistemico": "fato", "origem": "wiki", "rel": "especializa"}, {"alvo": "eng_sistemas_energia", "confianca": 1.0, "epistemico": "fato", "origem": "curriculo", "rel": "parte_de"}], "confianca": 1.0, "contraexemplos": [], "descricao": "Chavear grandes potências com semicondutores (MOSFET, IGBT, tiristor) para converter e controlar energia com alta eficiência. A tecnologia que casou a eletrônica com a energia — inversores, fontes, drives.", "epistemico": "fato", "exemplos": [], "id": "eletronica_de_potencia", "memoria": "semantica", "meta": {"dominio": "engenharia", "fonte": ""}, "origem": "wiki", "palavras_chave": [], "sinonimos": [], "tipo": "conceito", "titulo": "Eletrônica de Potência"}
\section{Eletrônica de Potência}
Chavear grandes potências com semicondutores (MOSFET, IGBT, tiristor) para converter e controlar energia com alta eficiência. A tecnologia que casou a eletrônica com a energia — inversores, fontes, drives.
\prov{relações: usa semicondutor; usa mosfet; especializa eletronica; parte\_de eng\_sistemas\_energia}
\prov{proveniência: wiki \textperiodcentered{} fato \textperiodcentered{} confiança 1.0 \textperiodcentered{} domínio engenharia \textperiodcentered{} história 5 versões no jornal}

%CRISTAL {"arestas": [{"alvo": "eletronica_analogica", "confianca": 1.0, "epistemico": "fato", "origem": "wiki", "rel": "especializa"}, {"alvo": "ondas_eletromagneticas", "confianca": 1.0, "epistemico": "fato", "origem": "wiki", "rel": "usa"}, {"alvo": "eng_eletronica", "confianca": 1.0, "epistemico": "fato", "origem": "curriculo", "rel": "parte_de"}], "confianca": 1.0, "contraexemplos": [], "descricao": "Circuitos que operam em rádio-frequência: LNA, mixer, oscilador, PLL, casamento de impedância. A camada física do transceptor — onde a eletrônica encontra a telecomunicação.", "epistemico": "fato", "exemplos": [], "id": "eletronica_de_rf", "memoria": "semantica", "meta": {"dominio": "engenharia", "fonte": ""}, "origem": "wiki", "palavras_chave": [], "sinonimos": [], "tipo": "conceito", "titulo": "Eletrônica de RF e Micro-ondas"}
\section{Eletrônica de RF e Micro-ondas}
Circuitos que operam em rádio-frequência: LNA, mixer, oscilador, PLL, casamento de impedância. A camada física do transceptor — onde a eletrônica encontra a telecomunicação.
\prov{relações: especializa eletronica\_analogica; usa ondas\_eletromagneticas; parte\_de eng\_eletronica}
\prov{proveniência: wiki \textperiodcentered{} fato \textperiodcentered{} confiança 1.0 \textperiodcentered{} domínio engenharia \textperiodcentered{} história 4 versões no jornal}

%CRISTAL {"arestas": [{"alvo": "logica_digital", "confianca": 1.0, "epistemico": "fato", "origem": "wiki", "rel": "usa"}, {"alvo": "eletronica", "confianca": 1.0, "epistemico": "fato", "origem": "wiki", "rel": "especializa"}, {"alvo": "eng_eletronica", "confianca": 1.0, "epistemico": "fato", "origem": "curriculo", "rel": "parte_de"}, {"alvo": "eng_telecomunicacoes", "confianca": 1.0, "epistemico": "fato", "origem": "curriculo", "rel": "parte_de"}, {"alvo": "eng_sistemas_energia", "confianca": 1.0, "epistemico": "fato", "origem": "curriculo", "rel": "parte_de"}], "confianca": 1.0, "contraexemplos": [], "descricao": "Construir lógica com 0 e 1: álgebra booleana, portas, flip-flops, máquinas de estado, memória. O andar onde o sinal vira símbolo — a base do computador e do processamento digital.", "epistemico": "fato", "exemplos": [], "id": "eletronica_digital", "memoria": "semantica", "meta": {"autor": "Boole", "dominio": "engenharia", "fonte": ""}, "origem": "wiki", "palavras_chave": [], "sinonimos": [], "tipo": "conceito", "titulo": "Eletrônica Digital / Sistemas Digitais"}
\section{Eletrônica Digital / Sistemas Digitais}
Construir lógica com 0 e 1: álgebra booleana, portas, flip-flops, máquinas de estado, memória. O andar onde o sinal vira símbolo — a base do computador e do processamento digital.
\prov{relações: usa logica\_digital; especializa eletronica; parte\_de eng\_eletronica; parte\_de eng\_telecomunicacoes; parte\_de eng\_sistemas\_energia}
\prov{proveniência: wiki \textperiodcentered{} fato \textperiodcentered{} confiança 1.0 \textperiodcentered{} domínio engenharia \textperiodcentered{} história 6 versões no jornal}

%CRISTAL {"arestas": [{"alvo": "eletronica", "confianca": 1.0, "epistemico": "fato", "origem": "wiki", "rel": "parte_de"}, {"alvo": "logica_digital", "confianca": 1.0, "epistemico": "fato", "origem": "wiki", "rel": "explica"}], "confianca": 1.0, "contraexemplos": [], "descricao": "Analógica lida com sinais contínuos (som, rádio); digital com níveis discretos (0/1) — a discretização que trouxe ruído-imunidade e o computador.", "epistemico": "inferencia", "exemplos": [], "id": "eletronica_digital_analogica", "memoria": "semantica", "meta": {"dominio": "engenharia", "fonte": ""}, "origem": "wiki", "palavras_chave": [], "sinonimos": [], "tipo": "conceito", "titulo": "Eletrônica Digital vs Analógica"}
\section{Eletrônica Digital vs Analógica}
Analógica lida com sinais contínuos (som, rádio); digital com níveis discretos (0/1) — a discretização que trouxe ruído-imunidade e o computador.
\prov{relações: parte\_de eletronica; explica logica\_digital}
\prov{proveniência: wiki \textperiodcentered{} inferencia \textperiodcentered{} confiança 1.0 \textperiodcentered{} domínio engenharia \textperiodcentered{} história 3 versões no jornal}

%CRISTAL {"arestas": [{"alvo": "protocolo_tcp", "confianca": 1.0, "epistemico": "fato", "origem": "wiki", "rel": "usa"}, {"alvo": "eng_telecomunicacoes", "confianca": 1.0, "epistemico": "fato", "origem": "curriculo", "rel": "parte_de"}], "confianca": 1.0, "contraexemplos": [], "descricao": "SMTP entrega a mensagem entre servidores; IMAP/POP a busca no cliente. O protocolo de aplicação mais antigo ainda de pé — store-and-forward, endereços @ e a origem do spam.", "epistemico": "fato", "exemplos": [], "id": "email_smtp", "memoria": "semantica", "meta": {"dominio": "engenharia", "fonte": ""}, "origem": "wiki", "palavras_chave": [], "sinonimos": [], "tipo": "conceito", "titulo": "E-mail (SMTP/IMAP)"}
\section{E-mail (SMTP/IMAP)}
SMTP entrega a mensagem entre servidores; IMAP/POP a busca no cliente. O protocolo de aplicação mais antigo ainda de pé — store-and-forward, endereços @ e a origem do spam.
\prov{relações: usa protocolo\_tcp; parte\_de eng\_telecomunicacoes}
\prov{proveniência: wiki \textperiodcentered{} fato \textperiodcentered{} confiança 1.0 \textperiodcentered{} domínio engenharia \textperiodcentered{} história 3 versões no jornal}

%CRISTAL {"arestas": [{"alvo": "modelo_osi", "confianca": 1.0, "epistemico": "fato", "origem": "wiki", "rel": "parte_de"}, {"alvo": "eng_telecomunicacoes", "confianca": 1.0, "epistemico": "fato", "origem": "curriculo", "rel": "parte_de"}], "confianca": 1.0, "contraexemplos": [], "descricao": "Cada camada embrulha o dado da de cima com seu próprio cabeçalho (e o desembrulha no destino) — quadro⊃pacote⊃segmento⊃dados. A cebola que faz as camadas coexistirem sem se conhecer.", "epistemico": "fato", "exemplos": [], "id": "encapsulamento", "memoria": "semantica", "meta": {"dominio": "engenharia", "fonte": ""}, "origem": "wiki", "palavras_chave": [], "sinonimos": [], "tipo": "conceito", "titulo": "Encapsulamento"}
\section{Encapsulamento}
Cada camada embrulha o dado da de cima com seu próprio cabeçalho (e o desembrulha no destino) — quadro\ensuremath{\supset}pacote\ensuremath{\supset}segmento\ensuremath{\supset}dados. A cebola que faz as camadas coexistirem sem se conhecer.
\prov{relações: parte\_de modelo\_osi; parte\_de eng\_telecomunicacoes}
\prov{proveniência: wiki \textperiodcentered{} fato \textperiodcentered{} confiança 1.0 \textperiodcentered{} domínio engenharia \textperiodcentered{} história 3 versões no jornal}

%CRISTAL {"arestas": [{"alvo": "camada_de_enlace", "confianca": 1.0, "epistemico": "fato", "origem": "wiki", "rel": "parte_de"}, {"alvo": "eng_telecomunicacoes", "confianca": 1.0, "epistemico": "fato", "origem": "curriculo", "rel": "parte_de"}], "confianca": 1.0, "contraexemplos": [], "descricao": "O identificador físico de 48 bits gravado na placa de rede — endereça o quadro dentro do segmento local. Plano e único de fábrica, ao contrário do IP, que é hierárquico e roteável.", "epistemico": "fato", "exemplos": [], "id": "endereco_mac", "memoria": "semantica", "meta": {"dominio": "engenharia", "fonte": ""}, "origem": "wiki", "palavras_chave": [], "sinonimos": [], "tipo": "conceito", "titulo": "Endereço MAC"}
\section{Endereço MAC}
O identificador físico de 48 bits gravado na placa de rede — endereça o quadro dentro do segmento local. Plano e único de fábrica, ao contrário do IP, que é hierárquico e roteável.
\prov{relações: parte\_de camada\_de\_enlace; parte\_de eng\_telecomunicacoes}
\prov{proveniência: wiki \textperiodcentered{} fato \textperiodcentered{} confiança 1.0 \textperiodcentered{} domínio engenharia \textperiodcentered{} história 3 versões no jornal}

%CRISTAL {"arestas": [{"alvo": "eletronica", "confianca": 1.0, "epistemico": "fato", "origem": "wiki", "rel": "relaciona"}, {"alvo": "termodinamica", "confianca": 1.0, "epistemico": "fato", "origem": "wiki", "rel": "usa"}, {"alvo": "prova_de_trabalho", "confianca": 1.0, "epistemico": "fato", "origem": "wiki", "rel": "relaciona"}], "confianca": 1.0, "contraexemplos": [], "descricao": "A forma de energia mais versátil e transportável — gerada, transmitida a longa distância e distribuída ao consumo. A espinha da civilização técnica.", "epistemico": "fato", "exemplos": [], "id": "energia_eletrica", "memoria": "semantica", "meta": {"dominio": "engenharia", "fonte": ""}, "origem": "wiki", "palavras_chave": [], "sinonimos": [], "tipo": "conceito", "titulo": "Energia Elétrica"}
\section{Energia Elétrica}
A forma de energia mais versátil e transportável — gerada, transmitida a longa distância e distribuída ao consumo. A espinha da civilização técnica.
\prov{relações: relaciona eletronica; usa termodinamica; relaciona prova\_de\_trabalho}
\prov{proveniência: wiki \textperiodcentered{} fato \textperiodcentered{} confiança 1.0 \textperiodcentered{} domínio engenharia \textperiodcentered{} história 4 versões no jornal}

%CRISTAL {"arestas": [{"alvo": "geracao_de_energia", "confianca": 1.0, "epistemico": "fato", "origem": "wiki", "rel": "especializa"}], "confianca": 1.0, "contraexemplos": [], "descricao": "O vento gira turbinas que acionam geradores — renovável e crescente, com o desafio da intermitência.", "epistemico": "inferencia", "exemplos": [], "id": "energia_eolica", "memoria": "semantica", "meta": {"dominio": "engenharia", "fonte": ""}, "origem": "wiki", "palavras_chave": [], "sinonimos": [], "tipo": "conceito", "titulo": "Energia Eólica"}
\section{Energia Eólica}
O vento gira turbinas que acionam geradores — renovável e crescente, com o desafio da intermitência.
\prov{relações: especializa geracao\_de\_energia}
\prov{proveniência: wiki \textperiodcentered{} inferencia \textperiodcentered{} confiança 1.0 \textperiodcentered{} domínio engenharia \textperiodcentered{} história 2 versões no jornal}

%CRISTAL {"arestas": [{"alvo": "geracao_de_energia", "confianca": 1.0, "epistemico": "fato", "origem": "wiki", "rel": "especializa"}], "confianca": 1.0, "contraexemplos": [], "descricao": "A queda d'água gira turbinas — a maior fonte renovável despachável; o modelo das grandes barragens (e do Brasil).", "epistemico": "inferencia", "exemplos": [], "id": "energia_hidreletrica", "memoria": "semantica", "meta": {"dominio": "engenharia", "fonte": ""}, "origem": "wiki", "palavras_chave": [], "sinonimos": [], "tipo": "conceito", "titulo": "Energia Hidrelétrica"}
\section{Energia Hidrelétrica}
A queda d'água gira turbinas — a maior fonte renovável despachável; o modelo das grandes barragens (e do Brasil).
\prov{relações: especializa geracao\_de\_energia}
\prov{proveniência: wiki \textperiodcentered{} inferencia \textperiodcentered{} confiança 1.0 \textperiodcentered{} domínio engenharia \textperiodcentered{} história 2 versões no jornal}

%CRISTAL {"arestas": [{"alvo": "geracao_de_energia", "confianca": 1.0, "epistemico": "fato", "origem": "wiki", "rel": "especializa"}, {"alvo": "celula_fotovoltaica", "confianca": 1.0, "epistemico": "fato", "origem": "wiki", "rel": "usa"}], "confianca": 1.0, "contraexemplos": [], "descricao": "Converter luz do sol em eletricidade — direto por células fotovoltaicas, ou por calor concentrado. Renovável e cada vez mais barata.", "epistemico": "inferencia", "exemplos": [], "id": "energia_solar", "memoria": "semantica", "meta": {"dominio": "engenharia", "fonte": ""}, "origem": "wiki", "palavras_chave": [], "sinonimos": [], "tipo": "conceito", "titulo": "Energia Solar"}
\section{Energia Solar}
Converter luz do sol em eletricidade — direto por células fotovoltaicas, ou por calor concentrado. Renovável e cada vez mais barata.
\prov{relações: especializa geracao\_de\_energia; usa celula\_fotovoltaica}
\prov{proveniência: wiki \textperiodcentered{} inferencia \textperiodcentered{} confiança 1.0 \textperiodcentered{} domínio engenharia \textperiodcentered{} história 3 versões no jornal}

%CRISTAL {"arestas": [{"alvo": "engenharia_eletrica", "confianca": 1.0, "epistemico": "fato", "origem": "wiki", "rel": "especializa"}], "confianca": 1.0, "contraexemplos": [], "descricao": "A ênfase do DISPOSITIVO e do CIRCUITO: semicondutores, circuitos integrados, sistemas embarcados, instrumentação e processamento de sinais. Onde a informação vira corrente controlada no silício.", "epistemico": "fato", "exemplos": [], "id": "eng_eletronica", "memoria": "semantica", "meta": {"dominio": "engenharia", "fonte": ""}, "origem": "wiki", "palavras_chave": [], "sinonimos": [], "tipo": "conceito", "titulo": "Engenharia Eletrônica"}
\section{Engenharia Eletrônica}
A ênfase do DISPOSITIVO e do CIRCUITO: semicondutores, circuitos integrados, sistemas embarcados, instrumentação e processamento de sinais. Onde a informação vira corrente controlada no silício.
\prov{relações: especializa engenharia\_eletrica}
\prov{proveniência: wiki \textperiodcentered{} fato \textperiodcentered{} confiança 1.0 \textperiodcentered{} domínio engenharia \textperiodcentered{} história 2 versões no jornal}

%CRISTAL {"arestas": [{"alvo": "engenharia_eletrica", "confianca": 1.0, "epistemico": "fato", "origem": "wiki", "rel": "especializa"}], "confianca": 1.0, "contraexemplos": [], "descricao": "A ênfase da POTÊNCIA: geração, transmissão, distribuição, máquinas elétricas, eletrônica de potência e proteção. Manter tensão, frequência e balanço numa rede que nunca pode parar.", "epistemico": "fato", "exemplos": [], "id": "eng_sistemas_energia", "memoria": "semantica", "meta": {"dominio": "engenharia", "fonte": ""}, "origem": "wiki", "palavras_chave": [], "sinonimos": [], "tipo": "conceito", "titulo": "Engenharia de Sistemas de Energia"}
\section{Engenharia de Sistemas de Energia}
A ênfase da POTÊNCIA: geração, transmissão, distribuição, máquinas elétricas, eletrônica de potência e proteção. Manter tensão, frequência e balanço numa rede que nunca pode parar.
\prov{relações: especializa engenharia\_eletrica}
\prov{proveniência: wiki \textperiodcentered{} fato \textperiodcentered{} confiança 1.0 \textperiodcentered{} domínio engenharia \textperiodcentered{} história 2 versões no jornal}

%CRISTAL {"arestas": [{"alvo": "engenharia_eletrica", "confianca": 1.0, "epistemico": "fato", "origem": "wiki", "rel": "especializa"}], "confianca": 1.0, "contraexemplos": [], "descricao": "A ênfase do SINAL QUE VIAJA: modulação, teoria da informação, propagação, antenas, fibra, redes e protocolos. Levar informação de A a B com o mínimo de erro sobre um canal ruidoso e compartilhado.", "epistemico": "fato", "exemplos": [], "id": "eng_telecomunicacoes", "memoria": "semantica", "meta": {"dominio": "engenharia", "fonte": ""}, "origem": "wiki", "palavras_chave": [], "sinonimos": [], "tipo": "conceito", "titulo": "Engenharia de Telecomunicações"}
\section{Engenharia de Telecomunicações}
A ênfase do SINAL QUE VIAJA: modulação, teoria da informação, propagação, antenas, fibra, redes e protocolos. Levar informação de A a B com o mínimo de erro sobre um canal ruidoso e compartilhado.
\prov{relações: especializa engenharia\_eletrica}
\prov{proveniência: wiki \textperiodcentered{} fato \textperiodcentered{} confiança 1.0 \textperiodcentered{} domínio engenharia \textperiodcentered{} história 2 versões no jornal}

%CRISTAL {"arestas": [{"alvo": "eletronica", "confianca": 1.0, "epistemico": "fato", "origem": "wiki", "rel": "usa"}, {"alvo": "fisica", "confianca": 1.0, "epistemico": "fato", "origem": "wiki", "rel": "usa"}, {"alvo": "matematica", "confianca": 1.0, "epistemico": "fato", "origem": "wiki", "rel": "usa"}], "confianca": 1.0, "contraexemplos": [], "descricao": "A grande área que domestica a eletricidade e o eletromagnetismo — do elétron no cristal ao sistema de potência de um país. Do seu tronco saem a eletrônica, as telecomunicações e a energia.", "epistemico": "fato", "exemplos": [], "id": "engenharia_eletrica", "memoria": "semantica", "meta": {"dominio": "engenharia", "fonte": ""}, "origem": "wiki", "palavras_chave": [], "sinonimos": [], "tipo": "conceito", "titulo": "Engenharia Elétrica"}
\section{Engenharia Elétrica}
A grande área que domestica a eletricidade e o eletromagnetismo — do elétron no cristal ao sistema de potência de um país. Do seu tronco saem a eletrônica, as telecomunicações e a energia.
\prov{relações: usa eletronica; usa fisica; usa matematica}
\prov{proveniência: wiki \textperiodcentered{} fato \textperiodcentered{} confiança 1.0 \textperiodcentered{} domínio engenharia \textperiodcentered{} história 4 versões no jornal}

%CRISTAL {"arestas": [{"alvo": "sistema_autonomo", "confianca": 1.0, "epistemico": "fato", "origem": "wiki", "rel": "usa"}, {"alvo": "rede_autonoma_broca", "confianca": 1.0, "epistemico": "fato", "origem": "wiki", "rel": "eh_um"}, {"alvo": "corpo_mineral", "confianca": 1.0, "epistemico": "fato", "origem": "wiki", "rel": "relaciona"}], "confianca": 1.0, "contraexemplos": [], "descricao": "Muitos robôs simples coordenando-se por regras locais fazem emergir comportamento global (formação, busca, construção). É um GRAFO de agentes — a rede autônoma da broca-so — vivo na borda: regras rígidas demais congelam (cristal), soltas demais viram caos, o enxame útil fica na borda.", "epistemico": "inferencia", "exemplos": [], "id": "enxame_robotico", "memoria": "semantica", "meta": {"dominio": "robotica", "fonte": "robótica e sistemas autônomos"}, "origem": "wiki", "palavras_chave": [], "sinonimos": [], "tipo": "conceito", "titulo": "Enxame Robótico (o grafo de agentes)"}
\section{Enxame Robótico (o grafo de agentes)}
Muitos robôs simples coordenando-se por regras locais fazem emergir comportamento global (formação, busca, construção). É um GRAFO de agentes — a rede autônoma da broca-so — vivo na borda: regras rígidas demais congelam (cristal), soltas demais viram caos, o enxame útil fica na borda.
\prov{relações: usa sistema\_autonomo; eh\_um rede\_autonoma\_broca; relaciona corpo\_mineral}
\prov{proveniência: wiki \textperiodcentered{} robótica e sistemas autônomos \textperiodcentered{} inferencia \textperiodcentered{} confiança 1.0 \textperiodcentered{} domínio robotica \textperiodcentered{} história 4 versões no jornal}

%CRISTAL {"arestas": [{"alvo": "calculo_engenharia", "confianca": 1.0, "epistemico": "fato", "origem": "wiki", "rel": "usa"}, {"alvo": "eng_eletronica", "confianca": 1.0, "epistemico": "fato", "origem": "curriculo", "rel": "parte_de"}, {"alvo": "eng_telecomunicacoes", "confianca": 1.0, "epistemico": "fato", "origem": "curriculo", "rel": "parte_de"}, {"alvo": "eng_sistemas_energia", "confianca": 1.0, "epistemico": "fato", "origem": "curriculo", "rel": "parte_de"}], "confianca": 1.0, "contraexemplos": [], "descricao": "Descrever um sistema pela lei de como ele muda — do circuito RLC ao decaimento e à onda. A EDO/EDP é o modelo padrão de todo sistema dinâmico da engenharia elétrica.", "epistemico": "fato", "exemplos": [], "id": "equacoes_diferenciais", "memoria": "semantica", "meta": {"dominio": "engenharia", "fonte": ""}, "origem": "wiki", "palavras_chave": [], "sinonimos": [], "tipo": "conceito", "titulo": "Equações Diferenciais"}
\section{Equações Diferenciais}
Descrever um sistema pela lei de como ele muda — do circuito RLC ao decaimento e à onda. A EDO/EDP é o modelo padrão de todo sistema dinâmico da engenharia elétrica.
\prov{relações: usa calculo\_engenharia; parte\_de eng\_eletronica; parte\_de eng\_telecomunicacoes; parte\_de eng\_sistemas\_energia}
\prov{proveniência: wiki \textperiodcentered{} fato \textperiodcentered{} confiança 1.0 \textperiodcentered{} domínio engenharia \textperiodcentered{} história 5 versões no jornal}

%CRISTAL {"arestas": [{"alvo": "mecanismo_de_la_hire", "confianca": 1.0, "epistemico": "fato", "origem": "wiki", "rel": "usa"}, {"alvo": "nucleossintese_dos_flocos", "confianca": 1.0, "epistemico": "fato", "origem": "wiki", "rel": "parte_de"}, {"alvo": "movimento_circular_dois_modos", "confianca": 1.0, "epistemico": "fato", "origem": "wiki", "rel": "usa"}, {"alvo": "vale_do_ferro", "confianca": 1.0, "epistemico": "fato", "origem": "wiki", "rel": "relaciona"}, {"alvo": "corpo_mineral", "confianca": 1.0, "epistemico": "fato", "origem": "wiki", "rel": "usa"}, {"alvo": "corpo_estelar", "confianca": 1.0, "epistemico": "fato", "origem": "wiki", "rel": "usa"}, {"alvo": "fusao_nuclear", "confianca": 1.0, "epistemico": "fato", "origem": "wiki", "rel": "relaciona"}, {"alvo": "familia_m-bonacci_gap_k", "confianca": 1.0, "epistemico": "fato", "origem": "wiki", "rel": "relaciona"}], "confianca": 1.0, "contraexemplos": [], "descricao": "Generalizar La Hire à nucleossíntese: o k_significativo de um floco é o nº de CÚSPIDES da sua figura de La Hire — a geometria do núcleo. A reta (k=2) é o PISO, o mais ligado (o ferro). Cada floco fundido sobe a escada k=2,3,4,… : acrescenta uma cúspide e ABRE a curva da reta ao disco (ρ→1) — o núcleo engorda e afrouxa, além do ferro. A nucleossíntese é subir cúspides; a reta mineral é o par de Tusi (k=2) dessa escada. E as duas rotações de cada degrau são os 2 modos da POD (órbita temporal).", "epistemico": "inferencia", "exemplos": [], "id": "escada_de_hipociclos", "memoria": "semantica", "meta": {"dominio": "imagem", "fonte": "la hire / nucleossíntese"}, "origem": "wiki", "palavras_chave": [], "sinonimos": [], "tipo": "conceito", "titulo": "A Escada de Hipociclóides (fundir sobe as cúspides)"}
\section{A Escada de Hipociclóides (fundir sobe as cúspides)}
Generalizar La Hire à nucleossíntese: o k\_significativo de um floco é o nº de CÚSPIDES da sua figura de La Hire — a geometria do núcleo. A reta (k=2) é o PISO, o mais ligado (o ferro). Cada floco fundido sobe a escada k=2,3,4,… : acrescenta uma cúspide e ABRE a curva da reta ao disco (\ensuremath{\rho}\ensuremath{\to}1) — o núcleo engorda e afrouxa, além do ferro. A nucleossíntese é subir cúspides; a reta mineral é o par de Tusi (k=2) dessa escada. E as duas rotações de cada degrau são os 2 modos da POD (órbita temporal).
\prov{relações: usa mecanismo\_de\_la\_hire; parte\_de nucleossintese\_dos\_flocos; usa movimento\_circular\_dois\_modos; relaciona vale\_do\_ferro; usa corpo\_mineral; usa corpo\_estelar; relaciona fusao\_nuclear; relaciona familia\_m-bonacci\_gap\_k}
\prov{proveniência: wiki \textperiodcentered{} la hire / nucleossíntese \textperiodcentered{} inferencia \textperiodcentered{} confiança 1.0 \textperiodcentered{} domínio imagem \textperiodcentered{} história 11 versões no jornal}

%CRISTAL {"arestas": [{"alvo": "balanceamento_de_carga", "confianca": 1.0, "epistemico": "fato", "origem": "wiki", "rel": "usa"}, {"alvo": "utilizacao_e_borda", "confianca": 1.0, "epistemico": "fato", "origem": "wiki", "rel": "usa"}], "confianca": 1.0, "contraexemplos": [], "descricao": "Adicionar ou remover servidores automaticamente conforme a carga sobe e desce, para manter ρ na faixa boa (nem ocioso/caro, nem saturado/lento). É o controle que segura a nuvem na borda servível — o autoscaling como regulador do ρ mineral.", "epistemico": "fato", "exemplos": [], "id": "escalabilidade_elastica", "memoria": "semantica", "meta": {"dominio": "engenharia", "fonte": "nuvem e data centers"}, "origem": "wiki", "palavras_chave": [], "sinonimos": [], "tipo": "conceito", "titulo": "Escalabilidade Elástica"}
\section{Escalabilidade Elástica}
Adicionar ou remover servidores automaticamente conforme a carga sobe e desce, para manter \ensuremath{\rho} na faixa boa (nem ocioso/caro, nem saturado/lento). É o controle que segura a nuvem na borda servível — o autoscaling como regulador do \ensuremath{\rho} mineral.
\prov{relações: usa balanceamento\_de\_carga; usa utilizacao\_e\_borda}
\prov{proveniência: wiki \textperiodcentered{} nuvem e data centers \textperiodcentered{} fato \textperiodcentered{} confiança 1.0 \textperiodcentered{} domínio engenharia \textperiodcentered{} história 3 versões no jornal}

%CRISTAL {"arestas": [{"alvo": "acesso_multiplo", "confianca": 1.0, "epistemico": "fato", "origem": "wiki", "rel": "usa"}, {"alvo": "eng_telecomunicacoes", "confianca": 1.0, "epistemico": "fato", "origem": "curriculo", "rel": "parte_de"}], "confianca": 1.0, "contraexemplos": [], "descricao": "Espalhar o sinal por banda muito maior que a necessária, com um código — resistência a interferência e sigilo, e acesso múltiplo por código (CDMA). Do GPS ao 3G.", "epistemico": "fato", "exemplos": [], "id": "espalhamento_espectral", "memoria": "semantica", "meta": {"autor": "Lamarr", "dominio": "engenharia", "fonte": ""}, "origem": "wiki", "palavras_chave": [], "sinonimos": [], "tipo": "conceito", "titulo": "Espalhamento Espectral (Spread Spectrum)"}
\section{Espalhamento Espectral (Spread Spectrum)}
Espalhar o sinal por banda muito maior que a necessária, com um código — resistência a interferência e sigilo, e acesso múltiplo por código (CDMA). Do GPS ao 3G.
\prov{relações: usa acesso\_multiplo; parte\_de eng\_telecomunicacoes}
\prov{proveniência: wiki \textperiodcentered{} fato \textperiodcentered{} confiança 1.0 \textperiodcentered{} domínio engenharia \textperiodcentered{} história 3 versões no jornal}

%CRISTAL {"arestas": [{"alvo": "decomposicao_espectral_imagem", "confianca": 1.0, "epistemico": "fato", "origem": "wiki", "rel": "parte_de"}, {"alvo": "numeros_de_pisot", "confianca": 1.0, "epistemico": "fato", "origem": "wiki", "rel": "relaciona"}, {"alvo": "vibracao_e_colapso_dois_tempos", "confianca": 1.0, "epistemico": "fato", "origem": "wiki", "rel": "relaciona"}, {"alvo": "orbita_de_partida", "confianca": 1.0, "epistemico": "fato", "origem": "wiki", "rel": "relaciona"}], "confianca": 1.0, "contraexemplos": [], "descricao": "O espectro {σ_i} DECAI (a 1ª camada domina): é uma órbita cristal/Pisot (λ=log(σ_2/σ_1)<0, contração). A iteração de potência (AᵀA)^m projeta na camada dominante — o FUNDAMENTAL — que é o mesmo A^k da mecânica quântica mineral projetando no estado fundamental (tempo imaginário). O espectro da imagem obedece a mesma parábola: 1 camada domina (cristal) ou o espectro é plano (caos).", "epistemico": "fato", "exemplos": [], "id": "espectro_e_orbita_cristal", "memoria": "semantica", "meta": {"dominio": "imagem", "fonte": "decomposição espectral"}, "origem": "wiki", "palavras_chave": [], "sinonimos": [], "tipo": "conceito", "titulo": "O Espectro é uma Órbita Cristal"}
\section{O Espectro é uma Órbita Cristal}
O espectro \{\ensuremath{\sigma}\_i\} DECAI (a 1ª camada domina): é uma órbita cristal/Pisot (\ensuremath{\lambda}=log(\ensuremath{\sigma}\_2/\ensuremath{\sigma}\_1)<0, contração). A iteração de potência (A\ensuremath{^{T}}A)\^{}m projeta na camada dominante — o FUNDAMENTAL — que é o mesmo A\^{}k da mecânica quântica mineral projetando no estado fundamental (tempo imaginário). O espectro da imagem obedece a mesma parábola: 1 camada domina (cristal) ou o espectro é plano (caos).
\prov{relações: parte\_de decomposicao\_espectral\_imagem; relaciona numeros\_de\_pisot; relaciona vibracao\_e\_colapso\_dois\_tempos; relaciona orbita\_de\_partida}
\prov{proveniência: wiki \textperiodcentered{} decomposição espectral \textperiodcentered{} fato \textperiodcentered{} confiança 1.0 \textperiodcentered{} domínio imagem \textperiodcentered{} história 5 versões no jornal}

%CRISTAL {"arestas": [{"alvo": "largura_de_banda", "confianca": 1.0, "epistemico": "fato", "origem": "wiki", "rel": "usa"}, {"alvo": "capacidade_de_shannon", "confianca": 1.0, "epistemico": "fato", "origem": "wiki", "rel": "relaciona"}], "confianca": 1.0, "contraexemplos": [], "descricao": "O rádio compartilha o mesmo espectro entre todos; reusá-lo no espaço (células) e no tempo (slots) e não pedir mais que Shannon é o que mantém a rede na borda servível em vez de afogá-la em interferência.", "epistemico": "fato", "exemplos": [], "id": "espectro_eletromagnetico_uso", "memoria": "semantica", "meta": {"dominio": "engenharia", "fonte": "conectividade e 5G"}, "origem": "wiki", "palavras_chave": [], "sinonimos": [], "tipo": "conceito", "titulo": "Uso do Espectro"}
\section{Uso do Espectro}
O rádio compartilha o mesmo espectro entre todos; reusá-lo no espaço (células) e no tempo (slots) e não pedir mais que Shannon é o que mantém a rede na borda servível em vez de afogá-la em interferência.
\prov{relações: usa largura\_de\_banda; relaciona capacidade\_de\_shannon}
\prov{proveniência: wiki \textperiodcentered{} conectividade e 5G \textperiodcentered{} fato \textperiodcentered{} confiança 1.0 \textperiodcentered{} domínio engenharia \textperiodcentered{} história 3 versões no jornal}

%CRISTAL {"arestas": [{"alvo": "decomposicao_espectral_imagem", "confianca": 1.0, "epistemico": "fato", "origem": "wiki", "rel": "parte_de"}, {"alvo": "o_floco_como_selo_do_canal", "confianca": 1.0, "epistemico": "fato", "origem": "wiki", "rel": "eh_um"}, {"alvo": "floco_phi", "confianca": 1.0, "epistemico": "fato", "origem": "wiki", "rel": "relaciona"}, {"alvo": "impressao_koch_compressao", "confianca": 1.0, "epistemico": "fato", "origem": "wiki", "rel": "relaciona"}, {"alvo": "reconstrucao_purificada", "confianca": 1.0, "epistemico": "fato", "origem": "wiki", "rel": "relaciona"}], "confianca": 1.0, "contraexemplos": [], "descricao": "A assinatura de uma imagem é o seu ESPECTRO IRREDUTÍVEL: as menores k camadas que guardam ~95% da energia. É o FLOCO da obra — compacto, guarda a essência, e reconstrói o resto a partir de si. Guardar o floco é guardar a assinatura. É o mesmo floco que sela o canal e nomeia a banda, agora o selo de uma imagem. Verif.: 95% de uma obra da Catedral (posto 300) em ~59 camadas (20% dos dados).", "epistemico": "fato", "exemplos": [], "id": "espectro_irredutivel_floco", "memoria": "semantica", "meta": {"dominio": "imagem", "fonte": "decomposição espectral"}, "origem": "wiki", "palavras_chave": [], "sinonimos": [], "tipo": "conceito", "titulo": "O Floco da Obra (o espectro irredutível)"}
\section{O Floco da Obra (o espectro irredutível)}
A assinatura de uma imagem é o seu ESPECTRO IRREDUTÍVEL: as menores k camadas que guardam \~{}95\% da energia. É o FLOCO da obra — compacto, guarda a essência, e reconstrói o resto a partir de si. Guardar o floco é guardar a assinatura. É o mesmo floco que sela o canal e nomeia a banda, agora o selo de uma imagem. Verif.: 95\% de uma obra da Catedral (posto 300) em \~{}59 camadas (20\% dos dados).
\prov{relações: parte\_de decomposicao\_espectral\_imagem; eh\_um o\_floco\_como\_selo\_do\_canal; relaciona floco\_phi; relaciona impressao\_koch\_compressao; relaciona reconstrucao\_purificada}
\prov{proveniência: wiki \textperiodcentered{} decomposição espectral \textperiodcentered{} fato \textperiodcentered{} confiança 1.0 \textperiodcentered{} domínio imagem \textperiodcentered{} história 6 versões no jornal}

%CRISTAL {"arestas": [{"alvo": "biblioteca_formas_minerais", "confianca": 1.0, "epistemico": "fato", "origem": "wiki", "rel": "parte_de"}, {"alvo": "numero_de_salem", "confianca": 1.0, "epistemico": "fato", "origem": "wiki", "rel": "relaciona"}, {"alvo": "corpo_mineral", "confianca": 1.0, "epistemico": "fato", "origem": "wiki", "rel": "usa"}], "confianca": 1.0, "contraexemplos": [], "descricao": "A forma gerada colocando sementes no ângulo 2π·frac(σ_n) e raio √i — o girassol. Cada metal dá uma espiral: a do φ é a filotaxia ótima (a mais irracional, empacota melhor); prata e bronze dão outras. A reta mineral virando forma no plano.", "epistemico": "inferencia", "exemplos": [], "id": "espiral_metalica", "memoria": "semantica", "meta": {"dominio": "imagem", "fonte": "formas minerais"}, "origem": "wiki", "palavras_chave": [], "sinonimos": [], "tipo": "conceito", "titulo": "Espiral Metálica (filotaxia de σ_n)"}
\section{Espiral Metálica (filotaxia de \ensuremath{\sigma}\_n)}
A forma gerada colocando sementes no ângulo 2\ensuremath{\pi}\textperiodcentered frac(\ensuremath{\sigma}\_n) e raio \ensuremath{\sqrt{\,}}i — o girassol. Cada metal dá uma espiral: a do \ensuremath{\varphi} é a filotaxia ótima (a mais irracional, empacota melhor); prata e bronze dão outras. A reta mineral virando forma no plano.
\prov{relações: parte\_de biblioteca\_formas\_minerais; relaciona numero\_de\_salem; usa corpo\_mineral}
\prov{proveniência: wiki \textperiodcentered{} formas minerais \textperiodcentered{} inferencia \textperiodcentered{} confiança 1.0 \textperiodcentered{} domínio imagem \textperiodcentered{} história 4 versões no jornal}

%CRISTAL {"arestas": [{"alvo": "rach_4_passos", "confianca": 1.0, "epistemico": "fato", "origem": "wiki", "rel": "continua"}, {"alvo": "camada_rrc_nr", "confianca": 1.0, "epistemico": "fato", "origem": "wiki", "rel": "usa"}, {"alvo": "estados_rrc_nr", "confianca": 1.0, "epistemico": "fato", "origem": "wiki", "rel": "usa"}], "confianca": 1.0, "contraexemplos": [], "descricao": "Após vencer o RACH, o UE pede e recebe a conexão RRC (RRC Setup / Setup Complete), saindo de IDLE para CONNECTED. É o berço em que a sinalização NAS do registro viaja rumo ao núcleo.", "epistemico": "fato", "exemplos": [], "id": "estabelecimento_conexao_rrc", "memoria": "semantica", "meta": {"dominio": "engenharia", "fonte": ""}, "origem": "wiki", "palavras_chave": [], "sinonimos": [], "tipo": "conceito", "titulo": "Estabelecimento de Conexão RRC"}
\section{Estabelecimento de Conexão RRC}
Após vencer o RACH, o UE pede e recebe a conexão RRC (RRC Setup / Setup Complete), saindo de IDLE para CONNECTED. É o berço em que a sinalização NAS do registro viaja rumo ao núcleo.
\prov{relações: continua rach\_4\_passos; usa camada\_rrc\_nr; usa estados\_rrc\_nr}
\prov{proveniência: wiki \textperiodcentered{} fato \textperiodcentered{} confiança 1.0 \textperiodcentered{} domínio engenharia \textperiodcentered{} história 4 versões no jornal}

%CRISTAL {"arestas": [{"alvo": "registro_5g", "confianca": 1.0, "epistemico": "fato", "origem": "wiki", "rel": "continua"}, {"alvo": "funcoes_de_rede_5gc", "confianca": 1.0, "epistemico": "fato", "origem": "wiki", "rel": "usa"}, {"alvo": "qos_fluxos_5g", "confianca": 1.0, "epistemico": "fato", "origem": "wiki", "rel": "usa"}, {"alvo": "protocolo_ip", "confianca": 1.0, "epistemico": "fato", "origem": "wiki", "rel": "usa"}], "confianca": 1.0, "contraexemplos": [], "descricao": "Registro dá presença; a SESSÃO PDU dá DADOS. O UE pede uma sessão, o SMF a estabelece, escolhe o UPF, atribui endereço IP e instala os fluxos de QoS — o túnel fim-a-fim por onde o pacote de internet finalmente flui.", "epistemico": "fato", "exemplos": [], "id": "estabelecimento_sessao_pdu", "memoria": "semantica", "meta": {"dominio": "engenharia", "fonte": ""}, "origem": "wiki", "palavras_chave": [], "sinonimos": [], "tipo": "conceito", "titulo": "Estabelecimento de Sessão PDU"}
\section{Estabelecimento de Sessão PDU}
Registro dá presença; a SESSÃO PDU dá DADOS. O UE pede uma sessão, o SMF a estabelece, escolhe o UPF, atribui endereço IP e instala os fluxos de QoS — o túnel fim-a-fim por onde o pacote de internet finalmente flui.
\prov{relações: continua registro\_5g; usa funcoes\_de\_rede\_5gc; usa qos\_fluxos\_5g; usa protocolo\_ip}
\prov{proveniência: wiki \textperiodcentered{} fato \textperiodcentered{} confiança 1.0 \textperiodcentered{} domínio engenharia \textperiodcentered{} história 5 versões no jornal}

%CRISTAL {"arestas": [{"alvo": "rede_eletrica", "confianca": 1.0, "epistemico": "fato", "origem": "wiki", "rel": "parte_de"}, {"alvo": "lyapunov_da_orbita", "confianca": 1.0, "epistemico": "fato", "origem": "wiki", "rel": "eh_um"}, {"alvo": "frequencia_da_rede", "confianca": 1.0, "epistemico": "fato", "origem": "wiki", "rel": "usa"}], "confianca": 1.0, "contraexemplos": [], "descricao": "Após uma falta, o ângulo de potência do rotor oscila; com amortecimento ele volta ao sincronismo (λ<0, CRISTAL), marginal (λ=0, BORDA), ou diverge e perde o sincronismo (λ>0, CAOS = blackout). Manter a rede é ficar na borda estável — o Lyapunov do ângulo de potência.", "epistemico": "inferencia", "exemplos": [], "id": "estabilidade_da_rede", "memoria": "semantica", "meta": {"classes": {"amortecida": "cristal (estável — volta ao sincronismo)", "marginal": "borda (marginal)", "negativa": "caos (perde sincronismo → blackout)"}, "dominio": "engenharia", "fonte": "usinas e geração"}, "origem": "wiki", "palavras_chave": [], "sinonimos": [], "tipo": "conceito", "titulo": "Estabilidade da Rede (o Lyapunov do sincronismo)"}
\section{Estabilidade da Rede (o Lyapunov do sincronismo)}
Após uma falta, o ângulo de potência do rotor oscila; com amortecimento ele volta ao sincronismo (\ensuremath{\lambda}<0, CRISTAL), marginal (\ensuremath{\lambda}=0, BORDA), ou diverge e perde o sincronismo (\ensuremath{\lambda}>0, CAOS = blackout). Manter a rede é ficar na borda estável — o Lyapunov do ângulo de potência.
\prov{relações: parte\_de rede\_eletrica; eh\_um lyapunov\_da\_orbita; usa frequencia\_da\_rede}
\prov{proveniência: wiki \textperiodcentered{} usinas e geração \textperiodcentered{} inferencia \textperiodcentered{} confiança 1.0 \textperiodcentered{} domínio engenharia \textperiodcentered{} história 8 versões no jornal}

%CRISTAL {"arestas": [{"alvo": "sistema_eletrico_potencia", "confianca": 1.0, "epistemico": "fato", "origem": "wiki", "rel": "parte_de"}, {"alvo": "sistemas_de_controle", "confianca": 1.0, "epistemico": "fato", "origem": "wiki", "rel": "usa"}, {"alvo": "eng_sistemas_energia", "confianca": 1.0, "epistemico": "fato", "origem": "curriculo", "rel": "parte_de"}], "confianca": 1.0, "contraexemplos": [], "descricao": "A rede permanece sincronizada após um distúrbio? Estabilidade de ângulo, tensão e frequência — perdê-la é o caminho do apagão em cascata. O SEP é um sistema dinâmico gigante à beira do caos.", "epistemico": "fato", "exemplos": [], "id": "estabilidade_potencia", "memoria": "semantica", "meta": {"dominio": "engenharia", "fonte": ""}, "origem": "wiki", "palavras_chave": [], "sinonimos": [], "tipo": "conceito", "titulo": "Estabilidade de Sistemas de Potência"}
\section{Estabilidade de Sistemas de Potência}
A rede permanece sincronizada após um distúrbio? Estabilidade de ângulo, tensão e frequência — perdê-la é o caminho do apagão em cascata. O SEP é um sistema dinâmico gigante à beira do caos.
\prov{relações: parte\_de sistema\_eletrico\_potencia; usa sistemas\_de\_controle; parte\_de eng\_sistemas\_energia}
\prov{proveniência: wiki \textperiodcentered{} fato \textperiodcentered{} confiança 1.0 \textperiodcentered{} domínio engenharia \textperiodcentered{} história 4 versões no jornal}

%CRISTAL {"arestas": [{"alvo": "camada_rrc_nr", "confianca": 1.0, "epistemico": "fato", "origem": "wiki", "rel": "parte_de"}], "confianca": 1.0, "contraexemplos": [], "descricao": "Além de ocioso e conectado, o NR criou o RRC_INACTIVE: o contexto do UE fica guardado, então religar é quase instantâneo e barato em bateria — feito para o tráfego esporádico de IoT e apps móveis.", "epistemico": "fato", "exemplos": [], "id": "estados_rrc_nr", "memoria": "semantica", "meta": {"dominio": "engenharia", "fonte": ""}, "origem": "wiki", "palavras_chave": [], "sinonimos": [], "tipo": "conceito", "titulo": "Estados RRC (IDLE, INACTIVE, CONNECTED)"}
\section{Estados RRC (IDLE, INACTIVE, CONNECTED)}
Além de ocioso e conectado, o NR criou o RRC\_INACTIVE: o contexto do UE fica guardado, então religar é quase instantâneo e barato em bateria — feito para o tráfego esporádico de IoT e apps móveis.
\prov{relações: parte\_de camada\_rrc\_nr}
\prov{proveniência: wiki \textperiodcentered{} fato \textperiodcentered{} confiança 1.0 \textperiodcentered{} domínio engenharia \textperiodcentered{} história 2 versões no jornal}

%CRISTAL {"arestas": [{"alvo": "estrategia", "confianca": 1.0, "epistemico": "fato", "origem": "wiki", "rel": "relaciona"}, {"alvo": "estrategia_posicional", "confianca": 1.0, "epistemico": "fato", "origem": "wiki", "rel": "relaciona"}, {"alvo": "xadrez_bai", "confianca": 1.0, "epistemico": "fato", "origem": "wiki", "rel": "parte_de"}], "confianca": 1.0, "contraexemplos": [], "descricao": "A estratégia posicional do xadrez É a medida Φ do BAI: a avaliação = material (valor das peças) + posição (tabelas peça-casa). É a régua que desempata ramos de mesma força tática — a organização estrutural que o colapso otimiza quando não há mate à vista.", "epistemico": "inferencia", "exemplos": [], "id": "estrategia_como_phi", "memoria": "semantica", "meta": {"dominio": "ponte", "fonte": "xadrez↔BAI"}, "origem": "wiki", "palavras_chave": [], "sinonimos": [], "tipo": "conceito", "titulo": "Estratégia = Φ (a organização estrutural)"}
\section{Estratégia = \ensuremath{\Phi} (a organização estrutural)}
A estratégia posicional do xadrez É a medida \ensuremath{\Phi} do BAI: a avaliação = material (valor das peças) + posição (tabelas peça-casa). É a régua que desempata ramos de mesma força tática — a organização estrutural que o colapso otimiza quando não há mate à vista.
\prov{relações: relaciona estrategia; relaciona estrategia\_posicional; parte\_de xadrez\_bai}
\prov{proveniência: wiki \textperiodcentered{} xadrez\ensuremath{\leftrightarrow}BAI \textperiodcentered{} inferencia \textperiodcentered{} confiança 1.0 \textperiodcentered{} domínio ponte \textperiodcentered{} história 4 versões no jornal}

%CRISTAL {"arestas": [{"alvo": "camada_rrc_nr", "confianca": 1.0, "epistemico": "fato", "origem": "wiki", "rel": "usa"}, {"alvo": "funcoes_de_rede_5gc", "confianca": 1.0, "epistemico": "fato", "origem": "wiki", "rel": "usa"}], "confianca": 1.0, "contraexemplos": [], "descricao": "A sinalização direta entre o celular e o núcleo (AMF), transparente ao rádio: registro na rede, autenticação, mobilidade e gestão de sessão. Viaja encapsulada dentro das mensagens RRC.", "epistemico": "fato", "exemplos": [], "id": "estrato_nas_5g", "memoria": "semantica", "meta": {"dominio": "engenharia", "fonte": ""}, "origem": "wiki", "palavras_chave": [], "sinonimos": [], "tipo": "conceito", "titulo": "Estrato de Não-Acesso (NAS)"}
\section{Estrato de Não-Acesso (NAS)}
A sinalização direta entre o celular e o núcleo (AMF), transparente ao rádio: registro na rede, autenticação, mobilidade e gestão de sessão. Viaja encapsulada dentro das mensagens RRC.
\prov{relações: usa camada\_rrc\_nr; usa funcoes\_de\_rede\_5gc}
\prov{proveniência: wiki \textperiodcentered{} fato \textperiodcentered{} confiança 1.0 \textperiodcentered{} domínio engenharia \textperiodcentered{} história 3 versões no jornal}

%CRISTAL {"arestas": [{"alvo": "painel_estelar_obs", "confianca": 1.0, "epistemico": "fato", "origem": "wiki", "rel": "confirma"}, {"alvo": "tres_bracos_estelar", "confianca": 1.0, "epistemico": "fato", "origem": "wiki", "rel": "usa"}], "confianca": 1.0, "contraexemplos": [], "descricao": "painel_estelar_atlas.py roda o binário C, lê o JSON e emite PAINEL_ESTELAR.md + um SVG por fecho (GATO no centro com Q e a+c²=1, três braços rotulados Geometria(τ)/Mecânica(c²)/Termo(a)). A figura em estrela é o mapa dos quadrantes projetado no plano observável.", "epistemico": "fato", "exemplos": [], "id": "estrela_svg_atlas", "memoria": "semantica", "meta": {"autor": "broca-so", "dominio": "cockpit", "fonte": "ula/painel_estelar_atlas.py"}, "origem": "wiki", "palavras_chave": [], "sinonimos": [], "tipo": "conceito", "titulo": "O Atlas — a Estrela Renderizada"}
\section{O Atlas — a Estrela Renderizada}
painel\_estelar\_atlas.py roda o binário C, lê o JSON e emite PAINEL\_ESTELAR.md + um SVG por fecho (GATO no centro com Q e a+c\ensuremath{^{2}}=1, três braços rotulados Geometria(\ensuremath{\tau})/Mecânica(c\ensuremath{^{2}})/Termo(a)). A figura em estrela é o mapa dos quadrantes projetado no plano observável.
\prov{relações: confirma painel\_estelar\_obs; usa tres\_bracos\_estelar}
\prov{proveniência: wiki \textperiodcentered{} ula/painel\_estelar\_atlas.py \textperiodcentered{} fato \textperiodcentered{} confiança 1.0 \textperiodcentered{} domínio cockpit \textperiodcentered{} história 3 versões no jornal}

%CRISTAL {"arestas": [{"alvo": "fisica_do_estado_solido", "confianca": 1.0, "epistemico": "fato", "origem": "wiki", "rel": "parte_de"}], "confianca": 1.0, "contraexemplos": [], "descricao": "Num sólido, os níveis de energia dos elétrons se agrupam em BANDAS permitidas separadas por lacunas — o que decide se o material conduz.", "epistemico": "fato", "exemplos": [], "id": "estrutura_de_bandas", "memoria": "semantica", "meta": {"dominio": "engenharia", "fonte": ""}, "origem": "wiki", "palavras_chave": [], "sinonimos": [], "tipo": "conceito", "titulo": "Estrutura de Bandas"}
\section{Estrutura de Bandas}
Num sólido, os níveis de energia dos elétrons se agrupam em BANDAS permitidas separadas por lacunas — o que decide se o material conduz.
\prov{relações: parte\_de fisica\_do\_estado\_solido}
\prov{proveniência: wiki \textperiodcentered{} fato \textperiodcentered{} confiança 1.0 \textperiodcentered{} domínio engenharia \textperiodcentered{} história 2 versões no jornal}

%CRISTAL {"arestas": [{"alvo": "camada_fisica_nr", "confianca": 1.0, "epistemico": "fato", "origem": "wiki", "rel": "parte_de"}, {"alvo": "numerologia_nr", "confianca": 1.0, "epistemico": "fato", "origem": "wiki", "rel": "usa"}], "confianca": 1.0, "contraexemplos": [], "descricao": "O tempo do NR: quadro de 10 ms, subquadro de 1 ms, e o SLOT (14 símbolos OFDM) cuja duração encolhe com a numerologia. O mini-slot permite começar a transmitir no meio do slot — a chave da latência URLLC.", "epistemico": "fato", "exemplos": [], "id": "estrutura_quadro_nr", "memoria": "semantica", "meta": {"dominio": "engenharia", "fonte": ""}, "origem": "wiki", "palavras_chave": [], "sinonimos": [], "tipo": "conceito", "titulo": "Estrutura de Quadro (Frame/Slot)"}
\section{Estrutura de Quadro (Frame/Slot)}
O tempo do NR: quadro de 10 ms, subquadro de 1 ms, e o SLOT (14 símbolos OFDM) cuja duração encolhe com a numerologia. O mini-slot permite começar a transmitir no meio do slot — a chave da latência URLLC.
\prov{relações: parte\_de camada\_fisica\_nr; usa numerologia\_nr}
\prov{proveniência: wiki \textperiodcentered{} fato \textperiodcentered{} confiança 1.0 \textperiodcentered{} domínio engenharia \textperiodcentered{} história 3 versões no jornal}

%CRISTAL {"arestas": [{"alvo": "camada_de_enlace", "confianca": 1.0, "epistemico": "fato", "origem": "wiki", "rel": "especializa"}, {"alvo": "endereco_mac", "confianca": 1.0, "epistemico": "fato", "origem": "wiki", "rel": "usa"}, {"alvo": "acesso_ao_meio", "confianca": 1.0, "epistemico": "fato", "origem": "wiki", "rel": "usa"}, {"alvo": "eng_telecomunicacoes", "confianca": 1.0, "epistemico": "fato", "origem": "curriculo", "rel": "parte_de"}], "confianca": 1.0, "contraexemplos": [], "descricao": "O protocolo de enlace dominante em redes locais: quadros com MAC de origem/destino, CRC, e switching. Do cabo coaxial ao switch de fibra, é a LAN de fato do planeta.", "epistemico": "fato", "exemplos": [], "id": "ethernet", "memoria": "semantica", "meta": {"autor": "Metcalfe", "dominio": "engenharia", "fonte": ""}, "origem": "wiki", "palavras_chave": [], "sinonimos": [], "tipo": "conceito", "titulo": "Ethernet (IEEE 802.3)"}
\section{Ethernet (IEEE 802.3)}
O protocolo de enlace dominante em redes locais: quadros com MAC de origem/destino, CRC, e switching. Do cabo coaxial ao switch de fibra, é a LAN de fato do planeta.
\prov{relações: especializa camada\_de\_enlace; usa endereco\_mac; usa acesso\_ao\_meio; parte\_de eng\_telecomunicacoes}
\prov{proveniência: wiki \textperiodcentered{} fato \textperiodcentered{} confiança 1.0 \textperiodcentered{} domínio engenharia \textperiodcentered{} história 5 versões no jornal}

%CRISTAL {"arestas": [{"alvo": "selecao_natural", "confianca": 1.0, "epistemico": "fato", "origem": "wiki", "rel": "explica"}, {"alvo": "treino_atrator_lyapunov", "confianca": 1.0, "epistemico": "fato", "origem": "wiki", "rel": "relaciona"}, {"alvo": "lyapunov_da_orbita", "confianca": 1.0, "epistemico": "fato", "origem": "wiki", "rel": "usa"}, {"alvo": "mutacao_borda_eigen", "confianca": 1.0, "epistemico": "fato", "origem": "wiki", "rel": "usa"}], "confianca": 1.0, "contraexemplos": [], "descricao": "A seleção natural escala a paisagem de aptidão (ou desce a energia livre) rumo a máximos locais — os ATRATORES, as espécies como pontos fixos. É a mesma descida do treino da IA: variação (o ruído da mutação) + seleção (o gradiente) convergindo a um cristal, com a mutação mantendo a busca na borda.", "epistemico": "inferencia", "exemplos": [], "id": "evolucao_atrator", "memoria": "semantica", "meta": {"dominio": "biotecnologia", "fonte": "biotecnologia e engenharia genética"}, "origem": "wiki", "palavras_chave": [], "sinonimos": [], "tipo": "conceito", "titulo": "Evolução = Descida a um Atrator (paisagem de aptidão)"}
\section{Evolução = Descida a um Atrator (paisagem de aptidão)}
A seleção natural escala a paisagem de aptidão (ou desce a energia livre) rumo a máximos locais — os ATRATORES, as espécies como pontos fixos. É a mesma descida do treino da IA: variação (o ruído da mutação) + seleção (o gradiente) convergindo a um cristal, com a mutação mantendo a busca na borda.
\prov{relações: explica selecao\_natural; relaciona treino\_atrator\_lyapunov; usa lyapunov\_da\_orbita; usa mutacao\_borda\_eigen}
\prov{proveniência: wiki \textperiodcentered{} biotecnologia e engenharia genética \textperiodcentered{} inferencia \textperiodcentered{} confiança 1.0 \textperiodcentered{} domínio biotecnologia \textperiodcentered{} história 5 versões no jornal}

%CRISTAL {"arestas": [{"alvo": "transicoes_reta_bai_newton", "confianca": 1.0, "epistemico": "fato", "origem": "wiki", "rel": "parte_de"}, {"alvo": "lyapunov_da_orbita", "confianca": 1.0, "epistemico": "fato", "origem": "wiki", "rel": "eh_um"}, {"alvo": "numero_de_salem", "confianca": 1.0, "epistemico": "fato", "origem": "wiki", "rel": "relaciona"}, {"alvo": "cristal_e_a_convergencia", "confianca": 1.0, "epistemico": "fato", "origem": "wiki", "rel": "usa"}, {"alvo": "borda_e_o_julia", "confianca": 1.0, "epistemico": "fato", "origem": "wiki", "rel": "usa"}], "confianca": 1.0, "contraexemplos": [], "descricao": "O que classifica a órbita é sempre o EXPOENTE: λ=log σ na reta mineral (o Lyapunov do metal), o MULTIPLICADOR |f''/f'| no Newton (a taxa de convergência), o ORÇAMENTO δ no BAI (a profundidade até o fecho). Um só número decide cristal, borda ou caos — a parábola é a mesma régua nas três leituras.", "epistemico": "inferencia", "exemplos": [], "id": "expoente_unico", "memoria": "semantica", "meta": {"dominio": "ponte", "fonte": "transições órbita"}, "origem": "wiki", "palavras_chave": [], "sinonimos": [], "tipo": "conceito", "titulo": "O Expoente é Único (log σ ↔ multiplicador ↔ orçamento δ)"}
\section{O Expoente é Único (log \ensuremath{\sigma} \ensuremath{\leftrightarrow} multiplicador \ensuremath{\leftrightarrow} orçamento \ensuremath{\delta})}
O que classifica a órbita é sempre o EXPOENTE: \ensuremath{\lambda}=log \ensuremath{\sigma} na reta mineral (o Lyapunov do metal), o MULTIPLICADOR |f''/f'| no Newton (a taxa de convergência), o ORÇAMENTO \ensuremath{\delta} no BAI (a profundidade até o fecho). Um só número decide cristal, borda ou caos — a parábola é a mesma régua nas três leituras.
\prov{relações: parte\_de transicoes\_reta\_bai\_newton; eh\_um lyapunov\_da\_orbita; relaciona numero\_de\_salem; usa cristal\_e\_a\_convergencia; usa borda\_e\_o\_julia}
\prov{proveniência: wiki \textperiodcentered{} transições órbita \textperiodcentered{} inferencia \textperiodcentered{} confiança 1.0 \textperiodcentered{} domínio ponte \textperiodcentered{} história 6 versões no jornal}

%CRISTAL {"arestas": [{"alvo": "camada_fisica_nr", "confianca": 1.0, "epistemico": "fato", "origem": "wiki", "rel": "usa"}, {"alvo": "propagacao_de_ondas", "confianca": 1.0, "epistemico": "fato", "origem": "wiki", "rel": "usa"}], "confianca": 1.0, "contraexemplos": [], "descricao": "FR1 (sub-6 GHz) dá alcance e penetração; FR2 (ondas milimétricas, 24–52 GHz) dá banda enorme mas mal atravessa parede ou chuva — só funciona com feixes estreitos e células densas.", "epistemico": "fato", "exemplos": [], "id": "faixas_fr1_fr2", "memoria": "semantica", "meta": {"dominio": "engenharia", "fonte": ""}, "origem": "wiki", "palavras_chave": [], "sinonimos": [], "tipo": "conceito", "titulo": "Faixas de Frequência (FR1 e FR2/mmWave)"}
\section{Faixas de Frequência (FR1 e FR2/mmWave)}
FR1 (sub-6 GHz) dá alcance e penetração; FR2 (ondas milimétricas, 24–52 GHz) dá banda enorme mas mal atravessa parede ou chuva — só funciona com feixes estreitos e células densas.
\prov{relações: usa camada\_fisica\_nr; usa propagacao\_de\_ondas}
\prov{proveniência: wiki \textperiodcentered{} fato \textperiodcentered{} confiança 1.0 \textperiodcentered{} domínio engenharia \textperiodcentered{} história 3 versões no jornal}

%CRISTAL {"arestas": [{"alvo": "handover_nr", "confianca": 1.0, "epistemico": "fato", "origem": "wiki", "rel": "relaciona"}, {"alvo": "estabelecimento_conexao_rrc", "confianca": 1.0, "epistemico": "fato", "origem": "wiki", "rel": "usa"}], "confianca": 1.0, "contraexemplos": [], "descricao": "Quando o rádio cai antes do handover completar (sombra, bloqueio), o UE declara RLF e tenta se RE-ESTABELECER numa célula adequada — reusando o contexto guardado, para não recomeçar do zero. A rede de segurança da mobilidade.", "epistemico": "fato", "exemplos": [], "id": "falha_enlace_radio_rlf", "memoria": "semantica", "meta": {"dominio": "engenharia", "fonte": ""}, "origem": "wiki", "palavras_chave": [], "sinonimos": [], "tipo": "conceito", "titulo": "Falha de Enlace e Re-estabelecimento (RLF)"}
\section{Falha de Enlace e Re-estabelecimento (RLF)}
Quando o rádio cai antes do handover completar (sombra, bloqueio), o UE declara RLF e tenta se RE-ESTABELECER numa célula adequada — reusando o contexto guardado, para não recomeçar do zero. A rede de segurança da mobilidade.
\prov{relações: relaciona handover\_nr; usa estabelecimento\_conexao\_rrc}
\prov{proveniência: wiki \textperiodcentered{} fato \textperiodcentered{} confiança 1.0 \textperiodcentered{} domínio engenharia \textperiodcentered{} história 3 versões no jornal}

%CRISTAL {"arestas": [{"alvo": "impressao_3d", "confianca": 1.0, "epistemico": "fato", "origem": "wiki", "rel": "usa"}, {"alvo": "amostragem", "confianca": 1.0, "epistemico": "fato", "origem": "wiki", "rel": "eh_um"}, {"alvo": "corpo_mineral", "confianca": 1.0, "epistemico": "fato", "origem": "wiki", "rel": "usa"}], "confianca": 1.0, "contraexemplos": [], "descricao": "O slicer corta o modelo 3D contínuo em camadas de altura h — a mesma discretização da amostragem de Nyquist. h é a RESOLUÇÃO; o degrau (a escada na superfície) é o erro de quantização da forma. h fino = fiel/liso (cristal); h grosso = degraus grosseiros (aliasing da forma, caos). Verif.: h=0.05 mm cristal, h=0.5 mm caos.", "epistemico": "inferencia", "exemplos": [], "id": "fatiamento_amostragem", "memoria": "semantica", "meta": {"classe_h005": "cristal (fiel, liso)", "classe_h05": "caos (degraus grosseiros)", "degrau_h005": 0.025, "dominio": "manufatura", "fonte": "impressão 3D e manufatura"}, "origem": "wiki", "palavras_chave": [], "sinonimos": [], "tipo": "conceito", "titulo": "Fatiar = Amostrar / Quantizar a Forma"}
\section{Fatiar = Amostrar / Quantizar a Forma}
O slicer corta o modelo 3D contínuo em camadas de altura h — a mesma discretização da amostragem de Nyquist. h é a RESOLUÇÃO; o degrau (a escada na superfície) é o erro de quantização da forma. h fino = fiel/liso (cristal); h grosso = degraus grosseiros (aliasing da forma, caos). Verif.: h=0.05 mm cristal, h=0.5 mm caos.
\prov{relações: usa impressao\_3d; eh\_um amostragem; usa corpo\_mineral}
\prov{proveniência: wiki \textperiodcentered{} impressão 3D e manufatura \textperiodcentered{} inferencia \textperiodcentered{} confiança 1.0 \textperiodcentered{} domínio manufatura \textperiodcentered{} história 4 versões no jornal}

%CRISTAL {"arestas": [{"alvo": "nucleo_estelar", "confianca": 1.0, "epistemico": "fato", "origem": "wiki", "rel": "usa"}, {"alvo": "olho_de_koch", "confianca": 1.0, "epistemico": "fato", "origem": "wiki", "rel": "usa"}], "confianca": 1.0, "contraexemplos": [], "descricao": "Habitantes cuja leitura fecha no atrator — os três braços convergem ao mesmo Q=1. Os nonces de referência 1572863 e 1966079 (habitante [295,143]) têm fecho=1 e núcleos idênticos (a coincidência é real, não erro de cópia). Scan de 2M citado no paper.", "epistemico": "fato", "exemplos": [], "id": "fecho_koch_painel", "memoria": "semantica", "meta": {"autor": "broca-so", "dominio": "cockpit", "fonte": "ula/painel_estelar_exp.c; ula/dados/painel_estelar.json"}, "origem": "wiki", "palavras_chave": [], "sinonimos": [], "tipo": "conceito", "titulo": "Fechos Koch no Painel (Q=1 convergente)"}
\section{Fechos Koch no Painel (Q=1 convergente)}
Habitantes cuja leitura fecha no atrator — os três braços convergem ao mesmo Q=1. Os nonces de referência 1572863 e 1966079 (habitante [295,143]) têm fecho=1 e núcleos idênticos (a coincidência é real, não erro de cópia). Scan de 2M citado no paper.
\prov{relações: usa nucleo\_estelar; usa olho\_de\_koch}
\prov{proveniência: wiki \textperiodcentered{} ula/painel\_estelar\_exp.c; ula/dados/painel\_estelar.json \textperiodcentered{} fato \textperiodcentered{} confiança 1.0 \textperiodcentered{} domínio cockpit \textperiodcentered{} história 3 versões no jornal}

%CRISTAL {"arestas": [{"alvo": "broca_rede_social_soberana", "confianca": 1.0, "epistemico": "fato", "origem": "wiki", "rel": "parte_de"}, {"alvo": "base_cliente_por_banda", "confianca": 1.0, "epistemico": "fato", "origem": "wiki", "rel": "usa"}, {"alvo": "software_arte_organismo", "confianca": 1.0, "epistemico": "fato", "origem": "wiki", "rel": "usa"}, {"alvo": "bolha_de_filtro", "confianca": 1.0, "epistemico": "fato", "origem": "wiki", "rel": "contradiz"}], "confianca": 1.0, "contraexemplos": [], "descricao": "O fluxo de apps e conteúdo que circula numa comunidade: quem sintoniza a banda vê, as sementes se propagam e evoluem entre comunidades. É o 'feed' das redes sociais modernas — mas ordenado por PERTENCIMENTO (o tecido), não por um algoritmo de engajamento; sem vigilância, sem bolha de filtro, sem economia da atenção. Onde o cliente vê e partilha o que a sua comunidade guarda.", "epistemico": "inferencia", "exemplos": [], "id": "feed_de_comunidade", "memoria": "semantica", "meta": {"dominio": "arquitetura", "fonte": "semilla_broca_rede_social"}, "origem": "wiki", "palavras_chave": [], "sinonimos": [], "tipo": "conceito", "titulo": "Feed de Comunidade (por banda, não por algoritmo)"}
\section{Feed de Comunidade (por banda, não por algoritmo)}
O fluxo de apps e conteúdo que circula numa comunidade: quem sintoniza a banda vê, as sementes se propagam e evoluem entre comunidades. É o 'feed' das redes sociais modernas — mas ordenado por PERTENCIMENTO (o tecido), não por um algoritmo de engajamento; sem vigilância, sem bolha de filtro, sem economia da atenção. Onde o cliente vê e partilha o que a sua comunidade guarda.
\prov{relações: parte\_de broca\_rede\_social\_soberana; usa base\_cliente\_por\_banda; usa software\_arte\_organismo; contradiz bolha\_de\_filtro}
\prov{proveniência: wiki \textperiodcentered{} semilla\_broca\_rede\_social \textperiodcentered{} inferencia \textperiodcentered{} confiança 1.0 \textperiodcentered{} domínio arquitetura \textperiodcentered{} história 10 versões no jornal}

%CRISTAL {"arestas": [{"alvo": "data_center", "confianca": 1.0, "epistemico": "fato", "origem": "wiki", "rel": "usa"}, {"alvo": "orbita_de_partida", "confianca": 1.0, "epistemico": "fato", "origem": "wiki", "rel": "relaciona"}], "confianca": 1.0, "contraexemplos": [], "descricao": "As requisições chegam a taxa λ e são servidas a taxa μ. A fila é uma órbita: cresce quando chegam mais do que saem. A teoria de filas (M/M/1) dá a latência e o tamanho médio a partir de λ e μ.", "epistemico": "fato", "exemplos": [], "id": "fila_de_requisicoes", "memoria": "semantica", "meta": {"dominio": "engenharia", "fonte": "nuvem e data centers"}, "origem": "wiki", "palavras_chave": [], "sinonimos": [], "tipo": "conceito", "titulo": "Fila de Requisições (M/M/1)"}
\section{Fila de Requisições (M/M/1)}
As requisições chegam a taxa \ensuremath{\lambda} e são servidas a taxa \ensuremath{\mu}. A fila é uma órbita: cresce quando chegam mais do que saem. A teoria de filas (M/M/1) dá a latência e o tamanho médio a partir de \ensuremath{\lambda} e \ensuremath{\mu}.
\prov{relações: usa data\_center; relaciona orbita\_de\_partida}
\prov{proveniência: wiki \textperiodcentered{} nuvem e data centers \textperiodcentered{} fato \textperiodcentered{} confiança 1.0 \textperiodcentered{} domínio engenharia \textperiodcentered{} história 3 versões no jornal}

%CRISTAL {"arestas": [{"alvo": "processamento_digital_sinais", "confianca": 1.0, "epistemico": "fato", "origem": "wiki", "rel": "usa"}, {"alvo": "transformada_de_fourier", "confianca": 1.0, "epistemico": "fato", "origem": "wiki", "rel": "usa"}, {"alvo": "eng_eletronica", "confianca": 1.0, "epistemico": "fato", "origem": "curriculo", "rel": "parte_de"}], "confianca": 1.0, "contraexemplos": [], "descricao": "Deixar passar certas frequências e barrar outras — passa-baixa, passa-faixa. O bisturi do domínio da frequência: separa sinal de ruído, banda de banda.", "epistemico": "fato", "exemplos": [], "id": "filtro", "memoria": "semantica", "meta": {"dominio": "engenharia", "fonte": ""}, "origem": "wiki", "palavras_chave": [], "sinonimos": [], "tipo": "conceito", "titulo": "Filtros (Analógicos e Digitais)"}
\section{Filtros (Analógicos e Digitais)}
Deixar passar certas frequências e barrar outras — passa-baixa, passa-faixa. O bisturi do domínio da frequência: separa sinal de ruído, banda de banda.
\prov{relações: usa processamento\_digital\_sinais; usa transformada\_de\_fourier; parte\_de eng\_eletronica}
\prov{proveniência: wiki \textperiodcentered{} fato \textperiodcentered{} confiança 1.0 \textperiodcentered{} domínio engenharia \textperiodcentered{} história 4 versões no jornal}

%CRISTAL {"arestas": [{"alvo": "seguranca_de_redes", "confianca": 1.0, "epistemico": "fato", "origem": "wiki", "rel": "parte_de"}, {"alvo": "eng_telecomunicacoes", "confianca": 1.0, "epistemico": "fato", "origem": "curriculo", "rel": "parte_de"}], "confianca": 1.0, "contraexemplos": [], "descricao": "O porteiro que filtra tráfego por regras (porta, IP, estado) — a fronteira entre o que se confia e o que não. Primeira linha de defesa perimetral da rede.", "epistemico": "inferencia", "exemplos": [], "id": "firewall", "memoria": "semantica", "meta": {"dominio": "engenharia", "fonte": ""}, "origem": "wiki", "palavras_chave": [], "sinonimos": [], "tipo": "conceito", "titulo": "Firewall"}
\section{Firewall}
O porteiro que filtra tráfego por regras (porta, IP, estado) — a fronteira entre o que se confia e o que não. Primeira linha de defesa perimetral da rede.
\prov{relações: parte\_de seguranca\_de\_redes; parte\_de eng\_telecomunicacoes}
\prov{proveniência: wiki \textperiodcentered{} inferencia \textperiodcentered{} confiança 1.0 \textperiodcentered{} domínio engenharia \textperiodcentered{} história 3 versões no jornal}

%CRISTAL {"arestas": [{"alvo": "cristalografia", "confianca": 1.0, "epistemico": "fato", "origem": "wiki", "rel": "usa"}, {"alvo": "fisica", "confianca": 1.0, "epistemico": "fato", "origem": "wiki", "rel": "parte_de"}], "confianca": 1.0, "contraexemplos": [], "descricao": "O estudo das propriedades da matéria sólida a partir de sua estrutura cristalina e dos elétrons — a base física da eletrônica.", "epistemico": "fato", "exemplos": [], "id": "fisica_do_estado_solido", "memoria": "semantica", "meta": {"dominio": "engenharia", "fonte": ""}, "origem": "wiki", "palavras_chave": [], "sinonimos": [], "tipo": "conceito", "titulo": "Física do Estado Sólido"}
\section{Física do Estado Sólido}
O estudo das propriedades da matéria sólida a partir de sua estrutura cristalina e dos elétrons — a base física da eletrônica.
\prov{relações: usa cristalografia; parte\_de fisica}
\prov{proveniência: wiki \textperiodcentered{} fato \textperiodcentered{} confiança 1.0 \textperiodcentered{} domínio engenharia \textperiodcentered{} história 3 versões no jornal}

%CRISTAL {"arestas": [{"alvo": "fisica", "confianca": 1.0, "epistemico": "fato", "origem": "wiki", "rel": "parte_de"}, {"alvo": "eng_eletronica", "confianca": 1.0, "epistemico": "fato", "origem": "curriculo", "rel": "parte_de"}, {"alvo": "eng_telecomunicacoes", "confianca": 1.0, "epistemico": "fato", "origem": "curriculo", "rel": "parte_de"}, {"alvo": "eng_sistemas_energia", "confianca": 1.0, "epistemico": "fato", "origem": "curriculo", "rel": "parte_de"}], "confianca": 1.0, "contraexemplos": [], "descricao": "Oscilação, onda, campo elétrico e magnético — a física que a engenharia elétrica encarna. Ponte entre a mecânica e as equações de Maxwell.", "epistemico": "fato", "exemplos": [], "id": "fisica_ondas_eng", "memoria": "semantica", "meta": {"dominio": "engenharia", "fonte": ""}, "origem": "wiki", "palavras_chave": [], "sinonimos": [], "tipo": "conceito", "titulo": "Física — Ondas e Eletromagnetismo"}
\section{Física — Ondas e Eletromagnetismo}
Oscilação, onda, campo elétrico e magnético — a física que a engenharia elétrica encarna. Ponte entre a mecânica e as equações de Maxwell.
\prov{relações: parte\_de fisica; parte\_de eng\_eletronica; parte\_de eng\_telecomunicacoes; parte\_de eng\_sistemas\_energia}
\prov{proveniência: wiki \textperiodcentered{} fato \textperiodcentered{} confiança 1.0 \textperiodcentered{} domínio engenharia \textperiodcentered{} história 5 versões no jornal}

%CRISTAL {"arestas": [{"alvo": "algebra_posicional_unificada", "confianca": 1.0, "epistemico": "fato", "origem": "wiki", "rel": "usa"}, {"alvo": "refino_continuo_floco_metais", "confianca": 1.0, "epistemico": "fato", "origem": "wiki", "rel": "parte_de"}, {"alvo": "beta_numeracao_metalica", "confianca": 1.0, "epistemico": "fato", "origem": "wiki", "rel": "usa"}, {"alvo": "tiffany", "confianca": 1.0, "epistemico": "fato", "origem": "wiki", "rel": "parte_de"}], "confianca": 1.0, "contraexemplos": [], "descricao": "O ganho de produção do M×M (agente produção, agora implementado): cada floco extraído fica ENDEREÇÁVEL por (metal n, posição k) — a coordenada da álgebra posicional. Busca por POSIÇÃO é salto direto O(1) (metalabem(n,k)); por VALOR é varredura O(k) (como buscaestrela). Cabeça de cadeia por metal em L['metais'][n]=emitidos,ab; ponte de continuidade: a prata (n=2) retoma a cadeia Pell antiga (L['pellab']) sem reiniciar. As 12 lanes viraram 12 metais reais (ouro n=1 … metal₁2), todos classificados 'cristal' (todo metálico grau 2 é Pisot). É a base para definir e PROCURAR um floco metálico específico. [I] o endereço (n,k) como chave do acervo metálico.", "epistemico": "inferencia", "exemplos": [], "id": "floco_enderecavel_metal_posicao", "memoria": "semantica", "meta": {"autor": "Lopes/broca-so", "dominio": "arquitetura", "fonte": "goldchain/minera.py:emite_metal (L['metais'][n]); metal_pos.metal_ab_em; symbolic.py:status"}, "origem": "wiki", "palavras_chave": [], "sinonimos": [], "tipo": "conceito", "titulo": "Todo floco endereçável por (metal n, posição) — a 'procura por um floco específico'"}
\section{Todo floco endereçável por (metal n, posição) — a 'procura por um floco específico'}
O ganho de produção do M$\times$M (agente produção, agora implementado): cada floco extraído fica ENDEREÇÁVEL por (metal n, posição k) — a coordenada da álgebra posicional. Busca por POSIÇÃO é salto direto O(1) (metalabem(n,k)); por VALOR é varredura O(k) (como buscaestrela). Cabeça de cadeia por metal em L['metais'][n]=emitidos,ab; ponte de continuidade: a prata (n=2) retoma a cadeia Pell antiga (L['pellab']) sem reiniciar. As 12 lanes viraram 12 metais reais (ouro n=1 … metal\ensuremath{_{1}}2), todos classificados 'cristal' (todo metálico grau 2 é Pisot). É a base para definir e PROCURAR um floco metálico específico. [I] o endereço (n,k) como chave do acervo metálico.
\prov{relações: usa algebra\_posicional\_unificada; parte\_de refino\_continuo\_floco\_metais; usa beta\_numeracao\_metalica; parte\_de tiffany}
\prov{proveniência: wiki \textperiodcentered{} goldchain/minera.py:emite\_metal (L['metais'][n]); metal\_pos.metal\_ab\_em; symbolic.py:status \textperiodcentered{} inferencia \textperiodcentered{} confiança 1.0 \textperiodcentered{} domínio arquitetura \textperiodcentered{} história 32 versões no jornal}

%CRISTAL {"arestas": [{"alvo": "espectro_irredutivel_floco", "confianca": 1.0, "epistemico": "fato", "origem": "wiki", "rel": "usa"}, {"alvo": "termodinamica_reta_mineral", "confianca": 1.0, "epistemico": "fato", "origem": "wiki", "rel": "usa"}, {"alvo": "o_floco_como_selo_do_canal", "confianca": 1.0, "epistemico": "fato", "origem": "wiki", "rel": "relaciona"}], "confianca": 1.0, "contraexemplos": [], "descricao": "A decomposição espectral faz do floco um RECIPIENTE de energia, e vale a 1ª lei: a energia total (Σσ_i²) divide-se em TRABALHO (a energia retida, recuperável pela reconstrução) + CALOR (a energia descartada na truncagem, dissipada). W(k) + Q(k) = 1 em toda compressão. O floco guarda o trabalho; a truncagem é o calor que escapa.", "epistemico": "fato", "exemplos": [], "id": "floco_recipiente_termico", "memoria": "semantica", "meta": {"dominio": "imagem", "fonte": "floco térmico"}, "origem": "wiki", "palavras_chave": [], "sinonimos": [], "tipo": "conceito", "titulo": "O Floco é um Recipiente Térmico de Trabalho"}
\section{O Floco é um Recipiente Térmico de Trabalho}
A decomposição espectral faz do floco um RECIPIENTE de energia, e vale a 1ª lei: a energia total (\ensuremath{\Sigma}\ensuremath{\sigma}\_i\ensuremath{^{2}}) divide-se em TRABALHO (a energia retida, recuperável pela reconstrução) + CALOR (a energia descartada na truncagem, dissipada). W(k) + Q(k) = 1 em toda compressão. O floco guarda o trabalho; a truncagem é o calor que escapa.
\prov{relações: usa espectro\_irredutivel\_floco; usa termodinamica\_reta\_mineral; relaciona o\_floco\_como\_selo\_do\_canal}
\prov{proveniência: wiki \textperiodcentered{} floco térmico \textperiodcentered{} fato \textperiodcentered{} confiança 1.0 \textperiodcentered{} domínio imagem \textperiodcentered{} história 8 versões no jornal}

%CRISTAL {"arestas": [{"alvo": "sistema_eletrico_potencia", "confianca": 1.0, "epistemico": "fato", "origem": "wiki", "rel": "parte_de"}, {"alvo": "potencia_ativa_reativa", "confianca": 1.0, "epistemico": "fato", "origem": "wiki", "rel": "usa"}, {"alvo": "eng_sistemas_energia", "confianca": 1.0, "epistemico": "fato", "origem": "curriculo", "rel": "parte_de"}], "confianca": 1.0, "contraexemplos": [], "descricao": "Calcular tensões, ângulos e fluxos de potência em regime permanente na malha — o estudo básico de operação e planejamento do SEP. Resolve-se por Newton-Raphson sobre equações não-lineares.", "epistemico": "fato", "exemplos": [], "id": "fluxo_de_carga", "memoria": "semantica", "meta": {"autor": "Newton-Raphson", "dominio": "engenharia", "fonte": ""}, "origem": "wiki", "palavras_chave": [], "sinonimos": [], "tipo": "conceito", "titulo": "Fluxo de Carga (Load Flow)"}
\section{Fluxo de Carga (Load Flow)}
Calcular tensões, ângulos e fluxos de potência em regime permanente na malha — o estudo básico de operação e planejamento do SEP. Resolve-se por Newton-Raphson sobre equações não-lineares.
\prov{relações: parte\_de sistema\_eletrico\_potencia; usa potencia\_ativa\_reativa; parte\_de eng\_sistemas\_energia}
\prov{proveniência: wiki \textperiodcentered{} fato \textperiodcentered{} confiança 1.0 \textperiodcentered{} domínio engenharia \textperiodcentered{} história 4 versões no jornal}

%CRISTAL {"arestas": [{"alvo": "ofdm", "confianca": 1.0, "epistemico": "fato", "origem": "wiki", "rel": "especializa"}, {"alvo": "camada_fisica_nr", "confianca": 1.0, "epistemico": "fato", "origem": "wiki", "rel": "parte_de"}], "confianca": 1.0, "contraexemplos": [], "descricao": "Descida e subida usam CP-OFDM; a subida pode ainda usar DFT-s-OFDM (SC-FDMA), que reduz o pico de potência — poupa a bateria e o amplificador do celular na borda da célula.", "epistemico": "fato", "exemplos": [], "id": "formas_de_onda_nr", "memoria": "semantica", "meta": {"dominio": "engenharia", "fonte": ""}, "origem": "wiki", "palavras_chave": [], "sinonimos": [], "tipo": "conceito", "titulo": "Formas de Onda (CP-OFDM e DFT-s-OFDM)"}
\section{Formas de Onda (CP-OFDM e DFT-s-OFDM)}
Descida e subida usam CP-OFDM; a subida pode ainda usar DFT-s-OFDM (SC-FDMA), que reduz o pico de potência — poupa a bateria e o amplificador do celular na borda da célula.
\prov{relações: especializa ofdm; parte\_de camada\_fisica\_nr}
\prov{proveniência: wiki \textperiodcentered{} fato \textperiodcentered{} confiança 1.0 \textperiodcentered{} domínio engenharia \textperiodcentered{} história 3 versões no jornal}

%CRISTAL {"arestas": [{"alvo": "porta_logica", "confianca": 1.0, "epistemico": "fato", "origem": "wiki", "rel": "usa"}, {"alvo": "linguagem_matriz_gatos", "confianca": 1.0, "epistemico": "fato", "origem": "wiki", "rel": "explica"}, {"alvo": "gato_2x2_e_fisico_nao_incidental", "confianca": 1.0, "epistemico": "fato", "origem": "wiki", "rel": "relaciona"}, {"alvo": "linguagem_descricao_hardware", "confianca": 1.0, "epistemico": "fato", "origem": "wiki", "rel": "usa"}, {"alvo": "circuito_integrado", "confianca": 1.0, "epistemico": "fato", "origem": "wiki", "rel": "especializa"}, {"alvo": "eng_eletronica", "confianca": 1.0, "epistemico": "fato", "origem": "curriculo", "rel": "parte_de"}], "confianca": 1.0, "contraexemplos": [], "descricao": "Um chip de lógica reconfigurável: blocos e interconexões que o HDL programa em campo. Prototipar hardware sem fabricar — o meio-termo entre software e ASIC.", "epistemico": "fato", "exemplos": [], "id": "fpga", "memoria": "semantica", "meta": {"dominio": "engenharia", "fonte": ""}, "origem": "wiki", "palavras_chave": [], "sinonimos": [], "tipo": "conceito", "titulo": "FPGA"}
\section{FPGA}
Um chip de lógica reconfigurável: blocos e interconexões que o HDL programa em campo. Prototipar hardware sem fabricar — o meio-termo entre software e ASIC.
\prov{relações: usa porta\_logica; explica linguagem\_matriz\_gatos; relaciona gato\_2x2\_e\_fisico\_nao\_incidental; usa linguagem\_descricao\_hardware; especializa circuito\_integrado; parte\_de eng\_eletronica}
\prov{proveniência: wiki \textperiodcentered{} fato \textperiodcentered{} confiança 1.0 \textperiodcentered{} domínio engenharia \textperiodcentered{} história 9 versões no jornal}

%CRISTAL {"arestas": [{"alvo": "motor_eletrico_tracao", "confianca": 1.0, "epistemico": "fato", "origem": "wiki", "rel": "usa"}, {"alvo": "bateria_cristal_movel", "confianca": 1.0, "epistemico": "fato", "origem": "wiki", "rel": "usa"}, {"alvo": "orbita_de_partida", "confianca": 1.0, "epistemico": "fato", "origem": "wiki", "rel": "relaciona"}], "confianca": 1.0, "contraexemplos": [], "descricao": "Ao frear, o motor funciona como GERADOR e devolve a energia cinética (½mv²) à bateria em vez de torrá-la em calor no freio. É a órbita REVERSÍVEL KE ⇄ cristal; a eficiência η<1 (o resto é atrito). Verificado: a 100 km/h recupera ~70% dos ~580 kJ cinéticos.", "epistemico": "inferencia", "exemplos": [], "id": "frenagem_regenerativa", "memoria": "semantica", "meta": {"dominio": "engenharia", "eta_regen": 0.7, "fonte": "mobilidade elétrica", "recuperado_frac": 0.7}, "origem": "wiki", "palavras_chave": [], "sinonimos": [], "tipo": "conceito", "titulo": "Frenagem Regenerativa (a órbita reversível)"}
\section{Frenagem Regenerativa (a órbita reversível)}
Ao frear, o motor funciona como GERADOR e devolve a energia cinética (\ensuremath{\tfrac12}mv\ensuremath{^{2}}) à bateria em vez de torrá-la em calor no freio. É a órbita REVERSÍVEL KE \ensuremath{\rightleftarrows} cristal; a eficiência \ensuremath{\eta}<1 (o resto é atrito). Verificado: a 100 km/h recupera \~{}70\% dos \~{}580 kJ cinéticos.
\prov{relações: usa motor\_eletrico\_tracao; usa bateria\_cristal\_movel; relaciona orbita\_de\_partida}
\prov{proveniência: wiki \textperiodcentered{} mobilidade elétrica \textperiodcentered{} inferencia \textperiodcentered{} confiança 1.0 \textperiodcentered{} domínio engenharia \textperiodcentered{} história 4 versões no jornal}

%CRISTAL {"arestas": [{"alvo": "corrente_alternada", "confianca": 1.0, "epistemico": "fato", "origem": "wiki", "rel": "usa"}, {"alvo": "invariante_estelar", "confianca": 1.0, "epistemico": "fato", "origem": "wiki", "rel": "relaciona"}], "confianca": 1.0, "contraexemplos": [], "descricao": "60 Hz (ou 50): todo gerador gira em SINCRONISMO com essa rotação. A frequência é a 'norma' conservada da rede — se a carga excede a geração, ela cai; se sobra, sobe. Manter 60 Hz é balancear geração e consumo em tempo real.", "epistemico": "inferencia", "exemplos": [], "id": "frequencia_da_rede", "memoria": "semantica", "meta": {"dominio": "engenharia", "fonte": "usinas e geração"}, "origem": "wiki", "palavras_chave": [], "sinonimos": [], "tipo": "conceito", "titulo": "Frequência da Rede (a norma)"}
\section{Frequência da Rede (a norma)}
60 Hz (ou 50): todo gerador gira em SINCRONISMO com essa rotação. A frequência é a 'norma' conservada da rede — se a carga excede a geração, ela cai; se sobra, sobe. Manter 60 Hz é balancear geração e consumo em tempo real.
\prov{relações: usa corrente\_alternada; relaciona invariante\_estelar}
\prov{proveniência: wiki \textperiodcentered{} usinas e geração \textperiodcentered{} inferencia \textperiodcentered{} confiança 1.0 \textperiodcentered{} domínio engenharia \textperiodcentered{} história 5 versões no jornal}

%CRISTAL {"arestas": [{"alvo": "imagem_bai", "confianca": 1.0, "epistemico": "fato", "origem": "wiki", "rel": "parte_de"}, {"alvo": "bordas_fronteira", "confianca": 1.0, "epistemico": "fato", "origem": "wiki", "rel": "usa"}, {"alvo": "horizonte_e_fail_closed", "confianca": 1.0, "epistemico": "fato", "origem": "wiki", "rel": "relaciona"}], "confianca": 1.0, "contraexemplos": [], "descricao": "A fronteira onde o crescimento PARA — o gradiente alto que barra a propagação — É o fail-closed do BAI (o orçamento δ esgota, nega por padrão) e o efeito horizonte do xadrez (não vê além, decide conservador). O limite da região, o limite da autoridade e o limite do cálculo são o mesmo limite.", "epistemico": "inferencia", "exemplos": [], "id": "fronteira_e_fail_closed", "memoria": "semantica", "meta": {"dominio": "ponte", "fonte": "imagem↔BAI"}, "origem": "wiki", "palavras_chave": [], "sinonimos": [], "tipo": "conceito", "titulo": "Fronteira = Fail-Closed (o horizonte)"}
\section{Fronteira = Fail-Closed (o horizonte)}
A fronteira onde o crescimento PARA — o gradiente alto que barra a propagação — É o fail-closed do BAI (o orçamento \ensuremath{\delta} esgota, nega por padrão) e o efeito horizonte do xadrez (não vê além, decide conservador). O limite da região, o limite da autoridade e o limite do cálculo são o mesmo limite.
\prov{relações: parte\_de imagem\_bai; usa bordas\_fronteira; relaciona horizonte\_e\_fail\_closed}
\prov{proveniência: wiki \textperiodcentered{} imagem\ensuremath{\leftrightarrow}BAI \textperiodcentered{} inferencia \textperiodcentered{} confiança 1.0 \textperiodcentered{} domínio ponte \textperiodcentered{} história 4 versões no jornal}

%CRISTAL {"arestas": [{"alvo": "nucleo_5gc_sba", "confianca": 1.0, "epistemico": "fato", "origem": "wiki", "rel": "parte_de"}, {"alvo": "protocolo_ip", "confianca": 1.0, "epistemico": "fato", "origem": "wiki", "rel": "usa"}], "confianca": 1.0, "contraexemplos": [], "descricao": "As peças do núcleo: AMF (acesso e mobilidade), SMF (sessão), UPF (plano de usuário, encaminha os pacotes IP), AUSF/UDM (autenticação e dados do assinante), PCF (política), NRF (repositório/descoberta), NSSF (seleção de slice). Cada uma um serviço independente e escalável.", "epistemico": "fato", "exemplos": [], "id": "funcoes_de_rede_5gc", "memoria": "semantica", "meta": {"dominio": "engenharia", "fonte": ""}, "origem": "wiki", "palavras_chave": [], "sinonimos": [], "tipo": "conceito", "titulo": "Funções de Rede do 5GC (AMF, SMF, UPF, …)"}
\section{Funções de Rede do 5GC (AMF, SMF, UPF, …)}
As peças do núcleo: AMF (acesso e mobilidade), SMF (sessão), UPF (plano de usuário, encaminha os pacotes IP), AUSF/UDM (autenticação e dados do assinante), PCF (política), NRF (repositório/descoberta), NSSF (seleção de slice). Cada uma um serviço independente e escalável.
\prov{relações: parte\_de nucleo\_5gc\_sba; usa protocolo\_ip}
\prov{proveniência: wiki \textperiodcentered{} fato \textperiodcentered{} confiança 1.0 \textperiodcentered{} domínio engenharia \textperiodcentered{} história 3 versões no jornal}

%CRISTAL {"arestas": [{"alvo": "imagem_bai", "confianca": 1.0, "epistemico": "fato", "origem": "wiki", "rel": "parte_de"}, {"alvo": "minimax_e_colapso", "confianca": 1.0, "epistemico": "fato", "origem": "wiki", "rel": "relaciona"}, {"alvo": "segmentacao_de_imagem", "confianca": 1.0, "epistemico": "fato", "origem": "wiki", "rel": "usa"}], "confianca": 1.0, "contraexemplos": [], "descricao": "Fundir ou descartar a região fraca/dominada (supressão não-máxima, merge de segmentos pequenos) É o deny-overrides do BAI (a proibição domina) e a poda alfa-beta do xadrez (o ramo dominado é cortado). A região que perde energia cede à vizinha, como a concessão dominada cede à proibição.", "epistemico": "inferencia", "exemplos": [], "id": "fusao_e_deny_overrides", "memoria": "semantica", "meta": {"dominio": "ponte", "fonte": "imagem↔BAI"}, "origem": "wiki", "palavras_chave": [], "sinonimos": [], "tipo": "conceito", "titulo": "Fusão de Regiões = Deny-Overrides"}
\section{Fusão de Regiões = Deny-Overrides}
Fundir ou descartar a região fraca/dominada (supressão não-máxima, merge de segmentos pequenos) É o deny-overrides do BAI (a proibição domina) e a poda alfa-beta do xadrez (o ramo dominado é cortado). A região que perde energia cede à vizinha, como a concessão dominada cede à proibição.
\prov{relações: parte\_de imagem\_bai; relaciona minimax\_e\_colapso; usa segmentacao\_de\_imagem}
\prov{proveniência: wiki \textperiodcentered{} imagem\ensuremath{\leftrightarrow}BAI \textperiodcentered{} inferencia \textperiodcentered{} confiança 1.0 \textperiodcentered{} domínio ponte \textperiodcentered{} história 4 versões no jornal}

%CRISTAL {"arestas": [{"alvo": "word_m_geral_fastpath_2x2", "confianca": 1.0, "epistemico": "fato", "origem": "wiki", "rel": "relaciona"}, {"alvo": "algebra_dupolinomial_espaco", "confianca": 1.0, "epistemico": "fato", "origem": "wiki", "rel": "relaciona"}, {"alvo": "tiffany", "confianca": 1.0, "epistemico": "fato", "origem": "wiki", "rel": "parte_de"}], "confianca": 1.0, "contraexemplos": [], "descricao": "Varredura (agente metal): o 2×2 do gato não é uma escolha arbitrária a generalizar — é físico. A Word[total,e] nasce das duas classes de bit do byte (par e=popcount(byte&0xAA), ímpar o=popcount(byte&0x55)); a lemniscata a+c²=1 é dim-2; o corpo é ℤ[i] de Gauss. Em C só existe GRAU 2 (não há cúbica/plástica no kernel — só em Python). Logo M×M colapsaria ZERO arquivos C e tocaria os ~48 que leem .total/.e, trocando struct fixa por tamanho variável. Conclusão: manter o 2×2 como caso especial OTIMIZADO (O(1) branch-free, o workhorse). M×M no C só se/quando um consumidor C precisar de cúbica no kernel — e aí num metal_pos.c à parte, sem tocar a Word. [D]", "epistemico": "fato", "exemplos": [], "id": "gato_2x2_e_fisico_nao_incidental", "memoria": "semantica", "meta": {"autor": "Lopes/broca-so", "dominio": "arquitetura", "fonte": "agente varredura: code/neuronio.c:10-18, ula/word.h:5-8, ula/metais.c:8-11, ula/parabola.c"}, "origem": "wiki", "palavras_chave": [], "sinonimos": [], "tipo": "conceito", "titulo": "O gato 2×2 é FÍSICO, não incidental — M×M no kernel C complica, não simplifica"}
\section{O gato 2$\times$2 é FÍSICO, não incidental — M$\times$M no kernel C complica, não simplifica}
Varredura (agente metal): o 2$\times$2 do gato não é uma escolha arbitrária a generalizar — é físico. A Word[total,e] nasce das duas classes de bit do byte (par e=popcount(byte\&0xAA), ímpar o=popcount(byte\&0x55)); a lemniscata a+c\ensuremath{^{2}}=1 é dim-2; o corpo é \ensuremath{\mathbb{Z}}[i] de Gauss. Em C só existe GRAU 2 (não há cúbica/plástica no kernel — só em Python). Logo M$\times$M colapsaria ZERO arquivos C e tocaria os \~{}48 que leem .total/.e, trocando struct fixa por tamanho variável. Conclusão: manter o 2$\times$2 como caso especial OTIMIZADO (O(1) branch-free, o workhorse). M$\times$M no C só se/quando um consumidor C precisar de cúbica no kernel — e aí num metal\_pos.c à parte, sem tocar a Word. [D]
\prov{relações: relaciona word\_m\_geral\_fastpath\_2x2; relaciona algebra\_dupolinomial\_espaco; parte\_de tiffany}
\prov{proveniência: wiki \textperiodcentered{} agente varredura: code/neuronio.c:10-18, ula/word.h:5-8, ula/metais.c:8-11, ula/parabola.c \textperiodcentered{} fato \textperiodcentered{} confiança 1.0 \textperiodcentered{} domínio arquitetura \textperiodcentered{} história 30 versões no jornal}

%CRISTAL {"arestas": [{"alvo": "imagem_como_orbita_gato", "confianca": 1.0, "epistemico": "fato", "origem": "wiki", "rel": "usa"}, {"alvo": "cifra", "confianca": 1.0, "epistemico": "fato", "origem": "wiki", "rel": "eh_um"}, {"alvo": "criptografia_simetrica", "confianca": 1.0, "epistemico": "fato", "origem": "wiki", "rel": "relaciona"}, {"alvo": "recorrencia_de_poincare_imagem", "confianca": 1.0, "epistemico": "fato", "origem": "wiki", "rel": "usa"}], "confianca": 1.0, "contraexemplos": [], "descricao": "Embaralhar a imagem em k passos do gato produz ruído (a mistura caótica). Decifrar é recuar pelo gato inverso A_n^{-1}=[[0,1],[1,-n]] — ou completar a órbita (T−k passos a mais). A chave é o número de passos; a recuperação é EXATA (bijeção). O gato é uma cifra simétrica: a ida e a volta da mesma órbita mineral. Verif.: k=64, decifra por inverso == original.", "epistemico": "inferencia", "exemplos": [], "id": "gato_como_cifra", "memoria": "semantica", "meta": {"dominio": "imagem", "fonte": "órbitas de imagem"}, "origem": "wiki", "palavras_chave": [], "sinonimos": [], "tipo": "conceito", "titulo": "O Gato como Cifra de Imagem"}
\section{O Gato como Cifra de Imagem}
Embaralhar a imagem em k passos do gato produz ruído (a mistura caótica). Decifrar é recuar pelo gato inverso A\_n\^{}\{-1\}=[[0,1],[1,-n]] — ou completar a órbita (T\ensuremath{-}k passos a mais). A chave é o número de passos; a recuperação é EXATA (bijeção). O gato é uma cifra simétrica: a ida e a volta da mesma órbita mineral. Verif.: k=64, decifra por inverso == original.
\prov{relações: usa imagem\_como\_orbita\_gato; eh\_um cifra; relaciona criptografia\_simetrica; usa recorrencia\_de\_poincare\_imagem}
\prov{proveniência: wiki \textperiodcentered{} órbitas de imagem \textperiodcentered{} inferencia \textperiodcentered{} confiança 1.0 \textperiodcentered{} domínio imagem \textperiodcentered{} história 5 versões no jornal}

%CRISTAL {"arestas": [{"alvo": "gato", "confianca": 1.0, "epistemico": "fato", "origem": "wiki", "rel": "explica"}, {"alvo": "numero_de_salem", "confianca": 1.0, "epistemico": "fato", "origem": "wiki", "rel": "relaciona"}, {"alvo": "lyapunov_da_orbita", "confianca": 1.0, "epistemico": "fato", "origem": "wiki", "rel": "usa"}, {"alvo": "imagem_como_orbita_gato", "confianca": 1.0, "epistemico": "fato", "origem": "wiki", "rel": "usa"}], "confianca": 1.0, "contraexemplos": [], "descricao": "A matriz A_n=[[n,1],[1,0]] é exatamente um mapa do gato de Arnold — o automorfismo hiperbólico do toro T². Para n=1 os autovalores são φ e −1/φ (o áureo); para n=2, 1±√2 (a prata). Cada metal σ_n é um gato de Arnold com sua taxa de mistura λ=log σ_n. O ícone da dinâmica caótica É a ULA mineral.", "epistemico": "inferencia", "exemplos": [], "id": "gato_de_arnold", "memoria": "semantica", "meta": {"dominio": "imagem", "fonte": "órbitas de imagem", "lyapunov_A1": 0.481}, "origem": "wiki", "palavras_chave": [], "sinonimos": [], "tipo": "conceito", "titulo": "O Gato é o Mapa de Arnold"}
\section{O Gato é o Mapa de Arnold}
A matriz A\_n=[[n,1],[1,0]] é exatamente um mapa do gato de Arnold — o automorfismo hiperbólico do toro T\ensuremath{^{2}}. Para n=1 os autovalores são \ensuremath{\varphi} e \ensuremath{-}1/\ensuremath{\varphi} (o áureo); para n=2, 1$\pm$\ensuremath{\sqrt{\,}}2 (a prata). Cada metal \ensuremath{\sigma}\_n é um gato de Arnold com sua taxa de mistura \ensuremath{\lambda}=log \ensuremath{\sigma}\_n. O ícone da dinâmica caótica É a ULA mineral.
\prov{relações: explica gato; relaciona numero\_de\_salem; usa lyapunov\_da\_orbita; usa imagem\_como\_orbita\_gato}
\prov{proveniência: wiki \textperiodcentered{} órbitas de imagem \textperiodcentered{} inferencia \textperiodcentered{} confiança 1.0 \textperiodcentered{} domínio imagem \textperiodcentered{} história 5 versões no jornal}

%CRISTAL {"arestas": [{"alvo": "energia_eletrica", "confianca": 1.0, "epistemico": "fato", "origem": "wiki", "rel": "parte_de"}, {"alvo": "inducao_eletromagnetica", "confianca": 1.0, "epistemico": "fato", "origem": "wiki", "rel": "usa"}], "confianca": 1.0, "contraexemplos": [], "descricao": "Converter uma fonte primária (queda d'água, calor, vento, sol, núcleo) em eletricidade — quase sempre girando um gerador por indução.", "epistemico": "fato", "exemplos": [], "id": "geracao_de_energia", "memoria": "semantica", "meta": {"dominio": "engenharia", "fonte": ""}, "origem": "wiki", "palavras_chave": [], "sinonimos": [], "tipo": "conceito", "titulo": "Geração de Energia"}
\section{Geração de Energia}
Converter uma fonte primária (queda d'água, calor, vento, sol, núcleo) em eletricidade — quase sempre girando um gerador por indução.
\prov{relações: parte\_de energia\_eletrica; usa inducao\_eletromagnetica}
\prov{proveniência: wiki \textperiodcentered{} fato \textperiodcentered{} confiança 1.0 \textperiodcentered{} domínio engenharia \textperiodcentered{} história 3 versões no jornal}

%CRISTAL {"arestas": [{"alvo": "sistema_eletrico_potencia", "confianca": 1.0, "epistemico": "fato", "origem": "wiki", "rel": "relaciona"}, {"alvo": "energia_solar", "confianca": 1.0, "epistemico": "fato", "origem": "wiki", "rel": "usa"}, {"alvo": "eng_sistemas_energia", "confianca": 1.0, "epistemico": "fato", "origem": "curriculo", "rel": "parte_de"}, {"alvo": "rede_eletrica", "confianca": 1.0, "epistemico": "fato", "origem": "wiki", "rel": "relaciona"}, {"alvo": "distribuicao_de_energia", "confianca": 1.0, "epistemico": "fato", "origem": "wiki", "rel": "relaciona"}], "confianca": 1.0, "contraexemplos": [], "descricao": "Gerar perto do consumo (painéis solares no telhado, pequenas eólicas) em vez de só em grandes usinas. Encurta a transmissão (menos perda), mas exige a rede coordenar muitas fontes intermitentes.", "epistemico": "fato", "exemplos": [], "id": "geracao_distribuida", "memoria": "semantica", "meta": {"dominio": "engenharia", "fonte": "distribuição e armazenamento"}, "origem": "wiki", "palavras_chave": [], "sinonimos": [], "tipo": "conceito", "titulo": "Geração Distribuída"}
\section{Geração Distribuída}
Gerar perto do consumo (painéis solares no telhado, pequenas eólicas) em vez de só em grandes usinas. Encurta a transmissão (menos perda), mas exige a rede coordenar muitas fontes intermitentes.
\prov{relações: relaciona sistema\_eletrico\_potencia; usa energia\_solar; parte\_de eng\_sistemas\_energia; relaciona rede\_eletrica; relaciona distribuicao\_de\_energia}
\prov{proveniência: wiki \textperiodcentered{} distribuição e armazenamento \textperiodcentered{} fato \textperiodcentered{} confiança 1.0 \textperiodcentered{} domínio engenharia \textperiodcentered{} história 10 versões no jornal}

%CRISTAL {"arestas": [{"alvo": "comunicacoes_moveis", "confianca": 1.0, "epistemico": "fato", "origem": "wiki", "rel": "parte_de"}, {"alvo": "ofdm", "confianca": 1.0, "epistemico": "fato", "origem": "wiki", "rel": "usa"}, {"alvo": "eng_telecomunicacoes", "confianca": 1.0, "epistemico": "fato", "origem": "curriculo", "rel": "parte_de"}], "confianca": 1.0, "contraexemplos": [], "descricao": "1G analógico → 2G digital/GSM (voz+SMS) → 3G dados/CDMA → 4G LTE tudo-IP/OFDM → 5G NR (banda larga móvel, baixa latência, IoT massiva). Cada geração é um salto de acesso múltiplo e codificação.", "epistemico": "fato", "exemplos": [], "id": "geracoes_celular", "memoria": "semantica", "meta": {"dominio": "engenharia", "fonte": ""}, "origem": "wiki", "palavras_chave": [], "sinonimos": [], "tipo": "conceito", "titulo": "Gerações do Celular (1G→5G)"}
\section{Gerações do Celular (1G\ensuremath{\to}5G)}
1G analógico \ensuremath{\to} 2G digital/GSM (voz+SMS) \ensuremath{\to} 3G dados/CDMA \ensuremath{\to} 4G LTE tudo-IP/OFDM \ensuremath{\to} 5G NR (banda larga móvel, baixa latência, IoT massiva). Cada geração é um salto de acesso múltiplo e codificação.
\prov{relações: parte\_de comunicacoes\_moveis; usa ofdm; parte\_de eng\_telecomunicacoes}
\prov{proveniência: wiki \textperiodcentered{} fato \textperiodcentered{} confiança 1.0 \textperiodcentered{} domínio engenharia \textperiodcentered{} história 4 versões no jornal}

%CRISTAL {"arestas": [{"alvo": "maquina_eletrica", "confianca": 1.0, "epistemico": "fato", "origem": "wiki", "rel": "especializa"}, {"alvo": "geracao_de_energia", "confianca": 1.0, "epistemico": "fato", "origem": "wiki", "rel": "usa"}, {"alvo": "eng_sistemas_energia", "confianca": 1.0, "epistemico": "fato", "origem": "curriculo", "rel": "parte_de"}, {"alvo": "usina_eletrica", "confianca": 1.0, "epistemico": "fato", "origem": "wiki", "rel": "parte_de"}, {"alvo": "conversao_eletromecanica", "confianca": 1.0, "epistemico": "fato", "origem": "wiki", "rel": "usa"}], "confianca": 1.0, "contraexemplos": [], "descricao": "A máquina rotativa que converte rotação em eletricidade por indução (Faraday): o campo girante induz tensão nas bobinas. A conversão eletromecânica no coração de quase toda usina.", "epistemico": "fato", "exemplos": [], "id": "gerador_eletrico", "memoria": "semantica", "meta": {"dominio": "engenharia", "fonte": "usinas e geração"}, "origem": "wiki", "palavras_chave": [], "sinonimos": [], "tipo": "conceito", "titulo": "Gerador Elétrico"}
\section{Gerador Elétrico}
A máquina rotativa que converte rotação em eletricidade por indução (Faraday): o campo girante induz tensão nas bobinas. A conversão eletromecânica no coração de quase toda usina.
\prov{relações: especializa maquina\_eletrica; usa geracao\_de\_energia; parte\_de eng\_sistemas\_energia; parte\_de usina\_eletrica; usa conversao\_eletromecanica}
\prov{proveniência: wiki \textperiodcentered{} usinas e geração \textperiodcentered{} fato \textperiodcentered{} confiança 1.0 \textperiodcentered{} domínio engenharia \textperiodcentered{} história 10 versões no jornal}

%CRISTAL {"arestas": [{"alvo": "massive_mimo_beamforming", "confianca": 1.0, "epistemico": "fato", "origem": "wiki", "rel": "usa"}, {"alvo": "redes_5g_nr", "confianca": 1.0, "epistemico": "fato", "origem": "wiki", "rel": "parte_de"}, {"alvo": "sinais_de_referencia_nr", "confianca": 1.0, "epistemico": "fato", "origem": "wiki", "rel": "usa"}], "confianca": 1.0, "contraexemplos": [], "descricao": "Em mmWave a conexão é um FEIXE estreito: a rede e o UE varrem, medem e escolhem o melhor par de feixes, e o trocam quando o UE gira ou é bloqueado (recuperação de falha de feixe). Um 'handover' dentro da própria célula, no nível da antena.", "epistemico": "fato", "exemplos": [], "id": "gerenciamento_de_feixe", "memoria": "semantica", "meta": {"dominio": "engenharia", "fonte": ""}, "origem": "wiki", "palavras_chave": [], "sinonimos": [], "tipo": "conceito", "titulo": "Gerenciamento de Feixe (Beam Management)"}
\section{Gerenciamento de Feixe (Beam Management)}
Em mmWave a conexão é um FEIXE estreito: a rede e o UE varrem, medem e escolhem o melhor par de feixes, e o trocam quando o UE gira ou é bloqueado (recuperação de falha de feixe). Um 'handover' dentro da própria célula, no nível da antena.
\prov{relações: usa massive\_mimo\_beamforming; parte\_de redes\_5g\_nr; usa sinais\_de\_referencia\_nr}
\prov{proveniência: wiki \textperiodcentered{} fato \textperiodcentered{} confiança 1.0 \textperiodcentered{} domínio engenharia \textperiodcentered{} história 4 versões no jornal}

%CRISTAL {"arestas": [{"alvo": "smart_city", "confianca": 1.0, "epistemico": "fato", "origem": "wiki", "rel": "parte_de"}, {"alvo": "carga_coordenada", "confianca": 1.0, "epistemico": "fato", "origem": "wiki", "rel": "usa"}, {"alvo": "vehicle_to_grid", "confianca": 1.0, "epistemico": "fato", "origem": "wiki", "rel": "usa"}, {"alvo": "armazenamento_estabiliza_rede", "confianca": 1.0, "epistemico": "fato", "origem": "wiki", "rel": "usa"}], "confianca": 1.0, "contraexemplos": [], "descricao": "O cérebro que orquestra geração distribuída, armazenamento, frota (V2G) e demanda numa cidade para manter a curva plana e a rede na borda. É o sistema de autorização/colapso aplicado à energia da cidade — decidir, a cada instante, quem carrega, quem descarrega, quem espera.", "epistemico": "fato", "exemplos": [], "id": "gestao_de_energia_urbana", "memoria": "semantica", "meta": {"dominio": "engenharia", "fonte": "rede de carregamento e smart cities"}, "origem": "wiki", "palavras_chave": [], "sinonimos": [], "tipo": "conceito", "titulo": "Gestão de Energia Urbana"}
\section{Gestão de Energia Urbana}
O cérebro que orquestra geração distribuída, armazenamento, frota (V2G) e demanda numa cidade para manter a curva plana e a rede na borda. É o sistema de autorização/colapso aplicado à energia da cidade — decidir, a cada instante, quem carrega, quem descarrega, quem espera.
\prov{relações: parte\_de smart\_city; usa carga\_coordenada; usa vehicle\_to\_grid; usa armazenamento\_estabiliza\_rede}
\prov{proveniência: wiki \textperiodcentered{} rede de carregamento e smart cities \textperiodcentered{} fato \textperiodcentered{} confiança 1.0 \textperiodcentered{} domínio engenharia \textperiodcentered{} história 10 versões no jornal}

%CRISTAL {"arestas": [{"alvo": "repositorio_broca_os", "confianca": 1.0, "epistemico": "fato", "origem": "wiki", "rel": "explica"}, {"alvo": "organismo_gex44", "confianca": 1.0, "epistemico": "fato", "origem": "wiki", "rel": "explica"}, {"alvo": "gex_minera_ouro_real", "confianca": 1.0, "epistemico": "fato", "origem": "wiki", "rel": "relaciona"}], "confianca": 1.0, "contraexemplos": [], "descricao": "O organismo GEX44 hospeda as bases dos clientes e recebe deploys pela própria rede (deploy por banda, túnel SSH até o olho). O repo central não é 'dono': é o lugar onde as comunidades soberanas guardam e atualizam suas apps, cada uma com sua autoridade.", "epistemico": "inferencia", "exemplos": [], "id": "gex44_repo_central", "memoria": "semantica", "meta": {"dominio": "arquitetura", "fonte": "semilla_repositorio_broca"}, "origem": "wiki", "palavras_chave": [], "sinonimos": [], "tipo": "conceito", "titulo": "gex44 como Repo Central"}
\section{gex44 como Repo Central}
O organismo GEX44 hospeda as bases dos clientes e recebe deploys pela própria rede (deploy por banda, túnel SSH até o olho). O repo central não é 'dono': é o lugar onde as comunidades soberanas guardam e atualizam suas apps, cada uma com sua autoridade.
\prov{relações: explica repositorio\_broca\_os; explica organismo\_gex44; relaciona gex\_minera\_ouro\_real}
\prov{proveniência: wiki \textperiodcentered{} semilla\_repositorio\_broca \textperiodcentered{} inferencia \textperiodcentered{} confiança 1.0 \textperiodcentered{} domínio arquitetura \textperiodcentered{} história 4 versões no jornal}

%CRISTAL {"arestas": [{"alvo": "p2p", "confianca": 1.0, "epistemico": "fato", "origem": "wiki", "rel": "especializa"}, {"alvo": "goldchain", "confianca": 1.0, "epistemico": "fato", "origem": "wiki", "rel": "explica"}, {"alvo": "consenso_distribuido", "confianca": 1.0, "epistemico": "fato", "origem": "wiki", "rel": "usa"}, {"alvo": "prova_de_trabalho", "confianca": 1.0, "epistemico": "fato", "origem": "wiki", "rel": "usa"}, {"alvo": "dois_generais", "confianca": 1.0, "epistemico": "fato", "origem": "wiki", "rel": "usa"}, {"alvo": "orbita_rede_dois_atratores", "confianca": 1.0, "epistemico": "fato", "origem": "wiki", "rel": "parte_de"}], "confianca": 1.0, "contraexemplos": [], "descricao": "O goldchain lido pela ótica das redes: uma rede PONTO A PONTO (sem servidor central) que atinge consenso por prova-de-trabalho — o consenso de Nakamoto. Cada nó é cliente e servidor, exatamente como o BitTorrent, e resolve o problema dos dois generais por custo computacional em vez de acordo.", "epistemico": "fato", "exemplos": [], "id": "goldchain_como_rede", "memoria": "semantica", "meta": {"dominio": "engenharia", "fonte": "goldchain/"}, "origem": "wiki", "palavras_chave": [], "sinonimos": [], "tipo": "conceito", "titulo": "Goldchain como Rede P2P"}
\section{Goldchain como Rede P2P}
O goldchain lido pela ótica das redes: uma rede PONTO A PONTO (sem servidor central) que atinge consenso por prova-de-trabalho — o consenso de Nakamoto. Cada nó é cliente e servidor, exatamente como o BitTorrent, e resolve o problema dos dois generais por custo computacional em vez de acordo.
\prov{relações: especializa p2p; explica goldchain; usa consenso\_distribuido; usa prova\_de\_trabalho; usa dois\_generais; parte\_de orbita\_rede\_dois\_atratores}
\prov{proveniência: wiki \textperiodcentered{} goldchain/ \textperiodcentered{} fato \textperiodcentered{} confiança 1.0 \textperiodcentered{} domínio engenharia \textperiodcentered{} história 7 versões no jornal}

%CRISTAL {"arestas": [{"alvo": "repositorio_broca_os", "confianca": 1.0, "epistemico": "fato", "origem": "wiki", "rel": "parte_de"}, {"alvo": "autoridade_intrinseca", "confianca": 1.0, "epistemico": "fato", "origem": "wiki", "rel": "usa"}, {"alvo": "cadeia_admins_comunidade", "confianca": 1.0, "epistemico": "fato", "origem": "wiki", "rel": "usa"}, {"alvo": "face_turno_fail_closed", "confianca": 1.0, "epistemico": "fato", "origem": "wiki", "rel": "usa"}], "confianca": 1.0, "contraexemplos": [], "descricao": "Escrever na base de um cliente exige AUTORIDADE: intrínseca (o app chegou decodável na banda — só quem tem o tecido publica) OU delegada (pode(peer, banda, 'deploy') pela cadeia de admins). Sem autoridade, recusado (fail-closed); nunca escreve fora de clientes/<banda>/apps/ (sem '..'/absoluto/dotfile).", "epistemico": "fato", "exemplos": [], "id": "governanca_deploy_banda", "memoria": "semantica", "meta": {"dominio": "arquitetura", "fonte": "semilla_repositorio_broca"}, "origem": "wiki", "palavras_chave": [], "sinonimos": [], "tipo": "conceito", "titulo": "Governança do Deploy por Banda"}
\section{Governança do Deploy por Banda}
Escrever na base de um cliente exige AUTORIDADE: intrínseca (o app chegou decodável na banda — só quem tem o tecido publica) OU delegada (pode(peer, banda, 'deploy') pela cadeia de admins). Sem autoridade, recusado (fail-closed); nunca escreve fora de clientes/<banda>/apps/ (sem '..'/absoluto/dotfile).
\prov{relações: parte\_de repositorio\_broca\_os; usa autoridade\_intrinseca; usa cadeia\_admins\_comunidade; usa face\_turno\_fail\_closed}
\prov{proveniência: wiki \textperiodcentered{} semilla\_repositorio\_broca \textperiodcentered{} fato \textperiodcentered{} confiança 1.0 \textperiodcentered{} domínio arquitetura \textperiodcentered{} história 5 versões no jornal}

%CRISTAL {"arestas": [{"alvo": "camada_fisica_nr", "confianca": 1.0, "epistemico": "fato", "origem": "wiki", "rel": "parte_de"}], "confianca": 1.0, "contraexemplos": [], "descricao": "O plano tempo×frequência dividido em Elementos de Recurso (1 subportadora × 1 símbolo) agrupados em Blocos de Recurso (12 subportadoras). O escalonador do gNB aloca esses blocos a cada usuário slot a slot.", "epistemico": "fato", "exemplos": [], "id": "grade_de_recursos_nr", "memoria": "semantica", "meta": {"dominio": "engenharia", "fonte": ""}, "origem": "wiki", "palavras_chave": [], "sinonimos": [], "tipo": "conceito", "titulo": "Grade de Recursos (RB, RE)"}
\section{Grade de Recursos (RB, RE)}
O plano tempo$\times$frequência dividido em Elementos de Recurso (1 subportadora $\times$ 1 símbolo) agrupados em Blocos de Recurso (12 subportadoras). O escalonador do gNB aloca esses blocos a cada usuário slot a slot.
\prov{relações: parte\_de camada\_fisica\_nr}
\prov{proveniência: wiki \textperiodcentered{} fato \textperiodcentered{} confiança 1.0 \textperiodcentered{} domínio engenharia \textperiodcentered{} história 2 versões no jornal}

%CRISTAL {"arestas": [{"alvo": "http2_http3", "confianca": 1.0, "epistemico": "fato", "origem": "wiki", "rel": "usa"}, {"alvo": "microservicos", "confianca": 1.0, "epistemico": "fato", "origem": "wiki", "rel": "usa"}, {"alvo": "eng_telecomunicacoes", "confianca": 1.0, "epistemico": "fato", "origem": "curriculo", "rel": "parte_de"}], "confianca": 1.0, "contraexemplos": [], "descricao": "Chamada de procedimento remoto sobre HTTP/2 com contrato binário (Protocol Buffers) — tipado, rápido e com streaming. A cola de alto desempenho entre microsserviços.", "epistemico": "fato", "exemplos": [], "id": "grpc", "memoria": "semantica", "meta": {"autor": "Google", "dominio": "engenharia", "fonte": ""}, "origem": "wiki", "palavras_chave": [], "sinonimos": [], "tipo": "conceito", "titulo": "gRPC"}
\section{gRPC}
Chamada de procedimento remoto sobre HTTP/2 com contrato binário (Protocol Buffers) — tipado, rápido e com streaming. A cola de alto desempenho entre microsserviços.
\prov{relações: usa http2\_http3; usa microservicos; parte\_de eng\_telecomunicacoes}
\prov{proveniência: wiki \textperiodcentered{} fato \textperiodcentered{} confiança 1.0 \textperiodcentered{} domínio engenharia \textperiodcentered{} história 4 versões no jornal}

%CRISTAL {"arestas": [{"alvo": "classificacao_de_formas", "confianca": 1.0, "epistemico": "fato", "origem": "wiki", "rel": "explica"}, {"alvo": "quasicristal_de_rotacoes", "confianca": 1.0, "epistemico": "fato", "origem": "wiki", "rel": "relaciona"}], "confianca": 1.0, "contraexemplos": [], "descricao": "Translação, rotação, escala e reflexão são matrizes que COMPÕEM: formam um grupo. rotação∘rotação soma ângulos; reflexão∘reflexão=identidade. A simetria de uma forma É um subgrupo diedral D_n (a tabela de Cayley fecha). Classificar formas e o grupo que as move são o mesmo objeto. A reflexão é uma rotação que inverte a orientação.", "epistemico": "fato", "exemplos": [], "id": "grupo_de_transformacoes", "memoria": "semantica", "meta": {"dominio": "imagem", "fonte": "formas e limite"}, "origem": "wiki", "palavras_chave": [], "sinonimos": [], "tipo": "conceito", "titulo": "O Grupo de Transformações (que compõe)"}
\section{O Grupo de Transformações (que compõe)}
Translação, rotação, escala e reflexão são matrizes que COMPÕEM: formam um grupo. rotação\ensuremath{\circ}rotação soma ângulos; reflexão\ensuremath{\circ}reflexão=identidade. A simetria de uma forma É um subgrupo diedral D\_n (a tabela de Cayley fecha). Classificar formas e o grupo que as move são o mesmo objeto. A reflexão é uma rotação que inverte a orientação.
\prov{relações: explica classificacao\_de\_formas; relaciona quasicristal\_de\_rotacoes}
\prov{proveniência: wiki \textperiodcentered{} formas e limite \textperiodcentered{} fato \textperiodcentered{} confiança 1.0 \textperiodcentered{} domínio imagem \textperiodcentered{} história 6 versões no jornal}

%CRISTAL {"arestas": [{"alvo": "linha_de_transmissao", "confianca": 1.0, "epistemico": "fato", "origem": "wiki", "rel": "especializa"}, {"alvo": "eletromagnetismo_eng", "confianca": 1.0, "epistemico": "fato", "origem": "wiki", "rel": "usa"}, {"alvo": "eng_telecomunicacoes", "confianca": 1.0, "epistemico": "fato", "origem": "curriculo", "rel": "parte_de"}], "confianca": 1.0, "contraexemplos": [], "descricao": "Em frequências altas, o sinal viaja confinado num tubo condutor ou dielétrico — modos, frequência de corte, cavidades ressonantes. A plataforma do radar e do enlace de micro-ondas.", "epistemico": "fato", "exemplos": [], "id": "guia_de_onda", "memoria": "semantica", "meta": {"dominio": "engenharia", "fonte": ""}, "origem": "wiki", "palavras_chave": [], "sinonimos": [], "tipo": "conceito", "titulo": "Guias de Onda e Micro-ondas"}
\section{Guias de Onda e Micro-ondas}
Em frequências altas, o sinal viaja confinado num tubo condutor ou dielétrico — modos, frequência de corte, cavidades ressonantes. A plataforma do radar e do enlace de micro-ondas.
\prov{relações: especializa linha\_de\_transmissao; usa eletromagnetismo\_eng; parte\_de eng\_telecomunicacoes}
\prov{proveniência: wiki \textperiodcentered{} fato \textperiodcentered{} confiança 1.0 \textperiodcentered{} domínio engenharia \textperiodcentered{} história 4 versões no jornal}

%CRISTAL {"arestas": [{"alvo": "handover_nr", "confianca": 1.0, "epistemico": "fato", "origem": "wiki", "rel": "especializa"}], "confianca": 1.0, "contraexemplos": [], "descricao": "A rede prepara VÁRIOS alvos com antecedência e dá ao UE a CONDIÇÃO para executar; o UE dispara sozinho quando a condição se cumpre. Robustez em alta velocidade — não depende do relatório chegar a tempo.", "epistemico": "fato", "exemplos": [], "id": "handover_condicional_cho", "memoria": "semantica", "meta": {"dominio": "engenharia", "fonte": ""}, "origem": "wiki", "palavras_chave": [], "sinonimos": [], "tipo": "conceito", "titulo": "Handover Condicional (CHO)"}
\section{Handover Condicional (CHO)}
A rede prepara VÁRIOS alvos com antecedência e dá ao UE a CONDIÇÃO para executar; o UE dispara sozinho quando a condição se cumpre. Robustez em alta velocidade — não depende do relatório chegar a tempo.
\prov{relações: especializa handover\_nr}
\prov{proveniência: wiki \textperiodcentered{} fato \textperiodcentered{} confiança 1.0 \textperiodcentered{} domínio engenharia \textperiodcentered{} história 2 versões no jornal}

%CRISTAL {"arestas": [{"alvo": "handover_nr", "confianca": 1.0, "epistemico": "fato", "origem": "wiki", "rel": "especializa"}, {"alvo": "dupla_conectividade_nr", "confianca": 1.0, "epistemico": "fato", "origem": "wiki", "rel": "usa"}], "confianca": 1.0, "contraexemplos": [], "descricao": "Dual Active Protocol Stack: o UE mantém origem E alvo ativos simultaneamente durante a troca, só largando a origem depois de firmar o alvo. Interrupção praticamente zero — para tráfego que não tolera corte.", "epistemico": "fato", "exemplos": [], "id": "handover_daps", "memoria": "semantica", "meta": {"dominio": "engenharia", "fonte": ""}, "origem": "wiki", "palavras_chave": [], "sinonimos": [], "tipo": "conceito", "titulo": "Handover DAPS (0 ms)"}
\section{Handover DAPS (0 ms)}
Dual Active Protocol Stack: o UE mantém origem E alvo ativos simultaneamente durante a troca, só largando a origem depois de firmar o alvo. Interrupção praticamente zero — para tráfego que não tolera corte.
\prov{relações: especializa handover\_nr; usa dupla\_conectividade\_nr}
\prov{proveniência: wiki \textperiodcentered{} fato \textperiodcentered{} confiança 1.0 \textperiodcentered{} domínio engenharia \textperiodcentered{} história 3 versões no jornal}

%CRISTAL {"arestas": [{"alvo": "handover_nr", "confianca": 1.0, "epistemico": "fato", "origem": "wiki", "rel": "especializa"}, {"alvo": "funcoes_de_rede_5gc", "confianca": 1.0, "epistemico": "fato", "origem": "wiki", "rel": "usa"}], "confianca": 1.0, "contraexemplos": [], "descricao": "Sem Xn entre os gNBs, a transferência sobe ao núcleo: o AMF coordena a preparação do alvo. Mais lento, usado entre áreas ou fabricantes que não têm enlace direto.", "epistemico": "inferencia", "exemplos": [], "id": "handover_n2", "memoria": "semantica", "meta": {"dominio": "engenharia", "fonte": ""}, "origem": "wiki", "palavras_chave": [], "sinonimos": [], "tipo": "conceito", "titulo": "Handover por N2 (via núcleo)"}
\section{Handover por N2 (via núcleo)}
Sem Xn entre os gNBs, a transferência sobe ao núcleo: o AMF coordena a preparação do alvo. Mais lento, usado entre áreas ou fabricantes que não têm enlace direto.
\prov{relações: especializa handover\_nr; usa funcoes\_de\_rede\_5gc}
\prov{proveniência: wiki \textperiodcentered{} inferencia \textperiodcentered{} confiança 1.0 \textperiodcentered{} domínio engenharia \textperiodcentered{} história 3 versões no jornal}

%CRISTAL {"arestas": [{"alvo": "redes_5g_nr", "confianca": 1.0, "epistemico": "fato", "origem": "wiki", "rel": "parte_de"}, {"alvo": "comunicacoes_moveis", "confianca": 1.0, "epistemico": "fato", "origem": "wiki", "rel": "especializa"}, {"alvo": "medicoes_mobilidade", "confianca": 1.0, "epistemico": "fato", "origem": "wiki", "rel": "usa"}, {"alvo": "rach_sem_contencao", "confianca": 1.0, "epistemico": "fato", "origem": "wiki", "rel": "usa"}, {"alvo": "eng_telecomunicacoes", "confianca": 1.0, "epistemico": "fato", "origem": "curriculo", "rel": "parte_de"}], "confianca": 1.0, "contraexemplos": [], "descricao": "Trocar a célula servidora sem derrubar a conexão, com o UE em movimento: a rede prepara o alvo, manda o comando, o UE faz acesso sem contenção no alvo e o núcleo comuta o caminho. Milissegundos de interrupção que o handover moderno tenta zerar.", "epistemico": "fato", "exemplos": [], "id": "handover_nr", "memoria": "semantica", "meta": {"dominio": "engenharia", "fonte": ""}, "origem": "wiki", "palavras_chave": [], "sinonimos": [], "tipo": "conceito", "titulo": "Handover (Transferência de Célula)"}
\section{Handover (Transferência de Célula)}
Trocar a célula servidora sem derrubar a conexão, com o UE em movimento: a rede prepara o alvo, manda o comando, o UE faz acesso sem contenção no alvo e o núcleo comuta o caminho. Milissegundos de interrupção que o handover moderno tenta zerar.
\prov{relações: parte\_de redes\_5g\_nr; especializa comunicacoes\_moveis; usa medicoes\_mobilidade; usa rach\_sem\_contencao; parte\_de eng\_telecomunicacoes}
\prov{proveniência: wiki \textperiodcentered{} fato \textperiodcentered{} confiança 1.0 \textperiodcentered{} domínio engenharia \textperiodcentered{} história 6 versões no jornal}

%CRISTAL {"arestas": [{"alvo": "handover_nr", "confianca": 1.0, "epistemico": "fato", "origem": "wiki", "rel": "especializa"}, {"alvo": "comutacao_de_caminho", "confianca": 1.0, "epistemico": "fato", "origem": "wiki", "rel": "usa"}], "confianca": 1.0, "contraexemplos": [], "descricao": "Quando os gNBs de origem e alvo se falam direto pela interface Xn: a preparação e a transferência de contexto vão por Xn, e só no fim o AMF/UPF é avisado (path switch). O caminho rápido e mais comum.", "epistemico": "fato", "exemplos": [], "id": "handover_xn", "memoria": "semantica", "meta": {"dominio": "engenharia", "fonte": ""}, "origem": "wiki", "palavras_chave": [], "sinonimos": [], "tipo": "conceito", "titulo": "Handover por Xn (entre gNBs)"}
\section{Handover por Xn (entre gNBs)}
Quando os gNBs de origem e alvo se falam direto pela interface Xn: a preparação e a transferência de contexto vão por Xn, e só no fim o AMF/UPF é avisado (path switch). O caminho rápido e mais comum.
\prov{relações: especializa handover\_nr; usa comutacao\_de\_caminho}
\prov{proveniência: wiki \textperiodcentered{} fato \textperiodcentered{} confiança 1.0 \textperiodcentered{} domínio engenharia \textperiodcentered{} história 3 versões no jornal}

%CRISTAL {"arestas": [{"alvo": "codificacao_de_canal", "confianca": 1.0, "epistemico": "fato", "origem": "wiki", "rel": "usa"}, {"alvo": "camada_mac_nr", "confianca": 1.0, "epistemico": "fato", "origem": "wiki", "rel": "parte_de"}], "confianca": 1.0, "contraexemplos": [], "descricao": "Combina correção de erro (FEC) com retransmissão (ARQ): o receptor guarda a tentativa falha e a COMBINA com a retransmissão (chase/incremental), em vez de descartá-la. Recupera erros com mínimo atraso — vital no URLLC.", "epistemico": "fato", "exemplos": [], "id": "harq_nr", "memoria": "semantica", "meta": {"dominio": "engenharia", "fonte": ""}, "origem": "wiki", "palavras_chave": [], "sinonimos": [], "tipo": "conceito", "titulo": "HARQ (ARQ Híbrido)"}
\section{HARQ (ARQ Híbrido)}
Combina correção de erro (FEC) com retransmissão (ARQ): o receptor guarda a tentativa falha e a COMBINA com a retransmissão (chase/incremental), em vez de descartá-la. Recupera erros com mínimo atraso — vital no URLLC.
\prov{relações: usa codificacao\_de\_canal; parte\_de camada\_mac\_nr}
\prov{proveniência: wiki \textperiodcentered{} fato \textperiodcentered{} confiança 1.0 \textperiodcentered{} domínio engenharia \textperiodcentered{} história 3 versões no jornal}

%CRISTAL {"arestas": [{"alvo": "armazenamento_o_cristal", "confianca": 1.0, "epistemico": "fato", "origem": "wiki", "rel": "eh_um"}, {"alvo": "energia_hidreletrica", "confianca": 1.0, "epistemico": "fato", "origem": "wiki", "rel": "relaciona"}], "confianca": 1.0, "contraexemplos": [], "descricao": "Bombeia água para um reservatório alto quando sobra energia e a solta pela turbina quando falta. É o maior armazenamento instalado do mundo (~90% da capacidade), η_rt~80% — a órbita entre dois níveis.", "epistemico": "fato", "exemplos": [], "id": "hidreletrica_reversivel", "memoria": "semantica", "meta": {"dominio": "engenharia", "fonte": "distribuição e armazenamento"}, "origem": "wiki", "palavras_chave": [], "sinonimos": [], "tipo": "conceito", "titulo": "Hidrelétrica Reversível (pumped hydro)"}
\section{Hidrelétrica Reversível (pumped hydro)}
Bombeia água para um reservatório alto quando sobra energia e a solta pela turbina quando falta. É o maior armazenamento instalado do mundo (\~{}90\% da capacidade), \ensuremath{\eta}\_rt\~{}80\% — a órbita entre dois níveis.
\prov{relações: eh\_um armazenamento\_o\_cristal; relaciona energia\_hidreletrica}
\prov{proveniência: wiki \textperiodcentered{} distribuição e armazenamento \textperiodcentered{} fato \textperiodcentered{} confiança 1.0 \textperiodcentered{} domínio engenharia \textperiodcentered{} história 6 versões no jornal}

%CRISTAL {"arestas": [{"alvo": "face_turno_fail_closed", "confianca": 1.0, "epistemico": "fato", "origem": "wiki", "rel": "relaciona"}, {"alvo": "xadrez_bai", "confianca": 1.0, "epistemico": "fato", "origem": "wiki", "rel": "parte_de"}], "confianca": 1.0, "contraexemplos": [], "descricao": "Quando o orçamento de busca esgota, o motor não vê além do horizonte e decide conservador pela avaliação — exatamente o fail-closed do BAI quando δ chega a zero (autoridade não propagada → nega). O limite do cálculo e o limite da autoridade são o mesmo limite.", "epistemico": "inferencia", "exemplos": [], "id": "horizonte_e_fail_closed", "memoria": "semantica", "meta": {"dominio": "ponte", "fonte": "xadrez↔BAI"}, "origem": "wiki", "palavras_chave": [], "sinonimos": [], "tipo": "conceito", "titulo": "Efeito Horizonte = Fail-Closed"}
\section{Efeito Horizonte = Fail-Closed}
Quando o orçamento de busca esgota, o motor não vê além do horizonte e decide conservador pela avaliação — exatamente o fail-closed do BAI quando \ensuremath{\delta} chega a zero (autoridade não propagada \ensuremath{\to} nega). O limite do cálculo e o limite da autoridade são o mesmo limite.
\prov{relações: relaciona face\_turno\_fail\_closed; parte\_de xadrez\_bai}
\prov{proveniência: wiki \textperiodcentered{} xadrez\ensuremath{\leftrightarrow}BAI \textperiodcentered{} inferencia \textperiodcentered{} confiança 1.0 \textperiodcentered{} domínio ponte \textperiodcentered{} história 3 versões no jornal}

%CRISTAL {"arestas": [{"alvo": "http", "confianca": 1.0, "epistemico": "fato", "origem": "wiki", "rel": "especializa"}, {"alvo": "quic", "confianca": 1.0, "epistemico": "fato", "origem": "wiki", "rel": "usa"}, {"alvo": "eng_telecomunicacoes", "confianca": 1.0, "epistemico": "fato", "origem": "curriculo", "rel": "parte_de"}], "confianca": 1.0, "contraexemplos": [], "descricao": "Evolução do HTTP: multiplexação de streams e compressão de cabeçalhos (H2 sobre TCP) e depois sobre QUIC (H3), matando o bloqueio de cabeça de fila. A web mais rápida sem mudar a semântica de GET/POST.", "epistemico": "fato", "exemplos": [], "id": "http2_http3", "memoria": "semantica", "meta": {"dominio": "engenharia", "fonte": ""}, "origem": "wiki", "palavras_chave": [], "sinonimos": [], "tipo": "conceito", "titulo": "HTTP/2 e HTTP/3"}
\section{HTTP/2 e HTTP/3}
Evolução do HTTP: multiplexação de streams e compressão de cabeçalhos (H2 sobre TCP) e depois sobre QUIC (H3), matando o bloqueio de cabeça de fila. A web mais rápida sem mudar a semântica de GET/POST.
\prov{relações: especializa http; usa quic; parte\_de eng\_telecomunicacoes}
\prov{proveniência: wiki \textperiodcentered{} fato \textperiodcentered{} confiança 1.0 \textperiodcentered{} domínio engenharia \textperiodcentered{} história 4 versões no jornal}

%CRISTAL {"arestas": [{"alvo": "registro_5g", "confianca": 1.0, "epistemico": "fato", "origem": "wiki", "rel": "usa"}, {"alvo": "seguranca_5g", "confianca": 1.0, "epistemico": "fato", "origem": "wiki", "rel": "relaciona"}], "confianca": 1.0, "contraexemplos": [], "descricao": "O identificador permanente (SUPI, o antigo IMSI) NUNCA vai claro no ar: viaja cifrado como SUCI. A rede então atribui um temporário (5G-GUTI) que troca com frequência — anonimato contra rastreio passivo.", "epistemico": "fato", "exemplos": [], "id": "identidades_5g", "memoria": "semantica", "meta": {"dominio": "engenharia", "fonte": ""}, "origem": "wiki", "palavras_chave": [], "sinonimos": [], "tipo": "conceito", "titulo": "Identidades do Assinante (SUPI/SUCI/5G-GUTI)"}
\section{Identidades do Assinante (SUPI/SUCI/5G-GUTI)}
O identificador permanente (SUPI, o antigo IMSI) NUNCA vai claro no ar: viaja cifrado como SUCI. A rede então atribui um temporário (5G-GUTI) que troca com frequência — anonimato contra rastreio passivo.
\prov{relações: usa registro\_5g; relaciona seguranca\_5g}
\prov{proveniência: wiki \textperiodcentered{} fato \textperiodcentered{} confiança 1.0 \textperiodcentered{} domínio engenharia \textperiodcentered{} história 3 versões no jornal}

%CRISTAL {"arestas": [{"alvo": "segmentacao_de_imagem", "confianca": 1.0, "epistemico": "fato", "origem": "wiki", "rel": "usa"}, {"alvo": "bai", "confianca": 1.0, "epistemico": "fato", "origem": "wiki", "rel": "usa"}, {"alvo": "xadrez_bai", "confianca": 1.0, "epistemico": "fato", "origem": "wiki", "rel": "relaciona"}, {"alvo": "laco_traducao_bai_peirce", "confianca": 1.0, "epistemico": "fato", "origem": "wiki", "rel": "relaciona"}, {"alvo": "teorema_geral_traducao", "confianca": 1.0, "epistemico": "fato", "origem": "wiki", "rel": "usa"}], "confianca": 1.0, "contraexemplos": [], "descricao": "Segmentar uma imagem e a Biblioteca de Autorização Isomórfica descrevem o MESMO objeto: uma busca que parte de uma semente, propaga, colapsa cada pixel num rótulo e fecha na fronteira quando o orçamento (a escala k) esgota. Terceira simbologia da mesma árvore (com xadrez e BAI); o dicionário traduz 11 construtos 1:1, a diferença (pixel, kernel, cor) é só bump. compor() dá imagem→BAI e imagem→xadrez de graça.", "epistemico": "fato", "exemplos": [], "id": "imagem_bai", "memoria": "semantica", "meta": {"dominio": "ponte", "fonte": "imagem↔BAI"}, "origem": "wiki", "palavras_chave": [], "sinonimos": [], "tipo": "conceito", "titulo": "Imagem ↔ BAI (a mesma árvore de decisão)"}
\section{Imagem \ensuremath{\leftrightarrow} BAI (a mesma árvore de decisão)}
Segmentar uma imagem e a Biblioteca de Autorização Isomórfica descrevem o MESMO objeto: uma busca que parte de uma semente, propaga, colapsa cada pixel num rótulo e fecha na fronteira quando o orçamento (a escala k) esgota. Terceira simbologia da mesma árvore (com xadrez e BAI); o dicionário traduz 11 construtos 1:1, a diferença (pixel, kernel, cor) é só bump. compor() dá imagem\ensuremath{\to}BAI e imagem\ensuremath{\to}xadrez de graça.
\prov{relações: usa segmentacao\_de\_imagem; usa bai; relaciona xadrez\_bai; relaciona laco\_traducao\_bai\_peirce; usa teorema\_geral\_traducao}
\prov{proveniência: wiki \textperiodcentered{} imagem\ensuremath{\leftrightarrow}BAI \textperiodcentered{} fato \textperiodcentered{} confiança 1.0 \textperiodcentered{} domínio ponte \textperiodcentered{} história 6 versões no jornal}

%CRISTAL {"arestas": [{"alvo": "gato_ula", "confianca": 1.0, "epistemico": "fato", "origem": "wiki", "rel": "usa"}, {"alvo": "gato", "confianca": 1.0, "epistemico": "fato", "origem": "wiki", "rel": "usa"}, {"alvo": "orbita_de_partida", "confianca": 1.0, "epistemico": "fato", "origem": "wiki", "rel": "eh_um"}, {"alvo": "computador_mineral", "confianca": 1.0, "epistemico": "fato", "origem": "wiki", "rel": "usa"}, {"alvo": "corpo_mineral", "confianca": 1.0, "epistemico": "fato", "origem": "wiki", "rel": "usa"}], "confianca": 1.0, "contraexemplos": [], "descricao": "Aplicar o gato A_n=[[n,1],[1,0]] a uma imagem N×N (mod N) é iterar uma órbita no toro discreto. O gato tem Lyapunov λ=log σ_n>0 — mistura caótica que embaralha a imagem — mas é uma BIJEÇÃO (det=±1). A imagem, portanto, percorre uma órbita finita. É o computador mineral operando sobre pixels: a mesma ULA (o gato) agora agita uma figura. Verif.: N=128, A_1 embaralha em poucos passos.", "epistemico": "inferencia", "exemplos": [], "id": "imagem_como_orbita_gato", "memoria": "semantica", "meta": {"dominio": "imagem", "fonte": "órbitas de imagem", "lyapunov_A1": 0.481}, "origem": "wiki", "palavras_chave": [], "sinonimos": [], "tipo": "conceito", "titulo": "A Imagem é uma Órbita do Gato"}
\section{A Imagem é uma Órbita do Gato}
Aplicar o gato A\_n=[[n,1],[1,0]] a uma imagem N$\times$N (mod N) é iterar uma órbita no toro discreto. O gato tem Lyapunov \ensuremath{\lambda}=log \ensuremath{\sigma}\_n>0 — mistura caótica que embaralha a imagem — mas é uma BIJEÇÃO (det=$\pm$1). A imagem, portanto, percorre uma órbita finita. É o computador mineral operando sobre pixels: a mesma ULA (o gato) agora agita uma figura. Verif.: N=128, A\_1 embaralha em poucos passos.
\prov{relações: usa gato\_ula; usa gato; eh\_um orbita\_de\_partida; usa computador\_mineral; usa corpo\_mineral}
\prov{proveniência: wiki \textperiodcentered{} órbitas de imagem \textperiodcentered{} inferencia \textperiodcentered{} confiança 1.0 \textperiodcentered{} domínio imagem \textperiodcentered{} história 6 versões no jornal}

%CRISTAL {"arestas": [{"alvo": "transformada_de_fourier", "confianca": 1.0, "epistemico": "fato", "origem": "wiki", "rel": "usa"}, {"alvo": "imagem_como_orbita_gato", "confianca": 1.0, "epistemico": "fato", "origem": "wiki", "rel": "relaciona"}, {"alvo": "numero_de_salem", "confianca": 1.0, "epistemico": "fato", "origem": "wiki", "rel": "relaciona"}, {"alvo": "corpo_mineral", "confianca": 1.0, "epistemico": "fato", "origem": "wiki", "rel": "usa"}], "confianca": 1.0, "contraexemplos": [], "descricao": "A transformada de Fourier 2D decompõe a imagem em rotações puras e^{iω·}, cada uma no círculo unitário (ρ=1 = Salem/borda). É o outro modo de ler a órbita: no espectro (frequências) em vez do espaço (pixels). Comprimir uma imagem é manter as rotações dominantes — o mesmo gesto do DMD nas órbitas minerais.", "epistemico": "inferencia", "exemplos": [], "id": "imagem_em_rotacoes_fourier", "memoria": "semantica", "meta": {"dominio": "imagem", "fonte": "órbitas de imagem"}, "origem": "wiki", "palavras_chave": [], "sinonimos": [], "tipo": "conceito", "titulo": "A Imagem como Soma de Rotações (Fourier)"}
\section{A Imagem como Soma de Rotações (Fourier)}
A transformada de Fourier 2D decompõe a imagem em rotações puras e\^{}\{i\ensuremath{\omega}\textperiodcentered \}, cada uma no círculo unitário (\ensuremath{\rho}=1 = Salem/borda). É o outro modo de ler a órbita: no espectro (frequências) em vez do espaço (pixels). Comprimir uma imagem é manter as rotações dominantes — o mesmo gesto do DMD nas órbitas minerais.
\prov{relações: usa transformada\_de\_fourier; relaciona imagem\_como\_orbita\_gato; relaciona numero\_de\_salem; usa corpo\_mineral}
\prov{proveniência: wiki \textperiodcentered{} órbitas de imagem \textperiodcentered{} inferencia \textperiodcentered{} confiança 1.0 \textperiodcentered{} domínio imagem \textperiodcentered{} história 5 versões no jornal}

%CRISTAL {"arestas": [{"alvo": "colapso", "confianca": 1.0, "epistemico": "fato", "origem": "wiki", "rel": "eh_um"}, {"alvo": "rotulagem_e_colapso", "confianca": 1.0, "epistemico": "fato", "origem": "wiki", "rel": "relaciona"}, {"alvo": "sensor_e_colapso", "confianca": 1.0, "epistemico": "fato", "origem": "wiki", "rel": "relaciona"}, {"alvo": "corpo_mineral", "confianca": 1.0, "epistemico": "fato", "origem": "wiki", "rel": "usa"}, {"alvo": "imagem_bai", "confianca": 1.0, "epistemico": "fato", "origem": "wiki", "rel": "relaciona"}], "confianca": 1.0, "contraexemplos": [], "descricao": "Ler um caractere de uma imagem (OCR) é um COLAPSO: o glifo, numa superposição de 'que letra sou?', desaba num rótulo. É o mesmo Ψ do sensor (mundo→bit), da segmentação (pixel→região) e do minimax (árvore→lance), agora glifo→caractere. Reconhecer é achar o cristal (template) mais próximo.", "epistemico": "fato", "exemplos": [], "id": "imagem_para_texto", "memoria": "semantica", "meta": {"dominio": "imagem", "fonte": "glifos minerais"}, "origem": "wiki", "palavras_chave": [], "sinonimos": [], "tipo": "conceito", "titulo": "Imagem → Texto (o reconhecimento como colapso)"}
\section{Imagem \ensuremath{\to} Texto (o reconhecimento como colapso)}
Ler um caractere de uma imagem (OCR) é um COLAPSO: o glifo, numa superposição de 'que letra sou?', desaba num rótulo. É o mesmo \ensuremath{\Psi} do sensor (mundo\ensuremath{\to}bit), da segmentação (pixel\ensuremath{\to}região) e do minimax (árvore\ensuremath{\to}lance), agora glifo\ensuremath{\to}caractere. Reconhecer é achar o cristal (template) mais próximo.
\prov{relações: eh\_um colapso; relaciona rotulagem\_e\_colapso; relaciona sensor\_e\_colapso; usa corpo\_mineral; relaciona imagem\_bai}
\prov{proveniência: wiki \textperiodcentered{} glifos minerais \textperiodcentered{} fato \textperiodcentered{} confiança 1.0 \textperiodcentered{} domínio imagem \textperiodcentered{} história 6 versões no jornal}

%CRISTAL {"arestas": [{"alvo": "manufatura", "confianca": 1.0, "epistemico": "fato", "origem": "wiki", "rel": "eh_um"}, {"alvo": "atuador", "confianca": 1.0, "epistemico": "fato", "origem": "wiki", "rel": "usa"}], "confianca": 1.0, "contraexemplos": [], "descricao": "Construir um objeto físico depositando material camada por camada, guiado por um arquivo digital (G-code). De protótipos a próteses e peças de foguete: a fabricação que parte direto do bit para a forma, sem molde.", "epistemico": "fato", "exemplos": [], "id": "impressao_3d", "memoria": "semantica", "meta": {"dominio": "manufatura", "fonte": "impressão 3D e manufatura"}, "origem": "wiki", "palavras_chave": [], "sinonimos": [], "tipo": "conceito", "titulo": "Impressão 3D"}
\section{Impressão 3D}
Construir um objeto físico depositando material camada por camada, guiado por um arquivo digital (G-code). De protótipos a próteses e peças de foguete: a fabricação que parte direto do bit para a forma, sem molde.
\prov{relações: eh\_um manufatura; usa atuador}
\prov{proveniência: wiki \textperiodcentered{} impressão 3D e manufatura \textperiodcentered{} fato \textperiodcentered{} confiança 1.0 \textperiodcentered{} domínio manufatura \textperiodcentered{} história 3 versões no jornal}

%CRISTAL {"arestas": [{"alvo": "impressao_3d", "confianca": 1.0, "epistemico": "fato", "origem": "wiki", "rel": "explica"}, {"alvo": "codificacao_decodificacao", "confianca": 1.0, "epistemico": "fato", "origem": "wiki", "rel": "usa"}, {"alvo": "transformada_mineral_7d", "confianca": 1.0, "epistemico": "fato", "origem": "wiki", "rel": "relaciona"}, {"alvo": "sensor_e_colapso", "confianca": 1.0, "epistemico": "fato", "origem": "wiki", "rel": "relaciona"}, {"alvo": "corpo_mineral", "confianca": 1.0, "epistemico": "fato", "origem": "wiki", "rel": "usa"}], "confianca": 1.0, "contraexemplos": [], "descricao": "Onde o sensor colapsava matéria→bits (codificar), a impressora faz bits→matéria (DECODIFICAR): o G-code é a órbita que a cabeça segue reconstruindo a forma a partir de sua descrição digital. São as duas metades da transformada mineral — codificar a informação e voltar a materializá-la.", "epistemico": "inferencia", "exemplos": [], "id": "impressao_decodifica", "memoria": "semantica", "meta": {"dominio": "manufatura", "fonte": "impressão 3D e manufatura"}, "origem": "wiki", "palavras_chave": [], "sinonimos": [], "tipo": "conceito", "titulo": "Impressão 3D = Decodificar (o inverso do sensor)"}
\section{Impressão 3D = Decodificar (o inverso do sensor)}
Onde o sensor colapsava matéria\ensuremath{\to}bits (codificar), a impressora faz bits\ensuremath{\to}matéria (DECODIFICAR): o G-code é a órbita que a cabeça segue reconstruindo a forma a partir de sua descrição digital. São as duas metades da transformada mineral — codificar a informação e voltar a materializá-la.
\prov{relações: explica impressao\_3d; usa codificacao\_decodificacao; relaciona transformada\_mineral\_7d; relaciona sensor\_e\_colapso; usa corpo\_mineral}
\prov{proveniência: wiki \textperiodcentered{} impressão 3D e manufatura \textperiodcentered{} inferencia \textperiodcentered{} confiança 1.0 \textperiodcentered{} domínio manufatura \textperiodcentered{} história 6 versões no jornal}

%CRISTAL {"arestas": [{"alvo": "energia_eletrica", "confianca": 1.0, "epistemico": "fato", "origem": "wiki", "rel": "explica"}], "confianca": 1.0, "contraexemplos": [], "descricao": "Faraday: um campo magnético variável induz tensão num condutor — o princípio do gerador, do motor e do transformador. A base da era elétrica.", "epistemico": "fato", "exemplos": [], "id": "inducao_eletromagnetica", "memoria": "semantica", "meta": {"autor": "Faraday", "dominio": "engenharia", "fonte": ""}, "origem": "wiki", "palavras_chave": [], "sinonimos": [], "tipo": "conceito", "titulo": "Indução Eletromagnética"}
\section{Indução Eletromagnética}
Faraday: um campo magnético variável induz tensão num condutor — o princípio do gerador, do motor e do transformador. A base da era elétrica.
\prov{relações: explica energia\_eletrica}
\prov{proveniência: wiki \textperiodcentered{} fato \textperiodcentered{} confiança 1.0 \textperiodcentered{} domínio engenharia \textperiodcentered{} história 2 versões no jornal}

%CRISTAL {"arestas": [{"alvo": "circuitos_eletricos", "confianca": 1.0, "epistemico": "fato", "origem": "wiki", "rel": "usa"}, {"alvo": "conversao_ad_da", "confianca": 1.0, "epistemico": "fato", "origem": "wiki", "rel": "usa"}, {"alvo": "eng_eletronica", "confianca": 1.0, "epistemico": "fato", "origem": "curriculo", "rel": "parte_de"}, {"alvo": "eng_telecomunicacoes", "confianca": 1.0, "epistemico": "fato", "origem": "curriculo", "rel": "parte_de"}, {"alvo": "eng_sistemas_energia", "confianca": 1.0, "epistemico": "fato", "origem": "curriculo", "rel": "parte_de"}], "confianca": 1.0, "contraexemplos": [], "descricao": "Medir tensão, corrente, potência e grandezas físicas com sensores, transdutores e incerteza controlada. Sem medida confiável não há projeto, controle nem diagnóstico.", "epistemico": "inferencia", "exemplos": [], "id": "instrumentacao_medidas", "memoria": "semantica", "meta": {"dominio": "engenharia", "fonte": ""}, "origem": "wiki", "palavras_chave": [], "sinonimos": [], "tipo": "conceito", "titulo": "Instrumentação e Medidas Elétricas"}
\section{Instrumentação e Medidas Elétricas}
Medir tensão, corrente, potência e grandezas físicas com sensores, transdutores e incerteza controlada. Sem medida confiável não há projeto, controle nem diagnóstico.
\prov{relações: usa circuitos\_eletricos; usa conversao\_ad\_da; parte\_de eng\_eletronica; parte\_de eng\_telecomunicacoes; parte\_de eng\_sistemas\_energia}
\prov{proveniência: wiki \textperiodcentered{} inferencia \textperiodcentered{} confiança 1.0 \textperiodcentered{} domínio engenharia \textperiodcentered{} história 6 versões no jornal}

%CRISTAL {"arestas": [{"alvo": "smart_city", "confianca": 1.0, "epistemico": "fato", "origem": "wiki", "rel": "relaciona"}], "confianca": 1.0, "contraexemplos": [], "descricao": "A malha de objetos físicos com sensores, processador e rede que medem o mundo e agem sobre ele — de medidores de energia a semáforos e geladeiras. São os OLHOS e as mãos da smart city: o mundo físico virando nós de um grafo de dados.", "epistemico": "fato", "exemplos": [], "id": "internet_das_coisas", "memoria": "semantica", "meta": {"dominio": "engenharia", "fonte": "IoT e sensores"}, "origem": "wiki", "palavras_chave": [], "sinonimos": [], "tipo": "conceito", "titulo": "Internet das Coisas (IoT)"}
\section{Internet das Coisas (IoT)}
A malha de objetos físicos com sensores, processador e rede que medem o mundo e agem sobre ele — de medidores de energia a semáforos e geladeiras. São os OLHOS e as mãos da smart city: o mundo físico virando nós de um grafo de dados.
\prov{relações: relaciona smart\_city}
\prov{proveniência: wiki \textperiodcentered{} IoT e sensores \textperiodcentered{} fato \textperiodcentered{} confiança 1.0 \textperiodcentered{} domínio engenharia \textperiodcentered{} história 4 versões no jornal}

%CRISTAL {"arestas": [{"alvo": "carro_eletrico", "confianca": 1.0, "epistemico": "fato", "origem": "wiki", "rel": "parte_de"}, {"alvo": "corrente_alternada", "confianca": 1.0, "epistemico": "fato", "origem": "wiki", "rel": "usa"}, {"alvo": "motor_eletrico_tracao", "confianca": 1.0, "epistemico": "fato", "origem": "wiki", "rel": "usa"}], "confianca": 1.0, "contraexemplos": [], "descricao": "Transforma a CC da bateria em CA trifásica de frequência variável para o motor (e o contrário na frenagem). É o que faz a ROTAÇÃO do motor a partir do cristal — o coração eletrônico do VE.", "epistemico": "fato", "exemplos": [], "id": "inversor_de_tracao", "memoria": "semantica", "meta": {"dominio": "engenharia", "fonte": "mobilidade elétrica"}, "origem": "wiki", "palavras_chave": [], "sinonimos": [], "tipo": "conceito", "titulo": "Inversor de Tração"}
\section{Inversor de Tração}
Transforma a CC da bateria em CA trifásica de frequência variável para o motor (e o contrário na frenagem). É o que faz a ROTAÇÃO do motor a partir do cristal — o coração eletrônico do VE.
\prov{relações: parte\_de carro\_eletrico; usa corrente\_alternada; usa motor\_eletrico\_tracao}
\prov{proveniência: wiki \textperiodcentered{} mobilidade elétrica \textperiodcentered{} fato \textperiodcentered{} confiança 1.0 \textperiodcentered{} domínio engenharia \textperiodcentered{} história 4 versões no jornal}

%CRISTAL {"arestas": [{"alvo": "seguranca_de_redes", "confianca": 1.0, "epistemico": "fato", "origem": "wiki", "rel": "especializa"}, {"alvo": "protocolo_ip", "confianca": 1.0, "epistemico": "fato", "origem": "wiki", "rel": "usa"}, {"alvo": "eng_telecomunicacoes", "confianca": 1.0, "epistemico": "fato", "origem": "curriculo", "rel": "parte_de"}], "confianca": 1.0, "contraexemplos": [], "descricao": "Cifrar e autenticar no nível do IP para criar um túnel seguro sobre a internet pública — a rede privada virtual. Estende a LAN confiável através de um meio hostil.", "epistemico": "fato", "exemplos": [], "id": "ipsec_vpn", "memoria": "semantica", "meta": {"dominio": "engenharia", "fonte": ""}, "origem": "wiki", "palavras_chave": [], "sinonimos": [], "tipo": "conceito", "titulo": "IPsec e VPN"}
\section{IPsec e VPN}
Cifrar e autenticar no nível do IP para criar um túnel seguro sobre a internet pública — a rede privada virtual. Estende a LAN confiável através de um meio hostil.
\prov{relações: especializa seguranca\_de\_redes; usa protocolo\_ip; parte\_de eng\_telecomunicacoes}
\prov{proveniência: wiki \textperiodcentered{} fato \textperiodcentered{} confiança 1.0 \textperiodcentered{} domínio engenharia \textperiodcentered{} história 4 versões no jornal}

%CRISTAL {"arestas": [{"alvo": "protocolo_ip", "confianca": 1.0, "epistemico": "fato", "origem": "wiki", "rel": "especializa"}, {"alvo": "eng_telecomunicacoes", "confianca": 1.0, "epistemico": "fato", "origem": "curriculo", "rel": "parte_de"}], "confianca": 1.0, "contraexemplos": [], "descricao": "IPv4 endereça com 32 bits (esgotados) e IPv6 com 128 — o endereçamento hierárquico e roteável que faz a internet ser uma só. Sub-redes, máscara e prefixo definem quem é local e quem exige roteamento.", "epistemico": "fato", "exemplos": [], "id": "ipv4_ipv6", "memoria": "semantica", "meta": {"dominio": "engenharia", "fonte": ""}, "origem": "wiki", "palavras_chave": [], "sinonimos": [], "tipo": "conceito", "titulo": "IPv4 e IPv6"}
\section{IPv4 e IPv6}
IPv4 endereça com 32 bits (esgotados) e IPv6 com 128 — o endereçamento hierárquico e roteável que faz a internet ser uma só. Sub-redes, máscara e prefixo definem quem é local e quem exige roteamento.
\prov{relações: especializa protocolo\_ip; parte\_de eng\_telecomunicacoes}
\prov{proveniência: wiki \textperiodcentered{} fato \textperiodcentered{} confiança 1.0 \textperiodcentered{} domínio engenharia \textperiodcentered{} história 3 versões no jornal}

%CRISTAL {"arestas": [{"alvo": "dopagem", "confianca": 1.0, "epistemico": "fato", "origem": "wiki", "rel": "usa"}], "confianca": 1.0, "contraexemplos": [], "descricao": "A fronteira entre silício tipo-P e tipo-N: deixa a corrente passar num sentido só. O coração do diodo, do transistor e da célula solar.", "epistemico": "fato", "exemplos": [], "id": "juncao_pn", "memoria": "semantica", "meta": {"dominio": "engenharia", "fonte": ""}, "origem": "wiki", "palavras_chave": [], "sinonimos": [], "tipo": "conceito", "titulo": "Junção P-N"}
\section{Junção P-N}
A fronteira entre silício tipo-P e tipo-N: deixa a corrente passar num sentido só. O coração do diodo, do transistor e da célula solar.
\prov{relações: usa dopagem}
\prov{proveniência: wiki \textperiodcentered{} fato \textperiodcentered{} confiança 1.0 \textperiodcentered{} domínio engenharia \textperiodcentered{} história 2 versões no jornal}

%CRISTAL {"arestas": [{"alvo": "segmentacao_de_imagem", "confianca": 1.0, "epistemico": "fato", "origem": "wiki", "rel": "parte_de"}, {"alvo": "gradiente_descendente", "confianca": 1.0, "epistemico": "fato", "origem": "wiki", "rel": "usa"}, {"alvo": "lyapunov_da_orbita", "confianca": 1.0, "epistemico": "fato", "origem": "wiki", "rel": "relaciona"}, {"alvo": "treino_atrator_lyapunov", "confianca": 1.0, "epistemico": "fato", "origem": "wiki", "rel": "relaciona"}], "confianca": 1.0, "contraexemplos": [], "descricao": "Agrupar os pixels em k centroides: cada pixel COLAPSA no centroide (cristal) mais próximo, e os centroides descem até o ponto fixo (Lloyd converge). Os k atratores são as cores/tons dominantes; a segmentação é o colapso de cada pixel ao seu cristal. A mesma descida a um mínimo do treino da IA e da evolução. Verif.: converge em poucas iterações aos k cristais.", "epistemico": "inferencia", "exemplos": [], "id": "kmeans_atratores", "memoria": "semantica", "meta": {"dominio": "imagem", "fonte": "segmentação"}, "origem": "wiki", "palavras_chave": [], "sinonimos": [], "tipo": "conceito", "titulo": "k-means = Descida aos Atratores (os cristais)"}
\section{k-means = Descida aos Atratores (os cristais)}
Agrupar os pixels em k centroides: cada pixel COLAPSA no centroide (cristal) mais próximo, e os centroides descem até o ponto fixo (Lloyd converge). Os k atratores são as cores/tons dominantes; a segmentação é o colapso de cada pixel ao seu cristal. A mesma descida a um mínimo do treino da IA e da evolução. Verif.: converge em poucas iterações aos k cristais.
\prov{relações: parte\_de segmentacao\_de\_imagem; usa gradiente\_descendente; relaciona lyapunov\_da\_orbita; relaciona treino\_atrator\_lyapunov}
\prov{proveniência: wiki \textperiodcentered{} segmentação \textperiodcentered{} inferencia \textperiodcentered{} confiança 1.0 \textperiodcentered{} domínio imagem \textperiodcentered{} história 10 versões no jornal}

%CRISTAL {"arestas": [{"alvo": "olho_de_cock_web", "confianca": 1.0, "epistemico": "fato", "origem": "wiki", "rel": "parte_de"}, {"alvo": "olho_de_koch", "confianca": 1.0, "epistemico": "fato", "origem": "wiki", "rel": "explica"}, {"alvo": "borda_critica_vazio", "confianca": 1.0, "epistemico": "fato", "origem": "wiki", "rel": "relaciona"}], "confianca": 1.0, "contraexemplos": [], "descricao": "Um <canvas> com a grade BSD n×k (sliders 4–128 / 1–64), dois atratores branco(c²) e preto(a) nos cantos, infusão só pelo mapa de borda z↦z² mod 1009; o front rasteriza a grade_bsd do servidor com difusão/dissipação Koopman. O olho de Koch desenhado.", "epistemico": "fato", "exemplos": [], "id": "koch_field_canvas", "memoria": "semantica", "meta": {"autor": "broca-so", "dominio": "cockpit", "fonte": "goldchain/cockpit/cockpit.js:KochField"}, "origem": "wiki", "palavras_chave": [], "sinonimos": [], "tipo": "conceito", "titulo": "Canvas BSD (KochField) — o olho na borda"}
\section{Canvas BSD (KochField) — o olho na borda}
Um <canvas> com a grade BSD n$\times$k (sliders 4–128 / 1–64), dois atratores branco(c\ensuremath{^{2}}) e preto(a) nos cantos, infusão só pelo mapa de borda z\ensuremath{\mapsto}z\ensuremath{^{2}} mod 1009; o front rasteriza a grade\_bsd do servidor com difusão/dissipação Koopman. O olho de Koch desenhado.
\prov{relações: parte\_de olho\_de\_cock\_web; explica olho\_de\_koch; relaciona borda\_critica\_vazio}
\prov{proveniência: wiki \textperiodcentered{} goldchain/cockpit/cockpit.js:KochField \textperiodcentered{} fato \textperiodcentered{} confiança 1.0 \textperiodcentered{} domínio cockpit \textperiodcentered{} história 4 versões no jornal}

%CRISTAL {"arestas": [{"alvo": "cotacao_parabola", "confianca": 1.0, "epistemico": "fato", "origem": "wiki", "rel": "usa"}, {"alvo": "observar_correlacao_hodge", "confianca": 1.0, "epistemico": "fato", "origem": "wiki", "rel": "relaciona"}, {"alvo": "ciclo_uno", "confianca": 1.0, "epistemico": "fato", "origem": "wiki", "rel": "relaciona"}], "confianca": 1.0, "contraexemplos": [], "descricao": "Constrói um operador Koopman reduzido k×k (contração por taxa de consumo + acoplamento via matriz das cúspides/Gram da grade BSD), tira autovalores e por ramo emite alerta (satura se |λ_dom|→1, esvazia se gradiente negativo) e score (ótimo quando |λ_dom|≈φ), escolhendo o ramo recomendado antes da próxima tx.", "epistemico": "fato", "exemplos": [], "id": "koopman_roteamento", "memoria": "semantica", "meta": {"autor": "broca-so", "dominio": "cockpit", "fonte": "goldchain/koopman_roteamento.py:predizer_ramos"}, "origem": "wiki", "palavras_chave": [], "sinonimos": [], "tipo": "conceito", "titulo": "Roteamento Koopman (autovalores da grade BSD)"}
\section{Roteamento Koopman (autovalores da grade BSD)}
Constrói um operador Koopman reduzido k$\times$k (contração por taxa de consumo + acoplamento via matriz das cúspides/Gram da grade BSD), tira autovalores e por ramo emite alerta (satura se |\ensuremath{\lambda}\_dom|\ensuremath{\to}1, esvazia se gradiente negativo) e score (ótimo quando |\ensuremath{\lambda}\_dom|\ensuremath{\approx}\ensuremath{\varphi}), escolhendo o ramo recomendado antes da próxima tx.
\prov{relações: usa cotacao\_parabola; relaciona observar\_correlacao\_hodge; relaciona ciclo\_uno}
\prov{proveniência: wiki \textperiodcentered{} goldchain/koopman\_roteamento.py:predizer\_ramos \textperiodcentered{} fato \textperiodcentered{} confiança 1.0 \textperiodcentered{} domínio cockpit \textperiodcentered{} história 4 versões no jornal}

%CRISTAL {"arestas": [{"alvo": "capacidade_de_shannon", "confianca": 1.0, "epistemico": "fato", "origem": "wiki", "rel": "usa"}, {"alvo": "rede_5g", "confianca": 1.0, "epistemico": "fato", "origem": "wiki", "rel": "usa"}], "confianca": 1.0, "contraexemplos": [], "descricao": "A faixa de frequências disponível B: quanto mais larga, mais dados (C cresce com B). 5G busca banda em ondas milimétricas (mmWave) — muita banda, mas alcance curto e bloqueada por paredes. O recurso escasso e leiloado da conectividade.", "epistemico": "fato", "exemplos": [], "id": "largura_de_banda", "memoria": "semantica", "meta": {"dominio": "engenharia", "fonte": "conectividade e 5G"}, "origem": "wiki", "palavras_chave": [], "sinonimos": [], "tipo": "conceito", "titulo": "Largura de Banda (espectro)"}
\section{Largura de Banda (espectro)}
A faixa de frequências disponível B: quanto mais larga, mais dados (C cresce com B). 5G busca banda em ondas milimétricas (mmWave) — muita banda, mas alcance curto e bloqueada por paredes. O recurso escasso e leiloado da conectividade.
\prov{relações: usa capacidade\_de\_shannon; usa rede\_5g}
\prov{proveniência: wiki \textperiodcentered{} conectividade e 5G \textperiodcentered{} fato \textperiodcentered{} confiança 1.0 \textperiodcentered{} domínio engenharia \textperiodcentered{} história 3 versões no jornal}

%CRISTAL {"arestas": [{"alvo": "rede_5g", "confianca": 1.0, "epistemico": "fato", "origem": "wiki", "rel": "parte_de"}, {"alvo": "conectividade", "confianca": 1.0, "epistemico": "fato", "origem": "wiki", "rel": "usa"}], "confianca": 1.0, "contraexemplos": [], "descricao": "O atraso entre pedir e receber. 5G a leva a ~1 ms, e é isso que fecha os LAÇOS de tempo real — carro autônomo, cirurgia remota, controle da rede. É o período do laço gerador+bit+loop: latência baixa deixa o colapso reagir ao mundo no ritmo do mundo.", "epistemico": "inferencia", "exemplos": [], "id": "latencia", "memoria": "semantica", "meta": {"dominio": "engenharia", "fonte": "conectividade e 5G"}, "origem": "wiki", "palavras_chave": [], "sinonimos": [], "tipo": "conceito", "titulo": "Latência (o tempo do laço)"}
\section{Latência (o tempo do laço)}
O atraso entre pedir e receber. 5G a leva a \~{}1 ms, e é isso que fecha os LAÇOS de tempo real — carro autônomo, cirurgia remota, controle da rede. É o período do laço gerador+bit+loop: latência baixa deixa o colapso reagir ao mundo no ritmo do mundo.
\prov{relações: parte\_de rede\_5g; usa conectividade}
\prov{proveniência: wiki \textperiodcentered{} conectividade e 5G \textperiodcentered{} inferencia \textperiodcentered{} confiança 1.0 \textperiodcentered{} domínio engenharia \textperiodcentered{} história 3 versões no jornal}

%CRISTAL {"arestas": [{"alvo": "juncao_pn", "confianca": 1.0, "epistemico": "fato", "origem": "wiki", "rel": "usa"}, {"alvo": "diodo", "confianca": 1.0, "epistemico": "fato", "origem": "wiki", "rel": "especializa"}], "confianca": 1.0, "contraexemplos": [], "descricao": "Diodo que emite luz quando a corrente recombina elétrons e lacunas na junção — iluminação eficiente e o pixel das telas.", "epistemico": "fato", "exemplos": [], "id": "led", "memoria": "semantica", "meta": {"dominio": "engenharia", "fonte": ""}, "origem": "wiki", "palavras_chave": [], "sinonimos": [], "tipo": "conceito", "titulo": "LED"}
\section{LED}
Diodo que emite luz quando a corrente recombina elétrons e lacunas na junção — iluminação eficiente e o pixel das telas.
\prov{relações: usa juncao\_pn; especializa diodo}
\prov{proveniência: wiki \textperiodcentered{} fato \textperiodcentered{} confiança 1.0 \textperiodcentered{} domínio engenharia \textperiodcentered{} história 3 versões no jornal}

%CRISTAL {"arestas": [{"alvo": "circuito_integrado", "confianca": 1.0, "epistemico": "fato", "origem": "wiki", "rel": "usa"}, {"alvo": "transistor", "confianca": 1.0, "epistemico": "fato", "origem": "wiki", "rel": "usa"}], "confianca": 1.0, "contraexemplos": [], "descricao": "Moore: o número de transistores num chip dobra a cada ~2 anos — a observação econômica que guiou 50 anos de progresso digital.", "epistemico": "inferencia", "exemplos": [], "id": "lei_de_moore", "memoria": "semantica", "meta": {"autor": "Moore", "dominio": "engenharia", "fonte": ""}, "origem": "wiki", "palavras_chave": [], "sinonimos": [], "tipo": "conceito", "titulo": "Lei de Moore"}
\section{Lei de Moore}
Moore: o número de transistores num chip dobra a cada \~{}2 anos — a observação econômica que guiou 50 anos de progresso digital.
\prov{relações: usa circuito\_integrado; usa transistor}
\prov{proveniência: wiki \textperiodcentered{} inferencia \textperiodcentered{} confiança 1.0 \textperiodcentered{} domínio engenharia \textperiodcentered{} história 3 versões no jornal}

%CRISTAL {"arestas": [{"alvo": "corrente_e_tensao", "confianca": 1.0, "epistemico": "fato", "origem": "wiki", "rel": "explica"}], "confianca": 1.0, "contraexemplos": [], "descricao": "V = R·I: a tensão através de um resistor é proporcional à corrente. A relação mais básica dos circuitos elétricos.", "epistemico": "fato", "exemplos": [], "id": "lei_de_ohm", "memoria": "semantica", "meta": {"autor": "Ohm", "dominio": "engenharia", "fonte": ""}, "origem": "wiki", "palavras_chave": [], "sinonimos": [], "tipo": "conceito", "titulo": "Lei de Ohm"}
\section{Lei de Ohm}
V = R\textperiodcentered I: a tensão através de um resistor é proporcional à corrente. A relação mais básica dos circuitos elétricos.
\prov{relações: explica corrente\_e\_tensao}
\prov{proveniência: wiki \textperiodcentered{} fato \textperiodcentered{} confiança 1.0 \textperiodcentered{} domínio engenharia \textperiodcentered{} história 2 versões no jornal}

%CRISTAL {"arestas": [{"alvo": "segmentacao_de_imagem", "confianca": 1.0, "epistemico": "fato", "origem": "wiki", "rel": "parte_de"}, {"alvo": "ligacao_por_overlap_limiar", "confianca": 1.0, "epistemico": "fato", "origem": "wiki", "rel": "relaciona"}, {"alvo": "ordem_borda_caos", "confianca": 1.0, "epistemico": "fato", "origem": "wiki", "rel": "relaciona"}, {"alvo": "corpo_mineral", "confianca": 1.0, "epistemico": "fato", "origem": "wiki", "rel": "usa"}], "confianca": 1.0, "contraexemplos": [], "descricao": "O histograma de intensidades tem picos — os cristais (fundo e objeto). O limiar de Otsu maximiza a variância entre-classes σ_b²=w0·w1·(μ0−μ1)²: o pico dela cai no VALE mais fundo entre os cristais — a borda que separa. Segmentar por limiar é colapsar cada pixel nessa fronteira. Verif.: dois cristais 60/200 → t na borda entre eles.", "epistemico": "inferencia", "exemplos": [], "id": "limiar_de_otsu", "memoria": "semantica", "meta": {"dominio": "imagem", "fonte": "segmentação"}, "origem": "wiki", "palavras_chave": [], "sinonimos": [], "tipo": "conceito", "titulo": "Limiar de Otsu (a borda entre dois cristais)"}
\section{Limiar de Otsu (a borda entre dois cristais)}
O histograma de intensidades tem picos — os cristais (fundo e objeto). O limiar de Otsu maximiza a variância entre-classes \ensuremath{\sigma}\_b\ensuremath{^{2}}=w0\textperiodcentered w1\textperiodcentered (\ensuremath{\mu}0\ensuremath{-}\ensuremath{\mu}1)\ensuremath{^{2}}: o pico dela cai no VALE mais fundo entre os cristais — a borda que separa. Segmentar por limiar é colapsar cada pixel nessa fronteira. Verif.: dois cristais 60/200 \ensuremath{\to} t na borda entre eles.
\prov{relações: parte\_de segmentacao\_de\_imagem; relaciona ligacao\_por\_overlap\_limiar; relaciona ordem\_borda\_caos; usa corpo\_mineral}
\prov{proveniência: wiki \textperiodcentered{} segmentação \textperiodcentered{} inferencia \textperiodcentered{} confiança 1.0 \textperiodcentered{} domínio imagem \textperiodcentered{} história 10 versões no jornal}

%CRISTAL {"arestas": [{"alvo": "pi_acima_da_reta_mineral", "confianca": 1.0, "epistemico": "fato", "origem": "wiki", "rel": "usa"}, {"alvo": "algebra_reta_mineral_7d", "confianca": 1.0, "epistemico": "fato", "origem": "wiki", "rel": "usa"}, {"alvo": "computador_mineral", "confianca": 1.0, "epistemico": "fato", "origem": "wiki", "rel": "relaciona"}], "confianca": 1.0, "contraexemplos": [], "descricao": "Em cada dimensão d a esfera tem volume π^{d/2}/Γ(d/2+1): π aparece à potência d/2. Os politopos algébricos (que a máquina computa) TENDEM à esfera — Arquimedes — mas π^{d/2} é transcendente. Em D2/D3 a máquina alcança; em D7 (o canto octoniônico, 𝔤₂=14) o politopo não preenche a bola com pontos tratáveis (maldição da dimensionalidade) e FICA AQUÉM. π^{7/2} recua além do alcance: o horizonte transcendente é o LIMITE DA MÁQUINA. Verif.: D3 π≈3.13, D7 π≈2.68 (aquém).", "epistemico": "fato", "exemplos": [], "id": "limite_da_maquina", "memoria": "semantica", "meta": {"dominio": "imagem", "fonte": "formas e limite"}, "origem": "wiki", "palavras_chave": [], "sinonimos": [], "tipo": "conceito", "titulo": "O Limite da Máquina (π^{d/2})"}
\section{O Limite da Máquina (\ensuremath{\pi}\^{}\{d/2\})}
Em cada dimensão d a esfera tem volume \ensuremath{\pi}\^{}\{d/2\}/\ensuremath{\Gamma}(d/2+1): \ensuremath{\pi} aparece à potência d/2. Os politopos algébricos (que a máquina computa) TENDEM à esfera — Arquimedes — mas \ensuremath{\pi}\^{}\{d/2\} é transcendente. Em D2/D3 a máquina alcança; em D7 (o canto octoniônico, \ensuremath{\mathfrak{g}}\ensuremath{_{2}}=14) o politopo não preenche a bola com pontos tratáveis (maldição da dimensionalidade) e FICA AQUÉM. \ensuremath{\pi}\^{}\{7/2\} recua além do alcance: o horizonte transcendente é o LIMITE DA MÁQUINA. Verif.: D3 \ensuremath{\pi}\ensuremath{\approx}3.13, D7 \ensuremath{\pi}\ensuremath{\approx}2.68 (aquém).
\prov{relações: usa pi\_acima\_da\_reta\_mineral; usa algebra\_reta\_mineral\_7d; relaciona computador\_mineral}
\prov{proveniência: wiki \textperiodcentered{} formas e limite \textperiodcentered{} fato \textperiodcentered{} confiança 1.0 \textperiodcentered{} domínio imagem \textperiodcentered{} história 7 versões no jornal}

%CRISTAL {"arestas": [{"alvo": "comprimir_aumenta_pressao", "confianca": 1.0, "epistemico": "fato", "origem": "wiki", "rel": "usa"}, {"alvo": "fusao_estelar", "confianca": 1.0, "epistemico": "fato", "origem": "wiki", "rel": "eh_um"}, {"alvo": "fusao_nuclear", "confianca": 1.0, "epistemico": "fato", "origem": "wiki", "rel": "usa"}, {"alvo": "vale_do_ferro", "confianca": 1.0, "epistemico": "fato", "origem": "wiki", "rel": "relaciona"}, {"alvo": "floco_de_neve", "confianca": 1.0, "epistemico": "fato", "origem": "wiki", "rel": "relaciona"}], "confianca": 1.0, "contraexemplos": [], "descricao": "Comprimir ao limite k=1 concentra ~toda a energia num só modo (σ_1²) — o ponto de IGNIÇÃO, a fusão. Como na fusão nuclear, comprimir além do limite libera calor. O FLOCO DE NEVE É UMA ESTRELA: côncavo (as pontas), e fundido sobe ao modo dominante — o 'ferro espectral' (Fe-56, o atrator nuclear onde a fusão para). Verif.: floco-estrela funde 99% em k=1; a curva de ligação sobe ao vale do ferro.", "epistemico": "fato", "exemplos": [], "id": "limite_de_fusao_do_floco", "memoria": "semantica", "meta": {"dominio": "imagem", "fonte": "floco térmico"}, "origem": "wiki", "palavras_chave": [], "sinonimos": [], "tipo": "conceito", "titulo": "O Limite de Fusão do Floco (a estrela)"}
\section{O Limite de Fusão do Floco (a estrela)}
Comprimir ao limite k=1 concentra \~{}toda a energia num só modo (\ensuremath{\sigma}\_1\ensuremath{^{2}}) — o ponto de IGNIÇÃO, a fusão. Como na fusão nuclear, comprimir além do limite libera calor. O FLOCO DE NEVE É UMA ESTRELA: côncavo (as pontas), e fundido sobe ao modo dominante — o 'ferro espectral' (Fe-56, o atrator nuclear onde a fusão para). Verif.: floco-estrela funde 99\% em k=1; a curva de ligação sobe ao vale do ferro.
\prov{relações: usa comprimir\_aumenta\_pressao; eh\_um fusao\_estelar; usa fusao\_nuclear; relaciona vale\_do\_ferro; relaciona floco\_de\_neve}
\prov{proveniência: wiki \textperiodcentered{} floco térmico \textperiodcentered{} fato \textperiodcentered{} confiança 1.0 \textperiodcentered{} domínio imagem \textperiodcentered{} história 11 versões no jornal}

%CRISTAL {"arestas": [{"alvo": "data_center", "confianca": 1.0, "epistemico": "fato", "origem": "wiki", "rel": "explica"}, {"alvo": "ciclo_termodinamico_orbita", "confianca": 1.0, "epistemico": "fato", "origem": "wiki", "rel": "relaciona"}, {"alvo": "computador_mineral", "confianca": 1.0, "epistemico": "fato", "origem": "wiki", "rel": "relaciona"}], "confianca": 1.0, "contraexemplos": [], "descricao": "Apagar 1 bit de informação dissipa no MÍNIMO kT·ln2 de calor (Landauer, ~2,9×10⁻²¹ J a 300K). A computação é física: a dissipação (o lado 'a' de a+c²=1) aparece no próprio ato de apagar. Só a computação reversível escaparia — por isso o data center é, no fundo, uma máquina térmica.", "epistemico": "inferencia", "exemplos": [], "id": "limite_de_landauer", "memoria": "semantica", "meta": {"dominio": "engenharia", "fonte": "nuvem e data centers", "joules_por_bit_300k": 2.870978885078724e-21}, "origem": "wiki", "palavras_chave": [], "sinonimos": [], "tipo": "conceito", "titulo": "Limite de Landauer (o chão termodinâmico)"}
\section{Limite de Landauer (o chão termodinâmico)}
Apagar 1 bit de informação dissipa no MÍNIMO kT\textperiodcentered ln2 de calor (Landauer, \~{}2,9$\times$10\ensuremath{^{-21}} J a 300K). A computação é física: a dissipação (o lado 'a' de a+c\ensuremath{^{2}}=1) aparece no próprio ato de apagar. Só a computação reversível escaparia — por isso o data center é, no fundo, uma máquina térmica.
\prov{relações: explica data\_center; relaciona ciclo\_termodinamico\_orbita; relaciona computador\_mineral}
\prov{proveniência: wiki \textperiodcentered{} nuvem e data centers \textperiodcentered{} inferencia \textperiodcentered{} confiança 1.0 \textperiodcentered{} domínio engenharia \textperiodcentered{} história 4 versões no jornal}

%CRISTAL {"arestas": [{"alvo": "logica_digital", "confianca": 1.0, "epistemico": "fato", "origem": "wiki", "rel": "usa"}, {"alvo": "linguagem_de_programacao", "confianca": 1.0, "epistemico": "fato", "origem": "wiki", "rel": "relaciona"}, {"alvo": "projeto_digital_vlsi", "confianca": 1.0, "epistemico": "fato", "origem": "wiki", "rel": "usa"}, {"alvo": "verilog", "confianca": 1.0, "epistemico": "fato", "origem": "wiki", "rel": "generaliza"}, {"alvo": "isa_da_broca", "confianca": 1.0, "epistemico": "fato", "origem": "wiki", "rel": "explica"}, {"alvo": "eng_eletronica", "confianca": 1.0, "epistemico": "fato", "origem": "curriculo", "rel": "parte_de"}], "confianca": 1.0, "contraexemplos": [], "descricao": "Descrever circuito como código concorrente (Verilog, VHDL) e sintetizar em portas — o hardware projetado como software. É a linguagem em que se escreveria a ULA-12 do sandbox.", "epistemico": "fato", "exemplos": [], "id": "linguagem_descricao_hardware", "memoria": "semantica", "meta": {"dominio": "engenharia", "fonte": ""}, "origem": "wiki", "palavras_chave": [], "sinonimos": [], "tipo": "conceito", "titulo": "Linguagens de Descrição de Hardware (HDL)"}
\section{Linguagens de Descrição de Hardware (HDL)}
Descrever circuito como código concorrente (Verilog, VHDL) e sintetizar em portas — o hardware projetado como software. É a linguagem em que se escreveria a ULA-12 do sandbox.
\prov{relações: usa logica\_digital; relaciona linguagem\_de\_programacao; usa projeto\_digital\_vlsi; generaliza verilog; explica isa\_da\_broca; parte\_de eng\_eletronica}
\prov{proveniência: wiki \textperiodcentered{} fato \textperiodcentered{} confiança 1.0 \textperiodcentered{} domínio engenharia \textperiodcentered{} história 8 versões no jornal}

%CRISTAL {"arestas": [{"alvo": "ondas_eletromagneticas", "confianca": 1.0, "epistemico": "fato", "origem": "wiki", "rel": "usa"}, {"alvo": "eletronica_de_rf", "confianca": 1.0, "epistemico": "fato", "origem": "wiki", "rel": "usa"}, {"alvo": "eng_telecomunicacoes", "confianca": 1.0, "epistemico": "fato", "origem": "curriculo", "rel": "parte_de"}], "confianca": 1.0, "contraexemplos": [], "descricao": "Conduzir sinal de alta frequência sem refletir: impedância característica, coeficiente de reflexão, ondas estacionárias, carta de Smith. Onde o fio deixa de ser fio e vira meio de onda.", "epistemico": "fato", "exemplos": [], "id": "linha_de_transmissao", "memoria": "semantica", "meta": {"dominio": "engenharia", "fonte": ""}, "origem": "wiki", "palavras_chave": [], "sinonimos": [], "tipo": "conceito", "titulo": "Linhas de Transmissão"}
\section{Linhas de Transmissão}
Conduzir sinal de alta frequência sem refletir: impedância característica, coeficiente de reflexão, ondas estacionárias, carta de Smith. Onde o fio deixa de ser fio e vira meio de onda.
\prov{relações: usa ondas\_eletromagneticas; usa eletronica\_de\_rf; parte\_de eng\_telecomunicacoes}
\prov{proveniência: wiki \textperiodcentered{} fato \textperiodcentered{} confiança 1.0 \textperiodcentered{} domínio engenharia \textperiodcentered{} história 4 versões no jornal}

%CRISTAL {"arestas": [{"alvo": "lyapunov_computacao", "confianca": 1.0, "epistemico": "fato", "origem": "wiki", "rel": "relaciona"}, {"alvo": "xadrez_bai", "confianca": 1.0, "epistemico": "fato", "origem": "wiki", "rel": "explica"}, {"alvo": "aprendizado_por_reforco", "confianca": 1.0, "epistemico": "fato", "origem": "wiki", "rel": "relaciona"}, {"alvo": "orbita_de_partida", "confianca": 1.0, "epistemico": "fato", "origem": "wiki", "rel": "explica"}, {"alvo": "transformada_metalica", "confianca": 1.0, "epistemico": "fato", "origem": "wiki", "rel": "usa"}, {"alvo": "minimax_e_colapso", "confianca": 1.0, "epistemico": "fato", "origem": "wiki", "rel": "relaciona"}], "confianca": 1.0, "contraexemplos": [], "descricao": "O expoente de Lyapunov mede a taxa em que órbitas vizinhas se separam: λ>0 = caos (um lance muda tudo — as partidas românticas de sacrifício, o material não prevê), λ<0 = contração (as estradas levam ao mesmo lugar — a vitória técnica, previsível). É a MESMA classificação da parábola (cristal=Pisot/contrai, borda=Salem, caos=fora) e o expoente da transformada metálica (o mapa x→σ_n·x tem expoentes ±log σ_n; o lado Pisot é −log σ_n<0). A escala é log φ≈0.481. Assim se PREVÊ a órbita de uma partida.", "epistemico": "inferencia", "exemplos": [], "id": "lyapunov_da_orbita", "memoria": "semantica", "meta": {"dominio": "ponte", "escala_log_phi": 0.4812, "fonte": "partidas↔órbitas"}, "origem": "wiki", "palavras_chave": [], "sinonimos": [], "tipo": "conceito", "titulo": "Lyapunov da Órbita"}
\section{Lyapunov da Órbita}
O expoente de Lyapunov mede a taxa em que órbitas vizinhas se separam: \ensuremath{\lambda}>0 = caos (um lance muda tudo — as partidas românticas de sacrifício, o material não prevê), \ensuremath{\lambda}<0 = contração (as estradas levam ao mesmo lugar — a vitória técnica, previsível). É a MESMA classificação da parábola (cristal=Pisot/contrai, borda=Salem, caos=fora) e o expoente da transformada metálica (o mapa x\ensuremath{\to}\ensuremath{\sigma}\_n\textperiodcentered x tem expoentes $\pm$log \ensuremath{\sigma}\_n; o lado Pisot é \ensuremath{-}log \ensuremath{\sigma}\_n<0). A escala é log \ensuremath{\varphi}\ensuremath{\approx}0.481. Assim se PREVÊ a órbita de uma partida.
\prov{relações: relaciona lyapunov\_computacao; explica xadrez\_bai; relaciona aprendizado\_por\_reforco; explica orbita\_de\_partida; usa transformada\_metalica; relaciona minimax\_e\_colapso}
\prov{proveniência: wiki \textperiodcentered{} partidas\ensuremath{\leftrightarrow}órbitas \textperiodcentered{} inferencia \textperiodcentered{} confiança 1.0 \textperiodcentered{} domínio ponte \textperiodcentered{} história 26 versões no jornal}

%CRISTAL {"arestas": [{"alvo": "sistemas_de_controle", "confianca": 1.0, "epistemico": "fato", "origem": "wiki", "rel": "explica"}, {"alvo": "sensor_e_colapso", "confianca": 1.0, "epistemico": "fato", "origem": "wiki", "rel": "usa"}, {"alvo": "neuronio", "confianca": 1.0, "epistemico": "fato", "origem": "wiki", "rel": "relaciona"}, {"alvo": "orbita_de_partida", "confianca": 1.0, "epistemico": "fato", "origem": "wiki", "rel": "relaciona"}, {"alvo": "corpo_mineral", "confianca": 1.0, "epistemico": "fato", "origem": "wiki", "rel": "usa"}], "confianca": 1.0, "contraexemplos": [], "descricao": "Sentir→decidir→agir→sentir: um laço fechado que é, ele mesmo, uma ÓRBITA no espaço de estados. É o cabeca.py com músculos — gerador (a política) + bit (a leitura do sensor) + loop. O robô não é uma coisa, é um laço que se sustenta. Verif.: o laço sensório-motor converge ao alvo (erro→0).", "epistemico": "inferencia", "exemplos": [], "id": "malha_de_controle_laco", "memoria": "semantica", "meta": {"dominio": "robotica", "erro_final_laco": 0.0022, "fonte": "robótica e sistemas autônomos"}, "origem": "wiki", "palavras_chave": [], "sinonimos": [], "tipo": "conceito", "titulo": "A Malha de Controle é o Laço"}
\section{A Malha de Controle é o Laço}
Sentir\ensuremath{\to}decidir\ensuremath{\to}agir\ensuremath{\to}sentir: um laço fechado que é, ele mesmo, uma ÓRBITA no espaço de estados. É o cabeca.py com músculos — gerador (a política) + bit (a leitura do sensor) + loop. O robô não é uma coisa, é um laço que se sustenta. Verif.: o laço sensório-motor converge ao alvo (erro\ensuremath{\to}0).
\prov{relações: explica sistemas\_de\_controle; usa sensor\_e\_colapso; relaciona neuronio; relaciona orbita\_de\_partida; usa corpo\_mineral}
\prov{proveniência: wiki \textperiodcentered{} robótica e sistemas autônomos \textperiodcentered{} inferencia \textperiodcentered{} confiança 1.0 \textperiodcentered{} domínio robotica \textperiodcentered{} história 6 versões no jornal}

%CRISTAL {"arestas": [{"alvo": "robotica", "confianca": 1.0, "epistemico": "fato", "origem": "wiki", "rel": "relaciona"}, {"alvo": "energia_eletrica", "confianca": 1.0, "epistemico": "fato", "origem": "wiki", "rel": "usa"}], "confianca": 1.0, "contraexemplos": [], "descricao": "Transformar matéria-prima em objetos úteis. A subtrativa remove material (torno, fresa); a ADITIVA acrescenta (impressão 3D). É onde a decisão da máquina vira coisa — o fim físico da cadeia de informação que começou no sensor.", "epistemico": "fato", "exemplos": [], "id": "manufatura", "memoria": "semantica", "meta": {"dominio": "manufatura", "fonte": "impressão 3D e manufatura"}, "origem": "wiki", "palavras_chave": [], "sinonimos": [], "tipo": "conceito", "titulo": "Manufatura"}
\section{Manufatura}
Transformar matéria-prima em objetos úteis. A subtrativa remove material (torno, fresa); a ADITIVA acrescenta (impressão 3D). É onde a decisão da máquina vira coisa — o fim físico da cadeia de informação que começou no sensor.
\prov{relações: relaciona robotica; usa energia\_eletrica}
\prov{proveniência: wiki \textperiodcentered{} impressão 3D e manufatura \textperiodcentered{} fato \textperiodcentered{} confiança 1.0 \textperiodcentered{} domínio manufatura \textperiodcentered{} história 3 versões no jornal}

%CRISTAL {"arestas": [{"alvo": "impressao_3d", "confianca": 1.0, "epistemico": "fato", "origem": "wiki", "rel": "explica"}, {"alvo": "orbita_de_partida", "confianca": 1.0, "epistemico": "fato", "origem": "wiki", "rel": "relaciona"}], "confianca": 1.0, "contraexemplos": [], "descricao": "Construir por ADIÇÃO iterativa: cada camada assenta sobre a anterior (camada_{n+1} sobre camada_n). É uma construção recursiva — uma órbita no espaço de build, empilhando a forma passo a passo em vez de esculpi-la.", "epistemico": "fato", "exemplos": [], "id": "manufatura_aditiva", "memoria": "semantica", "meta": {"dominio": "manufatura", "fonte": "impressão 3D e manufatura"}, "origem": "wiki", "palavras_chave": [], "sinonimos": [], "tipo": "conceito", "titulo": "Manufatura Aditiva (camada a camada)"}
\section{Manufatura Aditiva (camada a camada)}
Construir por ADIÇÃO iterativa: cada camada assenta sobre a anterior (camada\_\{n+1\} sobre camada\_n). É uma construção recursiva — uma órbita no espaço de build, empilhando a forma passo a passo em vez de esculpi-la.
\prov{relações: explica impressao\_3d; relaciona orbita\_de\_partida}
\prov{proveniência: wiki \textperiodcentered{} impressão 3D e manufatura \textperiodcentered{} fato \textperiodcentered{} confiança 1.0 \textperiodcentered{} domínio manufatura \textperiodcentered{} história 3 versões no jornal}

%CRISTAL {"arestas": [{"alvo": "conversao_eletromecanica", "confianca": 1.0, "epistemico": "fato", "origem": "wiki", "rel": "usa"}, {"alvo": "eng_sistemas_energia", "confianca": 1.0, "epistemico": "fato", "origem": "curriculo", "rel": "parte_de"}], "confianca": 1.0, "contraexemplos": [], "descricao": "Dispositivos de conversão eletromecânica rotativa — o mesmo objeto é motor (elétrico→mecânico) ou gerador (mecânico→elétrico). Enrolamentos, campo, torque e a curva de conjugado.", "epistemico": "fato", "exemplos": [], "id": "maquina_eletrica", "memoria": "semantica", "meta": {"dominio": "engenharia", "fonte": ""}, "origem": "wiki", "palavras_chave": [], "sinonimos": [], "tipo": "conceito", "titulo": "Máquinas Elétricas"}
\section{Máquinas Elétricas}
Dispositivos de conversão eletromecânica rotativa — o mesmo objeto é motor (elétrico\ensuremath{\to}mecânico) ou gerador (mecânico\ensuremath{\to}elétrico). Enrolamentos, campo, torque e a curva de conjugado.
\prov{relações: usa conversao\_eletromecanica; parte\_de eng\_sistemas\_energia}
\prov{proveniência: wiki \textperiodcentered{} fato \textperiodcentered{} confiança 1.0 \textperiodcentered{} domínio engenharia \textperiodcentered{} história 3 versões no jornal}

%CRISTAL {"arestas": [{"alvo": "mimo", "confianca": 1.0, "epistemico": "fato", "origem": "wiki", "rel": "especializa"}, {"alvo": "antena", "confianca": 1.0, "epistemico": "fato", "origem": "wiki", "rel": "usa"}, {"alvo": "faixas_fr1_fr2", "confianca": 1.0, "epistemico": "fato", "origem": "wiki", "rel": "usa"}], "confianca": 1.0, "contraexemplos": [], "descricao": "Dezenas a centenas de antenas no gNB formam FEIXES que apontam a energia a cada usuário (beamforming) e servem muitos em paralelo (MU-MIMO). É o que torna a mmWave viável e multiplica a capacidade.", "epistemico": "fato", "exemplos": [], "id": "massive_mimo_beamforming", "memoria": "semantica", "meta": {"dominio": "engenharia", "fonte": ""}, "origem": "wiki", "palavras_chave": [], "sinonimos": [], "tipo": "conceito", "titulo": "Massive MIMO e Beamforming"}
\section{Massive MIMO e Beamforming}
Dezenas a centenas de antenas no gNB formam FEIXES que apontam a energia a cada usuário (beamforming) e servem muitos em paralelo (MU-MIMO). É o que torna a mmWave viável e multiplica a capacidade.
\prov{relações: especializa mimo; usa antena; usa faixas\_fr1\_fr2}
\prov{proveniência: wiki \textperiodcentered{} fato \textperiodcentered{} confiança 1.0 \textperiodcentered{} domínio engenharia \textperiodcentered{} história 4 versões no jornal}

%CRISTAL {"arestas": [{"alvo": "semicondutor", "confianca": 1.0, "epistemico": "fato", "origem": "wiki", "rel": "explica"}, {"alvo": "cristalografia", "confianca": 1.0, "epistemico": "fato", "origem": "wiki", "rel": "usa"}, {"alvo": "eng_eletronica", "confianca": 1.0, "epistemico": "fato", "origem": "curriculo", "rel": "parte_de"}, {"alvo": "eng_telecomunicacoes", "confianca": 1.0, "epistemico": "fato", "origem": "curriculo", "rel": "parte_de"}, {"alvo": "eng_sistemas_energia", "confianca": 1.0, "epistemico": "fato", "origem": "curriculo", "rel": "parte_de"}], "confianca": 1.0, "contraexemplos": [], "descricao": "Condutores, isolantes, dielétricos, magnéticos e semicondutores — por que cada material conduz, isola ou chaveia. A base física que liga a cristalografia ao dispositivo.", "epistemico": "fato", "exemplos": [], "id": "materiais_eletricos", "memoria": "semantica", "meta": {"dominio": "engenharia", "fonte": ""}, "origem": "wiki", "palavras_chave": [], "sinonimos": [], "tipo": "conceito", "titulo": "Ciência dos Materiais Elétricos"}
\section{Ciência dos Materiais Elétricos}
Condutores, isolantes, dielétricos, magnéticos e semicondutores — por que cada material conduz, isola ou chaveia. A base física que liga a cristalografia ao dispositivo.
\prov{relações: explica semicondutor; usa cristalografia; parte\_de eng\_eletronica; parte\_de eng\_telecomunicacoes; parte\_de eng\_sistemas\_energia}
\prov{proveniência: wiki \textperiodcentered{} fato \textperiodcentered{} confiança 1.0 \textperiodcentered{} domínio engenharia \textperiodcentered{} história 6 versões no jornal}

%CRISTAL {"arestas": [{"alvo": "flip_e_rotacao_de_wick", "confianca": 1.0, "epistemico": "fato", "origem": "wiki", "rel": "eh_um"}, {"alvo": "flip_duopolinomios", "confianca": 1.0, "epistemico": "fato", "origem": "wiki", "rel": "relaciona"}, {"alvo": "mecanica_reta_mineral", "confianca": 1.0, "epistemico": "fato", "origem": "wiki", "rel": "usa"}, {"alvo": "cuspide_pisot", "confianca": 1.0, "epistemico": "fato", "origem": "wiki", "rel": "relaciona"}, {"alvo": "o_chicote_a_limpeza_nao_um_operador", "confianca": 1.0, "epistemico": "fato", "origem": "wiki", "rel": "relaciona"}, {"alvo": "flip_pg_pa", "confianca": 1.0, "epistemico": "fato", "origem": "wiki", "rel": "relaciona"}, {"alvo": "movimento_circular_dois_modos", "confianca": 1.0, "epistemico": "fato", "origem": "wiki", "rel": "relaciona"}, {"alvo": "parte_fria_t_zero", "confianca": 1.0, "epistemico": "fato", "origem": "wiki", "rel": "relaciona"}], "confianca": 1.0, "contraexemplos": [], "descricao": "O par de Tusi: um círculo de raio r rolando DENTRO de outro de raio R=k·r faz cada ponto traçar um HIPOCICLÓIDE de k cúspides. No caso k=2 (R=2r) o hipociclóide DEGENERA numa RETA — o diâmetro (|y|→0, ρ=σ_2/σ_1=0). Movimento circular vira retilíneo: é o FLIP feito mecânica — o que u=ln z faz às espirais (endireita), La Hire faz aos círculos. A reta é o cristal, o estado mais ligado. E a figura é sempre DUAS rotações contrárias (freq +1 e −(k−1)): em k=2 o par perfeito (+1,−1) cancela um eixo.", "epistemico": "fato", "exemplos": [], "id": "mecanismo_de_la_hire", "memoria": "semantica", "meta": {"dominio": "imagem", "fonte": "la hire / nucleossíntese"}, "origem": "wiki", "palavras_chave": [], "sinonimos": [], "tipo": "conceito", "titulo": "O Mecanismo de La Hire (o círculo que vira reta)"}
\section{O Mecanismo de La Hire (o círculo que vira reta)}
O par de Tusi: um círculo de raio r rolando DENTRO de outro de raio R=k\textperiodcentered r faz cada ponto traçar um HIPOCICLÓIDE de k cúspides. No caso k=2 (R=2r) o hipociclóide DEGENERA numa RETA — o diâmetro (|y|\ensuremath{\to}0, \ensuremath{\rho}=\ensuremath{\sigma}\_2/\ensuremath{\sigma}\_1=0). Movimento circular vira retilíneo: é o FLIP feito mecânica — o que u=ln z faz às espirais (endireita), La Hire faz aos círculos. A reta é o cristal, o estado mais ligado. E a figura é sempre DUAS rotações contrárias (freq +1 e \ensuremath{-}(k\ensuremath{-}1)): em k=2 o par perfeito (+1,\ensuremath{-}1) cancela um eixo.
\prov{relações: eh\_um flip\_e\_rotacao\_de\_wick; relaciona flip\_duopolinomios; usa mecanica\_reta\_mineral; relaciona cuspide\_pisot; relaciona o\_chicote\_a\_limpeza\_nao\_um\_operador; relaciona flip\_pg\_pa; relaciona movimento\_circular\_dois\_modos; relaciona parte\_fria\_t\_zero}
\prov{proveniência: wiki \textperiodcentered{} la hire / nucleossíntese \textperiodcentered{} fato \textperiodcentered{} confiança 1.0 \textperiodcentered{} domínio imagem \textperiodcentered{} história 12 versões no jornal}

%CRISTAL {"arestas": [{"alvo": "camada_rrc_nr", "confianca": 1.0, "epistemico": "fato", "origem": "wiki", "rel": "usa"}, {"alvo": "sinais_de_referencia_nr", "confianca": 1.0, "epistemico": "fato", "origem": "wiki", "rel": "usa"}], "confianca": 1.0, "contraexemplos": [], "descricao": "O gNB configura o UE para medir a potência/qualidade (RSRP/RSRQ/SINR) da célula servida e das vizinhas, e reportar quando um EVENTO dispara — o A3 ('vizinha ficou melhor que a servida por um offset') é o gatilho clássico do handover.", "epistemico": "fato", "exemplos": [], "id": "medicoes_mobilidade", "memoria": "semantica", "meta": {"dominio": "engenharia", "fonte": ""}, "origem": "wiki", "palavras_chave": [], "sinonimos": [], "tipo": "conceito", "titulo": "Medições e Relatório (Eventos A1–A6)"}
\section{Medições e Relatório (Eventos A1–A6)}
O gNB configura o UE para medir a potência/qualidade (RSRP/RSRQ/SINR) da célula servida e das vizinhas, e reportar quando um EVENTO dispara — o A3 ('vizinha ficou melhor que a servida por um offset') é o gatilho clássico do handover.
\prov{relações: usa camada\_rrc\_nr; usa sinais\_de\_referencia\_nr}
\prov{proveniência: wiki \textperiodcentered{} fato \textperiodcentered{} confiança 1.0 \textperiodcentered{} domínio engenharia \textperiodcentered{} história 3 versões no jornal}

%CRISTAL {"arestas": [{"alvo": "rede_autonoma_broca", "confianca": 1.0, "epistemico": "fato", "origem": "wiki", "rel": "parte_de"}, {"alvo": "meta_inducao_geral_simbolo", "confianca": 1.0, "epistemico": "fato", "origem": "wiki", "rel": "usa"}, {"alvo": "lei_unica_tres_dominios", "confianca": 1.0, "epistemico": "fato", "origem": "wiki", "rel": "usa"}], "confianca": 1.0, "contraexemplos": [], "descricao": "O mesmo operador (a parábola a+c², a álgebra de Peirce) rege comunicação, mineração, execução e fala — 'só o símbolo muda'. Por isso a rede é UMA álgebra atravessando domínios: o canal, o floco e o colapso são a mesma meta-indução em símbolos diferentes. É o que faz o organismo ser um, não um enxame.", "epistemico": "inferencia", "exemplos": [], "id": "meta_inducao_fio_da_rede", "memoria": "semantica", "meta": {"dominio": "arquitetura", "fonte": "semilla_rede_autonoma"}, "origem": "wiki", "palavras_chave": [], "sinonimos": [], "tipo": "conceito", "titulo": "A meta-indução como fio da rede: só o símbolo muda"}
\section{A meta-indução como fio da rede: só o símbolo muda}
O mesmo operador (a parábola a+c\ensuremath{^{2}}, a álgebra de Peirce) rege comunicação, mineração, execução e fala — 'só o símbolo muda'. Por isso a rede é UMA álgebra atravessando domínios: o canal, o floco e o colapso são a mesma meta-indução em símbolos diferentes. É o que faz o organismo ser um, não um enxame.
\prov{relações: parte\_de rede\_autonoma\_broca; usa meta\_inducao\_geral\_simbolo; usa lei\_unica\_tres\_dominios}
\prov{proveniência: wiki \textperiodcentered{} semilla\_rede\_autonoma \textperiodcentered{} inferencia \textperiodcentered{} confiança 1.0 \textperiodcentered{} domínio arquitetura \textperiodcentered{} história 4 versões no jornal}

%CRISTAL {"arestas": [{"alvo": "microprocessador_arquitetura", "confianca": 1.0, "epistemico": "fato", "origem": "wiki", "rel": "especializa"}, {"alvo": "conversor_ad_da_disc", "confianca": 1.0, "epistemico": "fato", "origem": "wiki", "rel": "usa"}, {"alvo": "eng_eletronica", "confianca": 1.0, "epistemico": "fato", "origem": "curriculo", "rel": "parte_de"}], "confianca": 1.0, "contraexemplos": [], "descricao": "Um computador inteiro num chip — CPU, memória, timers e periféricos (GPIO, ADC, UART). O cérebro barato do sistema embarcado; onde o software encontra o mundo físico.", "epistemico": "fato", "exemplos": [], "id": "microcontrolador", "memoria": "semantica", "meta": {"dominio": "engenharia", "fonte": ""}, "origem": "wiki", "palavras_chave": [], "sinonimos": [], "tipo": "conceito", "titulo": "Microcontrolador"}
\section{Microcontrolador}
Um computador inteiro num chip — CPU, memória, timers e periféricos (GPIO, ADC, UART). O cérebro barato do sistema embarcado; onde o software encontra o mundo físico.
\prov{relações: especializa microprocessador\_arquitetura; usa conversor\_ad\_da\_disc; parte\_de eng\_eletronica}
\prov{proveniência: wiki \textperiodcentered{} fato \textperiodcentered{} confiança 1.0 \textperiodcentered{} domínio engenharia \textperiodcentered{} história 4 versões no jornal}

%CRISTAL {"arestas": [{"alvo": "semicondutor", "confianca": 1.0, "epistemico": "fato", "origem": "wiki", "rel": "usa"}, {"alvo": "circuito_integrado", "confianca": 1.0, "epistemico": "fato", "origem": "wiki", "rel": "explica"}, {"alvo": "lei_de_moore", "confianca": 1.0, "epistemico": "fato", "origem": "wiki", "rel": "causa"}, {"alvo": "eng_eletronica", "confianca": 1.0, "epistemico": "fato", "origem": "curriculo", "rel": "parte_de"}], "confianca": 1.0, "contraexemplos": [], "descricao": "Fabricar milhões de transistores num chip: dopagem, litografia, oxidação, deposição — o processo CMOS que imprime a lógica no silício. A engenharia por trás da Lei de Moore.", "epistemico": "fato", "exemplos": [], "id": "microeletronica", "memoria": "semantica", "meta": {"dominio": "engenharia", "fonte": ""}, "origem": "wiki", "palavras_chave": [], "sinonimos": [], "tipo": "conceito", "titulo": "Microeletrônica"}
\section{Microeletrônica}
Fabricar milhões de transistores num chip: dopagem, litografia, oxidação, deposição — o processo CMOS que imprime a lógica no silício. A engenharia por trás da Lei de Moore.
\prov{relações: usa semicondutor; explica circuito\_integrado; causa lei\_de\_moore; parte\_de eng\_eletronica}
\prov{proveniência: wiki \textperiodcentered{} fato \textperiodcentered{} confiança 1.0 \textperiodcentered{} domínio engenharia \textperiodcentered{} história 5 versões no jornal}

%CRISTAL {"arestas": [{"alvo": "eletronica_digital", "confianca": 1.0, "epistemico": "fato", "origem": "wiki", "rel": "usa"}, {"alvo": "circuito_integrado", "confianca": 1.0, "epistemico": "fato", "origem": "wiki", "rel": "usa"}, {"alvo": "isa_da_broca", "confianca": 1.0, "epistemico": "fato", "origem": "wiki", "rel": "explica"}, {"alvo": "eng_eletronica", "confianca": 1.0, "epistemico": "fato", "origem": "curriculo", "rel": "parte_de"}], "confianca": 1.0, "contraexemplos": [], "descricao": "A ULA, os registradores, o contador de programa e o ciclo busca-decodifica-executa sobre um conjunto de instruções (ISA). É exatamente a ISA-12 que o sandbox do broca-so executa a partir de transistores.", "epistemico": "fato", "exemplos": [], "id": "microprocessador_arquitetura", "memoria": "semantica", "meta": {"autor": "von Neumann", "dominio": "engenharia", "fonte": ""}, "origem": "wiki", "palavras_chave": [], "sinonimos": [], "tipo": "conceito", "titulo": "Microprocessador e Arquitetura"}
\section{Microprocessador e Arquitetura}
A ULA, os registradores, o contador de programa e o ciclo busca-decodifica-executa sobre um conjunto de instruções (ISA). É exatamente a ISA-12 que o sandbox do broca-so executa a partir de transistores.
\prov{relações: usa eletronica\_digital; usa circuito\_integrado; explica isa\_da\_broca; parte\_de eng\_eletronica}
\prov{proveniência: wiki \textperiodcentered{} fato \textperiodcentered{} confiança 1.0 \textperiodcentered{} domínio engenharia \textperiodcentered{} história 5 versões no jornal}

%CRISTAL {"arestas": [{"alvo": "rede_eletrica", "confianca": 1.0, "epistemico": "fato", "origem": "wiki", "rel": "parte_de"}, {"alvo": "armazenamento_estabiliza_rede", "confianca": 1.0, "epistemico": "fato", "origem": "wiki", "rel": "usa"}, {"alvo": "numero_de_salem", "confianca": 1.0, "epistemico": "fato", "origem": "wiki", "rel": "relaciona"}, {"alvo": "estabilidade_da_rede", "confianca": 1.0, "epistemico": "fato", "origem": "wiki", "rel": "usa"}], "confianca": 1.0, "contraexemplos": [], "descricao": "Um pedaço da rede (bairro, campus, hospital) com geração e armazenamento próprios que pode operar conectado OU se ILHAR — destacar-se da rede e se manter estável sozinho numa falta. É uma sub-órbita fechada resiliente: um Salem local que resiste ao blackout geral. Verificado: 40+20 MW vs 50 de carga → ilha estável (+10).", "epistemico": "inferencia", "exemplos": [], "id": "microrrede", "memoria": "semantica", "meta": {"dominio": "engenharia", "fonte": "rede de carregamento e smart cities", "margem_exemplo": 10, "pode_ilhar": true}, "origem": "wiki", "palavras_chave": [], "sinonimos": [], "tipo": "conceito", "titulo": "Microrrede (o Salem local)"}
\section{Microrrede (o Salem local)}
Um pedaço da rede (bairro, campus, hospital) com geração e armazenamento próprios que pode operar conectado OU se ILHAR — destacar-se da rede e se manter estável sozinho numa falta. É uma sub-órbita fechada resiliente: um Salem local que resiste ao blackout geral. Verificado: 40+20 MW vs 50 de carga \ensuremath{\to} ilha estável (+10).
\prov{relações: parte\_de rede\_eletrica; usa armazenamento\_estabiliza\_rede; relaciona numero\_de\_salem; usa estabilidade\_da\_rede}
\prov{proveniência: wiki \textperiodcentered{} rede de carregamento e smart cities \textperiodcentered{} inferencia \textperiodcentered{} confiança 1.0 \textperiodcentered{} domínio engenharia \textperiodcentered{} história 10 versões no jornal}

%CRISTAL {"arestas": [{"alvo": "propagacao_de_ondas", "confianca": 1.0, "epistemico": "fato", "origem": "wiki", "rel": "usa"}, {"alvo": "capacidade_de_canal", "confianca": 1.0, "epistemico": "fato", "origem": "wiki", "rel": "confirma"}, {"alvo": "eng_telecomunicacoes", "confianca": 1.0, "epistemico": "fato", "origem": "curriculo", "rel": "parte_de"}], "confianca": 1.0, "contraexemplos": [], "descricao": "Usar várias antenas no transmissor e receptor para criar canais espaciais paralelos — multiplica a taxa sem mais banda, explorando o multipercurso. O salto de capacidade do Wi-Fi moderno e do 5G massivo.", "epistemico": "fato", "exemplos": [], "id": "mimo", "memoria": "semantica", "meta": {"dominio": "engenharia", "fonte": ""}, "origem": "wiki", "palavras_chave": [], "sinonimos": [], "tipo": "conceito", "titulo": "MIMO (Múltiplas Antenas)"}
\section{MIMO (Múltiplas Antenas)}
Usar várias antenas no transmissor e receptor para criar canais espaciais paralelos — multiplica a taxa sem mais banda, explorando o multipercurso. O salto de capacidade do Wi-Fi moderno e do 5G massivo.
\prov{relações: usa propagacao\_de\_ondas; confirma capacidade\_de\_canal; parte\_de eng\_telecomunicacoes}
\prov{proveniência: wiki \textperiodcentered{} fato \textperiodcentered{} confiança 1.0 \textperiodcentered{} domínio engenharia \textperiodcentered{} história 4 versões no jornal}

%CRISTAL {"arestas": [{"alvo": "capacidade_de_shannon", "confianca": 1.0, "epistemico": "fato", "origem": "wiki", "rel": "usa"}, {"alvo": "rede_5g", "confianca": 1.0, "epistemico": "fato", "origem": "wiki", "rel": "usa"}, {"alvo": "relacao_sinal_ruido", "confianca": 1.0, "epistemico": "fato", "origem": "wiki", "rel": "usa"}], "confianca": 1.0, "contraexemplos": [], "descricao": "Muitas antenas transmitindo juntas: por interferência construtiva elas FOCAM o feixe no usuário (beamforming) e abrem vários fluxos paralelos no mesmo espectro (MIMO). Sobe o SNR e a capacidade sem mais banda — moldar a rotação no espaço.", "epistemico": "fato", "exemplos": [], "id": "mimo_beamforming", "memoria": "semantica", "meta": {"dominio": "engenharia", "fonte": "conectividade e 5G"}, "origem": "wiki", "palavras_chave": [], "sinonimos": [], "tipo": "conceito", "titulo": "MIMO e Beamforming (focar a rotação)"}
\section{MIMO e Beamforming (focar a rotação)}
Muitas antenas transmitindo juntas: por interferência construtiva elas FOCAM o feixe no usuário (beamforming) e abrem vários fluxos paralelos no mesmo espectro (MIMO). Sobe o SNR e a capacidade sem mais banda — moldar a rotação no espaço.
\prov{relações: usa capacidade\_de\_shannon; usa rede\_5g; usa relacao\_sinal\_ruido}
\prov{proveniência: wiki \textperiodcentered{} conectividade e 5G \textperiodcentered{} fato \textperiodcentered{} confiança 1.0 \textperiodcentered{} domínio engenharia \textperiodcentered{} história 4 versões no jornal}

%CRISTAL {"arestas": [{"alvo": "minimax", "confianca": 1.0, "epistemico": "fato", "origem": "wiki", "rel": "relaciona"}, {"alvo": "minimax_alfabeta", "confianca": 1.0, "epistemico": "fato", "origem": "wiki", "rel": "usa"}, {"alvo": "xadrez_bai", "confianca": 1.0, "epistemico": "fato", "origem": "wiki", "rel": "parte_de"}, {"alvo": "valor_do_jogo_von_neumann", "confianca": 1.0, "epistemico": "fato", "origem": "wiki", "rel": "usa"}], "confianca": 1.0, "contraexemplos": [], "descricao": "O minimax do xadrez É o colapso Ψ do BAI: escolher, entre os ramos vigentes, o que resolve a árvore. A profundidade da busca É o orçamento δ (δ=0 → mede Φ e para); a poda alfa-beta É deny-overrides (o ramo dominado é cortado, a proibição vence). O mesmo algoritmo, dois nomes.", "epistemico": "inferencia", "exemplos": [], "id": "minimax_e_colapso", "memoria": "semantica", "meta": {"dominio": "ponte", "fonte": "xadrez↔BAI"}, "origem": "wiki", "palavras_chave": [], "sinonimos": [], "tipo": "conceito", "titulo": "Minimax = Colapso Ψ (a ponte mais forte)"}
\section{Minimax = Colapso \ensuremath{\Psi} (a ponte mais forte)}
O minimax do xadrez É o colapso \ensuremath{\Psi} do BAI: escolher, entre os ramos vigentes, o que resolve a árvore. A profundidade da busca É o orçamento \ensuremath{\delta} (\ensuremath{\delta}=0 \ensuremath{\to} mede \ensuremath{\Phi} e para); a poda alfa-beta É deny-overrides (o ramo dominado é cortado, a proibição vence). O mesmo algoritmo, dois nomes.
\prov{relações: relaciona minimax; usa minimax\_alfabeta; parte\_de xadrez\_bai; usa valor\_do\_jogo\_von\_neumann}
\prov{proveniência: wiki \textperiodcentered{} xadrez\ensuremath{\leftrightarrow}BAI \textperiodcentered{} inferencia \textperiodcentered{} confiança 1.0 \textperiodcentered{} domínio ponte \textperiodcentered{} história 5 versões no jornal}

%CRISTAL {"arestas": [{"alvo": "biblioteca_formas_minerais", "confianca": 1.0, "epistemico": "fato", "origem": "wiki", "rel": "parte_de"}, {"alvo": "gato", "confianca": 1.0, "epistemico": "fato", "origem": "wiki", "rel": "usa"}, {"alvo": "gato_como_cifra", "confianca": 1.0, "epistemico": "fato", "origem": "wiki", "rel": "relaciona"}, {"alvo": "imagem_como_orbita_gato", "confianca": 1.0, "epistemico": "fato", "origem": "wiki", "rel": "usa"}, {"alvo": "espiral_metalica", "confianca": 1.0, "epistemico": "fato", "origem": "wiki", "rel": "usa"}, {"alvo": "quasicristal_de_rotacoes", "confianca": 1.0, "epistemico": "fato", "origem": "wiki", "rel": "usa"}], "confianca": 1.0, "contraexemplos": [], "descricao": "Combinar duas formas: operadores diretos (multiplica, tela, xor, morph) ou o GATO COMO BATEDEIRA — agitar uma forma com o cat-map e injetar a outra a cada passo, deixando a mistura caótica fundir as duas. Misturar formas é operar suas órbitas; a mesma ida-e-volta do gato que embaralha e recupera.", "epistemico": "fato", "exemplos": [], "id": "mistura_de_formas", "memoria": "semantica", "meta": {"dominio": "imagem", "fonte": "formas minerais"}, "origem": "wiki", "palavras_chave": [], "sinonimos": [], "tipo": "conceito", "titulo": "Mistura de Formas (o gato como batedeira)"}
\section{Mistura de Formas (o gato como batedeira)}
Combinar duas formas: operadores diretos (multiplica, tela, xor, morph) ou o GATO COMO BATEDEIRA — agitar uma forma com o cat-map e injetar a outra a cada passo, deixando a mistura caótica fundir as duas. Misturar formas é operar suas órbitas; a mesma ida-e-volta do gato que embaralha e recupera.
\prov{relações: parte\_de biblioteca\_formas\_minerais; usa gato; relaciona gato\_como\_cifra; usa imagem\_como\_orbita\_gato; usa espiral\_metalica; usa quasicristal\_de\_rotacoes}
\prov{proveniência: wiki \textperiodcentered{} formas minerais \textperiodcentered{} fato \textperiodcentered{} confiança 1.0 \textperiodcentered{} domínio imagem \textperiodcentered{} história 7 versões no jornal}

%CRISTAL {"arestas": [{"alvo": "energia_eletrica", "confianca": 1.0, "epistemico": "fato", "origem": "wiki", "rel": "usa"}], "confianca": 1.0, "contraexemplos": [], "descricao": "Mover-se com eletricidade em vez de combustível: carros, ônibus, motos e bikes movidos a bateria e motor elétrico. Fecha a cadeia da energia — da usina à roda — e transfere o transporte para a rede.", "epistemico": "fato", "exemplos": [], "id": "mobilidade_eletrica", "memoria": "semantica", "meta": {"dominio": "engenharia", "fonte": "mobilidade elétrica"}, "origem": "wiki", "palavras_chave": [], "sinonimos": [], "tipo": "conceito", "titulo": "Mobilidade Elétrica"}
\section{Mobilidade Elétrica}
Mover-se com eletricidade em vez de combustível: carros, ônibus, motos e bikes movidos a bateria e motor elétrico. Fecha a cadeia da energia — da usina à roda — e transfere o transporte para a rede.
\prov{relações: usa energia\_eletrica}
\prov{proveniência: wiki \textperiodcentered{} mobilidade elétrica \textperiodcentered{} fato \textperiodcentered{} confiança 1.0 \textperiodcentered{} domínio engenharia \textperiodcentered{} história 2 versões no jornal}

%CRISTAL {"arestas": [{"alvo": "pilha_tcp_ip", "confianca": 1.0, "epistemico": "fato", "origem": "wiki", "rel": "generaliza"}, {"alvo": "rede_de_computadores", "confianca": 1.0, "epistemico": "fato", "origem": "wiki", "rel": "explica"}, {"alvo": "eng_telecomunicacoes", "confianca": 1.0, "epistemico": "fato", "origem": "curriculo", "rel": "parte_de"}], "confianca": 1.0, "contraexemplos": [], "descricao": "A referência que estratifica a comunicação em 7 camadas — física, enlace, rede, transporte, sessão, apresentação, aplicação. Cada camada serve a de cima e usa a de baixo, ignorando o resto. O mapa mental de toda rede; o TCP/IP a implementa em 4-5 camadas.", "epistemico": "fato", "exemplos": [], "id": "modelo_osi", "memoria": "semantica", "meta": {"autor": "ISO", "dominio": "engenharia", "fonte": ""}, "origem": "wiki", "palavras_chave": [], "sinonimos": [], "tipo": "conceito", "titulo": "Modelo OSI (7 Camadas)"}
\section{Modelo OSI (7 Camadas)}
A referência que estratifica a comunicação em 7 camadas — física, enlace, rede, transporte, sessão, apresentação, aplicação. Cada camada serve a de cima e usa a de baixo, ignorando o resto. O mapa mental de toda rede; o TCP/IP a implementa em 4-5 camadas.
\prov{relações: generaliza pilha\_tcp\_ip; explica rede\_de\_computadores; parte\_de eng\_telecomunicacoes}
\prov{proveniência: wiki \textperiodcentered{} fato \textperiodcentered{} confiança 1.0 \textperiodcentered{} domínio engenharia \textperiodcentered{} história 4 versões no jornal}

%CRISTAL {"arestas": [{"alvo": "autenticacao_5g_aka", "confianca": 1.0, "epistemico": "fato", "origem": "wiki", "rel": "continua"}, {"alvo": "camada_pdcp_nr", "confianca": 1.0, "epistemico": "fato", "origem": "wiki", "rel": "usa"}, {"alvo": "criptografia", "confianca": 1.0, "epistemico": "fato", "origem": "wiki", "rel": "usa"}], "confianca": 1.0, "contraexemplos": [], "descricao": "Depois de autenticar, ativam-se cifra e integridade em dois níveis: NAS (UE↔AMF) e AS (UE↔gNB, no PDCP). A partir daí toda sinalização e dado do ar vai protegido — o Security Mode Command liga o cadeado.", "epistemico": "fato", "exemplos": [], "id": "modo_seguranca_5g", "memoria": "semantica", "meta": {"dominio": "engenharia", "fonte": ""}, "origem": "wiki", "palavras_chave": [], "sinonimos": [], "tipo": "conceito", "titulo": "Ativação de Segurança (Security Mode)"}
\section{Ativação de Segurança (Security Mode)}
Depois de autenticar, ativam-se cifra e integridade em dois níveis: NAS (UE\ensuremath{\leftrightarrow}AMF) e AS (UE\ensuremath{\leftrightarrow}gNB, no PDCP). A partir daí toda sinalização e dado do ar vai protegido — o Security Mode Command liga o cadeado.
\prov{relações: continua autenticacao\_5g\_aka; usa camada\_pdcp\_nr; usa criptografia}
\prov{proveniência: wiki \textperiodcentered{} fato \textperiodcentered{} confiança 1.0 \textperiodcentered{} domínio engenharia \textperiodcentered{} história 4 versões no jornal}

%CRISTAL {"arestas": [{"alvo": "limite_de_fusao_do_floco", "confianca": 1.0, "epistemico": "fato", "origem": "wiki", "rel": "parte_de"}, {"alvo": "particulas_minerais", "confianca": 1.0, "epistemico": "fato", "origem": "wiki", "rel": "eh_um"}, {"alvo": "espectro_e_orbita_cristal", "confianca": 1.0, "epistemico": "fato", "origem": "wiki", "rel": "usa"}], "confianca": 1.0, "contraexemplos": [], "descricao": "Os modos espectrais {σ_i} do floco-estrela são as PARTÍCULAS da imagem: σ_1 (a maior massa) é o NÚCLEO, a cauda são os elétrons/fótons/ressonâncias — exatamente como particulas_minerais lê a raiz dominante como núcleo e a subdominante como elétron (por |raiz|). A física das partículas da máquina mineral, agora no espectro de uma imagem-estrela. Verif.: σ_1=57638 (núcleo), cauda elétrons.", "epistemico": "fato", "exemplos": [], "id": "modos_do_floco_sao_particulas", "memoria": "semantica", "meta": {"dominio": "imagem", "fonte": "floco térmico"}, "origem": "wiki", "palavras_chave": [], "sinonimos": [], "tipo": "conceito", "titulo": "Os Modos do Floco são Partículas"}
\section{Os Modos do Floco são Partículas}
Os modos espectrais \{\ensuremath{\sigma}\_i\} do floco-estrela são as PARTÍCULAS da imagem: \ensuremath{\sigma}\_1 (a maior massa) é o NÚCLEO, a cauda são os elétrons/fótons/ressonâncias — exatamente como particulas\_minerais lê a raiz dominante como núcleo e a subdominante como elétron (por |raiz|). A física das partículas da máquina mineral, agora no espectro de uma imagem-estrela. Verif.: \ensuremath{\sigma}\_1=57638 (núcleo), cauda elétrons.
\prov{relações: parte\_de limite\_de\_fusao\_do\_floco; eh\_um particulas\_minerais; usa espectro\_e\_orbita\_cristal}
\prov{proveniência: wiki \textperiodcentered{} floco térmico \textperiodcentered{} fato \textperiodcentered{} confiança 1.0 \textperiodcentered{} domínio imagem \textperiodcentered{} história 8 versões no jornal}

%CRISTAL {"arestas": [{"alvo": "poligono_tende_ao_circulo", "confianca": 1.0, "epistemico": "fato", "origem": "wiki", "rel": "usa"}, {"alvo": "algebra_reta_mineral_7d", "confianca": 1.0, "epistemico": "fato", "origem": "wiki", "rel": "relaciona"}, {"alvo": "pi_gerador_da_reta_acima", "confianca": 1.0, "epistemico": "fato", "origem": "wiki", "rel": "usa"}, {"alvo": "corpo_mineral", "confianca": 1.0, "epistemico": "fato", "origem": "wiki", "rel": "usa"}], "confianca": 1.0, "contraexemplos": [], "descricao": "Os modos do poliedro VIBRAM com os da esfera da reta acima: o Laplaciano do icosaedro tem modos de degenerescência 1, 3, 5 — exatamente os 2l+1 dos harmônicos esféricos (l=0,1,2: s, p, d). O poliedro discreto já carrega os primeiros modos da esfera contínua, e a subdivisão geodésica que o leva à esfera é uma PROGRESSÃO GEOMÉTRICA (faces ×4) — a assinatura mineral. A vibração é geométrica em PG.", "epistemico": "inferencia", "exemplos": [], "id": "modos_ressoam_em_pg", "memoria": "semantica", "meta": {"dominio": "imagem", "fonte": "formas e limite"}, "origem": "wiki", "palavras_chave": [], "sinonimos": [], "tipo": "conceito", "titulo": "Os Modos Vibram (em PG)"}
\section{Os Modos Vibram (em PG)}
Os modos do poliedro VIBRAM com os da esfera da reta acima: o Laplaciano do icosaedro tem modos de degenerescência 1, 3, 5 — exatamente os 2l+1 dos harmônicos esféricos (l=0,1,2: s, p, d). O poliedro discreto já carrega os primeiros modos da esfera contínua, e a subdivisão geodésica que o leva à esfera é uma PROGRESSÃO GEOMÉTRICA (faces $\times$4) — a assinatura mineral. A vibração é geométrica em PG.
\prov{relações: usa poligono\_tende\_ao\_circulo; relaciona algebra\_reta\_mineral\_7d; usa pi\_gerador\_da\_reta\_acima; usa corpo\_mineral}
\prov{proveniência: wiki \textperiodcentered{} formas e limite \textperiodcentered{} inferencia \textperiodcentered{} confiança 1.0 \textperiodcentered{} domínio imagem \textperiodcentered{} história 9 versões no jornal}

%CRISTAL {"arestas": [{"alvo": "sistema_de_comunicacao", "confianca": 1.0, "epistemico": "fato", "origem": "wiki", "rel": "parte_de"}, {"alvo": "transformada_de_fourier", "confianca": 1.0, "epistemico": "fato", "origem": "wiki", "rel": "usa"}, {"alvo": "eng_telecomunicacoes", "confianca": 1.0, "epistemico": "fato", "origem": "curriculo", "rel": "parte_de"}, {"alvo": "onda_portadora", "confianca": 1.0, "epistemico": "fato", "origem": "wiki", "rel": "usa"}, {"alvo": "portadora_rotacao_salem", "confianca": 1.0, "epistemico": "fato", "origem": "wiki", "rel": "usa"}], "confianca": 1.0, "contraexemplos": [], "descricao": "Gravar bits na portadora variando um parâmetro dela: amplitude (AM), frequência (FM), fase (PM) ou amplitude+fase juntas (QAM). É como a mensagem entra na rotação — a perturbação controlada do Salem.", "epistemico": "fato", "exemplos": [], "id": "modulacao", "memoria": "semantica", "meta": {"dominio": "engenharia", "fonte": "conectividade e 5G"}, "origem": "wiki", "palavras_chave": [], "sinonimos": [], "tipo": "conceito", "titulo": "Modulação"}
\section{Modulação}
Gravar bits na portadora variando um parâmetro dela: amplitude (AM), frequência (FM), fase (PM) ou amplitude+fase juntas (QAM). É como a mensagem entra na rotação — a perturbação controlada do Salem.
\prov{relações: parte\_de sistema\_de\_comunicacao; usa transformada\_de\_fourier; parte\_de eng\_telecomunicacoes; usa onda\_portadora; usa portadora\_rotacao\_salem}
\prov{proveniência: wiki \textperiodcentered{} conectividade e 5G \textperiodcentered{} fato \textperiodcentered{} confiança 1.0 \textperiodcentered{} domínio engenharia \textperiodcentered{} história 7 versões no jornal}

%CRISTAL {"arestas": [{"alvo": "modulacao", "confianca": 1.0, "epistemico": "fato", "origem": "wiki", "rel": "especializa"}, {"alvo": "eng_telecomunicacoes", "confianca": 1.0, "epistemico": "fato", "origem": "curriculo", "rel": "parte_de"}], "confianca": 1.0, "contraexemplos": [], "descricao": "Variar a AMPLITUDE da portadora com o sinal — simples e barata de demodular, mas frágil ao ruído (que é amplitude). O rádio AM e as bases da modulação analógica.", "epistemico": "fato", "exemplos": [], "id": "modulacao_am", "memoria": "semantica", "meta": {"dominio": "engenharia", "fonte": ""}, "origem": "wiki", "palavras_chave": [], "sinonimos": [], "tipo": "conceito", "titulo": "Modulação em Amplitude (AM)"}
\section{Modulação em Amplitude (AM)}
Variar a AMPLITUDE da portadora com o sinal — simples e barata de demodular, mas frágil ao ruído (que é amplitude). O rádio AM e as bases da modulação analógica.
\prov{relações: especializa modulacao; parte\_de eng\_telecomunicacoes}
\prov{proveniência: wiki \textperiodcentered{} fato \textperiodcentered{} confiança 1.0 \textperiodcentered{} domínio engenharia \textperiodcentered{} história 3 versões no jornal}

%CRISTAL {"arestas": [{"alvo": "modulacao", "confianca": 1.0, "epistemico": "fato", "origem": "wiki", "rel": "especializa"}, {"alvo": "eng_telecomunicacoes", "confianca": 1.0, "epistemico": "fato", "origem": "curriculo", "rel": "parte_de"}], "confianca": 1.0, "contraexemplos": [], "descricao": "Variar a FREQUÊNCIA ou a FASE da portadora — imune a ruído de amplitude, ao custo de banda (regra de Carson). O rádio FM e o antecessor das modulações digitais de fase.", "epistemico": "fato", "exemplos": [], "id": "modulacao_angular", "memoria": "semantica", "meta": {"autor": "Armstrong", "dominio": "engenharia", "fonte": ""}, "origem": "wiki", "palavras_chave": [], "sinonimos": [], "tipo": "conceito", "titulo": "Modulação em Frequência/Fase (FM/PM)"}
\section{Modulação em Frequência/Fase (FM/PM)}
Variar a FREQUÊNCIA ou a FASE da portadora — imune a ruído de amplitude, ao custo de banda (regra de Carson). O rádio FM e o antecessor das modulações digitais de fase.
\prov{relações: especializa modulacao; parte\_de eng\_telecomunicacoes}
\prov{proveniência: wiki \textperiodcentered{} fato \textperiodcentered{} confiança 1.0 \textperiodcentered{} domínio engenharia \textperiodcentered{} história 3 versões no jornal}

%CRISTAL {"arestas": [{"alvo": "modulacao", "confianca": 1.0, "epistemico": "fato", "origem": "wiki", "rel": "especializa"}, {"alvo": "eng_telecomunicacoes", "confianca": 1.0, "epistemico": "fato", "origem": "curriculo", "rel": "parte_de"}], "confianca": 1.0, "contraexemplos": [], "descricao": "Mapear bits em símbolos da portadora: chaveamento de amplitude (ASK), de frequência (FSK) ou de fase (PSK). A fase (BPSK/QPSK) é a mais robusta por energia — o alfabeto do rádio digital.", "epistemico": "fato", "exemplos": [], "id": "modulacao_digital", "memoria": "semantica", "meta": {"dominio": "engenharia", "fonte": ""}, "origem": "wiki", "palavras_chave": [], "sinonimos": [], "tipo": "conceito", "titulo": "Modulação Digital (ASK/FSK/PSK)"}
\section{Modulação Digital (ASK/FSK/PSK)}
Mapear bits em símbolos da portadora: chaveamento de amplitude (ASK), de frequência (FSK) ou de fase (PSK). A fase (BPSK/QPSK) é a mais robusta por energia — o alfabeto do rádio digital.
\prov{relações: especializa modulacao; parte\_de eng\_telecomunicacoes}
\prov{proveniência: wiki \textperiodcentered{} fato \textperiodcentered{} confiança 1.0 \textperiodcentered{} domínio engenharia \textperiodcentered{} história 3 versões no jornal}

%CRISTAL {"arestas": [{"alvo": "dtc_schema_ut", "confianca": 1.0, "epistemico": "fato", "origem": "wiki", "rel": "usa"}, {"alvo": "dtc_renderer", "confianca": 1.0, "epistemico": "fato", "origem": "wiki", "rel": "usa"}], "confianca": 1.0, "contraexemplos": [], "descricao": "Duas metades homónimas: no servidor embutir_dtc traduz o payload Hodge em U_t; no browser ModuladorVetorial.ingest recebe cada pulso SSE, deteta se é vivo/mecânico/tick novo e chama render ou renderPartial. A ponte payload→estado→pintura.", "epistemico": "fato", "exemplos": [], "id": "modulador_vetorial", "memoria": "semantica", "meta": {"autor": "broca-so", "dominio": "cockpit", "fonte": "goldchain/modulador_vetorial.py; cockpit/modulador_vetorial.js"}, "origem": "wiki", "palavras_chave": [], "sinonimos": [], "tipo": "conceito", "titulo": "O Modulador Vetorial (pulso→U_t→render)"}
\section{O Modulador Vetorial (pulso\ensuremath{\to}U\_t\ensuremath{\to}render)}
Duas metades homónimas: no servidor embutir\_dtc traduz o payload Hodge em U\_t; no browser ModuladorVetorial.ingest recebe cada pulso SSE, deteta se é vivo/mecânico/tick novo e chama render ou renderPartial. A ponte payload\ensuremath{\to}estado\ensuremath{\to}pintura.
\prov{relações: usa dtc\_schema\_ut; usa dtc\_renderer}
\prov{proveniência: wiki \textperiodcentered{} goldchain/modulador\_vetorial.py; cockpit/modulador\_vetorial.js \textperiodcentered{} fato \textperiodcentered{} confiança 1.0 \textperiodcentered{} domínio cockpit \textperiodcentered{} história 3 versões no jornal}

%CRISTAL {"arestas": [{"alvo": "transistor", "confianca": 1.0, "epistemico": "fato", "origem": "wiki", "rel": "especializa"}], "confianca": 1.0, "contraexemplos": [], "descricao": "O transistor de efeito de campo dominante: uma tensão no gate abre/fecha um canal — bilhões deles, minúsculos, formam cada processador.", "epistemico": "fato", "exemplos": [], "id": "mosfet", "memoria": "semantica", "meta": {"dominio": "engenharia", "fonte": ""}, "origem": "wiki", "palavras_chave": [], "sinonimos": [], "tipo": "conceito", "titulo": "MOSFET"}
\section{MOSFET}
O transistor de efeito de campo dominante: uma tensão no gate abre/fecha um canal — bilhões deles, minúsculos, formam cada processador.
\prov{relações: especializa transistor}
\prov{proveniência: wiki \textperiodcentered{} fato \textperiodcentered{} confiança 1.0 \textperiodcentered{} domínio engenharia \textperiodcentered{} história 2 versões no jornal}

%CRISTAL {"arestas": [{"alvo": "maquina_eletrica", "confianca": 1.0, "epistemico": "fato", "origem": "wiki", "rel": "especializa"}, {"alvo": "eng_sistemas_energia", "confianca": 1.0, "epistemico": "fato", "origem": "curriculo", "rel": "parte_de"}], "confianca": 1.0, "contraexemplos": [], "descricao": "Converte eletricidade em movimento: de indução (assíncrono, robusto, o cavalo de força da indústria), síncrono e de corrente contínua. Consome mais da metade da energia elétrica do mundo.", "epistemico": "fato", "exemplos": [], "id": "motor_eletrico", "memoria": "semantica", "meta": {"dominio": "engenharia", "fonte": ""}, "origem": "wiki", "palavras_chave": [], "sinonimos": [], "tipo": "conceito", "titulo": "Motor Elétrico"}
\section{Motor Elétrico}
Converte eletricidade em movimento: de indução (assíncrono, robusto, o cavalo de força da indústria), síncrono e de corrente contínua. Consome mais da metade da energia elétrica do mundo.
\prov{relações: especializa maquina\_eletrica; parte\_de eng\_sistemas\_energia}
\prov{proveniência: wiki \textperiodcentered{} fato \textperiodcentered{} confiança 1.0 \textperiodcentered{} domínio engenharia \textperiodcentered{} história 3 versões no jornal}

%CRISTAL {"arestas": [{"alvo": "carro_eletrico", "confianca": 1.0, "epistemico": "fato", "origem": "wiki", "rel": "parte_de"}, {"alvo": "gerador_eletrico", "confianca": 1.0, "epistemico": "fato", "origem": "wiki", "rel": "relaciona"}, {"alvo": "conversao_eletromecanica", "confianca": 1.0, "epistemico": "fato", "origem": "wiki", "rel": "usa"}], "confianca": 1.0, "contraexemplos": [], "descricao": "Converte a eletricidade da bateria em rotação (torque) direto no eixo — normalmente um motor CA trifásico. Rendimento ~90-95%, torque máximo desde zero rpm; a rotação que move a roda.", "epistemico": "fato", "exemplos": [], "id": "motor_eletrico_tracao", "memoria": "semantica", "meta": {"dominio": "engenharia", "fonte": "mobilidade elétrica"}, "origem": "wiki", "palavras_chave": [], "sinonimos": [], "tipo": "conceito", "titulo": "Motor Elétrico de Tração"}
\section{Motor Elétrico de Tração}
Converte a eletricidade da bateria em rotação (torque) direto no eixo — normalmente um motor CA trifásico. Rendimento \~{}90-95\%, torque máximo desde zero rpm; a rotação que move a roda.
\prov{relações: parte\_de carro\_eletrico; relaciona gerador\_eletrico; usa conversao\_eletromecanica}
\prov{proveniência: wiki \textperiodcentered{} mobilidade elétrica \textperiodcentered{} fato \textperiodcentered{} confiança 1.0 \textperiodcentered{} domínio engenharia \textperiodcentered{} história 4 versões no jornal}

%CRISTAL {"arestas": [{"alvo": "orbita_temporal_das_camadas", "confianca": 1.0, "epistemico": "fato", "origem": "wiki", "rel": "parte_de"}, {"alvo": "numero_de_salem", "confianca": 1.0, "epistemico": "fato", "origem": "wiki", "rel": "eh_um"}, {"alvo": "corrente_alternada", "confianca": 1.0, "epistemico": "fato", "origem": "wiki", "rel": "relaciona"}, {"alvo": "portadora_rotacao_salem", "confianca": 1.0, "epistemico": "fato", "origem": "wiki", "rel": "relaciona"}, {"alvo": "corpo_mineral", "confianca": 1.0, "epistemico": "fato", "origem": "wiki", "rel": "usa"}], "confianca": 1.0, "contraexemplos": [], "descricao": "Um objeto orbitando num círculo decompõe-se em 2 MODOS DOMINANTES cujos coeficientes temporais são cos e sen: no plano (modo1, modo2) eles traçam um CÍRCULO (ρ=1 = Salem/borda) — a mesma rotação pura da corrente alternada e da portadora de rádio. O movimento periódico é uma órbita Salem no espaço de modos. Verif.: raio constante (círculo) a partir da POD. A classe da órbita temporal segue a parábola: constante=cristal (fundo), oscila=borda (periódico), cresce=caos.", "epistemico": "fato", "exemplos": [], "id": "movimento_circular_dois_modos", "memoria": "semantica", "meta": {"dominio": "imagem", "fonte": "selo e órbita temporal"}, "origem": "wiki", "palavras_chave": [], "sinonimos": [], "tipo": "conceito", "titulo": "Movimento Circular = Dois Modos num Círculo"}
\section{Movimento Circular = Dois Modos num Círculo}
Um objeto orbitando num círculo decompõe-se em 2 MODOS DOMINANTES cujos coeficientes temporais são cos e sen: no plano (modo1, modo2) eles traçam um CÍRCULO (\ensuremath{\rho}=1 = Salem/borda) — a mesma rotação pura da corrente alternada e da portadora de rádio. O movimento periódico é uma órbita Salem no espaço de modos. Verif.: raio constante (círculo) a partir da POD. A classe da órbita temporal segue a parábola: constante=cristal (fundo), oscila=borda (periódico), cresce=caos.
\prov{relações: parte\_de orbita\_temporal\_das\_camadas; eh\_um numero\_de\_salem; relaciona corrente\_alternada; relaciona portadora\_rotacao\_salem; usa corpo\_mineral}
\prov{proveniência: wiki \textperiodcentered{} selo e órbita temporal \textperiodcentered{} fato \textperiodcentered{} confiança 1.0 \textperiodcentered{} domínio imagem \textperiodcentered{} história 6 versões no jornal}

%CRISTAL {"arestas": [{"alvo": "protocolo_tcp", "confianca": 1.0, "epistemico": "fato", "origem": "wiki", "rel": "usa"}, {"alvo": "sistemas_embarcados", "confianca": 1.0, "epistemico": "fato", "origem": "wiki", "rel": "aplicacao"}, {"alvo": "eng_telecomunicacoes", "confianca": 1.0, "epistemico": "fato", "origem": "curriculo", "rel": "parte_de"}], "confianca": 1.0, "contraexemplos": [], "descricao": "Protocolo leve de publicação/assinatura sobre TCP para IoT: sensores publicam tópicos, um broker distribui. Mínimo overhead para dispositivos restritos e redes instáveis.", "epistemico": "fato", "exemplos": [], "id": "mqtt", "memoria": "semantica", "meta": {"dominio": "engenharia", "fonte": ""}, "origem": "wiki", "palavras_chave": [], "sinonimos": [], "tipo": "conceito", "titulo": "MQTT"}
\section{MQTT}
Protocolo leve de publicação/assinatura sobre TCP para IoT: sensores publicam tópicos, um broker distribui. Mínimo overhead para dispositivos restritos e redes instáveis.
\prov{relações: usa protocolo\_tcp; aplicacao sistemas\_embarcados; parte\_de eng\_telecomunicacoes}
\prov{proveniência: wiki \textperiodcentered{} fato \textperiodcentered{} confiança 1.0 \textperiodcentered{} domínio engenharia \textperiodcentered{} história 4 versões no jornal}

%CRISTAL {"arestas": [{"alvo": "sistema_de_comunicacao", "confianca": 1.0, "epistemico": "fato", "origem": "wiki", "rel": "usa"}, {"alvo": "eng_telecomunicacoes", "confianca": 1.0, "epistemico": "fato", "origem": "curriculo", "rel": "parte_de"}], "confianca": 1.0, "contraexemplos": [], "descricao": "Vários sinais dividindo um mesmo meio: por frequência (FDM), por tempo (TDM) ou por comprimento de onda na fibra (WDM). Como uma via física carrega muitas conversas simultâneas.", "epistemico": "fato", "exemplos": [], "id": "multiplexacao", "memoria": "semantica", "meta": {"dominio": "engenharia", "fonte": ""}, "origem": "wiki", "palavras_chave": [], "sinonimos": [], "tipo": "conceito", "titulo": "Multiplexação (FDM/TDM/WDM)"}
\section{Multiplexação (FDM/TDM/WDM)}
Vários sinais dividindo um mesmo meio: por frequência (FDM), por tempo (TDM) ou por comprimento de onda na fibra (WDM). Como uma via física carrega muitas conversas simultâneas.
\prov{relações: usa sistema\_de\_comunicacao; parte\_de eng\_telecomunicacoes}
\prov{proveniência: wiki \textperiodcentered{} fato \textperiodcentered{} confiança 1.0 \textperiodcentered{} domínio engenharia \textperiodcentered{} história 3 versões no jornal}

%CRISTAL {"arestas": [{"alvo": "mutacao", "confianca": 1.0, "epistemico": "fato", "origem": "wiki", "rel": "explica"}, {"alvo": "numero_de_salem", "confianca": 1.0, "epistemico": "fato", "origem": "wiki", "rel": "eh_um"}, {"alvo": "criticidade", "confianca": 1.0, "epistemico": "fato", "origem": "wiki", "rel": "relaciona"}, {"alvo": "corpo_mineral", "confianca": 1.0, "epistemico": "fato", "origem": "wiki", "rel": "usa"}, {"alvo": "teoria_da_evolucao", "confianca": 1.0, "epistemico": "fato", "origem": "wiki", "rel": "usa"}], "confianca": 1.0, "contraexemplos": [], "descricao": "O limiar de erro de Eigen: um genoma de comprimento L com vantagem σ só preserva a informação se (1−μ)^L > 1/σ. μ baixo = genoma CONGELADO (cristal, sem evolução), no limiar = a BORDA (onde a evolução acontece), μ alto = a informação DERRETE (caos, catástrofe de erro). É uma transição de fase — o ponto Salem da vida, o mesmo da criticidade nuclear. Verif.: μ=0.001 cristal, μ=0.02 caos.", "epistemico": "inferencia", "exemplos": [], "id": "mutacao_borda_eigen", "memoria": "semantica", "meta": {"classe_mu001": "cristal (info preservada)", "classe_mu02": "caos (catástrofe de erro, info derrete)", "dominio": "biotecnologia", "fonte": "biotecnologia e engenharia genética", "mu_critico_L100": 0.0069}, "origem": "wiki", "palavras_chave": [], "sinonimos": [], "tipo": "conceito", "titulo": "Taxa de Mutação = a Borda (catástrofe de erro)"}
\section{Taxa de Mutação = a Borda (catástrofe de erro)}
O limiar de erro de Eigen: um genoma de comprimento L com vantagem \ensuremath{\sigma} só preserva a informação se (1\ensuremath{-}\ensuremath{\mu})\^{}L > 1/\ensuremath{\sigma}. \ensuremath{\mu} baixo = genoma CONGELADO (cristal, sem evolução), no limiar = a BORDA (onde a evolução acontece), \ensuremath{\mu} alto = a informação DERRETE (caos, catástrofe de erro). É uma transição de fase — o ponto Salem da vida, o mesmo da criticidade nuclear. Verif.: \ensuremath{\mu}=0.001 cristal, \ensuremath{\mu}=0.02 caos.
\prov{relações: explica mutacao; eh\_um numero\_de\_salem; relaciona criticidade; usa corpo\_mineral; usa teoria\_da\_evolucao}
\prov{proveniência: wiki \textperiodcentered{} biotecnologia e engenharia genética \textperiodcentered{} inferencia \textperiodcentered{} confiança 1.0 \textperiodcentered{} domínio biotecnologia \textperiodcentered{} história 6 versões no jornal}

%CRISTAL {"arestas": [{"alvo": "protocolo_ip", "confianca": 1.0, "epistemico": "fato", "origem": "wiki", "rel": "usa"}, {"alvo": "eng_telecomunicacoes", "confianca": 1.0, "epistemico": "fato", "origem": "curriculo", "rel": "parte_de"}], "confianca": 1.0, "contraexemplos": [], "descricao": "Reescrever endereços/portas para muitos hosts privados dividirem um IP público — o remendo que estendeu o IPv4 e, de quebra, isolou a LAN. Quebra a conectividade fim-a-fim.", "epistemico": "inferencia", "exemplos": [], "id": "nat", "memoria": "semantica", "meta": {"dominio": "engenharia", "fonte": ""}, "origem": "wiki", "palavras_chave": [], "sinonimos": [], "tipo": "conceito", "titulo": "NAT (Tradução de Endereços)"}
\section{NAT (Tradução de Endereços)}
Reescrever endereços/portas para muitos hosts privados dividirem um IP público — o remendo que estendeu o IPv4 e, de quebra, isolou a LAN. Quebra a conectividade fim-a-fim.
\prov{relações: usa protocolo\_ip; parte\_de eng\_telecomunicacoes}
\prov{proveniência: wiki \textperiodcentered{} inferencia \textperiodcentered{} confiança 1.0 \textperiodcentered{} domínio engenharia \textperiodcentered{} história 3 versões no jornal}

%CRISTAL {"arestas": [{"alvo": "nucleo_5gc_sba", "confianca": 1.0, "epistemico": "fato", "origem": "wiki", "rel": "usa"}, {"alvo": "virtualizacao", "confianca": 1.0, "epistemico": "fato", "origem": "wiki", "rel": "usa"}, {"alvo": "qos_fluxos_5g", "confianca": 1.0, "epistemico": "fato", "origem": "wiki", "rel": "usa"}], "confianca": 1.0, "contraexemplos": [], "descricao": "Criar redes lógicas fim-a-fim isoladas sobre a mesma infra física — uma fatia URLLC para a fábrica, uma eMBB para o estádio, cada uma com sua QoS e seu SLA. A virtualização levada ao operador de telecom.", "epistemico": "inferencia", "exemplos": [], "id": "network_slicing_5g", "memoria": "semantica", "meta": {"dominio": "engenharia", "fonte": ""}, "origem": "wiki", "palavras_chave": [], "sinonimos": [], "tipo": "conceito", "titulo": "Fatiamento de Rede (Network Slicing)"}
\section{Fatiamento de Rede (Network Slicing)}
Criar redes lógicas fim-a-fim isoladas sobre a mesma infra física — uma fatia URLLC para a fábrica, uma eMBB para o estádio, cada uma com sua QoS e seu SLA. A virtualização levada ao operador de telecom.
\prov{relações: usa nucleo\_5gc\_sba; usa virtualizacao; usa qos\_fluxos\_5g}
\prov{proveniência: wiki \textperiodcentered{} inferencia \textperiodcentered{} confiança 1.0 \textperiodcentered{} domínio engenharia \textperiodcentered{} história 4 versões no jornal}

%CRISTAL {"arestas": [{"alvo": "julia", "confianca": 1.0, "epistemico": "fato", "origem": "wiki", "rel": "usa"}, {"alvo": "flip_duopolinomios", "confianca": 1.0, "epistemico": "fato", "origem": "wiki", "rel": "usa"}, {"alvo": "pi_medida_da_borda", "confianca": 1.0, "epistemico": "fato", "origem": "wiki", "rel": "usa"}, {"alvo": "poligono_tende_ao_circulo", "confianca": 1.0, "epistemico": "fato", "origem": "wiki", "rel": "relaciona"}], "confianca": 1.0, "contraexemplos": [], "descricao": "A fundamentação para achar o grupo discreto: o método de Newton em z^n−1 converge, de sementes espalhadas, para as n raízes da unidade — as bacias de atração (o fractal de Newton) SÃO o grupo cíclico C_n, na borda ρ=1. Achar o limite (o círculo) e o fractal de Newton (o grupo discreto que o aproxima) são os dois lados da mesma ponte entre as retas. Verif.: z^n−1 → n raízes, todas ρ=1.", "epistemico": "fato", "exemplos": [], "id": "newton_acha_o_grupo_discreto", "memoria": "semantica", "meta": {"dominio": "imagem", "fonte": "formas e limite"}, "origem": "wiki", "palavras_chave": [], "sinonimos": [], "tipo": "conceito", "titulo": "Newton Acha o Grupo Discreto"}
\section{Newton Acha o Grupo Discreto}
A fundamentação para achar o grupo discreto: o método de Newton em z\^{}n\ensuremath{-}1 converge, de sementes espalhadas, para as n raízes da unidade — as bacias de atração (o fractal de Newton) SÃO o grupo cíclico C\_n, na borda \ensuremath{\rho}=1. Achar o limite (o círculo) e o fractal de Newton (o grupo discreto que o aproxima) são os dois lados da mesma ponte entre as retas. Verif.: z\^{}n\ensuremath{-}1 \ensuremath{\to} n raízes, todas \ensuremath{\rho}=1.
\prov{relações: usa julia; usa flip\_duopolinomios; usa pi\_medida\_da\_borda; relaciona poligono\_tende\_ao\_circulo}
\prov{proveniência: wiki \textperiodcentered{} formas e limite \textperiodcentered{} fato \textperiodcentered{} confiança 1.0 \textperiodcentered{} domínio imagem \textperiodcentered{} história 10 versões no jornal}

%CRISTAL {"arestas": [{"alvo": "redes_5g_nr", "confianca": 1.0, "epistemico": "fato", "origem": "wiki", "rel": "parte_de"}], "confianca": 1.0, "contraexemplos": [], "descricao": "A rede de acesso do 5G: as estações-base gNB conectadas ao núcleo pela interface NG e entre si por Xn. O gNB termina toda a pilha de rádio e agenda cada recurso do ar em tempo real.", "epistemico": "fato", "exemplos": [], "id": "ng_ran_gnb", "memoria": "semantica", "meta": {"dominio": "engenharia", "fonte": ""}, "origem": "wiki", "palavras_chave": [], "sinonimos": [], "tipo": "conceito", "titulo": "NG-RAN e gNB"}
\section{NG-RAN e gNB}
A rede de acesso do 5G: as estações-base gNB conectadas ao núcleo pela interface NG e entre si por Xn. O gNB termina toda a pilha de rádio e agenda cada recurso do ar em tempo real.
\prov{relações: parte\_de redes\_5g\_nr}
\prov{proveniência: wiki \textperiodcentered{} fato \textperiodcentered{} confiança 1.0 \textperiodcentered{} domínio engenharia \textperiodcentered{} história 2 versões no jornal}

%CRISTAL {"arestas": [{"alvo": "redes_5g_nr", "confianca": 1.0, "epistemico": "fato", "origem": "wiki", "rel": "parte_de"}, {"alvo": "http2_http3", "confianca": 1.0, "epistemico": "fato", "origem": "wiki", "rel": "usa"}, {"alvo": "microservicos", "confianca": 1.0, "epistemico": "fato", "origem": "wiki", "rel": "analogia"}, {"alvo": "rest", "confianca": 1.0, "epistemico": "fato", "origem": "wiki", "rel": "usa"}], "confianca": 1.0, "contraexemplos": [], "descricao": "O núcleo 5G abandona as caixas monolíticas por FUNÇÕES DE REDE que se descobrem e se chamam como microsserviços — a Service Based Architecture, com interfaces REST sobre HTTP/2 e JSON. O core de telecom virou nuvem de APIs.", "epistemico": "fato", "exemplos": [], "id": "nucleo_5gc_sba", "memoria": "semantica", "meta": {"dominio": "engenharia", "fonte": ""}, "origem": "wiki", "palavras_chave": [], "sinonimos": [], "tipo": "conceito", "titulo": "Núcleo 5G (5GC) e Arquitetura de Serviços (SBA)"}
\section{Núcleo 5G (5GC) e Arquitetura de Serviços (SBA)}
O núcleo 5G abandona as caixas monolíticas por FUNÇÕES DE REDE que se descobrem e se chamam como microsserviços — a Service Based Architecture, com interfaces REST sobre HTTP/2 e JSON. O core de telecom virou nuvem de APIs.
\prov{relações: parte\_de redes\_5g\_nr; usa http2\_http3; analogia microservicos; usa rest}
\prov{proveniência: wiki \textperiodcentered{} fato \textperiodcentered{} confiança 1.0 \textperiodcentered{} domínio engenharia \textperiodcentered{} história 5 versões no jornal}

%CRISTAL {"arestas": [{"alvo": "painel_estelar_obs", "confianca": 1.0, "epistemico": "fato", "origem": "wiki", "rel": "parte_de"}, {"alvo": "invariante_estelar", "confianca": 1.0, "epistemico": "fato", "origem": "wiki", "rel": "usa"}, {"alvo": "cotacao_parabola", "confianca": 1.0, "epistemico": "fato", "origem": "wiki", "rel": "usa"}], "confianca": 1.0, "contraexemplos": [], "descricao": "A zona central que registra Q (invariante quadrático de Lopes), a balança (a,c²) e os pesos (P_φ,P_ψ)≈(e/t, 1−e/t). Nos fechos de referência: Q=1.0, a=0.00623249, c²=0.99376751, P_φ=0.48474576. Lê-se de pow_metal_avaliar+gentil, não se recalcula.", "epistemico": "fato", "exemplos": [], "id": "nucleo_estelar", "memoria": "semantica", "meta": {"autor": "broca-so", "dominio": "cockpit", "fonte": "ula/painel_estelar.c; ula/dados/painel_estelar.json"}, "origem": "wiki", "palavras_chave": [], "sinonimos": [], "tipo": "conceito", "titulo": "O Núcleo Estelar (o invariante que não se recalcula)"}
\section{O Núcleo Estelar (o invariante que não se recalcula)}
A zona central que registra Q (invariante quadrático de Lopes), a balança (a,c\ensuremath{^{2}}) e os pesos (P\_\ensuremath{\varphi},P\_\ensuremath{\psi})\ensuremath{\approx}(e/t, 1\ensuremath{-}e/t). Nos fechos de referência: Q=1.0, a=0.00623249, c\ensuremath{^{2}}=0.99376751, P\_\ensuremath{\varphi}=0.48474576. Lê-se de pow\_metal\_avaliar+gentil, não se recalcula.
\prov{relações: parte\_de painel\_estelar\_obs; usa invariante\_estelar; usa cotacao\_parabola}
\prov{proveniência: wiki \textperiodcentered{} ula/painel\_estelar.c; ula/dados/painel\_estelar.json \textperiodcentered{} fato \textperiodcentered{} confiança 1.0 \textperiodcentered{} domínio cockpit \textperiodcentered{} história 5 versões no jornal}

%CRISTAL {"arestas": [{"alvo": "limite_de_fusao_do_floco", "confianca": 1.0, "epistemico": "fato", "origem": "wiki", "rel": "usa"}, {"alvo": "fusao_estelar", "confianca": 1.0, "epistemico": "fato", "origem": "wiki", "rel": "eh_um"}, {"alvo": "vale_do_ferro", "confianca": 1.0, "epistemico": "fato", "origem": "wiki", "rel": "relaciona"}, {"alvo": "espectro_irredutivel_floco", "confianca": 1.0, "epistemico": "fato", "origem": "wiki", "rel": "usa"}, {"alvo": "particulas_minerais", "confianca": 1.0, "epistemico": "fato", "origem": "wiki", "rel": "relaciona"}, {"alvo": "corpo_estelar", "confianca": 1.0, "epistemico": "fato", "origem": "wiki", "rel": "usa"}, {"alvo": "fusao_nuclear", "confianca": 1.0, "epistemico": "fato", "origem": "wiki", "rel": "relaciona"}, {"alvo": "nucleossintese_bbn", "confianca": 1.0, "epistemico": "fato", "origem": "wiki", "rel": "relaciona"}], "confianca": 1.0, "contraexemplos": [], "descricao": "Fundir dois flocos é SUPERPOR: C=norma(A+B), e refazer a SVD. A energia crua Σpx² SOMA (dois núcleos num só), e o espectro SE ESPALHA — o k_significativo do fundido ≥ o dos parciais: um núcleo mais pesado. A ENERGIA DE LIGAÇÃO POR NÚCLEON (σ_1²/Σσ², intensiva) sobe enquanto os flocos se alinham (rumo ao ferro) e cai quando são diversos (passou do ferro). Fundir a cadeia A→A+B→A+B+C é a nucleossíntese: constrói os elementos e acha o pico — o vale do ferro da imagem.", "epistemico": "fato", "exemplos": [], "id": "nucleossintese_dos_flocos", "memoria": "semantica", "meta": {"dominio": "imagem", "fonte": "la hire / nucleossíntese"}, "origem": "wiki", "palavras_chave": [], "sinonimos": [], "tipo": "conceito", "titulo": "A Nucleossíntese dos Flocos"}
\section{A Nucleossíntese dos Flocos}
Fundir dois flocos é SUPERPOR: C=norma(A+B), e refazer a SVD. A energia crua \ensuremath{\Sigma}px\ensuremath{^{2}} SOMA (dois núcleos num só), e o espectro SE ESPALHA — o k\_significativo do fundido \ensuremath{\ge} o dos parciais: um núcleo mais pesado. A ENERGIA DE LIGAÇÃO POR NÚCLEON (\ensuremath{\sigma}\_1\ensuremath{^{2}}/\ensuremath{\Sigma}\ensuremath{\sigma}\ensuremath{^{2}}, intensiva) sobe enquanto os flocos se alinham (rumo ao ferro) e cai quando são diversos (passou do ferro). Fundir a cadeia A\ensuremath{\to}A+B\ensuremath{\to}A+B+C é a nucleossíntese: constrói os elementos e acha o pico — o vale do ferro da imagem.
\prov{relações: usa limite\_de\_fusao\_do\_floco; eh\_um fusao\_estelar; relaciona vale\_do\_ferro; usa espectro\_irredutivel\_floco; relaciona particulas\_minerais; usa corpo\_estelar; relaciona fusao\_nuclear; relaciona nucleossintese\_bbn}
\prov{proveniência: wiki \textperiodcentered{} la hire / nucleossíntese \textperiodcentered{} fato \textperiodcentered{} confiança 1.0 \textperiodcentered{} domínio imagem \textperiodcentered{} história 11 versões no jornal}

%CRISTAL {"arestas": [{"alvo": "camada_fisica_nr", "confianca": 1.0, "epistemico": "fato", "origem": "wiki", "rel": "parte_de"}, {"alvo": "ofdm", "confianca": 1.0, "epistemico": "fato", "origem": "wiki", "rel": "usa"}], "confianca": 1.0, "contraexemplos": [], "descricao": "O NR escolhe o espaçamento entre subportadoras (15/30/60/120 kHz) conforme a banda e o uso: espaçamento largo encurta o símbolo (baixa latência, mmWave), estreito favorece cobertura. Um OFDM que muda de escala.", "epistemico": "fato", "exemplos": [], "id": "numerologia_nr", "memoria": "semantica", "meta": {"dominio": "engenharia", "fonte": ""}, "origem": "wiki", "palavras_chave": [], "sinonimos": [], "tipo": "conceito", "titulo": "Numerologia Flexível (Espaçamento de Subportadora)"}
\section{Numerologia Flexível (Espaçamento de Subportadora)}
O NR escolhe o espaçamento entre subportadoras (15/30/60/120 kHz) conforme a banda e o uso: espaçamento largo encurta o símbolo (baixa latência, mmWave), estreito favorece cobertura. Um OFDM que muda de escala.
\prov{relações: parte\_de camada\_fisica\_nr; usa ofdm}
\prov{proveniência: wiki \textperiodcentered{} fato \textperiodcentered{} confiança 1.0 \textperiodcentered{} domínio engenharia \textperiodcentered{} história 3 versões no jornal}

%CRISTAL {"arestas": [{"alvo": "porta_hodge_correlacao", "confianca": 1.0, "epistemico": "fato", "origem": "wiki", "rel": "usa"}, {"alvo": "hamiltoniano_quarentena", "confianca": 1.0, "epistemico": "fato", "origem": "wiki", "rel": "usa"}, {"alvo": "olho_invertido_koch", "confianca": 1.0, "epistemico": "fato", "origem": "wiki", "rel": "usa"}], "confianca": 1.0, "contraexemplos": [], "descricao": "Produz 'um instante Hodge': chama analisar_para_cockpit, acopla a predição Koopman com o olho invertido, calcula a assinatura e reporta porta_aberta/cúspides + os escalares (lastro, delta, fechou, parábola_ok, sync_ok). É o que o Hodge automatiza: a leitura do fluxo num único snapshot correlacionado.", "epistemico": "fato", "exemplos": [], "id": "observar_correlacao_hodge", "memoria": "semantica", "meta": {"autor": "broca-so", "dominio": "cockpit", "fonte": "goldchain/hodge_automatizar.py:observar_correlacao"}, "origem": "wiki", "palavras_chave": [], "sinonimos": [], "tipo": "conceito", "titulo": "O Instante Hodge (observar_correlacao)"}
\section{O Instante Hodge (observar\_correlacao)}
Produz 'um instante Hodge': chama analisar\_para\_cockpit, acopla a predição Koopman com o olho invertido, calcula a assinatura e reporta porta\_aberta/cúspides + os escalares (lastro, delta, fechou, parábola\_ok, sync\_ok). É o que o Hodge automatiza: a leitura do fluxo num único snapshot correlacionado.
\prov{relações: usa porta\_hodge\_correlacao; usa hamiltoniano\_quarentena; usa olho\_invertido\_koch}
\prov{proveniência: wiki \textperiodcentered{} goldchain/hodge\_automatizar.py:observar\_correlacao \textperiodcentered{} fato \textperiodcentered{} confiança 1.0 \textperiodcentered{} domínio cockpit \textperiodcentered{} história 4 versões no jornal}

%CRISTAL {"arestas": [{"alvo": "painel_estelar_obs", "confianca": 1.0, "epistemico": "fato", "origem": "wiki", "rel": "explica"}, {"alvo": "olho_de_koch", "confianca": 1.0, "epistemico": "fato", "origem": "wiki", "rel": "relaciona"}, {"alvo": "cockpit_observatorio", "confianca": 1.0, "epistemico": "fato", "origem": "wiki", "rel": "explica"}], "confianca": 1.0, "contraexemplos": [], "descricao": "O painel/cockpit NÃO otimiza mineração, não altera SHA-256, não escolhe nonces — OBSERVA a forma que o gato já impõe. Mandar o header completo preserva o associador; picotar destruiria a leitura. 'Um habitante, três braços, um invariante — nada a acrescentar.' A mesma vocação do olho de Koch que contempla e não minera.", "epistemico": "inferencia", "exemplos": [], "id": "observatorio_nao_fabrica", "memoria": "semantica", "meta": {"autor": "broca-so", "dominio": "cockpit", "fonte": "papers/painel_estelar.tex §V"}, "origem": "wiki", "palavras_chave": [], "sinonimos": [], "tipo": "conceito", "titulo": "Observatório, não Fábrica"}
\section{Observatório, não Fábrica}
O painel/cockpit NÃO otimiza mineração, não altera SHA-256, não escolhe nonces — OBSERVA a forma que o gato já impõe. Mandar o header completo preserva o associador; picotar destruiria a leitura. 'Um habitante, três braços, um invariante — nada a acrescentar.' A mesma vocação do olho de Koch que contempla e não minera.
\prov{relações: explica painel\_estelar\_obs; relaciona olho\_de\_koch; explica cockpit\_observatorio}
\prov{proveniência: wiki \textperiodcentered{} papers/painel\_estelar.tex §V \textperiodcentered{} inferencia \textperiodcentered{} confiança 1.0 \textperiodcentered{} domínio cockpit \textperiodcentered{} história 4 versões no jornal}

%CRISTAL {"arestas": [{"alvo": "qam", "confianca": 1.0, "epistemico": "fato", "origem": "wiki", "rel": "usa"}, {"alvo": "transformada_de_fourier", "confianca": 1.0, "epistemico": "fato", "origem": "wiki", "rel": "usa"}, {"alvo": "multiplexacao", "confianca": 1.0, "epistemico": "fato", "origem": "wiki", "rel": "especializa"}, {"alvo": "eng_telecomunicacoes", "confianca": 1.0, "epistemico": "fato", "origem": "curriculo", "rel": "parte_de"}], "confianca": 1.0, "contraexemplos": [], "descricao": "Dividir a banda em milhares de subportadoras ortogonais, cada uma com QAM lento — trocando um canal difícil de banda larga por muitos canais fáceis. A técnica do Wi-Fi, LTE, 5G e da TV digital.", "epistemico": "fato", "exemplos": [], "id": "ofdm", "memoria": "semantica", "meta": {"dominio": "engenharia", "fonte": ""}, "origem": "wiki", "palavras_chave": [], "sinonimos": [], "tipo": "conceito", "titulo": "OFDM (Multiplexação por Divisão de Frequências Ortogonais)"}
\section{OFDM (Multiplexação por Divisão de Frequências Ortogonais)}
Dividir a banda em milhares de subportadoras ortogonais, cada uma com QAM lento — trocando um canal difícil de banda larga por muitos canais fáceis. A técnica do Wi-Fi, LTE, 5G e da TV digital.
\prov{relações: usa qam; usa transformada\_de\_fourier; especializa multiplexacao; parte\_de eng\_telecomunicacoes}
\prov{proveniência: wiki \textperiodcentered{} fato \textperiodcentered{} confiança 1.0 \textperiodcentered{} domínio engenharia \textperiodcentered{} história 5 versões no jornal}

%CRISTAL {"arestas": [{"alvo": "cockpit_observatorio", "confianca": 1.0, "epistemico": "fato", "origem": "wiki", "rel": "parte_de"}, {"alvo": "cockpit_olho", "confianca": 1.0, "epistemico": "fato", "origem": "wiki", "rel": "explica"}, {"alvo": "goldchain", "confianca": 1.0, "epistemico": "fato", "origem": "wiki", "rel": "relaciona"}], "confianca": 1.0, "contraexemplos": [], "descricao": "A interface no navegador (<title>Olho de Cock · broca-so): observatório de um só gerador U_t, com seções cofre, engenharia de ouro, fluxo Navier-Stokes/voltas, gerador dinâmico, olho na borda, contabilidade em tempo real. Aprofunda o cockpit_olho (que estava raso). [WIP — precisa avançar.]", "epistemico": "inferencia", "exemplos": [], "id": "olho_de_cock_web", "memoria": "semantica", "meta": {"autor": "broca-so", "dominio": "cockpit", "fonte": "goldchain/cockpit/index.html", "wip": true}, "origem": "wiki", "palavras_chave": [], "sinonimos": [], "tipo": "conceito", "titulo": "O Olho de Cock (o cockpit web)"}
\section{O Olho de Cock (o cockpit web)}
A interface no navegador (<title>Olho de Cock \textperiodcentered  broca-so): observatório de um só gerador U\_t, com seções cofre, engenharia de ouro, fluxo Navier-Stokes/voltas, gerador dinâmico, olho na borda, contabilidade em tempo real. Aprofunda o cockpit\_olho (que estava raso). [WIP — precisa avançar.]
\prov{relações: parte\_de cockpit\_observatorio; explica cockpit\_olho; relaciona goldchain}
\prov{proveniência: wiki \textperiodcentered{} goldchain/cockpit/index.html \textperiodcentered{} inferencia \textperiodcentered{} confiança 1.0 \textperiodcentered{} domínio cockpit \textperiodcentered{} história 4 versões no jornal}

%CRISTAL {"arestas": [{"alvo": "olho_de_koch", "confianca": 1.0, "epistemico": "fato", "origem": "wiki", "rel": "explica"}, {"alvo": "colapso_koch_tiffany", "confianca": 1.0, "epistemico": "fato", "origem": "wiki", "rel": "usa"}, {"alvo": "broca_cockpit_engine", "confianca": 1.0, "epistemico": "fato", "origem": "wiki", "rel": "usa"}], "confianca": 1.0, "contraexemplos": [], "descricao": "Traz o mundo externo (mempool/B_out) de volta ao painel (W_in) atravessando o ribossomo Koch z²mod1009 — transformada 2-para-1 SEM inversa (o 'invertido' é o SENTIDO do fluxo mundo→painel, não uma inversa matemática). quantizar_mundo mistura fee/count/pendentes num atrator via XOR de três koch_z2; a janela abre por percepção do mundo.", "epistemico": "fato", "exemplos": [], "id": "olho_invertido_koch", "memoria": "semantica", "meta": {"autor": "broca-so", "dominio": "cockpit", "fonte": "goldchain/olho_invertido.py:koch_z2/quantizar_mundo/perceber_mundo"}, "origem": "wiki", "palavras_chave": [], "sinonimos": [], "tipo": "conceito", "titulo": "O Olho Invertido — B_out → z↦z² → W_in"}
\section{O Olho Invertido — B\_out \ensuremath{\to} z\ensuremath{\mapsto}z\ensuremath{^{2}} \ensuremath{\to} W\_in}
Traz o mundo externo (mempool/B\_out) de volta ao painel (W\_in) atravessando o ribossomo Koch z\ensuremath{^{2}}mod1009 — transformada 2-para-1 SEM inversa (o 'invertido' é o SENTIDO do fluxo mundo\ensuremath{\to}painel, não uma inversa matemática). quantizar\_mundo mistura fee/count/pendentes num atrator via XOR de três koch\_z2; a janela abre por percepção do mundo.
\prov{relações: explica olho\_de\_koch; usa colapso\_koch\_tiffany; usa broca\_cockpit\_engine}
\prov{proveniência: wiki \textperiodcentered{} goldchain/olho\_invertido.py:koch\_z2/quantizar\_mundo/perceber\_mundo \textperiodcentered{} fato \textperiodcentered{} confiança 1.0 \textperiodcentered{} domínio cockpit \textperiodcentered{} história 4 versões no jornal}

%CRISTAL {"arestas": [{"alvo": "conectividade", "confianca": 1.0, "epistemico": "fato", "origem": "wiki", "rel": "usa"}], "confianca": 1.0, "contraexemplos": [], "descricao": "A oscilação de rádio de alta frequência que CARREGA a informação. Sozinha não diz nada; modulá-la (mudar amplitude, fase ou frequência) grava os bits nela. O veículo da mensagem pelo ar.", "epistemico": "fato", "exemplos": [], "id": "onda_portadora", "memoria": "semantica", "meta": {"dominio": "engenharia", "fonte": "conectividade e 5G"}, "origem": "wiki", "palavras_chave": [], "sinonimos": [], "tipo": "conceito", "titulo": "Onda Portadora"}
\section{Onda Portadora}
A oscilação de rádio de alta frequência que CARREGA a informação. Sozinha não diz nada; modulá-la (mudar amplitude, fase ou frequência) grava os bits nela. O veículo da mensagem pelo ar.
\prov{relações: usa conectividade}
\prov{proveniência: wiki \textperiodcentered{} conectividade e 5G \textperiodcentered{} fato \textperiodcentered{} confiança 1.0 \textperiodcentered{} domínio engenharia \textperiodcentered{} história 2 versões no jornal}

%CRISTAL {"arestas": [{"alvo": "eletromagnetismo_eng", "confianca": 1.0, "epistemico": "fato", "origem": "wiki", "rel": "consequencia"}, {"alvo": "eng_eletronica", "confianca": 1.0, "epistemico": "fato", "origem": "curriculo", "rel": "parte_de"}, {"alvo": "eng_telecomunicacoes", "confianca": 1.0, "epistemico": "fato", "origem": "curriculo", "rel": "parte_de"}, {"alvo": "eng_sistemas_energia", "confianca": 1.0, "epistemico": "fato", "origem": "curriculo", "rel": "parte_de"}], "confianca": 1.0, "contraexemplos": [], "descricao": "A solução propagante de Maxwell: campo E e B que se sustentam viajando à velocidade da luz. Polarização, frequência, comprimento de onda — o veículo físico de toda radiocomunicação.", "epistemico": "fato", "exemplos": [], "id": "ondas_eletromagneticas", "memoria": "semantica", "meta": {"dominio": "engenharia", "fonte": ""}, "origem": "wiki", "palavras_chave": [], "sinonimos": [], "tipo": "conceito", "titulo": "Ondas Eletromagnéticas"}
\section{Ondas Eletromagnéticas}
A solução propagante de Maxwell: campo E e B que se sustentam viajando à velocidade da luz. Polarização, frequência, comprimento de onda — o veículo físico de toda radiocomunicação.
\prov{relações: consequencia eletromagnetismo\_eng; parte\_de eng\_eletronica; parte\_de eng\_telecomunicacoes; parte\_de eng\_sistemas\_energia}
\prov{proveniência: wiki \textperiodcentered{} fato \textperiodcentered{} confiança 1.0 \textperiodcentered{} domínio engenharia \textperiodcentered{} história 5 versões no jornal}

%CRISTAL {"arestas": [{"alvo": "led", "confianca": 1.0, "epistemico": "fato", "origem": "wiki", "rel": "generaliza"}, {"alvo": "semicondutor", "confianca": 1.0, "epistemico": "fato", "origem": "wiki", "rel": "usa"}, {"alvo": "eng_eletronica", "confianca": 1.0, "epistemico": "fato", "origem": "curriculo", "rel": "parte_de"}], "confianca": 1.0, "contraexemplos": [], "descricao": "Dispositivos que convertem elétron↔fóton: LED, laser semicondutor, fotodiodo. A base do display, do sensor de luz e — crucialmente — da comunicação por fibra óptica.", "epistemico": "fato", "exemplos": [], "id": "optoeletronica", "memoria": "semantica", "meta": {"dominio": "engenharia", "fonte": ""}, "origem": "wiki", "palavras_chave": [], "sinonimos": [], "tipo": "conceito", "titulo": "Optoeletrônica"}
\section{Optoeletrônica}
Dispositivos que convertem elétron\ensuremath{\leftrightarrow}fóton: LED, laser semicondutor, fotodiodo. A base do display, do sensor de luz e — crucialmente — da comunicação por fibra óptica.
\prov{relações: generaliza led; usa semicondutor; parte\_de eng\_eletronica}
\prov{proveniência: wiki \textperiodcentered{} fato \textperiodcentered{} confiança 1.0 \textperiodcentered{} domínio engenharia \textperiodcentered{} história 4 versões no jornal}

%CRISTAL {"arestas": [{"alvo": "classificacao_de_formas", "confianca": 1.0, "epistemico": "fato", "origem": "wiki", "rel": "especializa"}, {"alvo": "segmentacao_de_imagem", "confianca": 1.0, "epistemico": "fato", "origem": "wiki", "rel": "usa"}, {"alvo": "imagem_como_orbita_gato", "confianca": 1.0, "epistemico": "fato", "origem": "wiki", "rel": "relaciona"}, {"alvo": "corpo_mineral", "confianca": 1.0, "epistemico": "fato", "origem": "wiki", "rel": "usa"}], "confianca": 1.0, "contraexemplos": [], "descricao": "A órbita de um objeto é o seu CONTORNO como curva fechada. Segmenta-se o objeto da imagem (o colapso da segmentação), toma-se a maior componente e lê-se a silhueta em coordenadas polares: uma órbita periódica. Depois das formas sintéticas, os objetos reais entram na mesma máquina — o desenho e o mundo deixam a mesma assinatura de rotações fechadas.", "epistemico": "fato", "exemplos": [], "id": "orbita_de_objeto_real", "memoria": "semantica", "meta": {"dominio": "imagem", "fonte": "órbitas de objetos"}, "origem": "wiki", "palavras_chave": [], "sinonimos": [], "tipo": "conceito", "titulo": "A Órbita de um Objeto Real"}
\section{A Órbita de um Objeto Real}
A órbita de um objeto é o seu CONTORNO como curva fechada. Segmenta-se o objeto da imagem (o colapso da segmentação), toma-se a maior componente e lê-se a silhueta em coordenadas polares: uma órbita periódica. Depois das formas sintéticas, os objetos reais entram na mesma máquina — o desenho e o mundo deixam a mesma assinatura de rotações fechadas.
\prov{relações: especializa classificacao\_de\_formas; usa segmentacao\_de\_imagem; relaciona imagem\_como\_orbita\_gato; usa corpo\_mineral}
\prov{proveniência: wiki \textperiodcentered{} órbitas de objetos \textperiodcentered{} fato \textperiodcentered{} confiança 1.0 \textperiodcentered{} domínio imagem \textperiodcentered{} história 5 versões no jornal}

%CRISTAL {"arestas": [{"alvo": "xadrez", "confianca": 1.0, "epistemico": "fato", "origem": "wiki", "rel": "relaciona"}, {"alvo": "xadrez_bai", "confianca": 1.0, "epistemico": "fato", "origem": "wiki", "rel": "relaciona"}, {"alvo": "laco_traducao_bai_peirce", "confianca": 1.0, "epistemico": "fato", "origem": "wiki", "rel": "usa"}, {"alvo": "minimax_e_colapso", "confianca": 1.0, "epistemico": "fato", "origem": "wiki", "rel": "relaciona"}, {"alvo": "valor_do_jogo_von_neumann", "confianca": 1.0, "epistemico": "fato", "origem": "wiki", "rel": "relaciona"}], "confianca": 1.0, "contraexemplos": [], "descricao": "Uma partida de xadrez É uma ÓRBITA: a trajetória do tabuleiro sob o operador-lance, do inicial ao terminal (mate ou empate). É a mesma órbita do BAI (orbitar = colapsar = percorrer a árvore de decisão), agora percorrida de fato por dois jogadores. Cada posição é um ponto; a sequência de posições é a órbita; o mate é o ponto fixo.", "epistemico": "inferencia", "exemplos": [], "id": "orbita_de_partida", "memoria": "semantica", "meta": {"dominio": "ponte", "fonte": "partidas↔órbitas"}, "origem": "wiki", "palavras_chave": [], "sinonimos": [], "tipo": "conceito", "titulo": "Órbita de Partida"}
\section{Órbita de Partida}
Uma partida de xadrez É uma ÓRBITA: a trajetória do tabuleiro sob o operador-lance, do inicial ao terminal (mate ou empate). É a mesma órbita do BAI (orbitar = colapsar = percorrer a árvore de decisão), agora percorrida de fato por dois jogadores. Cada posição é um ponto; a sequência de posições é a órbita; o mate é o ponto fixo.
\prov{relações: relaciona xadrez; relaciona xadrez\_bai; usa laco\_traducao\_bai\_peirce; relaciona minimax\_e\_colapso; relaciona valor\_do\_jogo\_von\_neumann}
\prov{proveniência: wiki \textperiodcentered{} partidas\ensuremath{\leftrightarrow}órbitas \textperiodcentered{} inferencia \textperiodcentered{} confiança 1.0 \textperiodcentered{} domínio ponte \textperiodcentered{} história 54 versões no jornal}

%CRISTAL {"arestas": [{"alvo": "broca_rede_social_soberana", "confianca": 1.0, "epistemico": "fato", "origem": "wiki", "rel": "relaciona"}, {"alvo": "goldchain_como_rede", "confianca": 1.0, "epistemico": "fato", "origem": "wiki", "rel": "relaciona"}, {"alvo": "meta_inducao_geral_simbolo", "confianca": 1.0, "epistemico": "fato", "origem": "wiki", "rel": "usa"}, {"alvo": "rede_limpa_bump", "confianca": 1.0, "epistemico": "fato", "origem": "wiki", "rel": "usa"}, {"alvo": "banda_exclusiva_por_tecido", "confianca": 1.0, "epistemico": "fato", "origem": "wiki", "rel": "usa"}, {"alvo": "traducao_como_transformacao", "confianca": 1.0, "epistemico": "fato", "origem": "wiki", "rel": "usa"}, {"alvo": "lei_unica_tres_dominios", "confianca": 1.0, "epistemico": "fato", "origem": "wiki", "rel": "relaciona"}], "confianca": 1.0, "contraexemplos": [], "descricao": "A rede soberana (social/pertencimento) e o goldchain (valor/consenso) são a MESMA órbita — o substrato de pares numa banda trocando o bump com reciprocidade, sem dono. Traduzem 1:1 por um pivô ('rede'): membro⇄nó, comunidade⇄rede, post⇄transação, feed⇄broadcast, moderador⇄minerador. A diferença é SÓ interpretação, e essa diferença É o BUMP: goldchain adiciona consenso+valor+ledger; a soberana adiciona artefato-vivo+evolução+sentido. Meta-indução ('só o símbolo muda') + rede_limpa_bump ('só a diferença viaja') no mesmo gesto. Não são dois nós rivais — é um, com dois polos.", "epistemico": "inferencia", "exemplos": [], "id": "orbita_rede_dois_atratores", "memoria": "semantica", "meta": {"dominio": "arquitetura", "fonte": "semilla_orbita_dois_atratores"}, "origem": "wiki", "palavras_chave": ["rede soberana e goldchain", "a mesma rede duas leituras", "diferença entre rede social e goldchain", "por que rede soberana ressoa com goldchain", "uma órbita dois atratores", "o bump interpretativo"], "sinonimos": [], "tipo": "conceito", "titulo": "Uma Órbita, Dois Atratores: rede soberana e goldchain"}
\section{Uma Órbita, Dois Atratores: rede soberana e goldchain}
A rede soberana (social/pertencimento) e o goldchain (valor/consenso) são a MESMA órbita — o substrato de pares numa banda trocando o bump com reciprocidade, sem dono. Traduzem 1:1 por um pivô ('rede'): membro\ensuremath{\rightleftarrows}nó, comunidade\ensuremath{\rightleftarrows}rede, post\ensuremath{\rightleftarrows}transação, feed\ensuremath{\rightleftarrows}broadcast, moderador\ensuremath{\rightleftarrows}minerador. A diferença é SÓ interpretação, e essa diferença É o BUMP: goldchain adiciona consenso+valor+ledger; a soberana adiciona artefato-vivo+evolução+sentido. Meta-indução ('só o símbolo muda') + rede\_limpa\_bump ('só a diferença viaja') no mesmo gesto. Não são dois nós rivais — é um, com dois polos.
\prov{relações: relaciona broca\_rede\_social\_soberana; relaciona goldchain\_como\_rede; usa meta\_inducao\_geral\_simbolo; usa rede\_limpa\_bump; usa banda\_exclusiva\_por\_tecido; usa traducao\_como\_transformacao; relaciona lei\_unica\_tres\_dominios}
\prov{proveniência: wiki \textperiodcentered{} semilla\_orbita\_dois\_atratores \textperiodcentered{} inferencia \textperiodcentered{} confiança 1.0 \textperiodcentered{} domínio arquitetura \textperiodcentered{} história 8 versões no jornal}

%CRISTAL {"arestas": [{"alvo": "decomposicao_espectral_imagem", "confianca": 1.0, "epistemico": "fato", "origem": "wiki", "rel": "especializa"}, {"alvo": "orbita_de_partida", "confianca": 1.0, "epistemico": "fato", "origem": "wiki", "rel": "relaciona"}, {"alvo": "lyapunov_da_orbita", "confianca": 1.0, "epistemico": "fato", "origem": "wiki", "rel": "usa"}], "confianca": 1.0, "contraexemplos": [], "descricao": "A decomposição espectral de UMA imagem dá camadas estáticas; de uma SEQUÊNCIA de frames, dá camadas com DINÂMICA. A POD (SVD do vídeo empilhado, cada frame uma coluna) separa os MODOS ESPACIAIS (os eigen-frames) dos seus COEFICIENTES TEMPORAIS — e cada coeficiente é a ÓRBITA daquele modo no tempo. O movimento vira modos + órbitas: o vocabulário espacial e as trajetórias que o movem.", "epistemico": "fato", "exemplos": [], "id": "orbita_temporal_das_camadas", "memoria": "semantica", "meta": {"dominio": "imagem", "fonte": "selo e órbita temporal"}, "origem": "wiki", "palavras_chave": [], "sinonimos": [], "tipo": "conceito", "titulo": "A Órbita Temporal das Camadas"}
\section{A Órbita Temporal das Camadas}
A decomposição espectral de UMA imagem dá camadas estáticas; de uma SEQUÊNCIA de frames, dá camadas com DINÂMICA. A POD (SVD do vídeo empilhado, cada frame uma coluna) separa os MODOS ESPACIAIS (os eigen-frames) dos seus COEFICIENTES TEMPORAIS — e cada coeficiente é a ÓRBITA daquele modo no tempo. O movimento vira modos + órbitas: o vocabulário espacial e as trajetórias que o movem.
\prov{relações: especializa decomposicao\_espectral\_imagem; relaciona orbita\_de\_partida; usa lyapunov\_da\_orbita}
\prov{proveniência: wiki \textperiodcentered{} selo e órbita temporal \textperiodcentered{} fato \textperiodcentered{} confiança 1.0 \textperiodcentered{} domínio imagem \textperiodcentered{} história 4 versões no jornal}

%CRISTAL {"arestas": [{"alvo": "roteamento", "confianca": 1.0, "epistemico": "fato", "origem": "wiki", "rel": "especializa"}, {"alvo": "eng_telecomunicacoes", "confianca": 1.0, "epistemico": "fato", "origem": "curriculo", "rel": "parte_de"}], "confianca": 1.0, "contraexemplos": [], "descricao": "Protocolo de estado de enlace: cada roteador anuncia seus enlaces, todos montam o mapa e calculam o menor caminho (Dijkstra). O roteamento DENTRO de um domínio administrativo (IGP).", "epistemico": "fato", "exemplos": [], "id": "ospf", "memoria": "semantica", "meta": {"dominio": "engenharia", "fonte": ""}, "origem": "wiki", "palavras_chave": [], "sinonimos": [], "tipo": "conceito", "titulo": "OSPF (Roteamento Interno)"}
\section{OSPF (Roteamento Interno)}
Protocolo de estado de enlace: cada roteador anuncia seus enlaces, todos montam o mapa e calculam o menor caminho (Dijkstra). O roteamento DENTRO de um domínio administrativo (IGP).
\prov{relações: especializa roteamento; parte\_de eng\_telecomunicacoes}
\prov{proveniência: wiki \textperiodcentered{} fato \textperiodcentered{} confiança 1.0 \textperiodcentered{} domínio engenharia \textperiodcentered{} história 3 versões no jornal}

%CRISTAL {"arestas": [{"alvo": "olho_de_cock_web", "confianca": 1.0, "epistemico": "fato", "origem": "wiki", "rel": "parte_de"}, {"alvo": "goldchain", "confianca": 1.0, "epistemico": "fato", "origem": "wiki", "rel": "relaciona"}, {"alvo": "colapso_koch_tiffany", "confianca": 1.0, "epistemico": "fato", "origem": "wiki", "rel": "usa"}], "confianca": 1.0, "contraexemplos": [], "descricao": "O painel CRISTAL·RASTER: pipeline Koch→coord φ→atrator 3D→câmera 1/w→ρ→T=√(ρ/ρmax) — empacota o snapshot em bits, reconstrói um atrator por coordenadas-φ (razão áurea), projeta com câmera e rasteriza densidade→temperatura. Um frame por evento SSE, sem buffer.", "epistemico": "fato", "exemplos": [], "id": "ouro_raytracer_canvas", "memoria": "semantica", "meta": {"autor": "broca-so", "dominio": "cockpit", "fonte": "goldchain/cockpit/cockpit.js:OuroRaytracer"}, "origem": "wiki", "palavras_chave": [], "sinonimos": [], "tipo": "conceito", "titulo": "Canvas Cristal (OuroRaytracer)"}
\section{Canvas Cristal (OuroRaytracer)}
O painel CRISTAL\textperiodcentered RASTER: pipeline Koch\ensuremath{\to}coord \ensuremath{\varphi}\ensuremath{\to}atrator 3D\ensuremath{\to}câmera 1/w\ensuremath{\to}\ensuremath{\rho}\ensuremath{\to}T=\ensuremath{\sqrt{\,}}(\ensuremath{\rho}/\ensuremath{\rho}max) — empacota o snapshot em bits, reconstrói um atrator por coordenadas-\ensuremath{\varphi} (razão áurea), projeta com câmera e rasteriza densidade\ensuremath{\to}temperatura. Um frame por evento SSE, sem buffer.
\prov{relações: parte\_de olho\_de\_cock\_web; relaciona goldchain; usa colapso\_koch\_tiffany}
\prov{proveniência: wiki \textperiodcentered{} goldchain/cockpit/cockpit.js:OuroRaytracer \textperiodcentered{} fato \textperiodcentered{} confiança 1.0 \textperiodcentered{} domínio cockpit \textperiodcentered{} história 4 versões no jornal}

%CRISTAL {"arestas": [{"alvo": "estados_rrc_nr", "confianca": 1.0, "epistemico": "fato", "origem": "wiki", "rel": "usa"}, {"alvo": "funcoes_de_rede_5gc", "confianca": 1.0, "epistemico": "fato", "origem": "wiki", "rel": "usa"}, {"alvo": "acesso_aleatorio_nr", "confianca": 1.0, "epistemico": "fato", "origem": "wiki", "rel": "usa"}], "confianca": 1.0, "contraexemplos": [], "descricao": "Como a rede alcança um UE que está dormindo (IDLE/INACTIVE) quando chega tráfego ou uma chamada: envia um aviso de paging na área de registro/notificação, e o UE acorda e faz acesso aleatório para responder.", "epistemico": "fato", "exemplos": [], "id": "paginacao_5g", "memoria": "semantica", "meta": {"dominio": "engenharia", "fonte": ""}, "origem": "wiki", "palavras_chave": [], "sinonimos": [], "tipo": "conceito", "titulo": "Paginação (Paging)"}
\section{Paginação (Paging)}
Como a rede alcança um UE que está dormindo (IDLE/INACTIVE) quando chega tráfego ou uma chamada: envia um aviso de paging na área de registro/notificação, e o UE acorda e faz acesso aleatório para responder.
\prov{relações: usa estados\_rrc\_nr; usa funcoes\_de\_rede\_5gc; usa acesso\_aleatorio\_nr}
\prov{proveniência: wiki \textperiodcentered{} fato \textperiodcentered{} confiança 1.0 \textperiodcentered{} domínio engenharia \textperiodcentered{} história 4 versões no jornal}

%CRISTAL {"arestas": [{"alvo": "painel_estelar", "confianca": 1.0, "epistemico": "fato", "origem": "wiki", "rel": "explica"}, {"alvo": "cockpit_observatorio", "confianca": 1.0, "epistemico": "fato", "origem": "wiki", "rel": "parte_de"}, {"alvo": "matematica_estelar", "confianca": 1.0, "epistemico": "fato", "origem": "wiki", "rel": "usa"}, {"alvo": "invariante_estelar", "confianca": 1.0, "epistemico": "fato", "origem": "wiki", "rel": "usa"}], "confianca": 1.0, "contraexemplos": [], "descricao": "Instrumento com um núcleo (o invariante) e três braços (geometria, mecânica, termodinâmica) lidos sobre cada nonce PoW — 'não se fabricam três ciências, lê-se o mesmo gato em três braços'. Não altera consenso BTC nem SHA-256. Aprofunda/cristaliza o nó raso painel_estelar.", "epistemico": "fato", "exemplos": [], "id": "painel_estelar_obs", "memoria": "semantica", "meta": {"autor": "broca-so", "dominio": "cockpit", "fonte": "papers/painel_estelar.tex; ula/PAINEL_ESTELAR.md"}, "origem": "wiki", "palavras_chave": [], "sinonimos": [], "tipo": "conceito", "titulo": "Painel Estelar — o Observatório de Três Braços"}
\section{Painel Estelar — o Observatório de Três Braços}
Instrumento com um núcleo (o invariante) e três braços (geometria, mecânica, termodinâmica) lidos sobre cada nonce PoW — 'não se fabricam três ciências, lê-se o mesmo gato em três braços'. Não altera consenso BTC nem SHA-256. Aprofunda/cristaliza o nó raso painel\_estelar.
\prov{relações: explica painel\_estelar; parte\_de cockpit\_observatorio; usa matematica\_estelar; usa invariante\_estelar}
\prov{proveniência: wiki \textperiodcentered{} papers/painel\_estelar.tex; ula/PAINEL\_ESTELAR.md \textperiodcentered{} fato \textperiodcentered{} confiança 1.0 \textperiodcentered{} domínio cockpit \textperiodcentered{} história 5 versões no jornal}

%CRISTAL {"arestas": [{"alvo": "algebra_dupolinomial_espaco", "confianca": 1.0, "epistemico": "fato", "origem": "wiki", "rel": "usa"}, {"alvo": "symbolic_beta_pipe_so", "confianca": 1.0, "epistemico": "fato", "origem": "wiki", "rel": "relaciona"}, {"alvo": "tiffany", "confianca": 1.0, "epistemico": "fato", "origem": "wiki", "rel": "parte_de"}], "confianca": 1.0, "contraexemplos": [], "descricao": "A ÚNICA simplificação limpa que a varredura achou (agente arquitetura): pell.py reimplementava a recorrência que já é metal_pos.passo_metal(2,·). Feito: pell.py agora DELEGA a metal_pos (passo_pell=passo_metal(2), pell_ab_em=metal_ab_em(2)), mantendo só a POLÍTICA do goldchain (TETO, proximo_pell que pula P0,P1). Verificado idêntico ao Pell antigo em i=0..59; emite_pell ok; 62/62 testes. Zero churn de import, zero migração de ledger — a prata é o caso n=2 da álgebra posicional, e agora o código diz isso. [D] a duplicação silenciosa fechada.", "epistemico": "fato", "exemplos": [], "id": "pell_shim_fonte_unica", "memoria": "semantica", "meta": {"autor": "Lopes/broca-so", "dominio": "arquitetura", "fonte": "goldchain/pell.py (shim); goldchain/metal_pos.py (passo_metal, metal_ab_em)"}, "origem": "wiki", "palavras_chave": [], "sinonimos": [], "tipo": "conceito", "titulo": "Dedup: pell.py virou shim sobre metal_pos (n=2) — fonte única da recorrência"}
\section{Dedup: pell.py virou shim sobre metal\_pos (n=2) — fonte única da recorrência}
A ÚNICA simplificação limpa que a varredura achou (agente arquitetura): pell.py reimplementava a recorrência que já é metal\_pos.passo\_metal(2,\textperiodcentered ). Feito: pell.py agora DELEGA a metal\_pos (passo\_pell=passo\_metal(2), pell\_ab\_em=metal\_ab\_em(2)), mantendo só a POLÍTICA do goldchain (TETO, proximo\_pell que pula P0,P1). Verificado idêntico ao Pell antigo em i=0..59; emite\_pell ok; 62/62 testes. Zero churn de import, zero migração de ledger — a prata é o caso n=2 da álgebra posicional, e agora o código diz isso. [D] a duplicação silenciosa fechada.
\prov{relações: usa algebra\_dupolinomial\_espaco; relaciona symbolic\_beta\_pipe\_so; parte\_de tiffany}
\prov{proveniência: wiki \textperiodcentered{} goldchain/pell.py (shim); goldchain/metal\_pos.py (passo\_metal, metal\_ab\_em) \textperiodcentered{} fato \textperiodcentered{} confiança 1.0 \textperiodcentered{} domínio arquitetura \textperiodcentered{} história 30 versões no jornal}

%CRISTAL {"arestas": [{"alvo": "realimentacao_lyapunov", "confianca": 1.0, "epistemico": "fato", "origem": "wiki", "rel": "usa"}, {"alvo": "numero_de_salem", "confianca": 1.0, "epistemico": "fato", "origem": "wiki", "rel": "relaciona"}, {"alvo": "corpo_mineral", "confianca": 1.0, "epistemico": "fato", "origem": "wiki", "rel": "usa"}], "confianca": 1.0, "contraexemplos": [], "descricao": "Um pêndulo invertido é instável: em malha aberta cai (λ>0, ponto de sela, caos). O controle o segura no equilíbrio INSTÁVEL (λ<0) — o robô cavalga o ponto Salem entre tombar para um lado e para o outro. Andar é cair controlado; ficar de pé é viver na borda. Verif.: aberta λ=+0.69 (cai), controle λ=−0.69 (de pé), K na borda λ=0.", "epistemico": "inferencia", "exemplos": [], "id": "pendulo_invertido_borda", "memoria": "semantica", "meta": {"de_pe_controle": true, "dominio": "robotica", "fonte": "robótica e sistemas autônomos", "lambda_aberta": 0.693, "lambda_controle": -0.693}, "origem": "wiki", "palavras_chave": [], "sinonimos": [], "tipo": "conceito", "titulo": "Pêndulo Invertido = Equilibrar-se na Borda"}
\section{Pêndulo Invertido = Equilibrar-se na Borda}
Um pêndulo invertido é instável: em malha aberta cai (\ensuremath{\lambda}>0, ponto de sela, caos). O controle o segura no equilíbrio INSTÁVEL (\ensuremath{\lambda}<0) — o robô cavalga o ponto Salem entre tombar para um lado e para o outro. Andar é cair controlado; ficar de pé é viver na borda. Verif.: aberta \ensuremath{\lambda}=+0.69 (cai), controle \ensuremath{\lambda}=\ensuremath{-}0.69 (de pé), K na borda \ensuremath{\lambda}=0.
\prov{relações: usa realimentacao\_lyapunov; relaciona numero\_de\_salem; usa corpo\_mineral}
\prov{proveniência: wiki \textperiodcentered{} robótica e sistemas autônomos \textperiodcentered{} inferencia \textperiodcentered{} confiança 1.0 \textperiodcentered{} domínio robotica \textperiodcentered{} história 4 versões no jornal}

%CRISTAL {"arestas": [{"alvo": "transmissao_de_energia", "confianca": 1.0, "epistemico": "fato", "origem": "wiki", "rel": "explica"}, {"alvo": "corpo_mineral", "confianca": 1.0, "epistemico": "fato", "origem": "wiki", "rel": "relaciona"}], "confianca": 1.0, "contraexemplos": [], "descricao": "Todo condutor com resistência R dissipa P=I²R em calor. É a fricção da rede — o lado dissipativo (o 'a' de a+c²=1); transmitir em alta tensão existe justamente para minimizá-la.", "epistemico": "inferencia", "exemplos": [], "id": "perda_joule", "memoria": "semantica", "meta": {"dominio": "engenharia", "fonte": "distribuição e armazenamento"}, "origem": "wiki", "palavras_chave": [], "sinonimos": [], "tipo": "conceito", "titulo": "Perda Joule (I²R — a dissipação)"}
\section{Perda Joule (I\ensuremath{^{2}}R — a dissipação)}
Todo condutor com resistência R dissipa P=I\ensuremath{^{2}}R em calor. É a fricção da rede — o lado dissipativo (o 'a' de a+c\ensuremath{^{2}}=1); transmitir em alta tensão existe justamente para minimizá-la.
\prov{relações: explica transmissao\_de\_energia; relaciona corpo\_mineral}
\prov{proveniência: wiki \textperiodcentered{} distribuição e armazenamento \textperiodcentered{} inferencia \textperiodcentered{} confiança 1.0 \textperiodcentered{} domínio engenharia \textperiodcentered{} história 6 versões no jornal}

%CRISTAL {"arestas": [{"alvo": "pi_emerge_do_limite", "confianca": 1.0, "epistemico": "fato", "origem": "wiki", "rel": "usa"}, {"alvo": "numeros_de_pisot", "confianca": 1.0, "epistemico": "fato", "origem": "wiki", "rel": "relaciona"}, {"alvo": "numero_de_salem", "confianca": 1.0, "epistemico": "fato", "origem": "wiki", "rel": "relaciona"}, {"alvo": "corpo_mineral", "confianca": 1.0, "epistemico": "fato", "origem": "wiki", "rel": "contradiz"}], "confianca": 1.0, "contraexemplos": [], "descricao": "A reta mineral é ALGÉBRICA (os metais σ_n: Pisot=cristal, Salem=borda, caos). π NÃO está nela: sua fração contínua é aperiódica — TRANSCENDENTE, não quadrática como φ e √2. π não é metal nem Pisot nem Salem; ressoa com a reta mineral (é o limite dos polígonos) mas vive ACIMA dela. Verif.: classifica_real(π)='superior (não-quadrático)'.", "epistemico": "fato", "exemplos": [], "id": "pi_acima_da_reta_mineral", "memoria": "semantica", "meta": {"dominio": "imagem", "fonte": "formas e limite"}, "origem": "wiki", "palavras_chave": [], "sinonimos": [], "tipo": "conceito", "titulo": "π Está Acima da Reta Mineral"}
\section{\ensuremath{\pi} Está Acima da Reta Mineral}
A reta mineral é ALGÉBRICA (os metais \ensuremath{\sigma}\_n: Pisot=cristal, Salem=borda, caos). \ensuremath{\pi} NÃO está nela: sua fração contínua é aperiódica — TRANSCENDENTE, não quadrática como \ensuremath{\varphi} e \ensuremath{\sqrt{\,}}2. \ensuremath{\pi} não é metal nem Pisot nem Salem; ressoa com a reta mineral (é o limite dos polígonos) mas vive ACIMA dela. Verif.: classifica\_real(\ensuremath{\pi})='superior (não-quadrático)'.
\prov{relações: usa pi\_emerge\_do\_limite; relaciona numeros\_de\_pisot; relaciona numero\_de\_salem; contradiz corpo\_mineral}
\prov{proveniência: wiki \textperiodcentered{} formas e limite \textperiodcentered{} fato \textperiodcentered{} confiança 1.0 \textperiodcentered{} domínio imagem \textperiodcentered{} história 10 versões no jornal}

%CRISTAL {"arestas": [{"alvo": "poligono_tende_ao_circulo", "confianca": 1.0, "epistemico": "fato", "origem": "wiki", "rel": "usa"}, {"alvo": "corpo_mineral", "confianca": 1.0, "epistemico": "fato", "origem": "wiki", "rel": "usa"}], "confianca": 1.0, "contraexemplos": [], "descricao": "No limite polígono→círculo, a razão área/R² — medida direto nos pixels — CONVERGE para π. É o método de Arquimedes: os polígonos são os aproximantes, o círculo é o limite, π é a sua assinatura. π é o INVARIANTE do atrator circular, como cada metal σ_n é a assinatura da sua cadeia. Verif.: (n/2)·sin(2π/n) → π.", "epistemico": "fato", "exemplos": [], "id": "pi_emerge_do_limite", "memoria": "semantica", "meta": {"dominio": "imagem", "fonte": "formas e limite"}, "origem": "wiki", "palavras_chave": [], "sinonimos": [], "tipo": "conceito", "titulo": "π Emerge do Limite (Arquimedes)"}
\section{\ensuremath{\pi} Emerge do Limite (Arquimedes)}
No limite polígono\ensuremath{\to}círculo, a razão área/R\ensuremath{^{2}} — medida direto nos pixels — CONVERGE para \ensuremath{\pi}. É o método de Arquimedes: os polígonos são os aproximantes, o círculo é o limite, \ensuremath{\pi} é a sua assinatura. \ensuremath{\pi} é o INVARIANTE do atrator circular, como cada metal \ensuremath{\sigma}\_n é a assinatura da sua cadeia. Verif.: (n/2)\textperiodcentered sin(2\ensuremath{\pi}/n) \ensuremath{\to} \ensuremath{\pi}.
\prov{relações: usa poligono\_tende\_ao\_circulo; usa corpo\_mineral}
\prov{proveniência: wiki \textperiodcentered{} formas e limite \textperiodcentered{} fato \textperiodcentered{} confiança 1.0 \textperiodcentered{} domínio imagem \textperiodcentered{} história 6 versões no jornal}

%CRISTAL {"arestas": [{"alvo": "pi_medida_da_borda", "confianca": 1.0, "epistemico": "fato", "origem": "wiki", "rel": "usa"}, {"alvo": "quasicristal_de_rotacoes", "confianca": 1.0, "epistemico": "fato", "origem": "wiki", "rel": "relaciona"}, {"alvo": "transformada_de_fourier", "confianca": 1.0, "epistemico": "fato", "origem": "wiki", "rel": "relaciona"}], "confianca": 1.0, "contraexemplos": [], "descricao": "As simetrias da reta mineral são DISCRETAS (os grupos D_n, as raízes da unidade e^{2πik/n} — pontos racionais no círculo). π gera a simetria CONTÍNUA (o grupo U(1), o círculo inteiro e^{2πi·x}). Assim como σ_n gera a cadeia metálica, π gera o círculo — a RETA TRANSCENDENTE ACIMA. Achar o limite (π) é achar o gancho para a próxima reta. Ponte honesta, não prova.", "epistemico": "hipotese", "exemplos": [], "id": "pi_gerador_da_reta_acima", "memoria": "semantica", "meta": {"dominio": "imagem", "fonte": "formas e limite"}, "origem": "wiki", "palavras_chave": [], "sinonimos": [], "tipo": "conceito", "titulo": "π é o Gerador da Reta Acima"}
\section{\ensuremath{\pi} é o Gerador da Reta Acima}
As simetrias da reta mineral são DISCRETAS (os grupos D\_n, as raízes da unidade e\^{}\{2\ensuremath{\pi}ik/n\} — pontos racionais no círculo). \ensuremath{\pi} gera a simetria CONTÍNUA (o grupo U(1), o círculo inteiro e\^{}\{2\ensuremath{\pi}i\textperiodcentered x\}). Assim como \ensuremath{\sigma}\_n gera a cadeia metálica, \ensuremath{\pi} gera o círculo — a RETA TRANSCENDENTE ACIMA. Achar o limite (\ensuremath{\pi}) é achar o gancho para a próxima reta. Ponte honesta, não prova.
\prov{relações: usa pi\_medida\_da\_borda; relaciona quasicristal\_de\_rotacoes; relaciona transformada\_de\_fourier}
\prov{proveniência: wiki \textperiodcentered{} formas e limite \textperiodcentered{} hipotese \textperiodcentered{} confiança 1.0 \textperiodcentered{} domínio imagem \textperiodcentered{} história 8 versões no jornal}

%CRISTAL {"arestas": [{"alvo": "pi_acima_da_reta_mineral", "confianca": 1.0, "epistemico": "fato", "origem": "wiki", "rel": "usa"}, {"alvo": "numero_de_salem", "confianca": 1.0, "epistemico": "fato", "origem": "wiki", "rel": "relaciona"}, {"alvo": "ordem_borda_caos", "confianca": 1.0, "epistemico": "fato", "origem": "wiki", "rel": "relaciona"}], "confianca": 1.0, "contraexemplos": [], "descricao": "A borda da parábola mineral é o círculo ρ=1 (onde vivem os Salem, λ=0). As raízes da unidade (os vértices do polígono) têm todas |z|=1 → estão na borda. π é a MEDIDA dessa borda: a circunferência do círculo ρ=1. π não é um ponto da parábola — é o comprimento da sua crista, o número da borda entre o cristal (|z|<1) e o caos (|z|>1).", "epistemico": "inferencia", "exemplos": [], "id": "pi_medida_da_borda", "memoria": "semantica", "meta": {"dominio": "imagem", "fonte": "formas e limite"}, "origem": "wiki", "palavras_chave": [], "sinonimos": [], "tipo": "conceito", "titulo": "π é a Medida da Borda (ρ=1)"}
\section{\ensuremath{\pi} é a Medida da Borda (\ensuremath{\rho}=1)}
A borda da parábola mineral é o círculo \ensuremath{\rho}=1 (onde vivem os Salem, \ensuremath{\lambda}=0). As raízes da unidade (os vértices do polígono) têm todas |z|=1 \ensuremath{\to} estão na borda. \ensuremath{\pi} é a MEDIDA dessa borda: a circunferência do círculo \ensuremath{\rho}=1. \ensuremath{\pi} não é um ponto da parábola — é o comprimento da sua crista, o número da borda entre o cristal (|z|<1) e o caos (|z|>1).
\prov{relações: usa pi\_acima\_da\_reta\_mineral; relaciona numero\_de\_salem; relaciona ordem\_borda\_caos}
\prov{proveniência: wiki \textperiodcentered{} formas e limite \textperiodcentered{} inferencia \textperiodcentered{} confiança 1.0 \textperiodcentered{} domínio imagem \textperiodcentered{} história 8 versões no jornal}

%CRISTAL {"arestas": [{"alvo": "rede_autonoma_broca", "confianca": 1.0, "epistemico": "fato", "origem": "wiki", "rel": "usa"}, {"alvo": "repositorio_broca_os", "confianca": 1.0, "epistemico": "fato", "origem": "wiki", "rel": "usa"}, {"alvo": "gex44_repo_central", "confianca": 1.0, "epistemico": "fato", "origem": "wiki", "rel": "usa"}, {"alvo": "ci_cd", "confianca": 1.0, "epistemico": "fato", "origem": "wiki", "rel": "especializa"}, {"alvo": "cristal_pipeline_runner", "confianca": 1.0, "epistemico": "fato", "origem": "wiki", "rel": "relaciona"}], "confianca": 1.0, "contraexemplos": [], "descricao": "O código se atualiza no gex44 (repo central) pela PRÓPRIA rede que ele define — dogfooding. Fluxo scripts/gex44/deploy_banda.sh: bootstrap (sync seguro) → up (olho+agente) → deploy (por banda) → verify (checksum) → down. O deploy não é comandado de fora: o código flui pela antena.", "epistemico": "fato", "exemplos": [], "id": "pipeline_deploy_broca", "memoria": "semantica", "meta": {"dominio": "arquitetura", "fonte": "scripts/gex44/deploy_banda.sh"}, "origem": "wiki", "palavras_chave": [], "sinonimos": [], "tipo": "conceito", "titulo": "O Pipeline de Deploy do broca-os"}
\section{O Pipeline de Deploy do broca-os}
O código se atualiza no gex44 (repo central) pela PRÓPRIA rede que ele define — dogfooding. Fluxo scripts/gex44/deploy\_banda.sh: bootstrap (sync seguro) \ensuremath{\to} up (olho+agente) \ensuremath{\to} deploy (por banda) \ensuremath{\to} verify (checksum) \ensuremath{\to} down. O deploy não é comandado de fora: o código flui pela antena.
\prov{relações: usa rede\_autonoma\_broca; usa repositorio\_broca\_os; usa gex44\_repo\_central; especializa ci\_cd; relaciona cristal\_pipeline\_runner}
\prov{proveniência: wiki \textperiodcentered{} scripts/gex44/deploy\_banda.sh \textperiodcentered{} fato \textperiodcentered{} confiança 1.0 \textperiodcentered{} domínio arquitetura \textperiodcentered{} história 6 versões no jornal}

%CRISTAL {"arestas": [{"alvo": "circuito_integrado", "confianca": 1.0, "epistemico": "fato", "origem": "wiki", "rel": "usa"}, {"alvo": "eng_eletronica", "confianca": 1.0, "epistemico": "fato", "origem": "curriculo", "rel": "parte_de"}], "confianca": 1.0, "contraexemplos": [], "descricao": "O substrato que interliga os componentes: trilhas, camadas, integridade de sinal e de potência. Onde o esquemático vira objeto físico manufaturável.", "epistemico": "inferencia", "exemplos": [], "id": "placa_circuito_impresso", "memoria": "semantica", "meta": {"dominio": "engenharia", "fonte": ""}, "origem": "wiki", "palavras_chave": [], "sinonimos": [], "tipo": "conceito", "titulo": "Placa de Circuito Impresso (PCB)"}
\section{Placa de Circuito Impresso (PCB)}
O substrato que interliga os componentes: trilhas, camadas, integridade de sinal e de potência. Onde o esquemático vira objeto físico manufaturável.
\prov{relações: usa circuito\_integrado; parte\_de eng\_eletronica}
\prov{proveniência: wiki \textperiodcentered{} inferencia \textperiodcentered{} confiança 1.0 \textperiodcentered{} domínio engenharia \textperiodcentered{} história 3 versões no jornal}

%CRISTAL {"arestas": [{"alvo": "minimax_e_colapso", "confianca": 1.0, "epistemico": "fato", "origem": "wiki", "rel": "eh_um"}, {"alvo": "sistema_autonomo", "confianca": 1.0, "epistemico": "fato", "origem": "wiki", "rel": "usa"}, {"alvo": "bai", "confianca": 1.0, "epistemico": "fato", "origem": "wiki", "rel": "usa"}], "confianca": 1.0, "contraexemplos": [], "descricao": "Planejar uma rota ou sequência de ações é buscar sobre futuros possíveis com um ORÇAMENTO de profundidade (MPC, A*, árvore de busca) e escolher o melhor ramo — o mesmo colapso do xadrez (minimax + poda) e o Ψ do BAI com seu δ. O robô que planeja está jogando xadrez contra o mundo.", "epistemico": "inferencia", "exemplos": [], "id": "planejamento_minimax_delta", "memoria": "semantica", "meta": {"dominio": "robotica", "fonte": "robótica e sistemas autônomos"}, "origem": "wiki", "palavras_chave": [], "sinonimos": [], "tipo": "conceito", "titulo": "Planejar = o Colapso Minimax com Orçamento δ"}
\section{Planejar = o Colapso Minimax com Orçamento \ensuremath{\delta}}
Planejar uma rota ou sequência de ações é buscar sobre futuros possíveis com um ORÇAMENTO de profundidade (MPC, A*, árvore de busca) e escolher o melhor ramo — o mesmo colapso do xadrez (minimax + poda) e o \ensuremath{\Psi} do BAI com seu \ensuremath{\delta}. O robô que planeja está jogando xadrez contra o mundo.
\prov{relações: eh\_um minimax\_e\_colapso; usa sistema\_autonomo; usa bai}
\prov{proveniência: wiki \textperiodcentered{} robótica e sistemas autônomos \textperiodcentered{} inferencia \textperiodcentered{} confiança 1.0 \textperiodcentered{} domínio robotica \textperiodcentered{} história 4 versões no jornal}

%CRISTAL {"arestas": [{"alvo": "eletronica_de_rf", "confianca": 1.0, "epistemico": "fato", "origem": "wiki", "rel": "usa"}, {"alvo": "realimentacao_eletronica", "confianca": 1.0, "epistemico": "fato", "origem": "wiki", "rel": "usa"}, {"alvo": "eng_eletronica", "confianca": 1.0, "epistemico": "fato", "origem": "curriculo", "rel": "parte_de"}], "confianca": 1.0, "contraexemplos": [], "descricao": "Uma malha que trava um oscilador na fase de uma referência — sintetiza frequências, recupera relógio e demodula. Onipresente em rádio, clock e comunicação digital.", "epistemico": "fato", "exemplos": [], "id": "pll", "memoria": "semantica", "meta": {"dominio": "engenharia", "fonte": ""}, "origem": "wiki", "palavras_chave": [], "sinonimos": [], "tipo": "conceito", "titulo": "PLL (Malha de Captura de Fase)"}
\section{PLL (Malha de Captura de Fase)}
Uma malha que trava um oscilador na fase de uma referência — sintetiza frequências, recupera relógio e demodula. Onipresente em rádio, clock e comunicação digital.
\prov{relações: usa eletronica\_de\_rf; usa realimentacao\_eletronica; parte\_de eng\_eletronica}
\prov{proveniência: wiki \textperiodcentered{} fato \textperiodcentered{} confiança 1.0 \textperiodcentered{} domínio engenharia \textperiodcentered{} história 4 versões no jornal}

%CRISTAL {"arestas": [{"alvo": "belo_platonico", "confianca": 1.0, "epistemico": "fato", "origem": "wiki", "rel": "relaciona"}, {"alvo": "razao_aurea", "confianca": 1.0, "epistemico": "fato", "origem": "wiki", "rel": "usa"}, {"alvo": "poligono_tende_ao_circulo", "confianca": 1.0, "epistemico": "fato", "origem": "wiki", "rel": "relaciona"}, {"alvo": "corpo_estelar", "confianca": 1.0, "epistemico": "fato", "origem": "wiki", "rel": "usa"}], "confianca": 1.0, "contraexemplos": [], "descricao": "Os cinco sólidos platônicos, e a máquina os lê como minerais. O ICOSAEDRO (vértices = permutações cíclicas de (0,±1,±φ)) e o DODECAEDRO (o cubo + (0,±1/φ,±φ)) são construídos a partir de φ = σ_1, o metal OURO: a razão áurea que os funda é o primeiro metal, e a geometria dos sólidos é a reta mineral. Verif. em transformacoes_geometricas.py: V/A corretos (tetra 4/6, cubo 8/12, octa 6/12, icosa 12/30, dodeca 20/30) e φ nos vértices do icosaedro (0,1,φ).", "epistemico": "fato", "exemplos": [], "id": "poliedros_de_ouro", "memoria": "semantica", "meta": {"dominio": "imagem", "fonte": "poliedros minerais"}, "origem": "wiki", "palavras_chave": [], "sinonimos": [], "tipo": "conceito", "titulo": "Os Poliedros são de Ouro (φ)"}
\section{Os Poliedros são de Ouro (\ensuremath{\varphi})}
Os cinco sólidos platônicos, e a máquina os lê como minerais. O ICOSAEDRO (vértices = permutações cíclicas de (0,$\pm$1,$\pm$\ensuremath{\varphi})) e o DODECAEDRO (o cubo + (0,$\pm$1/\ensuremath{\varphi},$\pm$\ensuremath{\varphi})) são construídos a partir de \ensuremath{\varphi} = \ensuremath{\sigma}\_1, o metal OURO: a razão áurea que os funda é o primeiro metal, e a geometria dos sólidos é a reta mineral. Verif. em transformacoes\_geometricas.py: V/A corretos (tetra 4/6, cubo 8/12, octa 6/12, icosa 12/30, dodeca 20/30) e \ensuremath{\varphi} nos vértices do icosaedro (0,1,\ensuremath{\varphi}).
\prov{relações: relaciona belo\_platonico; usa razao\_aurea; relaciona poligono\_tende\_ao\_circulo; usa corpo\_estelar}
\prov{proveniência: wiki \textperiodcentered{} poliedros minerais \textperiodcentered{} fato \textperiodcentered{} confiança 1.0 \textperiodcentered{} domínio imagem \textperiodcentered{} história 5 versões no jornal}

%CRISTAL {"arestas": [{"alvo": "grupo_de_transformacoes", "confianca": 1.0, "epistemico": "fato", "origem": "wiki", "rel": "usa"}, {"alvo": "numero_de_salem", "confianca": 1.0, "epistemico": "fato", "origem": "wiki", "rel": "relaciona"}], "confianca": 1.0, "contraexemplos": [], "descricao": "Aumentando o número de lados, a ordem de simetria → ∞ e a irregularidade → 0: o polígono vira CÍRCULO. O círculo é o ATRATOR da sequência de polígonos — a forma de simetria máxima, o cristal circular (ρ=1). Verif.: irregularidade 0.22 (triângulo) → 0.001 (200-ágono).", "epistemico": "inferencia", "exemplos": [], "id": "poligono_tende_ao_circulo", "memoria": "semantica", "meta": {"dominio": "imagem", "fonte": "formas e limite"}, "origem": "wiki", "palavras_chave": [], "sinonimos": [], "tipo": "conceito", "titulo": "O Polígono Tende ao Círculo"}
\section{O Polígono Tende ao Círculo}
Aumentando o número de lados, a ordem de simetria \ensuremath{\to} \ensuremath{\infty} e a irregularidade \ensuremath{\to} 0: o polígono vira CÍRCULO. O círculo é o ATRATOR da sequência de polígonos — a forma de simetria máxima, o cristal circular (\ensuremath{\rho}=1). Verif.: irregularidade 0.22 (triângulo) \ensuremath{\to} 0.001 (200-ágono).
\prov{relações: usa grupo\_de\_transformacoes; relaciona numero\_de\_salem}
\prov{proveniência: wiki \textperiodcentered{} formas e limite \textperiodcentered{} inferencia \textperiodcentered{} confiança 1.0 \textperiodcentered{} domínio imagem \textperiodcentered{} história 6 versões no jornal}

%CRISTAL {"arestas": [{"alvo": "porta_mecanica", "confianca": 1.0, "epistemico": "fato", "origem": "wiki", "rel": "explica"}, {"alvo": "cotacao_parabola", "confianca": 1.0, "epistemico": "fato", "origem": "wiki", "rel": "usa"}, {"alvo": "broca_cockpit_engine", "confianca": 1.0, "epistemico": "fato", "origem": "wiki", "rel": "usa"}, {"alvo": "colapso_koch_tiffany", "confianca": 1.0, "epistemico": "fato", "origem": "wiki", "rel": "relaciona"}], "confianca": 1.0, "contraexemplos": [], "descricao": "porta_aberta(payload) devolve True só quando a costura fecha por dentro simultaneamente: fechamento.fechou E parabola_global.ok (a+c²=1) E sincronizacao.sincronizado_ok E (se houver) hodge_quantico.ok. O poll de 0.8s é só carga de CPU, não é tick — a porta abre por CORRELAÇÃO mecânica, não por tempo. Formaliza a porta_mecanica.", "epistemico": "fato", "exemplos": [], "id": "porta_hodge_correlacao", "memoria": "semantica", "meta": {"autor": "broca-so", "dominio": "cockpit", "fonte": "goldchain/hodge_automatizar.py:porta_aberta"}, "origem": "wiki", "palavras_chave": [], "sinonimos": [], "tipo": "conceito", "titulo": "A Porta Hodge — Colapso por Correlação (sem relógio)"}
\section{A Porta Hodge — Colapso por Correlação (sem relógio)}
porta\_aberta(payload) devolve True só quando a costura fecha por dentro simultaneamente: fechamento.fechou E parabola\_global.ok (a+c\ensuremath{^{2}}=1) E sincronizacao.sincronizado\_ok E (se houver) hodge\_quantico.ok. O poll de 0.8s é só carga de CPU, não é tick — a porta abre por CORRELAÇÃO mecânica, não por tempo. Formaliza a porta\_mecanica.
\prov{relações: explica porta\_mecanica; usa cotacao\_parabola; usa broca\_cockpit\_engine; relaciona colapso\_koch\_tiffany}
\prov{proveniência: wiki \textperiodcentered{} goldchain/hodge\_automatizar.py:porta\_aberta \textperiodcentered{} fato \textperiodcentered{} confiança 1.0 \textperiodcentered{} domínio cockpit \textperiodcentered{} história 5 versões no jornal}

%CRISTAL {"arestas": [{"alvo": "semicondutor", "confianca": 1.0, "epistemico": "fato", "origem": "wiki", "rel": "parte_de"}], "confianca": 1.0, "contraexemplos": [], "descricao": "O que carrega corrente no semicondutor: o elétron (negativo) e a LACUNA (a ausência de elétron, que se move como carga positiva).", "epistemico": "fato", "exemplos": [], "id": "portador_de_carga", "memoria": "semantica", "meta": {"dominio": "engenharia", "fonte": ""}, "origem": "wiki", "palavras_chave": [], "sinonimos": [], "tipo": "conceito", "titulo": "Portador de Carga"}
\section{Portador de Carga}
O que carrega corrente no semicondutor: o elétron (negativo) e a LACUNA (a ausência de elétron, que se move como carga positiva).
\prov{relações: parte\_de semicondutor}
\prov{proveniência: wiki \textperiodcentered{} fato \textperiodcentered{} confiança 1.0 \textperiodcentered{} domínio engenharia \textperiodcentered{} história 2 versões no jornal}

%CRISTAL {"arestas": [{"alvo": "onda_portadora", "confianca": 1.0, "epistemico": "fato", "origem": "wiki", "rel": "explica"}, {"alvo": "numero_de_salem", "confianca": 1.0, "epistemico": "fato", "origem": "wiki", "rel": "eh_um"}, {"alvo": "corrente_alternada", "confianca": 1.0, "epistemico": "fato", "origem": "wiki", "rel": "relaciona"}, {"alvo": "corpo_mineral", "confianca": 1.0, "epistemico": "fato", "origem": "wiki", "rel": "usa"}], "confianca": 1.0, "contraexemplos": [], "descricao": "A onda portadora e^{iωt} percorre o CÍRCULO UNITÁRIO (ρ=1): é a rotação pura — o mesmo Salem/borda da corrente alternada e do neurônio-rádio. Modular é PERTURBAR essa rotação para gravar bits; a informação viaja como um floco na rotação. Verif.: raio=1.000, no círculo.", "epistemico": "inferencia", "exemplos": [], "id": "portadora_rotacao_salem", "memoria": "semantica", "meta": {"dominio": "engenharia", "fonte": "conectividade e 5G", "no_circulo": true, "raio": 1.0}, "origem": "wiki", "palavras_chave": [], "sinonimos": [], "tipo": "conceito", "titulo": "A Portadora é uma Rotação (Salem)"}
\section{A Portadora é uma Rotação (Salem)}
A onda portadora e\^{}\{i\ensuremath{\omega}t\} percorre o CÍRCULO UNITÁRIO (\ensuremath{\rho}=1): é a rotação pura — o mesmo Salem/borda da corrente alternada e do neurônio-rádio. Modular é PERTURBAR essa rotação para gravar bits; a informação viaja como um floco na rotação. Verif.: raio=1.000, no círculo.
\prov{relações: explica onda\_portadora; eh\_um numero\_de\_salem; relaciona corrente\_alternada; usa corpo\_mineral}
\prov{proveniência: wiki \textperiodcentered{} conectividade e 5G \textperiodcentered{} inferencia \textperiodcentered{} confiança 1.0 \textperiodcentered{} domínio engenharia \textperiodcentered{} história 5 versões no jornal}

%CRISTAL {"arestas": [{"alvo": "corrente_alternada_continua", "confianca": 1.0, "epistemico": "fato", "origem": "wiki", "rel": "usa"}, {"alvo": "circuitos_eletricos", "confianca": 1.0, "epistemico": "fato", "origem": "wiki", "rel": "usa"}, {"alvo": "eng_sistemas_energia", "confianca": 1.0, "epistemico": "fato", "origem": "curriculo", "rel": "parte_de"}], "confianca": 1.0, "contraexemplos": [], "descricao": "Em CA, a potência tem parte que faz trabalho (ativa, W) e parte que só oscila com o campo (reativa, VAr). O fator de potência mede o desperdício; corrigi-lo é chave na eficiência do SEP.", "epistemico": "fato", "exemplos": [], "id": "potencia_ativa_reativa", "memoria": "semantica", "meta": {"dominio": "engenharia", "fonte": ""}, "origem": "wiki", "palavras_chave": [], "sinonimos": [], "tipo": "conceito", "titulo": "Potência Ativa, Reativa e Fator de Potência"}
\section{Potência Ativa, Reativa e Fator de Potência}
Em CA, a potência tem parte que faz trabalho (ativa, W) e parte que só oscila com o campo (reativa, VAr). O fator de potência mede o desperdício; corrigi-lo é chave na eficiência do SEP.
\prov{relações: usa corrente\_alternada\_continua; usa circuitos\_eletricos; parte\_de eng\_sistemas\_energia}
\prov{proveniência: wiki \textperiodcentered{} fato \textperiodcentered{} confiança 1.0 \textperiodcentered{} domínio engenharia \textperiodcentered{} história 4 versões no jornal}

%CRISTAL {"arestas": [{"alvo": "acesso_aleatorio_nr", "confianca": 1.0, "epistemico": "fato", "origem": "wiki", "rel": "parte_de"}, {"alvo": "ssb_sincronizacao_nr", "confianca": 1.0, "epistemico": "fato", "origem": "wiki", "rel": "usa"}], "confianca": 1.0, "contraexemplos": [], "descricao": "A assinatura Zadoff-Chu que o UE escolhe e transmite no PRACH; o instante/frequência define o RA-RNTI com que a rede endereça a resposta. Preâmbulos são recurso escasso — colisão é dois UEs escolhendo o mesmo.", "epistemico": "fato", "exemplos": [], "id": "preambulo_rach", "memoria": "semantica", "meta": {"dominio": "engenharia", "fonte": ""}, "origem": "wiki", "palavras_chave": [], "sinonimos": [], "tipo": "conceito", "titulo": "Preâmbulo e RA-RNTI"}
\section{Preâmbulo e RA-RNTI}
A assinatura Zadoff-Chu que o UE escolhe e transmite no PRACH; o instante/frequência define o RA-RNTI com que a rede endereça a resposta. Preâmbulos são recurso escasso — colisão é dois UEs escolhendo o mesmo.
\prov{relações: parte\_de acesso\_aleatorio\_nr; usa ssb\_sincronizacao\_nr}
\prov{proveniência: wiki \textperiodcentered{} fato \textperiodcentered{} confiança 1.0 \textperiodcentered{} domínio engenharia \textperiodcentered{} história 3 versões no jornal}

%CRISTAL {"arestas": [{"alvo": "matematica", "confianca": 1.0, "epistemico": "fato", "origem": "wiki", "rel": "parte_de"}, {"alvo": "eng_eletronica", "confianca": 1.0, "epistemico": "fato", "origem": "curriculo", "rel": "parte_de"}, {"alvo": "eng_telecomunicacoes", "confianca": 1.0, "epistemico": "fato", "origem": "curriculo", "rel": "parte_de"}, {"alvo": "eng_sistemas_energia", "confianca": 1.0, "epistemico": "fato", "origem": "curriculo", "rel": "parte_de"}], "confianca": 1.0, "contraexemplos": [], "descricao": "O cálculo do acaso: variáveis aleatórias, distribuições, esperança e processos no tempo. Base do RUÍDO, da teoria da informação e da estimação — sem ela não há telecomunicação real.", "epistemico": "fato", "exemplos": [], "id": "probabilidade_e_estatistica", "memoria": "semantica", "meta": {"dominio": "engenharia", "fonte": ""}, "origem": "wiki", "palavras_chave": [], "sinonimos": [], "tipo": "conceito", "titulo": "Probabilidade e Processos Estocásticos"}
\section{Probabilidade e Processos Estocásticos}
O cálculo do acaso: variáveis aleatórias, distribuições, esperança e processos no tempo. Base do RUÍDO, da teoria da informação e da estimação — sem ela não há telecomunicação real.
\prov{relações: parte\_de matematica; parte\_de eng\_eletronica; parte\_de eng\_telecomunicacoes; parte\_de eng\_sistemas\_energia}
\prov{proveniência: wiki \textperiodcentered{} fato \textperiodcentered{} confiança 1.0 \textperiodcentered{} domínio engenharia \textperiodcentered{} história 5 versões no jornal}

%CRISTAL {"arestas": [{"alvo": "sinais_e_sistemas", "confianca": 1.0, "epistemico": "fato", "origem": "wiki", "rel": "usa"}, {"alvo": "transformada_de_fourier", "confianca": 1.0, "epistemico": "fato", "origem": "wiki", "rel": "usa"}, {"alvo": "eng_eletronica", "confianca": 1.0, "epistemico": "fato", "origem": "curriculo", "rel": "parte_de"}], "confianca": 1.0, "contraexemplos": [], "descricao": "Amostrar o sinal e processá-lo com números: filtros FIR/IIR, FFT, convolução. Barato, exato e programável — substituiu o analógico em áudio, imagem e comunicação.", "epistemico": "fato", "exemplos": [], "id": "processamento_digital_sinais", "memoria": "semantica", "meta": {"autor": "Shannon", "dominio": "engenharia", "fonte": ""}, "origem": "wiki", "palavras_chave": [], "sinonimos": [], "tipo": "conceito", "titulo": "Processamento Digital de Sinais (DSP)"}
\section{Processamento Digital de Sinais (DSP)}
Amostrar o sinal e processá-lo com números: filtros FIR/IIR, FFT, convolução. Barato, exato e programável — substituiu o analógico em áudio, imagem e comunicação.
\prov{relações: usa sinais\_e\_sistemas; usa transformada\_de\_fourier; parte\_de eng\_eletronica}
\prov{proveniência: wiki \textperiodcentered{} fato \textperiodcentered{} confiança 1.0 \textperiodcentered{} domínio engenharia \textperiodcentered{} história 4 versões no jornal}

%CRISTAL {"arestas": [{"alvo": "mosfet", "confianca": 1.0, "epistemico": "fato", "origem": "wiki", "rel": "usa"}, {"alvo": "microeletronica", "confianca": 1.0, "epistemico": "fato", "origem": "wiki", "rel": "parte_de"}, {"alvo": "isa_da_broca", "confianca": 1.0, "epistemico": "fato", "origem": "wiki", "rel": "explica"}, {"alvo": "eng_eletronica", "confianca": 1.0, "epistemico": "fato", "origem": "curriculo", "rel": "parte_de"}], "confianca": 1.0, "contraexemplos": [], "descricao": "A tecnologia dominante: pares complementares nMOS+pMOS que só consomem ao chavear. Baixa potência estática — por isso todo processador e memória é CMOS. É o NAND que o sandbox do broca-so realiza.", "epistemico": "fato", "exemplos": [], "id": "processo_cmos", "memoria": "semantica", "meta": {"dominio": "engenharia", "fonte": ""}, "origem": "wiki", "palavras_chave": [], "sinonimos": [], "tipo": "conceito", "titulo": "Processo CMOS"}
\section{Processo CMOS}
A tecnologia dominante: pares complementares nMOS+pMOS que só consomem ao chavear. Baixa potência estática — por isso todo processador e memória é CMOS. É o NAND que o sandbox do broca-so realiza.
\prov{relações: usa mosfet; parte\_de microeletronica; explica isa\_da\_broca; parte\_de eng\_eletronica}
\prov{proveniência: wiki \textperiodcentered{} fato \textperiodcentered{} confiança 1.0 \textperiodcentered{} domínio engenharia \textperiodcentered{} história 5 versões no jornal}

%CRISTAL {"arestas": [{"alvo": "circuito_integrado", "confianca": 1.0, "epistemico": "fato", "origem": "wiki", "rel": "usa"}, {"alvo": "eletronica_digital", "confianca": 1.0, "epistemico": "fato", "origem": "wiki", "rel": "especializa"}, {"alvo": "eng_eletronica", "confianca": 1.0, "epistemico": "fato", "origem": "curriculo", "rel": "parte_de"}], "confianca": 1.0, "contraexemplos": [], "descricao": "Integração em altíssima escala: descrever, sintetizar e verificar milhões de portas. Do RTL ao layout, com temporização, potência e área como restrições. Onde a arquitetura vira silício.", "epistemico": "fato", "exemplos": [], "id": "projeto_digital_vlsi", "memoria": "semantica", "meta": {"dominio": "engenharia", "fonte": ""}, "origem": "wiki", "palavras_chave": [], "sinonimos": [], "tipo": "conceito", "titulo": "Projeto Digital e VLSI"}
\section{Projeto Digital e VLSI}
Integração em altíssima escala: descrever, sintetizar e verificar milhões de portas. Do RTL ao layout, com temporização, potência e área como restrições. Onde a arquitetura vira silício.
\prov{relações: usa circuito\_integrado; especializa eletronica\_digital; parte\_de eng\_eletronica}
\prov{proveniência: wiki \textperiodcentered{} fato \textperiodcentered{} confiança 1.0 \textperiodcentered{} domínio engenharia \textperiodcentered{} história 4 versões no jornal}

%CRISTAL {"arestas": [{"alvo": "ondas_eletromagneticas", "confianca": 1.0, "epistemico": "fato", "origem": "wiki", "rel": "usa"}, {"alvo": "eng_telecomunicacoes", "confianca": 1.0, "epistemico": "fato", "origem": "curriculo", "rel": "parte_de"}], "confianca": 1.0, "contraexemplos": [], "descricao": "Como a onda perde potência e se distorce no caminho: perda no espaço livre, difração, reflexão, sombreamento e DESVANECIMENTO (fading) por multipercurso. O canal real, hostil e variável.", "epistemico": "fato", "exemplos": [], "id": "propagacao_de_ondas", "memoria": "semantica", "meta": {"dominio": "engenharia", "fonte": ""}, "origem": "wiki", "palavras_chave": [], "sinonimos": [], "tipo": "conceito", "titulo": "Propagação de Ondas de Rádio"}
\section{Propagação de Ondas de Rádio}
Como a onda perde potência e se distorce no caminho: perda no espaço livre, difração, reflexão, sombreamento e DESVANECIMENTO (fading) por multipercurso. O canal real, hostil e variável.
\prov{relações: usa ondas\_eletromagneticas; parte\_de eng\_telecomunicacoes}
\prov{proveniência: wiki \textperiodcentered{} fato \textperiodcentered{} confiança 1.0 \textperiodcentered{} domínio engenharia \textperiodcentered{} história 3 versões no jornal}

%CRISTAL {"arestas": [{"alvo": "sistema_eletrico_potencia", "confianca": 1.0, "epistemico": "fato", "origem": "wiki", "rel": "parte_de"}, {"alvo": "curto_circuito", "confianca": 1.0, "epistemico": "fato", "origem": "wiki", "rel": "usa"}, {"alvo": "eng_sistemas_energia", "confianca": 1.0, "epistemico": "fato", "origem": "curriculo", "rel": "parte_de"}], "confianca": 1.0, "contraexemplos": [], "descricao": "Detectar a falta e isolar só o trecho defeituoso em milissegundos — relés, disjuntores, coordenação e seletividade. O sistema imunológico da rede: rápido, seletivo e confiável.", "epistemico": "fato", "exemplos": [], "id": "protecao_de_sistemas", "memoria": "semantica", "meta": {"dominio": "engenharia", "fonte": ""}, "origem": "wiki", "palavras_chave": [], "sinonimos": [], "tipo": "conceito", "titulo": "Proteção de Sistemas Elétricos"}
\section{Proteção de Sistemas Elétricos}
Detectar a falta e isolar só o trecho defeituoso em milissegundos — relés, disjuntores, coordenação e seletividade. O sistema imunológico da rede: rápido, seletivo e confiável.
\prov{relações: parte\_de sistema\_eletrico\_potencia; usa curto\_circuito; parte\_de eng\_sistemas\_energia}
\prov{proveniência: wiki \textperiodcentered{} fato \textperiodcentered{} confiança 1.0 \textperiodcentered{} domínio engenharia \textperiodcentered{} história 4 versões no jornal}

%CRISTAL {"arestas": [{"alvo": "camada_de_enlace", "confianca": 1.0, "epistemico": "fato", "origem": "wiki", "rel": "usa"}, {"alvo": "protocolo_ip", "confianca": 1.0, "epistemico": "fato", "origem": "wiki", "rel": "usa"}, {"alvo": "eng_telecomunicacoes", "confianca": 1.0, "epistemico": "fato", "origem": "curriculo", "rel": "parte_de"}], "confianca": 1.0, "contraexemplos": [], "descricao": "Descobre qual MAC corresponde a um IP na rede local — a cola entre a camada de rede (IP) e a de enlace (MAC). Sem ARP, o pacote sabe o IP do vizinho mas não como lhe entregar o quadro.", "epistemico": "fato", "exemplos": [], "id": "protocolo_arp", "memoria": "semantica", "meta": {"dominio": "engenharia", "fonte": ""}, "origem": "wiki", "palavras_chave": [], "sinonimos": [], "tipo": "conceito", "titulo": "ARP"}
\section{ARP}
Descobre qual MAC corresponde a um IP na rede local — a cola entre a camada de rede (IP) e a de enlace (MAC). Sem ARP, o pacote sabe o IP do vizinho mas não como lhe entregar o quadro.
\prov{relações: usa camada\_de\_enlace; usa protocolo\_ip; parte\_de eng\_telecomunicacoes}
\prov{proveniência: wiki \textperiodcentered{} fato \textperiodcentered{} confiança 1.0 \textperiodcentered{} domínio engenharia \textperiodcentered{} história 4 versões no jornal}

%CRISTAL {"arestas": [{"alvo": "modelo_osi", "confianca": 1.0, "epistemico": "fato", "origem": "wiki", "rel": "analogia"}, {"alvo": "bai_psi_colapso", "confianca": 1.0, "epistemico": "fato", "origem": "wiki", "rel": "explica"}, {"alvo": "websocket", "confianca": 1.0, "epistemico": "fato", "origem": "wiki", "rel": "usa"}, {"alvo": "capacidade_de_canal", "confianca": 1.0, "epistemico": "fato", "origem": "wiki", "rel": "analogia"}], "confianca": 1.0, "contraexemplos": [], "descricao": "A conversa da Tiffany é, ela mesma, um protocolo de camada de aplicação: pergunta→COLAPSO(Ψ) sobre o grafo→resposta, com fail-closed (recusa se c²<0.5, como um checksum que não fecha). Roda sobre a pilha (HTTP/WebSocket no cockpit) mas define sua própria semântica de requisição-resposta epistêmica.", "epistemico": "inferencia", "exemplos": [], "id": "protocolo_colapso_psi", "memoria": "semantica", "meta": {"dominio": "engenharia", "fonte": "conversa/colapso.py; conversa/parabola_algebra.py"}, "origem": "wiki", "palavras_chave": [], "sinonimos": [], "tipo": "conceito", "titulo": "Colapso Ψ como Protocolo (nosso)"}
\section{Colapso \ensuremath{\Psi} como Protocolo (nosso)}
A conversa da Tiffany é, ela mesma, um protocolo de camada de aplicação: pergunta\ensuremath{\to}COLAPSO(\ensuremath{\Psi}) sobre o grafo\ensuremath{\to}resposta, com fail-closed (recusa se c\ensuremath{^{2}}<0.5, como um checksum que não fecha). Roda sobre a pilha (HTTP/WebSocket no cockpit) mas define sua própria semântica de requisição-resposta epistêmica.
\prov{relações: analogia modelo\_osi; explica bai\_psi\_colapso; usa websocket; analogia capacidade\_de\_canal}
\prov{proveniência: wiki \textperiodcentered{} conversa/colapso.py; conversa/parabola\_algebra.py \textperiodcentered{} inferencia \textperiodcentered{} confiança 1.0 \textperiodcentered{} domínio engenharia \textperiodcentered{} história 5 versões no jornal}

%CRISTAL {"arestas": [{"alvo": "protocolo_ip", "confianca": 1.0, "epistemico": "fato", "origem": "wiki", "rel": "usa"}, {"alvo": "eng_telecomunicacoes", "confianca": 1.0, "epistemico": "fato", "origem": "curriculo", "rel": "parte_de"}], "confianca": 1.0, "contraexemplos": [], "descricao": "O canal de controle e diagnóstico do IP: mensagens de erro, inalcançável, e o eco do ping/traceroute. A rede reportando sobre si mesma — não carrega dados de usuário.", "epistemico": "fato", "exemplos": [], "id": "protocolo_icmp", "memoria": "semantica", "meta": {"dominio": "engenharia", "fonte": ""}, "origem": "wiki", "palavras_chave": [], "sinonimos": [], "tipo": "conceito", "titulo": "ICMP"}
\section{ICMP}
O canal de controle e diagnóstico do IP: mensagens de erro, inalcançável, e o eco do ping/traceroute. A rede reportando sobre si mesma — não carrega dados de usuário.
\prov{relações: usa protocolo\_ip; parte\_de eng\_telecomunicacoes}
\prov{proveniência: wiki \textperiodcentered{} fato \textperiodcentered{} confiança 1.0 \textperiodcentered{} domínio engenharia \textperiodcentered{} história 3 versões no jornal}

%CRISTAL {"arestas": [{"alvo": "modelo_osi", "confianca": 1.0, "epistemico": "fato", "origem": "wiki", "rel": "especializa"}, {"alvo": "protocolo_tcp", "confianca": 1.0, "epistemico": "fato", "origem": "wiki", "rel": "usa"}, {"alvo": "prova_de_trabalho", "confianca": 1.0, "epistemico": "fato", "origem": "wiki", "rel": "usa"}, {"alvo": "symbolic_beta_pipe_so", "confianca": 1.0, "epistemico": "fato", "origem": "wiki", "rel": "explica"}, {"alvo": "goldchain_como_rede", "confianca": 1.0, "epistemico": "fato", "origem": "wiki", "rel": "parte_de"}], "confianca": 1.0, "contraexemplos": [], "descricao": "O protocolo de aplicação que distribui TRABALHO de mineração entre nós: o servidor envia o cabeçalho e a meta, o minerador devolve a prova. No broca-so, o broca-stratum do gex44 fala este protocolo minerando os 12 metais ao vivo. É um protocolo de camada de aplicação, como o HTTP, mas para prova-de-trabalho.", "epistemico": "fato", "exemplos": [], "id": "protocolo_stratum", "memoria": "semantica", "meta": {"dominio": "engenharia", "fonte": "scripts/gex44/run_symbolic.sh; goldchain/symbolic.py"}, "origem": "wiki", "palavras_chave": [], "sinonimos": [], "tipo": "conceito", "titulo": "Protocolo Stratum (nosso, mineração)"}
\section{Protocolo Stratum (nosso, mineração)}
O protocolo de aplicação que distribui TRABALHO de mineração entre nós: o servidor envia o cabeçalho e a meta, o minerador devolve a prova. No broca-so, o broca-stratum do gex44 fala este protocolo minerando os 12 metais ao vivo. É um protocolo de camada de aplicação, como o HTTP, mas para prova-de-trabalho.
\prov{relações: especializa modelo\_osi; usa protocolo\_tcp; usa prova\_de\_trabalho; explica symbolic\_beta\_pipe\_so; parte\_de goldchain\_como\_rede}
\prov{proveniência: wiki \textperiodcentered{} scripts/gex44/run\_symbolic.sh; goldchain/symbolic.py \textperiodcentered{} fato \textperiodcentered{} confiança 1.0 \textperiodcentered{} domínio engenharia \textperiodcentered{} história 6 versões no jornal}

%CRISTAL {"arestas": [{"alvo": "modulacao_digital", "confianca": 1.0, "epistemico": "fato", "origem": "wiki", "rel": "especializa"}, {"alvo": "capacidade_de_canal", "confianca": 1.0, "epistemico": "fato", "origem": "wiki", "rel": "usa"}, {"alvo": "eng_telecomunicacoes", "confianca": 1.0, "epistemico": "fato", "origem": "curriculo", "rel": "parte_de"}, {"alvo": "modulacao", "confianca": 1.0, "epistemico": "fato", "origem": "wiki", "rel": "eh_um"}], "confianca": 1.0, "contraexemplos": [], "descricao": "Quadrature Amplitude Modulation: cada símbolo é um PONTO no plano complexo (amplitude+fase). M-QAM tem M pontos e carrega log₂(M) bits/símbolo. Mais pontos = mais bits, mas a distância mínima entre eles encolhe — a constelação sobe à borda do ruído.", "epistemico": "fato", "exemplos": [], "id": "qam", "memoria": "semantica", "meta": {"dominio": "engenharia", "fonte": "conectividade e 5G"}, "origem": "wiki", "palavras_chave": [], "sinonimos": [], "tipo": "conceito", "titulo": "QAM (constelação de símbolos)"}
\section{QAM (constelação de símbolos)}
Quadrature Amplitude Modulation: cada símbolo é um PONTO no plano complexo (amplitude+fase). M-QAM tem M pontos e carrega log\ensuremath{_{2}}(M) bits/símbolo. Mais pontos = mais bits, mas a distância mínima entre eles encolhe — a constelação sobe à borda do ruído.
\prov{relações: especializa modulacao\_digital; usa capacidade\_de\_canal; parte\_de eng\_telecomunicacoes; eh\_um modulacao}
\prov{proveniência: wiki \textperiodcentered{} conectividade e 5G \textperiodcentered{} fato \textperiodcentered{} confiança 1.0 \textperiodcentered{} domínio engenharia \textperiodcentered{} história 6 versões no jornal}

%CRISTAL {"arestas": [{"alvo": "funcoes_de_rede_5gc", "confianca": 1.0, "epistemico": "fato", "origem": "wiki", "rel": "usa"}], "confianca": 1.0, "contraexemplos": [], "descricao": "No 5G a QoS é por FLUXO, identificado pelo 5QI, que fixa prioridade, atraso e perda toleráveis. O SDAP o mapeia no rádio e o UPF o aplica no núcleo — QoS fim-a-fim granular, base do slicing e do URLLC.", "epistemico": "fato", "exemplos": [], "id": "qos_fluxos_5g", "memoria": "semantica", "meta": {"dominio": "engenharia", "fonte": ""}, "origem": "wiki", "palavras_chave": [], "sinonimos": [], "tipo": "conceito", "titulo": "QoS por Fluxos (5QI)"}
\section{QoS por Fluxos (5QI)}
No 5G a QoS é por FLUXO, identificado pelo 5QI, que fixa prioridade, atraso e perda toleráveis. O SDAP o mapeia no rádio e o UPF o aplica no núcleo — QoS fim-a-fim granular, base do slicing e do URLLC.
\prov{relações: usa funcoes\_de\_rede\_5gc}
\prov{proveniência: wiki \textperiodcentered{} fato \textperiodcentered{} confiança 1.0 \textperiodcentered{} domínio engenharia \textperiodcentered{} história 2 versões no jornal}

%CRISTAL {"arestas": [{"alvo": "eletronica_de_potencia", "confianca": 1.0, "epistemico": "fato", "origem": "wiki", "rel": "causa"}, {"alvo": "sistema_eletrico_potencia", "confianca": 1.0, "epistemico": "fato", "origem": "wiki", "rel": "parte_de"}, {"alvo": "eng_sistemas_energia", "confianca": 1.0, "epistemico": "fato", "origem": "curriculo", "rel": "parte_de"}], "confianca": 1.0, "contraexemplos": [], "descricao": "Manter a onda limpa: harmônicos, afundamentos, flutuações e desequilíbrio que a eletrônica moderna injeta e sofre. O preço oculto de uma rede cheia de conversores chaveados.", "epistemico": "inferencia", "exemplos": [], "id": "qualidade_de_energia", "memoria": "semantica", "meta": {"dominio": "engenharia", "fonte": ""}, "origem": "wiki", "palavras_chave": [], "sinonimos": [], "tipo": "conceito", "titulo": "Qualidade de Energia"}
\section{Qualidade de Energia}
Manter a onda limpa: harmônicos, afundamentos, flutuações e desequilíbrio que a eletrônica moderna injeta e sofre. O preço oculto de uma rede cheia de conversores chaveados.
\prov{relações: causa eletronica\_de\_potencia; parte\_de sistema\_eletrico\_potencia; parte\_de eng\_sistemas\_energia}
\prov{proveniência: wiki \textperiodcentered{} inferencia \textperiodcentered{} confiança 1.0 \textperiodcentered{} domínio engenharia \textperiodcentered{} história 4 versões no jornal}

%CRISTAL {"arestas": [{"alvo": "biblioteca_formas_minerais", "confianca": 1.0, "epistemico": "fato", "origem": "wiki", "rel": "parte_de"}, {"alvo": "numero_de_salem", "confianca": 1.0, "epistemico": "fato", "origem": "wiki", "rel": "eh_um"}, {"alvo": "numeros_de_pisot", "confianca": 1.0, "epistemico": "fato", "origem": "wiki", "rel": "relaciona"}, {"alvo": "transformada_de_fourier", "confianca": 1.0, "epistemico": "fato", "origem": "wiki", "rel": "usa"}, {"alvo": "imagem_em_rotacoes_fourier", "confianca": 1.0, "epistemico": "fato", "origem": "wiki", "rel": "usa"}], "confianca": 1.0, "contraexemplos": [], "descricao": "Somar rotações puras giradas de 2π/dobras produz uma interferência quasiperiódica; 5 dobras dá simetria pentagonal (o padrão de Penrose), proibida em cristais periódicos. É a forma-Salem: muitas rotações ρ=1 superpostas — a mesma estrutura que liga quasicristais aos números de Pisot.", "epistemico": "inferencia", "exemplos": [], "id": "quasicristal_de_rotacoes", "memoria": "semantica", "meta": {"dominio": "imagem", "fonte": "formas minerais"}, "origem": "wiki", "palavras_chave": [], "sinonimos": [], "tipo": "conceito", "titulo": "Quasicristal de Rotações (soma de Salem)"}
\section{Quasicristal de Rotações (soma de Salem)}
Somar rotações puras giradas de 2\ensuremath{\pi}/dobras produz uma interferência quasiperiódica; 5 dobras dá simetria pentagonal (o padrão de Penrose), proibida em cristais periódicos. É a forma-Salem: muitas rotações \ensuremath{\rho}=1 superpostas — a mesma estrutura que liga quasicristais aos números de Pisot.
\prov{relações: parte\_de biblioteca\_formas\_minerais; eh\_um numero\_de\_salem; relaciona numeros\_de\_pisot; usa transformada\_de\_fourier; usa imagem\_em\_rotacoes\_fourier}
\prov{proveniência: wiki \textperiodcentered{} formas minerais \textperiodcentered{} inferencia \textperiodcentered{} confiança 1.0 \textperiodcentered{} domínio imagem \textperiodcentered{} história 6 versões no jornal}

%CRISTAL {"arestas": [{"alvo": "protocolo_udp", "confianca": 1.0, "epistemico": "fato", "origem": "wiki", "rel": "usa"}, {"alvo": "tls", "confianca": 1.0, "epistemico": "fato", "origem": "wiki", "rel": "usa"}, {"alvo": "protocolo_tcp", "confianca": 1.0, "epistemico": "fato", "origem": "wiki", "rel": "contradiz"}, {"alvo": "eng_telecomunicacoes", "confianca": 1.0, "epistemico": "fato", "origem": "curriculo", "rel": "parte_de"}], "confianca": 1.0, "contraexemplos": [], "descricao": "Transporte moderno sobre UDP com cifra (TLS 1.3) embutida, streams sem bloqueio de cabeça de fila e conexão em 0-1 RTT — a base do HTTP/3. O TCP reinventado no espaço do usuário.", "epistemico": "fato", "exemplos": [], "id": "quic", "memoria": "semantica", "meta": {"autor": "Google", "dominio": "engenharia", "fonte": ""}, "origem": "wiki", "palavras_chave": [], "sinonimos": [], "tipo": "conceito", "titulo": "QUIC"}
\section{QUIC}
Transporte moderno sobre UDP com cifra (TLS 1.3) embutida, streams sem bloqueio de cabeça de fila e conexão em 0-1 RTT — a base do HTTP/3. O TCP reinventado no espaço do usuário.
\prov{relações: usa protocolo\_udp; usa tls; contradiz protocolo\_tcp; parte\_de eng\_telecomunicacoes}
\prov{proveniência: wiki \textperiodcentered{} fato \textperiodcentered{} confiança 1.0 \textperiodcentered{} domínio engenharia \textperiodcentered{} história 5 versões no jornal}

%CRISTAL {"arestas": [{"alvo": "acesso_aleatorio_nr", "confianca": 1.0, "epistemico": "fato", "origem": "wiki", "rel": "especializa"}], "confianca": 1.0, "contraexemplos": [], "descricao": "Novidade do NR: MsgA junta preâmbulo E dados numa só transmissão; MsgB responde com resolução e concessão. Corta a latência do acesso pela metade — importante no URLLC e na banda não-licenciada.", "epistemico": "fato", "exemplos": [], "id": "rach_2_passos", "memoria": "semantica", "meta": {"dominio": "engenharia", "fonte": ""}, "origem": "wiki", "palavras_chave": [], "sinonimos": [], "tipo": "conceito", "titulo": "RACH de 2 Passos (5G)"}
\section{RACH de 2 Passos (5G)}
Novidade do NR: MsgA junta preâmbulo E dados numa só transmissão; MsgB responde com resolução e concessão. Corta a latência do acesso pela metade — importante no URLLC e na banda não-licenciada.
\prov{relações: especializa acesso\_aleatorio\_nr}
\prov{proveniência: wiki \textperiodcentered{} fato \textperiodcentered{} confiança 1.0 \textperiodcentered{} domínio engenharia \textperiodcentered{} história 2 versões no jornal}

%CRISTAL {"arestas": [{"alvo": "acesso_aleatorio_nr", "confianca": 1.0, "epistemico": "fato", "origem": "wiki", "rel": "especializa"}], "confianca": 1.0, "contraexemplos": [], "descricao": "Msg1: UE envia preâmbulo no PRACH → Msg2: gNB responde com a RAR (concede recurso e adianta o tempo) → Msg3: UE se identifica no recurso concedido → Msg4: gNB resolve a contenção (confirma qual UE venceu). O aperto de mão de quem chega sem ser esperado.", "epistemico": "fato", "exemplos": [], "id": "rach_4_passos", "memoria": "semantica", "meta": {"dominio": "engenharia", "fonte": ""}, "origem": "wiki", "palavras_chave": [], "sinonimos": [], "tipo": "conceito", "titulo": "RACH de 4 Passos (com contenção)"}
\section{RACH de 4 Passos (com contenção)}
Msg1: UE envia preâmbulo no PRACH \ensuremath{\to} Msg2: gNB responde com a RAR (concede recurso e adianta o tempo) \ensuremath{\to} Msg3: UE se identifica no recurso concedido \ensuremath{\to} Msg4: gNB resolve a contenção (confirma qual UE venceu). O aperto de mão de quem chega sem ser esperado.
\prov{relações: especializa acesso\_aleatorio\_nr}
\prov{proveniência: wiki \textperiodcentered{} fato \textperiodcentered{} confiança 1.0 \textperiodcentered{} domínio engenharia \textperiodcentered{} história 2 versões no jornal}

%CRISTAL {"arestas": [{"alvo": "acesso_aleatorio_nr", "confianca": 1.0, "epistemico": "fato", "origem": "wiki", "rel": "especializa"}], "confianca": 1.0, "contraexemplos": [], "descricao": "Quando a rede JÁ sabe que o UE vai chegar (handover, recuperação de feixe), ela reserva um preâmbulo dedicado — sem disputa, sem Msg4. Acesso rápido e determinístico, usado no handover.", "epistemico": "fato", "exemplos": [], "id": "rach_sem_contencao", "memoria": "semantica", "meta": {"dominio": "engenharia", "fonte": ""}, "origem": "wiki", "palavras_chave": [], "sinonimos": [], "tipo": "conceito", "titulo": "RACH sem Contenção"}
\section{RACH sem Contenção}
Quando a rede JÁ sabe que o UE vai chegar (handover, recuperação de feixe), ela reserva um preâmbulo dedicado — sem disputa, sem Msg4. Acesso rápido e determinístico, usado no handover.
\prov{relações: especializa acesso\_aleatorio\_nr}
\prov{proveniência: wiki \textperiodcentered{} fato \textperiodcentered{} confiança 1.0 \textperiodcentered{} domínio engenharia \textperiodcentered{} história 2 versões no jornal}

%CRISTAL {"arestas": [{"alvo": "ruido_eletronico", "confianca": 1.0, "epistemico": "fato", "origem": "wiki", "rel": "usa"}, {"alvo": "eng_telecomunicacoes", "confianca": 1.0, "epistemico": "fato", "origem": "curriculo", "rel": "parte_de"}], "confianca": 1.0, "contraexemplos": [], "descricao": "Quanto o sinal se ergue acima do ruído — o parâmetro que decide taxa e taxa de erro. Eb/N0 normaliza por bit; a curva BER×Eb/N0 é o boletim de desempenho de todo esquema digital.", "epistemico": "fato", "exemplos": [], "id": "razao_sinal_ruido", "memoria": "semantica", "meta": {"dominio": "engenharia", "fonte": ""}, "origem": "wiki", "palavras_chave": [], "sinonimos": [], "tipo": "conceito", "titulo": "Relação Sinal-Ruído (SNR) e Eb/N0"}
\section{Relação Sinal-Ruído (SNR) e Eb/N0}
Quanto o sinal se ergue acima do ruído — o parâmetro que decide taxa e taxa de erro. Eb/N0 normaliza por bit; a curva BER$\times$Eb/N0 é o boletim de desempenho de todo esquema digital.
\prov{relações: usa ruido\_eletronico; parte\_de eng\_telecomunicacoes}
\prov{proveniência: wiki \textperiodcentered{} fato \textperiodcentered{} confiança 1.0 \textperiodcentered{} domínio engenharia \textperiodcentered{} história 3 versões no jornal}

%CRISTAL {"arestas": [{"alvo": "sistemas_de_controle", "confianca": 1.0, "epistemico": "fato", "origem": "wiki", "rel": "relaciona"}, {"alvo": "eng_eletronica", "confianca": 1.0, "epistemico": "fato", "origem": "curriculo", "rel": "parte_de"}], "confianca": 1.0, "contraexemplos": [], "descricao": "Devolver parte da saída à entrada — negativa estabiliza ganho e linearidade; positiva oscila. O princípio que une o amp-op, o oscilador e o controle. Barkhausen dita a oscilação.", "epistemico": "fato", "exemplos": [], "id": "realimentacao_eletronica", "memoria": "semantica", "meta": {"autor": "Black", "dominio": "engenharia", "fonte": ""}, "origem": "wiki", "palavras_chave": [], "sinonimos": [], "tipo": "conceito", "titulo": "Realimentação (Feedback)"}
\section{Realimentação (Feedback)}
Devolver parte da saída à entrada — negativa estabiliza ganho e linearidade; positiva oscila. O princípio que une o amp-op, o oscilador e o controle. Barkhausen dita a oscilação.
\prov{relações: relaciona sistemas\_de\_controle; parte\_de eng\_eletronica}
\prov{proveniência: wiki \textperiodcentered{} fato \textperiodcentered{} confiança 1.0 \textperiodcentered{} domínio engenharia \textperiodcentered{} história 3 versões no jornal}

%CRISTAL {"arestas": [{"alvo": "feedback_retroalimentacao", "confianca": 1.0, "epistemico": "fato", "origem": "wiki", "rel": "explica"}, {"alvo": "realimentacao_eletronica", "confianca": 1.0, "epistemico": "fato", "origem": "wiki", "rel": "relaciona"}, {"alvo": "lyapunov_da_orbita", "confianca": 1.0, "epistemico": "fato", "origem": "wiki", "rel": "eh_um"}, {"alvo": "estabilidade_da_rede", "confianca": 1.0, "epistemico": "fato", "origem": "wiki", "rel": "relaciona"}, {"alvo": "corpo_mineral", "confianca": 1.0, "epistemico": "fato", "origem": "wiki", "rel": "usa"}], "confianca": 1.0, "contraexemplos": [], "descricao": "Um controlador leva o erro a zero (o atrator). O ganho K é o Lyapunov: pequeno converge devagar (cristal), alto demais oscila e diverge (caos), no meio a borda rápida. É exatamente o λ da estabilidade da rede elétrica e do treino da IA, agora no corpo do robô. Verif.: K=0.3 λ=−0.36 (cristal), K=2.3 λ=+0.26 (caos, treme e cai).", "epistemico": "inferencia", "exemplos": [], "id": "realimentacao_lyapunov", "memoria": "semantica", "meta": {"dominio": "robotica", "fonte": "robótica e sistemas autônomos", "lambda_K03": -0.357, "lambda_K23": 0.262}, "origem": "wiki", "palavras_chave": [], "sinonimos": [], "tipo": "conceito", "titulo": "Realimentação = o Lyapunov da Estabilidade"}
\section{Realimentação = o Lyapunov da Estabilidade}
Um controlador leva o erro a zero (o atrator). O ganho K é o Lyapunov: pequeno converge devagar (cristal), alto demais oscila e diverge (caos), no meio a borda rápida. É exatamente o \ensuremath{\lambda} da estabilidade da rede elétrica e do treino da IA, agora no corpo do robô. Verif.: K=0.3 \ensuremath{\lambda}=\ensuremath{-}0.36 (cristal), K=2.3 \ensuremath{\lambda}=+0.26 (caos, treme e cai).
\prov{relações: explica feedback\_retroalimentacao; relaciona realimentacao\_eletronica; eh\_um lyapunov\_da\_orbita; relaciona estabilidade\_da\_rede; usa corpo\_mineral}
\prov{proveniência: wiki \textperiodcentered{} robótica e sistemas autônomos \textperiodcentered{} inferencia \textperiodcentered{} confiança 1.0 \textperiodcentered{} domínio robotica \textperiodcentered{} história 6 versões no jornal}

%CRISTAL {"arestas": [{"alvo": "bateria_de_tracao", "confianca": 1.0, "epistemico": "fato", "origem": "wiki", "rel": "usa"}, {"alvo": "distribuicao_de_energia", "confianca": 1.0, "epistemico": "fato", "origem": "wiki", "rel": "relaciona"}], "confianca": 1.0, "contraexemplos": [], "descricao": "Reabastecer o cristal: CA lenta (casa/destino) ou CC rápida (rota, centenas de kW). O SoC sobe — a órbita de carga do reservatório, agora num carro. A infraestrutura de recarga é a distribuição estendida aos nós móveis.", "epistemico": "fato", "exemplos": [], "id": "recarga_veicular", "memoria": "semantica", "meta": {"dominio": "engenharia", "fonte": "mobilidade elétrica"}, "origem": "wiki", "palavras_chave": [], "sinonimos": [], "tipo": "conceito", "titulo": "Recarga Veicular"}
\section{Recarga Veicular}
Reabastecer o cristal: CA lenta (casa/destino) ou CC rápida (rota, centenas de kW). O SoC sobe — a órbita de carga do reservatório, agora num carro. A infraestrutura de recarga é a distribuição estendida aos nós móveis.
\prov{relações: usa bateria\_de\_tracao; relaciona distribuicao\_de\_energia}
\prov{proveniência: wiki \textperiodcentered{} mobilidade elétrica \textperiodcentered{} fato \textperiodcentered{} confiança 1.0 \textperiodcentered{} domínio engenharia \textperiodcentered{} história 3 versões no jornal}

%CRISTAL {"arestas": [{"alvo": "colapso_fatorado_glifo", "confianca": 1.0, "epistemico": "fato", "origem": "wiki", "rel": "parte_de"}, {"alvo": "reconhecer_caractere_forma", "confianca": 1.0, "epistemico": "fato", "origem": "wiki", "rel": "usa"}, {"alvo": "kmeans_atratores", "confianca": 1.0, "epistemico": "fato", "origem": "wiki", "rel": "usa"}], "confianca": 1.0, "contraexemplos": [], "descricao": "O estilo colapsa em DOIS ESTÁGIOS: sabido o caractere, compara-se a APARÊNCIA da consulta com as seis versões daquela MESMA letra (a mesma letra em fontes diferentes é bem distinta) e desaba na mais próxima. Condicionar ao conteúdo torna o sinal de estilo forte — ~97%. Densidade lê o negrito, inclinação o itálico, a aparência da letra separa serif/sans/mono/casual.", "epistemico": "inferencia", "exemplos": [], "id": "reconhecer_caligrafia_estilo", "memoria": "semantica", "meta": {"dominio": "imagem", "fonte": "glifos minerais"}, "origem": "wiki", "palavras_chave": [], "sinonimos": [], "tipo": "conceito", "titulo": "Reconhecer a Caligrafia (colapso do estilo)"}
\section{Reconhecer a Caligrafia (colapso do estilo)}
O estilo colapsa em DOIS ESTÁGIOS: sabido o caractere, compara-se a APARÊNCIA da consulta com as seis versões daquela MESMA letra (a mesma letra em fontes diferentes é bem distinta) e desaba na mais próxima. Condicionar ao conteúdo torna o sinal de estilo forte — \~{}97\%. Densidade lê o negrito, inclinação o itálico, a aparência da letra separa serif/sans/mono/casual.
\prov{relações: parte\_de colapso\_fatorado\_glifo; usa reconhecer\_caractere\_forma; usa kmeans\_atratores}
\prov{proveniência: wiki \textperiodcentered{} glifos minerais \textperiodcentered{} inferencia \textperiodcentered{} confiança 1.0 \textperiodcentered{} domínio imagem \textperiodcentered{} história 4 versões no jornal}

%CRISTAL {"arestas": [{"alvo": "colapso_fatorado_glifo", "confianca": 1.0, "epistemico": "fato", "origem": "wiki", "rel": "parte_de"}, {"alvo": "kmeans_atratores", "confianca": 1.0, "epistemico": "fato", "origem": "wiki", "rel": "usa"}, {"alvo": "segmentacao_de_imagem", "confianca": 1.0, "epistemico": "fato", "origem": "wiki", "rel": "relaciona"}], "confianca": 1.0, "contraexemplos": [], "descricao": "O conteúdo colapsa pela FORMA feita invariante ao estilo: desinclina (tira o itálico), normaliza à caixa (tira escala/posição) e amostra numa grade grossa (alisa o traço). Assim 'A' serif e 'A' sans caem no mesmo cristal. É o colapso ao caractere mais próximo — ~96% em consultas perturbadas.", "epistemico": "inferencia", "exemplos": [], "id": "reconhecer_caractere_forma", "memoria": "semantica", "meta": {"dominio": "imagem", "fonte": "glifos minerais"}, "origem": "wiki", "palavras_chave": [], "sinonimos": [], "tipo": "conceito", "titulo": "Reconhecer o Caractere (colapso da forma)"}
\section{Reconhecer o Caractere (colapso da forma)}
O conteúdo colapsa pela FORMA feita invariante ao estilo: desinclina (tira o itálico), normaliza à caixa (tira escala/posição) e amostra numa grade grossa (alisa o traço). Assim 'A' serif e 'A' sans caem no mesmo cristal. É o colapso ao caractere mais próximo — \~{}96\% em consultas perturbadas.
\prov{relações: parte\_de colapso\_fatorado\_glifo; usa kmeans\_atratores; relaciona segmentacao\_de\_imagem}
\prov{proveniência: wiki \textperiodcentered{} glifos minerais \textperiodcentered{} inferencia \textperiodcentered{} confiança 1.0 \textperiodcentered{} domínio imagem \textperiodcentered{} história 4 versões no jornal}

%CRISTAL {"arestas": [{"alvo": "imagem_como_orbita_gato", "confianca": 1.0, "epistemico": "fato", "origem": "wiki", "rel": "explica"}, {"alvo": "experimento_2_recorrencias_e_metais", "confianca": 1.0, "epistemico": "fato", "origem": "wiki", "rel": "relaciona"}, {"alvo": "neuronio_e_toro_124", "confianca": 1.0, "epistemico": "fato", "origem": "wiki", "rel": "relaciona"}, {"alvo": "conjectura_de_poincare", "confianca": 1.0, "epistemico": "fato", "origem": "wiki", "rel": "relaciona"}, {"alvo": "lyapunov_da_orbita", "confianca": 1.0, "epistemico": "fato", "origem": "wiki", "rel": "usa"}], "confianca": 1.0, "contraexemplos": [], "descricao": "Num espaço FINITO, uma bijeção determinística sempre recorre: depois de T passos a imagem RETORNA idêntica (A_n^T ≡ I mod N). T é a ordem do gato em GL(2, Z_N) — depende só de N e n, não do conteúdo. O caos (λ>0) que mistura é o mesmo que, num toro discreto, garante o retorno exato. Verif.: N=128, A_1 → T=192; A^192(img)==img.", "epistemico": "inferencia", "exemplos": [], "id": "recorrencia_de_poincare_imagem", "memoria": "semantica", "meta": {"T_N128_n1": 192, "dominio": "imagem", "fonte": "órbitas de imagem", "retorna_exata": true}, "origem": "wiki", "palavras_chave": [], "sinonimos": [], "tipo": "conceito", "titulo": "Recorrência de Poincaré (o caos que volta)"}
\section{Recorrência de Poincaré (o caos que volta)}
Num espaço FINITO, uma bijeção determinística sempre recorre: depois de T passos a imagem RETORNA idêntica (A\_n\^{}T \ensuremath{\equiv} I mod N). T é a ordem do gato em GL(2, Z\_N) — depende só de N e n, não do conteúdo. O caos (\ensuremath{\lambda}>0) que mistura é o mesmo que, num toro discreto, garante o retorno exato. Verif.: N=128, A\_1 \ensuremath{\to} T=192; A\^{}192(img)==img.
\prov{relações: explica imagem\_como\_orbita\_gato; relaciona experimento\_2\_recorrencias\_e\_metais; relaciona neuronio\_e\_toro\_124; relaciona conjectura\_de\_poincare; usa lyapunov\_da\_orbita}
\prov{proveniência: wiki \textperiodcentered{} órbitas de imagem \textperiodcentered{} inferencia \textperiodcentered{} confiança 1.0 \textperiodcentered{} domínio imagem \textperiodcentered{} história 6 versões no jornal}

%CRISTAL {"arestas": [{"alvo": "conectividade", "confianca": 1.0, "epistemico": "fato", "origem": "wiki", "rel": "eh_um"}, {"alvo": "internet_das_coisas", "confianca": 1.0, "epistemico": "fato", "origem": "wiki", "rel": "relaciona"}, {"alvo": "smart_city", "confianca": 1.0, "epistemico": "fato", "origem": "wiki", "rel": "relaciona"}], "confianca": 1.0, "contraexemplos": [], "descricao": "A quinta geração de rede móvel: banda larga (até Gbps), latência baixa (~1 ms) e densidade enorme de dispositivos. Feita para a IoT e o tempo real — ampliando a banda (mmWave) e a eficiência espectral (MIMO massivo, beamforming).", "epistemico": "fato", "exemplos": [], "id": "rede_5g", "memoria": "semantica", "meta": {"dominio": "engenharia", "fonte": "conectividade e 5G"}, "origem": "wiki", "palavras_chave": [], "sinonimos": [], "tipo": "conceito", "titulo": "Rede 5G"}
\section{Rede 5G}
A quinta geração de rede móvel: banda larga (até Gbps), latência baixa (\~{}1 ms) e densidade enorme de dispositivos. Feita para a IoT e o tempo real — ampliando a banda (mmWave) e a eficiência espectral (MIMO massivo, beamforming).
\prov{relações: eh\_um conectividade; relaciona internet\_das\_coisas; relaciona smart\_city}
\prov{proveniência: wiki \textperiodcentered{} conectividade e 5G \textperiodcentered{} fato \textperiodcentered{} confiança 1.0 \textperiodcentered{} domínio engenharia \textperiodcentered{} história 4 versões no jornal}

%CRISTAL {"arestas": [{"alvo": "sistema_e_lei_nao_dono", "confianca": 1.0, "epistemico": "fato", "origem": "wiki", "rel": "usa"}, {"alvo": "organismo_gex44", "confianca": 1.0, "epistemico": "fato", "origem": "wiki", "rel": "explica"}, {"alvo": "broca_os", "confianca": 1.0, "epistemico": "fato", "origem": "wiki", "rel": "parte_de"}, {"alvo": "bai_tese_mae", "confianca": 1.0, "epistemico": "fato", "origem": "wiki", "rel": "usa"}], "confianca": 1.0, "contraexemplos": [], "descricao": "A visão consolidada: um organismo distribuído que se auto-gerencia sem autoridade central — 'é a lei, não o dono'. Cada software é uma célula que fala por canais numa banda; comunidades emergem de pares que compartilham um tecido; a ordem vem do recall, não de um scheduler. A comunicação viva já é fato; a rede autônoma é o design que a fecha num só corpo.", "epistemico": "inferencia", "exemplos": [], "id": "rede_autonoma_broca", "memoria": "semantica", "meta": {"dominio": "arquitetura", "fonte": "semilla_rede_autonoma"}, "origem": "wiki", "palavras_chave": [], "sinonimos": [], "tipo": "conceito", "titulo": "A Rede Autônoma do broca-so (organismo sem dono)"}
\section{A Rede Autônoma do broca-so (organismo sem dono)}
A visão consolidada: um organismo distribuído que se auto-gerencia sem autoridade central — 'é a lei, não o dono'. Cada software é uma célula que fala por canais numa banda; comunidades emergem de pares que compartilham um tecido; a ordem vem do recall, não de um scheduler. A comunicação viva já é fato; a rede autônoma é o design que a fecha num só corpo.
\prov{relações: usa sistema\_e\_lei\_nao\_dono; explica organismo\_gex44; parte\_de broca\_os; usa bai\_tese\_mae}
\prov{proveniência: wiki \textperiodcentered{} semilla\_rede\_autonoma \textperiodcentered{} inferencia \textperiodcentered{} confiança 1.0 \textperiodcentered{} domínio arquitetura \textperiodcentered{} história 5 versões no jornal}

%CRISTAL {"arestas": [{"alvo": "cristalografia", "confianca": 1.0, "epistemico": "fato", "origem": "wiki", "rel": "parte_de"}, {"alvo": "reticulado", "confianca": 1.0, "epistemico": "fato", "origem": "wiki", "rel": "relaciona"}, {"alvo": "bases_pisot", "confianca": 1.0, "epistemico": "fato", "origem": "wiki", "rel": "relaciona"}], "confianca": 1.0, "contraexemplos": [], "descricao": "O padrão periódico de átomos num cristal — as 14 redes de Bravais; a repetição infinita de uma célula unitária. A ordem tornada matéria.", "epistemico": "fato", "exemplos": [], "id": "rede_cristalina", "memoria": "semantica", "meta": {"autor": "Bravais", "dominio": "engenharia", "fonte": ""}, "origem": "wiki", "palavras_chave": [], "sinonimos": [], "tipo": "conceito", "titulo": "Rede Cristalina"}
\section{Rede Cristalina}
O padrão periódico de átomos num cristal — as 14 redes de Bravais; a repetição infinita de uma célula unitária. A ordem tornada matéria.
\prov{relações: parte\_de cristalografia; relaciona reticulado; relaciona bases\_pisot}
\prov{proveniência: wiki \textperiodcentered{} fato \textperiodcentered{} confiança 1.0 \textperiodcentered{} domínio engenharia \textperiodcentered{} história 4 versões no jornal}

%CRISTAL {"arestas": [{"alvo": "carregador_publico", "confianca": 1.0, "epistemico": "fato", "origem": "wiki", "rel": "usa"}, {"alvo": "distribuicao_de_energia", "confianca": 1.0, "epistemico": "fato", "origem": "wiki", "rel": "relaciona"}, {"alvo": "recarga_veicular", "confianca": 1.0, "epistemico": "fato", "origem": "wiki", "rel": "usa"}], "confianca": 1.0, "contraexemplos": [], "descricao": "A malha de eletropostos (CA lento, CC rápido) que abastece a frota elétrica. É a distribuição estendida a nós MÓVEIS e estocásticos: a carga não fica mais parada num prédio, ela circula e pluga onde e quando quiser.", "epistemico": "fato", "exemplos": [], "id": "rede_de_carregamento", "memoria": "semantica", "meta": {"dominio": "engenharia", "fonte": "rede de carregamento e smart cities"}, "origem": "wiki", "palavras_chave": [], "sinonimos": [], "tipo": "conceito", "titulo": "Rede de Carregamento"}
\section{Rede de Carregamento}
A malha de eletropostos (CA lento, CC rápido) que abastece a frota elétrica. É a distribuição estendida a nós MÓVEIS e estocásticos: a carga não fica mais parada num prédio, ela circula e pluga onde e quando quiser.
\prov{relações: usa carregador\_publico; relaciona distribuicao\_de\_energia; usa recarga\_veicular}
\prov{proveniência: wiki \textperiodcentered{} rede de carregamento e smart cities \textperiodcentered{} fato \textperiodcentered{} confiança 1.0 \textperiodcentered{} domínio engenharia \textperiodcentered{} história 8 versões no jornal}

%CRISTAL {"arestas": [{"alvo": "internet_das_coisas", "confianca": 1.0, "epistemico": "fato", "origem": "wiki", "rel": "parte_de"}, {"alvo": "sensor", "confianca": 1.0, "epistemico": "fato", "origem": "wiki", "rel": "usa"}, {"alvo": "banco_de_dados_grafo", "confianca": 1.0, "epistemico": "fato", "origem": "wiki", "rel": "relaciona"}, {"alvo": "smart_city", "confianca": 1.0, "epistemico": "fato", "origem": "wiki", "rel": "relaciona"}], "confianca": 1.0, "contraexemplos": [], "descricao": "Muitos sensores conectados, cada um um colapso, alimentando um sistema comum: uma malha de olhos que é, ela mesma, um GRAFO — nós que medem, arestas que comunicam. A percepção distribuída da cidade, a mesma estrutura da broca-so.", "epistemico": "fato", "exemplos": [], "id": "rede_de_sensores", "memoria": "semantica", "meta": {"dominio": "engenharia", "fonte": "IoT e sensores"}, "origem": "wiki", "palavras_chave": [], "sinonimos": [], "tipo": "conceito", "titulo": "Rede de Sensores (o grafo de olhos)"}
\section{Rede de Sensores (o grafo de olhos)}
Muitos sensores conectados, cada um um colapso, alimentando um sistema comum: uma malha de olhos que é, ela mesma, um GRAFO — nós que medem, arestas que comunicam. A percepção distribuída da cidade, a mesma estrutura da broca-so.
\prov{relações: parte\_de internet\_das\_coisas; usa sensor; relaciona banco\_de\_dados\_grafo; relaciona smart\_city}
\prov{proveniência: wiki \textperiodcentered{} IoT e sensores \textperiodcentered{} fato \textperiodcentered{} confiança 1.0 \textperiodcentered{} domínio engenharia \textperiodcentered{} história 10 versões no jornal}

%CRISTAL {"arestas": [{"alvo": "geracao_de_energia", "confianca": 1.0, "epistemico": "fato", "origem": "wiki", "rel": "usa"}, {"alvo": "transmissao_de_energia", "confianca": 1.0, "epistemico": "fato", "origem": "wiki", "rel": "usa"}, {"alvo": "mao_invisivel", "confianca": 1.0, "epistemico": "fato", "origem": "wiki", "rel": "relaciona"}, {"alvo": "distribuicao_de_energia", "confianca": 1.0, "epistemico": "fato", "origem": "wiki", "rel": "relaciona"}, {"alvo": "energia_eletrica", "confianca": 1.0, "epistemico": "fato", "origem": "wiki", "rel": "usa"}], "confianca": 1.0, "contraexemplos": [], "descricao": "A malha que liga usinas a consumidores: transmissão em alta tensão (menos perda), transformadores que sobem/descem a tensão, e distribuição. A maior máquina já construída, sincronizada numa só rotação.", "epistemico": "fato", "exemplos": [], "id": "rede_eletrica", "memoria": "semantica", "meta": {"dominio": "engenharia", "fonte": "usinas e geração"}, "origem": "wiki", "palavras_chave": [], "sinonimos": [], "tipo": "conceito", "titulo": "Rede Elétrica (o grid)"}
\section{Rede Elétrica (o grid)}
A malha que liga usinas a consumidores: transmissão em alta tensão (menos perda), transformadores que sobem/descem a tensão, e distribuição. A maior máquina já construída, sincronizada numa só rotação.
\prov{relações: usa geracao\_de\_energia; usa transmissao\_de\_energia; relaciona mao\_invisivel; relaciona distribuicao\_de\_energia; usa energia\_eletrica}
\prov{proveniência: wiki \textperiodcentered{} usinas e geração \textperiodcentered{} fato \textperiodcentered{} confiança 1.0 \textperiodcentered{} domínio engenharia \textperiodcentered{} história 11 versões no jornal}

%CRISTAL {"fusao": [{"arestas": [{"alvo": "rede_eletrica", "confianca": 1.0, "epistemico": "fato", "origem": "wiki", "rel": "usa"}, {"alvo": "estabilidade_da_rede", "confianca": 1.0, "epistemico": "fato", "origem": "wiki", "rel": "usa"}], "confianca": 1.0, "contraexemplos": [], "descricao": "A rede com sensores, medição e controle que balanceia geração e carga em tempo real — integra renováveis intermitentes e armazenamento distribuído, e reage a faltas para não propagar o blackout.", "epistemico": "fato", "exemplos": [], "id": "rede_inteligente", "memoria": "semantica", "meta": {"dominio": "engenharia", "fonte": "distribuição e armazenamento"}, "origem": "wiki", "palavras_chave": [], "sinonimos": [], "tipo": "conceito", "titulo": "Rede Inteligente (smart grid)"}, {"arestas": [{"alvo": "sistema_eletrico_potencia", "confianca": 1.0, "epistemico": "fato", "origem": "wiki", "rel": "especializa"}, {"alvo": "rede_de_computadores", "confianca": 1.0, "epistemico": "fato", "origem": "wiki", "rel": "usa"}, {"alvo": "geracao_distribuida", "confianca": 1.0, "epistemico": "fato", "origem": "wiki", "rel": "usa"}, {"alvo": "eng_sistemas_energia", "confianca": 1.0, "epistemico": "fato", "origem": "curriculo", "rel": "parte_de"}], "confianca": 1.0, "contraexemplos": [], "descricao": "Instrumentar a rede com medição, comunicação e controle bidirecionais — para integrar renováveis intermitentes, resposta à demanda e geração distribuída. O SEP encontrando as telecomunicações.", "epistemico": "inferencia", "exemplos": [], "id": "smart_grid", "memoria": "semantica", "meta": {"dominio": "engenharia", "fonte": ""}, "origem": "wiki", "palavras_chave": [], "sinonimos": [], "tipo": "conceito", "titulo": "Rede Inteligente (Smart Grid)"}], "id": "rede_inteligente", "tipo": "conceito"}
\section{Rede Inteligente (smart grid)}
A rede com sensores, medição e controle que balanceia geração e carga em tempo real — integra renováveis intermitentes e armazenamento distribuído, e reage a faltas para não propagar o blackout.
\prov{relações: usa rede\_eletrica; usa estabilidade\_da\_rede}
\prov{proveniência: wiki \textperiodcentered{} distribuição e armazenamento \textperiodcentered{} fato \textperiodcentered{} confiança 1.0 \textperiodcentered{} domínio engenharia \textperiodcentered{} história 6 versões no jornal}
\prov{a segunda face --- smart\_grid:}
Instrumentar a rede com medição, comunicação e controle bidirecionais — para integrar renováveis intermitentes, resposta à demanda e geração distribuída. O SEP encontrando as telecomunicações.
\prov{fusão de curadoria (julgada): absorve smart\_grid --- as duas faces acima, intactas no registo}

%CRISTAL {"arestas": [{"alvo": "comunicacoes_moveis", "confianca": 1.0, "epistemico": "fato", "origem": "wiki", "rel": "usa"}, {"alvo": "eng_telecomunicacoes", "confianca": 1.0, "epistemico": "fato", "origem": "curriculo", "rel": "parte_de"}, {"alvo": "geracoes_celular", "confianca": 1.0, "epistemico": "fato", "origem": "wiki", "rel": "continua"}], "confianca": 1.0, "contraexemplos": [], "descricao": "A interface de rádio de 5ª geração (3GPP): banda flexível de sub-1GHz a ondas milimétricas, latência de milissegundos e densidade massiva. Não é 'internet mais rápida' — é uma plataforma reconfigurável para três usos radicalmente diferentes no mesmo espectro.", "epistemico": "fato", "exemplos": [], "id": "redes_5g_nr", "memoria": "semantica", "meta": {"autor": "3GPP", "dominio": "engenharia", "fonte": ""}, "origem": "wiki", "palavras_chave": [], "sinonimos": [], "tipo": "conceito", "titulo": "5G NR (New Radio)"}
\section{5G NR (New Radio)}
A interface de rádio de 5ª geração (3GPP): banda flexível de sub-1GHz a ondas milimétricas, latência de milissegundos e densidade massiva. Não é 'internet mais rápida' — é uma plataforma reconfigurável para três usos radicalmente diferentes no mesmo espectro.
\prov{relações: usa comunicacoes\_moveis; parte\_de eng\_telecomunicacoes; continua geracoes\_celular}
\prov{proveniência: wiki \textperiodcentered{} fato \textperiodcentered{} confiança 1.0 \textperiodcentered{} domínio engenharia \textperiodcentered{} história 5 versões no jornal}

%CRISTAL {"arestas": [{"alvo": "codificacao_de_canal", "confianca": 1.0, "epistemico": "fato", "origem": "wiki", "rel": "especializa"}, {"alvo": "cifra_metais_an", "confianca": 1.0, "epistemico": "fato", "origem": "wiki", "rel": "relaciona"}, {"alvo": "eng_telecomunicacoes", "confianca": 1.0, "epistemico": "fato", "origem": "curriculo", "rel": "parte_de"}], "confianca": 1.0, "contraexemplos": [], "descricao": "Código de bloco sobre corpos finitos que corrige ERROS EM RAJADA — o que arranha um CD, um QR ou um setor de disco. Álgebra de 𝔽_q aplicada: opera em símbolos, não bits.", "epistemico": "fato", "exemplos": [], "id": "reed_solomon", "memoria": "semantica", "meta": {"autor": "Reed-Solomon", "dominio": "engenharia", "fonte": ""}, "origem": "wiki", "palavras_chave": [], "sinonimos": [], "tipo": "conceito", "titulo": "Reed-Solomon"}
\section{Reed-Solomon}
Código de bloco sobre corpos finitos que corrige ERROS EM RAJADA — o que arranha um CD, um QR ou um setor de disco. Álgebra de \ensuremath{\mathbb{F}}\_q aplicada: opera em símbolos, não bits.
\prov{relações: especializa codificacao\_de\_canal; relaciona cifra\_metais\_an; parte\_de eng\_telecomunicacoes}
\prov{proveniência: wiki \textperiodcentered{} fato \textperiodcentered{} confiança 1.0 \textperiodcentered{} domínio engenharia \textperiodcentered{} história 4 versões no jornal}

%CRISTAL {"arestas": [{"alvo": "symbolic_beta_pipe_so", "confianca": 1.0, "epistemico": "fato", "origem": "wiki", "rel": "especializa"}, {"alvo": "algebra_dupolinomial_espaco", "confianca": 1.0, "epistemico": "fato", "origem": "wiki", "rel": "usa"}, {"alvo": "word_m_geral_fastpath_2x2", "confianca": 1.0, "epistemico": "fato", "origem": "wiki", "rel": "relaciona"}, {"alvo": "floco_enderecavel_metal_posicao", "confianca": 1.0, "epistemico": "fato", "origem": "wiki", "rel": "relaciona"}, {"alvo": "tiffany", "confianca": 1.0, "epistemico": "fato", "origem": "wiki", "rel": "parte_de"}], "confianca": 1.0, "contraexemplos": [], "descricao": "Correção do Aarão (2 jul 2026) e o desenho implementado: cada bloco do stratum é o FLOCO (pacote seed = PoW nativo, 1×/job, com floco Koch + c² + ν) — e esse pacote contém VÁRIOS metais, não só prata. Antes cada lane já minerava todos os metais mas rotulava tudo como prata; não era problema porque a TAXA DE CONVERSÃO entre metais SE CONSERVA (N(σ_n)=det A_n=−1 p/ todo n — o invariante do refino). Agora minerar = EXTRAIR/REFINAR: dividir o floco em metais reaplicando a reta de cada um (x²−n·x−1, ×A_n). O refino é CONTÍNUO: enquanto não chega outro bloco, o pipe segue extraindo (um passo por chunk minerado, no tempo morto), cada vez mais fundo; quando chega OUTRO bloco, PARA e RECOMEÇA com o novo floco. Mais tempo de bloco = refino mais profundo (a profundidade acompanha o trabalho real). Implementado em goldchain/symbolic.py (classe Refinador: nova_floco reseta no bloco novo, passo extrai em rodízio pelos metais, flush em lote — I/O contido, a lição do olho). Cada floco vira um registro compacto FLOCO ↦ ranges por metal; a POSIÇÃO é a moeda (recomputável, sem inchar o ledger). [D] o pacote refinado em metais, contínuo entre blocos.", "epistemico": "fato", "exemplos": [], "id": "refino_continuo_floco_metais", "memoria": "semantica", "meta": {"autor": "Lopes/broca-so", "dominio": "arquitetura", "fonte": "goldchain/symbolic.py (Refinador, live/simular/recuperar); minera.py:emite_metal; metal_pos.passo_metal(n)"}, "origem": "wiki", "palavras_chave": [], "sinonimos": [], "tipo": "conceito", "titulo": "O FLOCO é um pacote multimetal; minerar = refinar contínuo, dividindo-o em metais (implementado)"}
\section{O FLOCO é um pacote multimetal; minerar = refinar contínuo, dividindo-o em metais (implementado)}
Correção do Aarão (2 jul 2026) e o desenho implementado: cada bloco do stratum é o FLOCO (pacote seed = PoW nativo, 1$\times$/job, com floco Koch + c\ensuremath{^{2}} + \ensuremath{\nu}) — e esse pacote contém VÁRIOS metais, não só prata. Antes cada lane já minerava todos os metais mas rotulava tudo como prata; não era problema porque a TAXA DE CONVERSÃO entre metais SE CONSERVA (N(\ensuremath{\sigma}\_n)=det A\_n=\ensuremath{-}1 p/ todo n — o invariante do refino). Agora minerar = EXTRAIR/REFINAR: dividir o floco em metais reaplicando a reta de cada um (x\ensuremath{^{2}}\ensuremath{-}n\textperiodcentered x\ensuremath{-}1, $\times$A\_n). O refino é CONTÍNUO: enquanto não chega outro bloco, o pipe segue extraindo (um passo por chunk minerado, no tempo morto), cada vez mais fundo; quando chega OUTRO bloco, PARA e RECOMEÇA com o novo floco. Mais tempo de bloco = refino mais profundo (a profundidade acompanha o trabalho real). Implementado em goldchain/symbolic.py (classe Refinador: nova\_floco reseta no bloco novo, passo extrai em rodízio pelos metais, flush em lote — I/O contido, a lição do olho). Cada floco vira um registro compacto FLOCO \ensuremath{\mapsto} ranges por metal; a POSIÇÃO é a moeda (recomputável, sem inchar o ledger). [D] o pacote refinado em metais, contínuo entre blocos.
\prov{relações: especializa symbolic\_beta\_pipe\_so; usa algebra\_dupolinomial\_espaco; relaciona word\_m\_geral\_fastpath\_2x2; relaciona floco\_enderecavel\_metal\_posicao; parte\_de tiffany}
\prov{proveniência: wiki \textperiodcentered{} goldchain/symbolic.py (Refinador, live/simular/recuperar); minera.py:emite\_metal; metal\_pos.passo\_metal(n) \textperiodcentered{} fato \textperiodcentered{} confiança 1.0 \textperiodcentered{} domínio arquitetura \textperiodcentered{} história 32 versões no jornal}

%CRISTAL {"arestas": [{"alvo": "data_center", "confianca": 1.0, "epistemico": "fato", "origem": "wiki", "rel": "parte_de"}, {"alvo": "limite_de_landauer", "confianca": 1.0, "epistemico": "fato", "origem": "wiki", "rel": "usa"}, {"alvo": "ciclo_termodinamico_orbita", "confianca": 1.0, "epistemico": "fato", "origem": "wiki", "rel": "usa"}], "confianca": 1.0, "contraexemplos": [], "descricao": "Tirar o calor dos servidores é metade do problema de um data center. O PUE (energia total / energia de TI) mede o desperdício; muito dele é o calor de Landauer amplificado 10⁹× pela eletrônica real. O ciclo térmico da geração reaparece aqui, agora para RESFRIAR — a órbita fechada de novo.", "epistemico": "fato", "exemplos": [], "id": "refrigeracao_data_center", "memoria": "semantica", "meta": {"dominio": "engenharia", "fonte": "nuvem e data centers"}, "origem": "wiki", "palavras_chave": [], "sinonimos": [], "tipo": "conceito", "titulo": "Refrigeração e PUE"}
\section{Refrigeração e PUE}
Tirar o calor dos servidores é metade do problema de um data center. O PUE (energia total / energia de TI) mede o desperdício; muito dele é o calor de Landauer amplificado 10\ensuremath{^{9}}$\times$ pela eletrônica real. O ciclo térmico da geração reaparece aqui, agora para RESFRIAR — a órbita fechada de novo.
\prov{relações: parte\_de data\_center; usa limite\_de\_landauer; usa ciclo\_termodinamico\_orbita}
\prov{proveniência: wiki \textperiodcentered{} nuvem e data centers \textperiodcentered{} fato \textperiodcentered{} confiança 1.0 \textperiodcentered{} domínio engenharia \textperiodcentered{} história 4 versões no jornal}

%CRISTAL {"arestas": [{"alvo": "redes_5g_nr", "confianca": 1.0, "epistemico": "fato", "origem": "wiki", "rel": "parte_de"}, {"alvo": "estabelecimento_conexao_rrc", "confianca": 1.0, "epistemico": "fato", "origem": "wiki", "rel": "usa"}, {"alvo": "estrato_nas_5g", "confianca": 1.0, "epistemico": "fato", "origem": "wiki", "rel": "usa"}, {"alvo": "comunicacoes_moveis", "confianca": 1.0, "epistemico": "fato", "origem": "wiki", "rel": "especializa"}, {"alvo": "eng_telecomunicacoes", "confianca": 1.0, "epistemico": "fato", "origem": "curriculo", "rel": "parte_de"}], "confianca": 1.0, "contraexemplos": [], "descricao": "O procedimento que põe o UE conhecido e alcançável na rede 5G (substitui o 'attach' do 4G): sobre a conexão RRC, o UE manda Registration Request (NAS) ao AMF, passa por autenticação e segurança, e recebe Registration Accept. Só depois dele há mobilidade e sessão de dados.", "epistemico": "fato", "exemplos": [], "id": "registro_5g", "memoria": "semantica", "meta": {"dominio": "engenharia", "fonte": ""}, "origem": "wiki", "palavras_chave": [], "sinonimos": [], "tipo": "conceito", "titulo": "Registro na Rede (Registration)"}
\section{Registro na Rede (Registration)}
O procedimento que põe o UE conhecido e alcançável na rede 5G (substitui o 'attach' do 4G): sobre a conexão RRC, o UE manda Registration Request (NAS) ao AMF, passa por autenticação e segurança, e recebe Registration Accept. Só depois dele há mobilidade e sessão de dados.
\prov{relações: parte\_de redes\_5g\_nr; usa estabelecimento\_conexao\_rrc; usa estrato\_nas\_5g; especializa comunicacoes\_moveis; parte\_de eng\_telecomunicacoes}
\prov{proveniência: wiki \textperiodcentered{} fato \textperiodcentered{} confiança 1.0 \textperiodcentered{} domínio engenharia \textperiodcentered{} história 6 versões no jornal}

%CRISTAL {"arestas": [{"alvo": "sensor", "confianca": 1.0, "epistemico": "fato", "origem": "wiki", "rel": "usa"}, {"alvo": "corpo_mineral", "confianca": 1.0, "epistemico": "fato", "origem": "wiki", "rel": "relaciona"}], "confianca": 1.0, "contraexemplos": [], "descricao": "SNR: quão forte o sinal está acima do ruído de fundo. Alto = sinal legível (cristal); no limiar = borda; afogado = ruído (caos). Filtrar (média, Kalman) é extrair a órbita do ruído — separar o floco do chiado. Verif.: 20 dB cristal, 5 dB borda, −3 dB caos.", "epistemico": "inferencia", "exemplos": [], "id": "relacao_sinal_ruido", "memoria": "semantica", "meta": {"classes_snr": {"-3": "caos (afogado no ruído)", "20": "cristal (sinal legível)", "5": "borda (no limiar)"}, "dominio": "engenharia", "fonte": "IoT e sensores"}, "origem": "wiki", "palavras_chave": [], "sinonimos": [], "tipo": "conceito", "titulo": "Relação Sinal-Ruído (o floco contra o ruído)"}
\section{Relação Sinal-Ruído (o floco contra o ruído)}
SNR: quão forte o sinal está acima do ruído de fundo. Alto = sinal legível (cristal); no limiar = borda; afogado = ruído (caos). Filtrar (média, Kalman) é extrair a órbita do ruído — separar o floco do chiado. Verif.: 20 dB cristal, 5 dB borda, \ensuremath{-}3 dB caos.
\prov{relações: usa sensor; relaciona corpo\_mineral}
\prov{proveniência: wiki \textperiodcentered{} IoT e sensores \textperiodcentered{} inferencia \textperiodcentered{} confiança 1.0 \textperiodcentered{} domínio engenharia \textperiodcentered{} história 6 versões no jornal}

%CRISTAL {"arestas": [{"alvo": "replicacao_pcr", "confianca": 1.0, "epistemico": "fato", "origem": "wiki", "rel": "relaciona"}, {"alvo": "auto_replicacao_reprap", "confianca": 1.0, "epistemico": "fato", "origem": "wiki", "rel": "relaciona"}, {"alvo": "reproducao_tiffany", "confianca": 1.0, "epistemico": "fato", "origem": "wiki", "rel": "relaciona"}, {"alvo": "genoma", "confianca": 1.0, "epistemico": "fato", "origem": "wiki", "rel": "usa"}], "confianca": 1.0, "contraexemplos": [], "descricao": "O DNA se copia (semiconservativa) — o auto-replicador que a manufatura só imita (RepRap, von Neumann). E a REVISÃO (proofreading 3'→5') + o reparo de mismatch derrubam a taxa de erro em ordens de grandeza: o checksum da vida, a norma que sobrevive à cópia. Verif.: 4 → 9 noves com revisão+reparo.", "epistemico": "inferencia", "exemplos": [], "id": "replicacao_checksum", "memoria": "semantica", "meta": {"dominio": "biotecnologia", "fonte": "biotecnologia e engenharia genética", "noves_completo": 9.0, "noves_crua": 4.0}, "origem": "wiki", "palavras_chave": [], "sinonimos": [], "tipo": "conceito", "titulo": "Replicação = Auto-Replicação + Checksum"}
\section{Replicação = Auto-Replicação + Checksum}
O DNA se copia (semiconservativa) — o auto-replicador que a manufatura só imita (RepRap, von Neumann). E a REVISÃO (proofreading 3'\ensuremath{\to}5') + o reparo de mismatch derrubam a taxa de erro em ordens de grandeza: o checksum da vida, a norma que sobrevive à cópia. Verif.: 4 \ensuremath{\to} 9 noves com revisão+reparo.
\prov{relações: relaciona replicacao\_pcr; relaciona auto\_replicacao\_reprap; relaciona reproducao\_tiffany; usa genoma}
\prov{proveniência: wiki \textperiodcentered{} biotecnologia e engenharia genética \textperiodcentered{} inferencia \textperiodcentered{} confiança 1.0 \textperiodcentered{} domínio biotecnologia \textperiodcentered{} história 5 versões no jornal}

%CRISTAL {"arestas": [{"alvo": "data_center", "confianca": 1.0, "epistemico": "fato", "origem": "wiki", "rel": "usa"}, {"alvo": "armazenamento_o_cristal", "confianca": 1.0, "epistemico": "fato", "origem": "wiki", "rel": "eh_um"}, {"alvo": "corpo_mineral", "confianca": 1.0, "epistemico": "fato", "origem": "wiki", "rel": "usa"}], "confianca": 1.0, "contraexemplos": [], "descricao": "Guardar N cópias do dado em máquinas independentes: perde-se só se TODAS falham (p^N). A disponibilidade 1−p^N sobe rápido com N (os 'noves'). A redundância é a norma/checksum que sobrevive à falha — o cristal da nuvem, o mesmo princípio de integridade do armazenamento. Verif.: 3 cópias a 1% de falha → 6 noves.", "epistemico": "inferencia", "exemplos": [], "id": "replicacao_dados", "memoria": "semantica", "meta": {"disponibilidade_3x_1pct": 0.999999, "dominio": "engenharia", "fonte": "nuvem e data centers", "noves": 6.0}, "origem": "wiki", "palavras_chave": [], "sinonimos": [], "tipo": "conceito", "titulo": "Replicação (o cristal da nuvem)"}
\section{Replicação (o cristal da nuvem)}
Guardar N cópias do dado em máquinas independentes: perde-se só se TODAS falham (p\^{}N). A disponibilidade 1\ensuremath{-}p\^{}N sobe rápido com N (os 'noves'). A redundância é a norma/checksum que sobrevive à falha — o cristal da nuvem, o mesmo princípio de integridade do armazenamento. Verif.: 3 cópias a 1\% de falha \ensuremath{\to} 6 noves.
\prov{relações: usa data\_center; eh\_um armazenamento\_o\_cristal; usa corpo\_mineral}
\prov{proveniência: wiki \textperiodcentered{} nuvem e data centers \textperiodcentered{} inferencia \textperiodcentered{} confiança 1.0 \textperiodcentered{} domínio engenharia \textperiodcentered{} história 4 versões no jornal}

%CRISTAL {"arestas": [{"alvo": "rede_autonoma_broca", "confianca": 1.0, "epistemico": "fato", "origem": "wiki", "rel": "parte_de"}, {"alvo": "roadmap_banda_no_olho", "confianca": 1.0, "epistemico": "fato", "origem": "wiki", "rel": "usa"}], "confianca": 1.0, "contraexemplos": [], "descricao": "Onde cada cliente salva suas aplicações no gex44 — uma base isolada por comunidade (banda), servida pela rede viva (band_bus) e governada por autoridade BAI. O deploy e o armazenamento acontecem POR BANDA: a comunidade é a unidade de posse. goldchain/repo_agente.py o realiza.", "epistemico": "fato", "exemplos": [], "id": "repositorio_broca_os", "memoria": "semantica", "meta": {"dominio": "arquitetura", "fonte": "goldchain/repo_agente.py"}, "origem": "wiki", "palavras_chave": [], "sinonimos": [], "tipo": "conceito", "titulo": "O Repositório do broca-os"}
\section{O Repositório do broca-os}
Onde cada cliente salva suas aplicações no gex44 — uma base isolada por comunidade (banda), servida pela rede viva (band\_bus) e governada por autoridade BAI. O deploy e o armazenamento acontecem POR BANDA: a comunidade é a unidade de posse. goldchain/repo\_agente.py o realiza.
\prov{relações: parte\_de rede\_autonoma\_broca; usa roadmap\_banda\_no\_olho}
\prov{proveniência: wiki \textperiodcentered{} goldchain/repo\_agente.py \textperiodcentered{} fato \textperiodcentered{} confiança 1.0 \textperiodcentered{} domínio arquitetura \textperiodcentered{} história 3 versões no jornal}

%CRISTAL {"arestas": [{"alvo": "curva_de_carga", "confianca": 1.0, "epistemico": "fato", "origem": "wiki", "rel": "usa"}, {"alvo": "carga_coordenada", "confianca": 1.0, "epistemico": "fato", "origem": "wiki", "rel": "relaciona"}], "confianca": 1.0, "contraexemplos": [], "descricao": "A rede pede aos consumidores para deslocar ou cortar carga nos picos, em troca de tarifa menor. Reshapa a curva de carga por incentivo em vez de por mais geração — o lado da demanda virando flexível, um amortecedor da órbita do dia.", "epistemico": "fato", "exemplos": [], "id": "resposta_a_demanda", "memoria": "semantica", "meta": {"dominio": "engenharia", "fonte": "rede de carregamento e smart cities"}, "origem": "wiki", "palavras_chave": [], "sinonimos": [], "tipo": "conceito", "titulo": "Resposta à Demanda (demand response)"}
\section{Resposta à Demanda (demand response)}
A rede pede aos consumidores para deslocar ou cortar carga nos picos, em troca de tarifa menor. Reshapa a curva de carga por incentivo em vez de por mais geração — o lado da demanda virando flexível, um amortecedor da órbita do dia.
\prov{relações: usa curva\_de\_carga; relaciona carga\_coordenada}
\prov{proveniência: wiki \textperiodcentered{} rede de carregamento e smart cities \textperiodcentered{} fato \textperiodcentered{} confiança 1.0 \textperiodcentered{} domínio engenharia \textperiodcentered{} história 6 versões no jornal}

%CRISTAL {"arestas": [{"alvo": "floco_recipiente_termico", "confianca": 1.0, "epistemico": "fato", "origem": "wiki", "rel": "parte_de"}, {"alvo": "parte_fria_t_zero", "confianca": 1.0, "epistemico": "fato", "origem": "wiki", "rel": "relaciona"}, {"alvo": "decomposicao_isomorfa_ao_bai", "confianca": 1.0, "epistemico": "fato", "origem": "wiki", "rel": "usa"}, {"alvo": "corpo_mineral", "confianca": 1.0, "epistemico": "fato", "origem": "wiki", "rel": "usa"}], "confianca": 1.0, "contraexemplos": [], "descricao": "Duas metades do recipiente térmico. AS RETAS SÃO O FRIO: o espectro {σ_i} (a reta mineral discreta, os 'metais' da imagem) é a parte FRIA — T→0, cristalizada, o idempotente. O CAMINHO PELO BAI das órbitas — reconstruir a imagem somando as camadas (o colapso Ψ) — é a máquina DISSIPANDO o calor descartado: a parte quente. O frio guarda o espectro; o quente é o caminho da reconstrução.", "epistemico": "inferencia", "exemplos": [], "id": "retas_frio_bai_dissipa", "memoria": "semantica", "meta": {"dominio": "imagem", "fonte": "floco térmico"}, "origem": "wiki", "palavras_chave": [], "sinonimos": [], "tipo": "conceito", "titulo": "As Retas são o Frio; o BAI Dissipa"}
\section{As Retas são o Frio; o BAI Dissipa}
Duas metades do recipiente térmico. AS RETAS SÃO O FRIO: o espectro \{\ensuremath{\sigma}\_i\} (a reta mineral discreta, os 'metais' da imagem) é a parte FRIA — T\ensuremath{\to}0, cristalizada, o idempotente. O CAMINHO PELO BAI das órbitas — reconstruir a imagem somando as camadas (o colapso \ensuremath{\Psi}) — é a máquina DISSIPANDO o calor descartado: a parte quente. O frio guarda o espectro; o quente é o caminho da reconstrução.
\prov{relações: parte\_de floco\_recipiente\_termico; relaciona parte\_fria\_t\_zero; usa decomposicao\_isomorfa\_ao\_bai; usa corpo\_mineral}
\prov{proveniência: wiki \textperiodcentered{} floco térmico \textperiodcentered{} inferencia \textperiodcentered{} confiança 1.0 \textperiodcentered{} domínio imagem \textperiodcentered{} história 10 versões no jornal}

%CRISTAL {"arestas": [{"alvo": "estados_rrc_nr", "confianca": 1.0, "epistemico": "fato", "origem": "wiki", "rel": "usa"}, {"alvo": "acesso_aleatorio_nr", "confianca": 1.0, "epistemico": "fato", "origem": "wiki", "rel": "usa"}], "confianca": 1.0, "contraexemplos": [], "descricao": "O truque de eficiência do 5G: de RRC_INACTIVE o UE RETOMA a conexão com pouca sinalização (o contexto ficou guardado no gNB), em vez de registrar de novo. Religa quase instantâneo — feito para o tráfego em rajadas de apps e IoT.", "epistemico": "fato", "exemplos": [], "id": "retomada_rrc_inactive", "memoria": "semantica", "meta": {"dominio": "engenharia", "fonte": ""}, "origem": "wiki", "palavras_chave": [], "sinonimos": [], "tipo": "conceito", "titulo": "Retomada de INACTIVE (RRC Resume)"}
\section{Retomada de INACTIVE (RRC Resume)}
O truque de eficiência do 5G: de RRC\_INACTIVE o UE RETOMA a conexão com pouca sinalização (o contexto ficou guardado no gNB), em vez de registrar de novo. Religa quase instantâneo — feito para o tráfego em rajadas de apps e IoT.
\prov{relações: usa estados\_rrc\_nr; usa acesso\_aleatorio\_nr}
\prov{proveniência: wiki \textperiodcentered{} fato \textperiodcentered{} confiança 1.0 \textperiodcentered{} domínio engenharia \textperiodcentered{} história 3 versões no jornal}

%CRISTAL {"arestas": [{"alvo": "roadmap_comunidade_por_banda", "confianca": 1.0, "epistemico": "fato", "origem": "wiki", "rel": "continua"}, {"alvo": "autonomia_subtrativa", "confianca": 1.0, "epistemico": "fato", "origem": "wiki", "rel": "usa"}], "confianca": 1.0, "contraexemplos": [], "descricao": "CONSTRUÍDO: goldchain/celula.py — uma célula reativa SEM scheduler/timer/tick: só age quando uma diferença chega na sua banda (evento no /bus); a ordem emerge do recall (plugável: passo convergente ou colapso da Tiffany). Verificado ao vivo: duas células (A→B, B→A) sem maestro; SEM semente nada flui (não há tick central); UMA semente dispara uma cascata (16→8→4→2→1) que converge e PARA SOZINHA. Prova de autonomia_subtrativa. Caveat: o ENGINE do olho ainda tem seu pulso interno (metabolismo do olho, não scheduler da rede); a camada de rede (bus+células) é 100% reativa. recall=colapso VERIFICADO ao vivo: passeio associativo pelo grafo (OFDM→QAM→Capacidade de Canal→Entropia de Shannon), célula a célula, sem maestro, disparado por uma semente. A 1ª recall carrega o grafo (~8.5s) por processo de célula; depois anda rápido.", "epistemico": "fato", "exemplos": [], "id": "roadmap_auto_gestao_recall", "memoria": "semantica", "meta": {"dominio": "arquitetura", "fonte": "semilla_rede_autonoma"}, "origem": "wiki", "palavras_chave": [], "sinonimos": [], "tipo": "conceito", "titulo": "Roadmap 4 — auto-gestão pelo recall (tirar o scheduler)"}
\section{Roadmap 4 — auto-gestão pelo recall (tirar o scheduler)}
CONSTRUÍDO: goldchain/celula.py — uma célula reativa SEM scheduler/timer/tick: só age quando uma diferença chega na sua banda (evento no /bus); a ordem emerge do recall (plugável: passo convergente ou colapso da Tiffany). Verificado ao vivo: duas células (A\ensuremath{\to}B, B\ensuremath{\to}A) sem maestro; SEM semente nada flui (não há tick central); UMA semente dispara uma cascata (16\ensuremath{\to}8\ensuremath{\to}4\ensuremath{\to}2\ensuremath{\to}1) que converge e PARA SOZINHA. Prova de autonomia\_subtrativa. Caveat: o ENGINE do olho ainda tem seu pulso interno (metabolismo do olho, não scheduler da rede); a camada de rede (bus+células) é 100\% reativa. recall=colapso VERIFICADO ao vivo: passeio associativo pelo grafo (OFDM\ensuremath{\to}QAM\ensuremath{\to}Capacidade de Canal\ensuremath{\to}Entropia de Shannon), célula a célula, sem maestro, disparado por uma semente. A 1ª recall carrega o grafo (\~{}8.5s) por processo de célula; depois anda rápido.
\prov{relações: continua roadmap\_comunidade\_por\_banda; usa autonomia\_subtrativa}
\prov{proveniência: wiki \textperiodcentered{} semilla\_rede\_autonoma \textperiodcentered{} fato \textperiodcentered{} confiança 1.0 \textperiodcentered{} domínio arquitetura \textperiodcentered{} história 5 versões no jornal}

%CRISTAL {"arestas": [{"alvo": "rede_autonoma_broca", "confianca": 1.0, "epistemico": "fato", "origem": "wiki", "rel": "parte_de"}, {"alvo": "banda_exclusiva_por_tecido", "confianca": 1.0, "epistemico": "fato", "origem": "wiki", "rel": "usa"}, {"alvo": "cockpit_olho", "confianca": 1.0, "epistemico": "fato", "origem": "wiki", "rel": "usa"}, {"alvo": "roteamento_multi_tenant", "confianca": 1.0, "epistemico": "fato", "origem": "wiki", "rel": "usa"}], "confianca": 1.0, "contraexemplos": [], "descricao": "CONSTRUÍDO: o cockpit (goldchain/cockpit_server.py) agora transmite POR BANDA — '/sse/broca?tenant=<id>' ou '?banda=<hex>' emoldura cada evento com dado ⊕ keystream(banda_do_tecido) via reuso de neuronio/antena_canal; quem tem a banda sintoniza, quem não tem lê ruído base64. Sem parâmetro = texto puro (retrocompatível). O 1º tijolo: o olho deixou de ser um canal global single-tenant. Frontend WIRED: goldchain/cockpit/banda.js (Web Crypto) desmoldura o fluxo quando a página abre com ?banda=<hex>; sem parâmetro, UI padrão intocada. Decode verificado em node byte-a-byte contra o Python; integração no browser guardada mas não testada em navegador real.", "epistemico": "fato", "exemplos": [], "id": "roadmap_banda_no_olho", "memoria": "semantica", "meta": {"dominio": "arquitetura", "fonte": "semilla_rede_autonoma"}, "origem": "wiki", "palavras_chave": [], "sinonimos": [], "tipo": "conceito", "titulo": "Roadmap 1 — banda exclusiva no olho"}
\section{Roadmap 1 — banda exclusiva no olho}
CONSTRUÍDO: o cockpit (goldchain/cockpit\_server.py) agora transmite POR BANDA — '/sse/broca?tenant=<id>' ou '?banda=<hex>' emoldura cada evento com dado \ensuremath{\oplus} keystream(banda\_do\_tecido) via reuso de neuronio/antena\_canal; quem tem a banda sintoniza, quem não tem lê ruído base64. Sem parâmetro = texto puro (retrocompatível). O 1º tijolo: o olho deixou de ser um canal global single-tenant. Frontend WIRED: goldchain/cockpit/banda.js (Web Crypto) desmoldura o fluxo quando a página abre com ?banda=<hex>; sem parâmetro, UI padrão intocada. Decode verificado em node byte-a-byte contra o Python; integração no browser guardada mas não testada em navegador real.
\prov{relações: parte\_de rede\_autonoma\_broca; usa banda\_exclusiva\_por\_tecido; usa cockpit\_olho; usa roteamento\_multi\_tenant}
\prov{proveniência: wiki \textperiodcentered{} semilla\_rede\_autonoma \textperiodcentered{} fato \textperiodcentered{} confiança 1.0 \textperiodcentered{} domínio arquitetura \textperiodcentered{} história 7 versões no jornal}

%CRISTAL {"arestas": [{"alvo": "roadmap_olho_multibanda", "confianca": 1.0, "epistemico": "fato", "origem": "wiki", "rel": "continua"}, {"alvo": "comunidade_organica_por_banda", "confianca": 1.0, "epistemico": "fato", "origem": "wiki", "rel": "usa"}], "confianca": 1.0, "contraexemplos": [], "descricao": "CONSTRUÍDO: presença viva por banda (goldchain/band_bus.py roster/mapa + cockpit 'GET /comunidade?banda=' e 'GET /comunidades'). Assinar o /bus com ?peer=&nome=&tipo= ENTRA na comunidade da banda; ao cair a conexão, some (presença orgânica). O mapa mostra as PONTES — pares presentes em ≥2 bandas = comunidades que se cruzam. Verificado: software E arte coexistem numa banda, e um par-ponte liga duas. O central só APRESENTA (diretório), o conteúdo segue no bus — nunca é correio. SEMENTES QUE CIRCULAM construído (goldchain/semente.py): um artefato (software/arte) nasce numa banda e se reproduz cruzando comunidades via pares-PONTE (replanta nas outras bandas), com linhagem e dedup+TTL (não reflooda). Verificado ao vivo: semente de arte A→B→C, geração 2, linhagem [semeador,ponte1,ponte2].", "epistemico": "fato", "exemplos": [], "id": "roadmap_comunidade_por_banda", "memoria": "semantica", "meta": {"dominio": "arquitetura", "fonte": "semilla_rede_autonoma"}, "origem": "wiki", "palavras_chave": [], "sinonimos": [], "tipo": "conceito", "titulo": "Roadmap 3 — comunidades por banda compartilhada"}
\section{Roadmap 3 — comunidades por banda compartilhada}
CONSTRUÍDO: presença viva por banda (goldchain/band\_bus.py roster/mapa + cockpit 'GET /comunidade?banda=' e 'GET /comunidades'). Assinar o /bus com ?peer=\&nome=\&tipo= ENTRA na comunidade da banda; ao cair a conexão, some (presença orgânica). O mapa mostra as PONTES — pares presentes em \ensuremath{\ge}2 bandas = comunidades que se cruzam. Verificado: software E arte coexistem numa banda, e um par-ponte liga duas. O central só APRESENTA (diretório), o conteúdo segue no bus — nunca é correio. SEMENTES QUE CIRCULAM construído (goldchain/semente.py): um artefato (software/arte) nasce numa banda e se reproduz cruzando comunidades via pares-PONTE (replanta nas outras bandas), com linhagem e dedup+TTL (não reflooda). Verificado ao vivo: semente de arte A\ensuremath{\to}B\ensuremath{\to}C, geração 2, linhagem [semeador,ponte1,ponte2].
\prov{relações: continua roadmap\_olho\_multibanda; usa comunidade\_organica\_por\_banda}
\prov{proveniência: wiki \textperiodcentered{} semilla\_rede\_autonoma \textperiodcentered{} fato \textperiodcentered{} confiança 1.0 \textperiodcentered{} domínio arquitetura \textperiodcentered{} história 5 versões no jornal}

%CRISTAL {"arestas": [{"alvo": "roadmap_banda_no_olho", "confianca": 1.0, "epistemico": "fato", "origem": "wiki", "rel": "continua"}, {"alvo": "a_antena_transmite_em_tempo_real", "confianca": 1.0, "epistemico": "fato", "origem": "wiki", "rel": "usa"}], "confianca": 1.0, "contraexemplos": [], "descricao": "CONSTRUÍDO: goldchain/band_bus.py (BandBus/BUS) + cockpit 'GET /bus?banda=' (assina) e 'POST /bus?banda=' (publica). Fan-out por banda, SEM histórico (antena, não estante), e SÓ A DIFERENÇA viaja (payload idêntico consecutivo não reenvia). Verificado ao vivo: isolamento por banda, dedup, concorrência (ThreadingHTTPServer). Refino pendente: delta/diff estrutural fino (hoje 'só a diferença' = dedup de idênticos) e wire do frontend ao /bus.", "epistemico": "fato", "exemplos": [], "id": "roadmap_olho_multibanda", "memoria": "semantica", "meta": {"dominio": "arquitetura", "fonte": "semilla_rede_autonoma"}, "origem": "wiki", "palavras_chave": [], "sinonimos": [], "tipo": "conceito", "titulo": "Roadmap 2 — o olho como barramento multi-banda"}
\section{Roadmap 2 — o olho como barramento multi-banda}
CONSTRUÍDO: goldchain/band\_bus.py (BandBus/BUS) + cockpit 'GET /bus?banda=' (assina) e 'POST /bus?banda=' (publica). Fan-out por banda, SEM histórico (antena, não estante), e SÓ A DIFERENÇA viaja (payload idêntico consecutivo não reenvia). Verificado ao vivo: isolamento por banda, dedup, concorrência (ThreadingHTTPServer). Refino pendente: delta/diff estrutural fino (hoje 'só a diferença' = dedup de idênticos) e wire do frontend ao /bus.
\prov{relações: continua roadmap\_banda\_no\_olho; usa a\_antena\_transmite\_em\_tempo\_real}
\prov{proveniência: wiki \textperiodcentered{} semilla\_rede\_autonoma \textperiodcentered{} fato \textperiodcentered{} confiança 1.0 \textperiodcentered{} domínio arquitetura \textperiodcentered{} história 4 versões no jornal}

%CRISTAL {"arestas": [{"alvo": "cibernetica", "confianca": 1.0, "epistemico": "fato", "origem": "wiki", "rel": "eh_um"}], "confianca": 1.0, "contraexemplos": [], "descricao": "A engenharia de máquinas que percebem, decidem e agem no mundo físico: sensores, processador e atuadores fechados num laço. É a IA ganhando corpo — onde a decisão vira movimento e o movimento muda o mundo que se volta a medir.", "epistemico": "fato", "exemplos": [], "id": "robotica", "memoria": "semantica", "meta": {"dominio": "robotica", "fonte": "robótica e sistemas autônomos"}, "origem": "wiki", "palavras_chave": [], "sinonimos": [], "tipo": "conceito", "titulo": "Robótica"}
\section{Robótica}
A engenharia de máquinas que percebem, decidem e agem no mundo físico: sensores, processador e atuadores fechados num laço. É a IA ganhando corpo — onde a decisão vira movimento e o movimento muda o mundo que se volta a medir.
\prov{relações: eh\_um cibernetica}
\prov{proveniência: wiki \textperiodcentered{} robótica e sistemas autônomos \textperiodcentered{} fato \textperiodcentered{} confiança 1.0 \textperiodcentered{} domínio robotica \textperiodcentered{} história 2 versões no jornal}

%CRISTAL {"arestas": [{"alvo": "protocolo_ip", "confianca": 1.0, "epistemico": "fato", "origem": "wiki", "rel": "usa"}, {"alvo": "comutacao_de_pacotes", "confianca": 1.0, "epistemico": "fato", "origem": "wiki", "rel": "usa"}, {"alvo": "eng_telecomunicacoes", "confianca": 1.0, "epistemico": "fato", "origem": "curriculo", "rel": "parte_de"}], "confianca": 1.0, "contraexemplos": [], "descricao": "Decidir por qual salto um pacote segue rumo ao destino — tabelas construídas por algoritmos de menor caminho. O que transforma muitas redes locais numa internet navegável.", "epistemico": "fato", "exemplos": [], "id": "roteamento", "memoria": "semantica", "meta": {"autor": "Dijkstra", "dominio": "engenharia", "fonte": ""}, "origem": "wiki", "palavras_chave": [], "sinonimos": [], "tipo": "conceito", "titulo": "Roteamento"}
\section{Roteamento}
Decidir por qual salto um pacote segue rumo ao destino — tabelas construídas por algoritmos de menor caminho. O que transforma muitas redes locais numa internet navegável.
\prov{relações: usa protocolo\_ip; usa comutacao\_de\_pacotes; parte\_de eng\_telecomunicacoes}
\prov{proveniência: wiki \textperiodcentered{} fato \textperiodcentered{} confiança 1.0 \textperiodcentered{} domínio engenharia \textperiodcentered{} história 4 versões no jornal}

%CRISTAL {"arestas": [{"alvo": "imagem_bai", "confianca": 1.0, "epistemico": "fato", "origem": "wiki", "rel": "parte_de"}, {"alvo": "colapso", "confianca": 1.0, "epistemico": "fato", "origem": "wiki", "rel": "eh_um"}, {"alvo": "minimax_e_colapso", "confianca": 1.0, "epistemico": "fato", "origem": "wiki", "rel": "relaciona"}, {"alvo": "kmeans_atratores", "confianca": 1.0, "epistemico": "fato", "origem": "wiki", "rel": "usa"}, {"alvo": "limiar_de_otsu", "confianca": 1.0, "epistemico": "fato", "origem": "wiki", "rel": "usa"}], "confianca": 1.0, "contraexemplos": [], "descricao": "Atribuir um rótulo a cada pixel — a superposição 'a que região pertenço?' desabando numa resposta — É o colapso Ψ do BAI. O k-means colapsa cada pixel no cristal mais próximo; Otsu colapsa na borda do histograma. O mesmo Ψ do minimax do xadrez, agora por pixel. É a tradução central do dicionário: rotulagem ↔ colapso_psi ↔ minimax.", "epistemico": "inferencia", "exemplos": [], "id": "rotulagem_e_colapso", "memoria": "semantica", "meta": {"dominio": "ponte", "fonte": "imagem↔BAI"}, "origem": "wiki", "palavras_chave": [], "sinonimos": [], "tipo": "conceito", "titulo": "Rotulagem = Colapso Ψ (a ponte mais forte)"}
\section{Rotulagem = Colapso \ensuremath{\Psi} (a ponte mais forte)}
Atribuir um rótulo a cada pixel — a superposição 'a que região pertenço?' desabando numa resposta — É o colapso \ensuremath{\Psi} do BAI. O k-means colapsa cada pixel no cristal mais próximo; Otsu colapsa na borda do histograma. O mesmo \ensuremath{\Psi} do minimax do xadrez, agora por pixel. É a tradução central do dicionário: rotulagem \ensuremath{\leftrightarrow} colapso\_psi \ensuremath{\leftrightarrow} minimax.
\prov{relações: parte\_de imagem\_bai; eh\_um colapso; relaciona minimax\_e\_colapso; usa kmeans\_atratores; usa limiar\_de\_otsu}
\prov{proveniência: wiki \textperiodcentered{} imagem\ensuremath{\leftrightarrow}BAI \textperiodcentered{} inferencia \textperiodcentered{} confiança 1.0 \textperiodcentered{} domínio ponte \textperiodcentered{} história 6 versões no jornal}

%CRISTAL {"arestas": [{"alvo": "eletronica_analogica", "confianca": 1.0, "epistemico": "fato", "origem": "wiki", "rel": "relaciona"}, {"alvo": "probabilidade_e_estatistica", "confianca": 1.0, "epistemico": "fato", "origem": "wiki", "rel": "usa"}, {"alvo": "eng_eletronica", "confianca": 1.0, "epistemico": "fato", "origem": "curriculo", "rel": "parte_de"}], "confianca": 1.0, "contraexemplos": [], "descricao": "A flutuação inevitável — térmico (Johnson), de disparo (shot), 1/f — que põe piso ao menor sinal detectável. O inimigo comum da eletrônica de precisão e da telecomunicação.", "epistemico": "fato", "exemplos": [], "id": "ruido_eletronico", "memoria": "semantica", "meta": {"dominio": "engenharia", "fonte": ""}, "origem": "wiki", "palavras_chave": [], "sinonimos": [], "tipo": "conceito", "titulo": "Ruído Eletrônico"}
\section{Ruído Eletrônico}
A flutuação inevitável — térmico (Johnson), de disparo (shot), 1/f — que põe piso ao menor sinal detectável. O inimigo comum da eletrônica de precisão e da telecomunicação.
\prov{relações: relaciona eletronica\_analogica; usa probabilidade\_e\_estatistica; parte\_de eng\_eletronica}
\prov{proveniência: wiki \textperiodcentered{} fato \textperiodcentered{} confiança 1.0 \textperiodcentered{} domínio engenharia \textperiodcentered{} história 4 versões no jornal}

%CRISTAL {"arestas": [{"alvo": "colapso", "confianca": 1.0, "epistemico": "fato", "origem": "wiki", "rel": "eh_um"}, {"alvo": "classificacao", "confianca": 1.0, "epistemico": "fato", "origem": "wiki", "rel": "relaciona"}, {"alvo": "sensor_e_colapso", "confianca": 1.0, "epistemico": "fato", "origem": "wiki", "rel": "relaciona"}, {"alvo": "corpo_mineral", "confianca": 1.0, "epistemico": "fato", "origem": "wiki", "rel": "usa"}], "confianca": 1.0, "contraexemplos": [], "descricao": "Partir a imagem em regiões (objeto/fundo, coisas). Cada pixel, numa superposição de 'a que região pertenço?', DESABA num rótulo — é o colapso do BAI por pixel. Os métodos (limiar, k-means, região, grafo) são o mesmo gesto em roupas diferentes: colapsar o campo de pixels em partes.", "epistemico": "fato", "exemplos": [], "id": "segmentacao_de_imagem", "memoria": "semantica", "meta": {"dominio": "imagem", "fonte": "segmentação"}, "origem": "wiki", "palavras_chave": [], "sinonimos": [], "tipo": "conceito", "titulo": "Segmentação de Imagem"}
\section{Segmentação de Imagem}
Partir a imagem em regiões (objeto/fundo, coisas). Cada pixel, numa superposição de 'a que região pertenço?', DESABA num rótulo — é o colapso do BAI por pixel. Os métodos (limiar, k-means, região, grafo) são o mesmo gesto em roupas diferentes: colapsar o campo de pixels em partes.
\prov{relações: eh\_um colapso; relaciona classificacao; relaciona sensor\_e\_colapso; usa corpo\_mineral}
\prov{proveniência: wiki \textperiodcentered{} segmentação \textperiodcentered{} fato \textperiodcentered{} confiança 1.0 \textperiodcentered{} domínio imagem \textperiodcentered{} história 10 versões no jornal}

%CRISTAL {"arestas": [{"alvo": "camada_pdcp_nr", "confianca": 1.0, "epistemico": "fato", "origem": "wiki", "rel": "usa"}, {"alvo": "seguranca_de_redes", "confianca": 1.0, "epistemico": "fato", "origem": "wiki", "rel": "especializa"}, {"alvo": "criptografia", "confianca": 1.0, "epistemico": "fato", "origem": "wiki", "rel": "usa"}], "confianca": 1.0, "contraexemplos": [], "descricao": "Autenticação MÚTUA rede↔dispositivo (5G-AKA), cifragem e integridade em PDCP e NAS, e o ocultamento do identificador permanente (SUPI→SUCI cifrado) — o 5G fechou o rastreio por IMSI que assombrava o 4G.", "epistemico": "fato", "exemplos": [], "id": "seguranca_5g", "memoria": "semantica", "meta": {"dominio": "engenharia", "fonte": ""}, "origem": "wiki", "palavras_chave": [], "sinonimos": [], "tipo": "conceito", "titulo": "Segurança 5G (5G-AKA)"}
\section{Segurança 5G (5G-AKA)}
Autenticação MÚTUA rede\ensuremath{\leftrightarrow}dispositivo (5G-AKA), cifragem e integridade em PDCP e NAS, e o ocultamento do identificador permanente (SUPI\ensuremath{\to}SUCI cifrado) — o 5G fechou o rastreio por IMSI que assombrava o 4G.
\prov{relações: usa camada\_pdcp\_nr; especializa seguranca\_de\_redes; usa criptografia}
\prov{proveniência: wiki \textperiodcentered{} fato \textperiodcentered{} confiança 1.0 \textperiodcentered{} domínio engenharia \textperiodcentered{} história 4 versões no jornal}

%CRISTAL {"arestas": [{"alvo": "criptografia", "confianca": 1.0, "epistemico": "fato", "origem": "wiki", "rel": "usa"}, {"alvo": "rede_de_computadores", "confianca": 1.0, "epistemico": "fato", "origem": "wiki", "rel": "parte_de"}, {"alvo": "eng_telecomunicacoes", "confianca": 1.0, "epistemico": "fato", "origem": "curriculo", "rel": "parte_de"}], "confianca": 1.0, "contraexemplos": [], "descricao": "Proteger a comunicação: sigilo, integridade, autenticação e disponibilidade — contra escuta, adulteração e personificação. A camada transversal que a criptografia realiza na prática.", "epistemico": "fato", "exemplos": [], "id": "seguranca_de_redes", "memoria": "semantica", "meta": {"dominio": "engenharia", "fonte": ""}, "origem": "wiki", "palavras_chave": [], "sinonimos": [], "tipo": "conceito", "titulo": "Segurança de Redes"}
\section{Segurança de Redes}
Proteger a comunicação: sigilo, integridade, autenticação e disponibilidade — contra escuta, adulteração e personificação. A camada transversal que a criptografia realiza na prática.
\prov{relações: usa criptografia; parte\_de rede\_de\_computadores; parte\_de eng\_telecomunicacoes}
\prov{proveniência: wiki \textperiodcentered{} fato \textperiodcentered{} confiança 1.0 \textperiodcentered{} domínio engenharia \textperiodcentered{} história 4 versões no jornal}

%CRISTAL {"arestas": [{"alvo": "registro_5g", "confianca": 1.0, "epistemico": "fato", "origem": "wiki", "rel": "parte_de"}, {"alvo": "funcoes_de_rede_5gc", "confianca": 1.0, "epistemico": "fato", "origem": "wiki", "rel": "usa"}], "confianca": 1.0, "contraexemplos": [], "descricao": "A rede escolhe qual AMF servirá o UE (por slice, carga, região) e busca o perfil de assinatura no UDM — quais serviços, QoS e slices o cliente pode usar. O NRF ajuda a descobrir a função certa.", "epistemico": "inferencia", "exemplos": [], "id": "selecao_amf_registro", "memoria": "semantica", "meta": {"dominio": "engenharia", "fonte": ""}, "origem": "wiki", "palavras_chave": [], "sinonimos": [], "tipo": "conceito", "titulo": "Seleção de AMF e Assinatura"}
\section{Seleção de AMF e Assinatura}
A rede escolhe qual AMF servirá o UE (por slice, carga, região) e busca o perfil de assinatura no UDM — quais serviços, QoS e slices o cliente pode usar. O NRF ajuda a descobrir a função certa.
\prov{relações: parte\_de registro\_5g; usa funcoes\_de\_rede\_5gc}
\prov{proveniência: wiki \textperiodcentered{} inferencia \textperiodcentered{} confiança 1.0 \textperiodcentered{} domínio engenharia \textperiodcentered{} história 3 versões no jornal}

%CRISTAL {"arestas": [{"alvo": "estados_rrc_nr", "confianca": 1.0, "epistemico": "fato", "origem": "wiki", "rel": "usa"}, {"alvo": "ssb_sincronizacao_nr", "confianca": 1.0, "epistemico": "fato", "origem": "wiki", "rel": "usa"}], "confianca": 1.0, "contraexemplos": [], "descricao": "Em IDLE não há handover: o próprio UE escolhe em qual célula acampar (seleção) e migra por conta própria quando outra fica melhor (reseleção), lendo o SSB e a informação do sistema. Mobilidade barata, sem a rede no laço.", "epistemico": "fato", "exemplos": [], "id": "selecao_reselecao_celula", "memoria": "semantica", "meta": {"dominio": "engenharia", "fonte": ""}, "origem": "wiki", "palavras_chave": [], "sinonimos": [], "tipo": "conceito", "titulo": "Seleção e Reseleção de Célula (IDLE)"}
\section{Seleção e Reseleção de Célula (IDLE)}
Em IDLE não há handover: o próprio UE escolhe em qual célula acampar (seleção) e migra por conta própria quando outra fica melhor (reseleção), lendo o SSB e a informação do sistema. Mobilidade barata, sem a rede no laço.
\prov{relações: usa estados\_rrc\_nr; usa ssb\_sincronizacao\_nr}
\prov{proveniência: wiki \textperiodcentered{} fato \textperiodcentered{} confiança 1.0 \textperiodcentered{} domínio engenharia \textperiodcentered{} história 3 versões no jornal}

%CRISTAL {"arestas": [{"alvo": "espectro_irredutivel_floco", "confianca": 1.0, "epistemico": "fato", "origem": "wiki", "rel": "usa"}, {"alvo": "floco_selo_16bytes", "confianca": 1.0, "epistemico": "fato", "origem": "wiki", "rel": "eh_um"}, {"alvo": "gato_como_cifra", "confianca": 1.0, "epistemico": "fato", "origem": "wiki", "rel": "usa"}, {"alvo": "cifra_de_cristal", "confianca": 1.0, "epistemico": "fato", "origem": "wiki", "rel": "usa"}, {"alvo": "impressao_koch_compressao", "confianca": 1.0, "epistemico": "fato", "origem": "wiki", "rel": "relaciona"}], "confianca": 1.0, "contraexemplos": [], "descricao": "O floco (espectro irredutível) é a assinatura da imagem, e serve para três coisas: COMPRIMIR (manda-se só as k camadas SVD, não a imagem inteira), SELAR (o sha256 do floco quantizado, 16 bytes, é a identidade — o mesmo floco_selo do canal, agora de uma imagem) e CIFRAR (o floco viaja embaralhado pelo gato metálico, a Cifra de Cristal, e decifra exato). Adulterar um bit do floco muda o selo. Verif.: reconstrói, sela em 16 bytes, cifra e decifra exato.", "epistemico": "fato", "exemplos": [], "id": "selo_de_imagem", "memoria": "semantica", "meta": {"dominio": "imagem", "fonte": "selo e órbita temporal"}, "origem": "wiki", "palavras_chave": [], "sinonimos": [], "tipo": "conceito", "titulo": "O Selo de Imagem (o floco como cifra)"}
\section{O Selo de Imagem (o floco como cifra)}
O floco (espectro irredutível) é a assinatura da imagem, e serve para três coisas: COMPRIMIR (manda-se só as k camadas SVD, não a imagem inteira), SELAR (o sha256 do floco quantizado, 16 bytes, é a identidade — o mesmo floco\_selo do canal, agora de uma imagem) e CIFRAR (o floco viaja embaralhado pelo gato metálico, a Cifra de Cristal, e decifra exato). Adulterar um bit do floco muda o selo. Verif.: reconstrói, sela em 16 bytes, cifra e decifra exato.
\prov{relações: usa espectro\_irredutivel\_floco; eh\_um floco\_selo\_16bytes; usa gato\_como\_cifra; usa cifra\_de\_cristal; relaciona impressao\_koch\_compressao}
\prov{proveniência: wiki \textperiodcentered{} selo e órbita temporal \textperiodcentered{} fato \textperiodcentered{} confiança 1.0 \textperiodcentered{} domínio imagem \textperiodcentered{} história 6 versões no jornal}

%CRISTAL {"arestas": [{"alvo": "imagem_bai", "confianca": 1.0, "epistemico": "fato", "origem": "wiki", "rel": "parte_de"}, {"alvo": "crescer_regiao", "confianca": 1.0, "epistemico": "fato", "origem": "wiki", "rel": "usa"}], "confianca": 1.0, "contraexemplos": [], "descricao": "O pixel-semente (ou marcador) de onde a região cresce É a raiz do BAI — a autoridade intrínseca de onde a propagação parte, o rei do xadrez a proteger. Toda a segmentação por crescimento herda a identidade da semente, como a cadeia de autorização herda da raiz.", "epistemico": "inferencia", "exemplos": [], "id": "semente_e_raiz", "memoria": "semantica", "meta": {"dominio": "ponte", "fonte": "imagem↔BAI"}, "origem": "wiki", "palavras_chave": [], "sinonimos": [], "tipo": "conceito", "titulo": "Semente = Raiz (a autoridade intrínseca)"}
\section{Semente = Raiz (a autoridade intrínseca)}
O pixel-semente (ou marcador) de onde a região cresce É a raiz do BAI — a autoridade intrínseca de onde a propagação parte, o rei do xadrez a proteger. Toda a segmentação por crescimento herda a identidade da semente, como a cadeia de autorização herda da raiz.
\prov{relações: parte\_de imagem\_bai; usa crescer\_regiao}
\prov{proveniência: wiki \textperiodcentered{} imagem\ensuremath{\leftrightarrow}BAI \textperiodcentered{} inferencia \textperiodcentered{} confiança 1.0 \textperiodcentered{} domínio ponte \textperiodcentered{} história 3 versões no jornal}

%CRISTAL {"arestas": [{"alvo": "condutor_isolante_semicondutor", "confianca": 1.0, "epistemico": "fato", "origem": "wiki", "rel": "especializa"}, {"alvo": "fisica_do_estado_solido", "confianca": 1.0, "epistemico": "fato", "origem": "wiki", "rel": "usa"}], "confianca": 1.0, "contraexemplos": [], "descricao": "Material (silício, germânio) cuja condução é controlável por dopagem, luz ou tensão — o meio-termo que tornou possível o transistor e o chip.", "epistemico": "fato", "exemplos": [], "id": "semicondutor", "memoria": "semantica", "meta": {"dominio": "engenharia", "fonte": ""}, "origem": "wiki", "palavras_chave": [], "sinonimos": [], "tipo": "conceito", "titulo": "Semicondutor"}
\section{Semicondutor}
Material (silício, germânio) cuja condução é controlável por dopagem, luz ou tensão — o meio-termo que tornou possível o transistor e o chip.
\prov{relações: especializa condutor\_isolante\_semicondutor; usa fisica\_do\_estado\_solido}
\prov{proveniência: wiki \textperiodcentered{} fato \textperiodcentered{} confiança 1.0 \textperiodcentered{} domínio engenharia \textperiodcentered{} história 3 versões no jornal}

%CRISTAL {"arestas": [{"alvo": "internet_das_coisas", "confianca": 1.0, "epistemico": "fato", "origem": "wiki", "rel": "parte_de"}], "confianca": 1.0, "contraexemplos": [], "descricao": "O dispositivo que transforma uma grandeza física (temperatura, luz, corrente, posição) num sinal elétrico legível. É a fronteira entre o mundo contínuo e a máquina discreta — o ponto onde a realidade entra no grafo.", "epistemico": "fato", "exemplos": [], "id": "sensor", "memoria": "semantica", "meta": {"dominio": "engenharia", "fonte": "IoT e sensores"}, "origem": "wiki", "palavras_chave": [], "sinonimos": [], "tipo": "conceito", "titulo": "Sensor"}
\section{Sensor}
O dispositivo que transforma uma grandeza física (temperatura, luz, corrente, posição) num sinal elétrico legível. É a fronteira entre o mundo contínuo e a máquina discreta — o ponto onde a realidade entra no grafo.
\prov{relações: parte\_de internet\_das\_coisas}
\prov{proveniência: wiki \textperiodcentered{} IoT e sensores \textperiodcentered{} fato \textperiodcentered{} confiança 1.0 \textperiodcentered{} domínio engenharia \textperiodcentered{} história 4 versões no jornal}

%CRISTAL {"arestas": [{"alvo": "sensor", "confianca": 1.0, "epistemico": "fato", "origem": "wiki", "rel": "explica"}, {"alvo": "colapso", "confianca": 1.0, "epistemico": "fato", "origem": "wiki", "rel": "eh_um"}, {"alvo": "bai", "confianca": 1.0, "epistemico": "fato", "origem": "wiki", "rel": "usa"}, {"alvo": "neuronio", "confianca": 1.0, "epistemico": "fato", "origem": "wiki", "rel": "relaciona"}, {"alvo": "olho_gex44_stream", "confianca": 1.0, "epistemico": "fato", "origem": "wiki", "rel": "relaciona"}, {"alvo": "decoerencia_colapso_medida", "confianca": 1.0, "epistemico": "fato", "origem": "wiki", "rel": "relaciona"}, {"alvo": "corpo_mineral", "confianca": 1.0, "epistemico": "fato", "origem": "wiki", "rel": "usa"}], "confianca": 1.0, "contraexemplos": [], "descricao": "Medir é DESABAR o contínuo do mundo num bit/número: o conversor A/D quantiza uma grandeza em N bits. É o mesmo colapso do BAI (a superposição do real → uma leitura) e é o OLHO da máquina — como o olho lê o floco do gex44. A lição do neurônio: 1 bit + 1 operador, sem buffer. A resolução (o erro de quantização) é o preço do colapso. Verif.: 8 bits → erro ~0.0002.", "epistemico": "inferencia", "exemplos": [], "id": "sensor_e_colapso", "memoria": "semantica", "meta": {"dominio": "engenharia", "erro_8bits": 0.00015, "fonte": "IoT e sensores", "resolucao_8bits": 0.00392}, "origem": "wiki", "palavras_chave": [], "sinonimos": [], "tipo": "conceito", "titulo": "O Sensor é um Colapso (o olho da máquina)"}
\section{O Sensor é um Colapso (o olho da máquina)}
Medir é DESABAR o contínuo do mundo num bit/número: o conversor A/D quantiza uma grandeza em N bits. É o mesmo colapso do BAI (a superposição do real \ensuremath{\to} uma leitura) e é o OLHO da máquina — como o olho lê o floco do gex44. A lição do neurônio: 1 bit + 1 operador, sem buffer. A resolução (o erro de quantização) é o preço do colapso. Verif.: 8 bits \ensuremath{\to} erro \~{}0.0002.
\prov{relações: explica sensor; eh\_um colapso; usa bai; relaciona neuronio; relaciona olho\_gex44\_stream; relaciona decoerencia\_colapso\_medida; usa corpo\_mineral}
\prov{proveniência: wiki \textperiodcentered{} IoT e sensores \textperiodcentered{} inferencia \textperiodcentered{} confiança 1.0 \textperiodcentered{} domínio engenharia \textperiodcentered{} história 11 versões no jornal}

%CRISTAL {"arestas": [{"alvo": "instrumentacao_medidas", "confianca": 1.0, "epistemico": "fato", "origem": "wiki", "rel": "usa"}, {"alvo": "eng_eletronica", "confianca": 1.0, "epistemico": "fato", "origem": "curriculo", "rel": "parte_de"}], "confianca": 1.0, "contraexemplos": [], "descricao": "Converter grandeza física (temperatura, luz, força) em sinal elétrico e vice-versa. Os órgãos dos sentidos do sistema eletrônico; a origem de todo dado do mundo real.", "epistemico": "fato", "exemplos": [], "id": "sensor_transdutor", "memoria": "semantica", "meta": {"dominio": "engenharia", "fonte": ""}, "origem": "wiki", "palavras_chave": [], "sinonimos": [], "tipo": "conceito", "titulo": "Sensores e Transdutores"}
\section{Sensores e Transdutores}
Converter grandeza física (temperatura, luz, força) em sinal elétrico e vice-versa. Os órgãos dos sentidos do sistema eletrônico; a origem de todo dado do mundo real.
\prov{relações: usa instrumentacao\_medidas; parte\_de eng\_eletronica}
\prov{proveniência: wiki \textperiodcentered{} fato \textperiodcentered{} confiança 1.0 \textperiodcentered{} domínio engenharia \textperiodcentered{} história 3 versões no jornal}

%CRISTAL {"arestas": [{"alvo": "semicondutor", "confianca": 1.0, "epistemico": "fato", "origem": "wiki", "rel": "especializa"}, {"alvo": "rede_cristalina", "confianca": 1.0, "epistemico": "fato", "origem": "wiki", "rel": "usa"}], "confianca": 1.0, "contraexemplos": [], "descricao": "O semicondutor do século: abundante, estável, com óxido isolante nativo — o elemento sobre o qual toda a era digital foi construída.", "epistemico": "fato", "exemplos": [], "id": "silicio", "memoria": "semantica", "meta": {"dominio": "engenharia", "fonte": ""}, "origem": "wiki", "palavras_chave": [], "sinonimos": [], "tipo": "conceito", "titulo": "Silício"}
\section{Silício}
O semicondutor do século: abundante, estável, com óxido isolante nativo — o elemento sobre o qual toda a era digital foi construída.
\prov{relações: especializa semicondutor; usa rede\_cristalina}
\prov{proveniência: wiki \textperiodcentered{} fato \textperiodcentered{} confiança 1.0 \textperiodcentered{} domínio engenharia \textperiodcentered{} história 3 versões no jornal}

%CRISTAL {"arestas": [{"alvo": "cristalografia", "confianca": 1.0, "epistemico": "fato", "origem": "wiki", "rel": "parte_de"}, {"alvo": "teoria_de_grupos", "confianca": 1.0, "epistemico": "fato", "origem": "wiki", "rel": "usa"}, {"alvo": "programa_de_erlangen", "confianca": 1.0, "epistemico": "fato", "origem": "wiki", "rel": "relaciona"}], "confianca": 1.0, "contraexemplos": [], "descricao": "As operações (rotação, reflexão, translação) que deixam o cristal idêntico — os 230 grupos espaciais; a teoria de grupos feita pedra.", "epistemico": "fato", "exemplos": [], "id": "simetria_cristalina", "memoria": "semantica", "meta": {"autor": "Fedorov", "dominio": "engenharia", "fonte": ""}, "origem": "wiki", "palavras_chave": [], "sinonimos": [], "tipo": "conceito", "titulo": "Simetria Cristalina"}
\section{Simetria Cristalina}
As operações (rotação, reflexão, translação) que deixam o cristal idêntico — os 230 grupos espaciais; a teoria de grupos feita pedra.
\prov{relações: parte\_de cristalografia; usa teoria\_de\_grupos; relaciona programa\_de\_erlangen}
\prov{proveniência: wiki \textperiodcentered{} fato \textperiodcentered{} confiança 1.0 \textperiodcentered{} domínio engenharia \textperiodcentered{} história 4 versões no jornal}

%CRISTAL {"arestas": [{"alvo": "camada_fisica_nr", "confianca": 1.0, "epistemico": "fato", "origem": "wiki", "rel": "parte_de"}, {"alvo": "massive_mimo_beamforming", "confianca": 1.0, "epistemico": "fato", "origem": "wiki", "rel": "usa"}], "confianca": 1.0, "contraexemplos": [], "descricao": "Pilotos conhecidos que medem o canal: DMRS para demodular, CSI-RS para o UE reportar a qualidade e escolher o feixe, SRS para o gNB sondar a subida, PT-RS para rastrear ruído de fase em mmWave.", "epistemico": "fato", "exemplos": [], "id": "sinais_de_referencia_nr", "memoria": "semantica", "meta": {"dominio": "engenharia", "fonte": ""}, "origem": "wiki", "palavras_chave": [], "sinonimos": [], "tipo": "conceito", "titulo": "Sinais de Referência (DMRS, CSI-RS, SRS, PT-RS)"}
\section{Sinais de Referência (DMRS, CSI-RS, SRS, PT-RS)}
Pilotos conhecidos que medem o canal: DMRS para demodular, CSI-RS para o UE reportar a qualidade e escolher o feixe, SRS para o gNB sondar a subida, PT-RS para rastrear ruído de fase em mmWave.
\prov{relações: parte\_de camada\_fisica\_nr; usa massive\_mimo\_beamforming}
\prov{proveniência: wiki \textperiodcentered{} fato \textperiodcentered{} confiança 1.0 \textperiodcentered{} domínio engenharia \textperiodcentered{} história 3 versões no jornal}

%CRISTAL {"arestas": [{"alvo": "transformada_de_fourier", "confianca": 1.0, "epistemico": "fato", "origem": "wiki", "rel": "usa"}, {"alvo": "equacoes_diferenciais", "confianca": 1.0, "epistemico": "fato", "origem": "wiki", "rel": "usa"}, {"alvo": "eng_eletronica", "confianca": 1.0, "epistemico": "fato", "origem": "curriculo", "rel": "parte_de"}, {"alvo": "eng_telecomunicacoes", "confianca": 1.0, "epistemico": "fato", "origem": "curriculo", "rel": "parte_de"}, {"alvo": "eng_sistemas_energia", "confianca": 1.0, "epistemico": "fato", "origem": "curriculo", "rel": "parte_de"}], "confianca": 1.0, "contraexemplos": [], "descricao": "Tratar o sinal (função do tempo) e o sistema (o que o transforma) como objetos matemáticos: linearidade, invariância no tempo, convolução, resposta ao impulso. A gramática do processamento de sinais.", "epistemico": "fato", "exemplos": [], "id": "sinais_e_sistemas", "memoria": "semantica", "meta": {"dominio": "engenharia", "fonte": ""}, "origem": "wiki", "palavras_chave": [], "sinonimos": [], "tipo": "conceito", "titulo": "Sinais e Sistemas"}
\section{Sinais e Sistemas}
Tratar o sinal (função do tempo) e o sistema (o que o transforma) como objetos matemáticos: linearidade, invariância no tempo, convolução, resposta ao impulso. A gramática do processamento de sinais.
\prov{relações: usa transformada\_de\_fourier; usa equacoes\_diferenciais; parte\_de eng\_eletronica; parte\_de eng\_telecomunicacoes; parte\_de eng\_sistemas\_energia}
\prov{proveniência: wiki \textperiodcentered{} fato \textperiodcentered{} confiança 1.0 \textperiodcentered{} domínio engenharia \textperiodcentered{} história 6 versões no jornal}

%CRISTAL {"arestas": [{"alvo": "robotica", "confianca": 1.0, "epistemico": "fato", "origem": "wiki", "rel": "relaciona"}, {"alvo": "rede_autonoma_broca", "confianca": 1.0, "epistemico": "fato", "origem": "wiki", "rel": "relaciona"}, {"alvo": "bai", "confianca": 1.0, "epistemico": "fato", "origem": "wiki", "rel": "usa"}], "confianca": 1.0, "contraexemplos": [], "descricao": "Um sistema que decide as próprias ações sem controle humano contínuo: percebe o ambiente, avalia opções e age sob sua própria autorização. Do drone ao carro sem motorista — autonomia é fechar o laço de decisão dentro da própria máquina.", "epistemico": "fato", "exemplos": [], "id": "sistema_autonomo", "memoria": "semantica", "meta": {"dominio": "robotica", "fonte": "robótica e sistemas autônomos"}, "origem": "wiki", "palavras_chave": [], "sinonimos": [], "tipo": "conceito", "titulo": "Sistema Autônomo"}
\section{Sistema Autônomo}
Um sistema que decide as próprias ações sem controle humano contínuo: percebe o ambiente, avalia opções e age sob sua própria autorização. Do drone ao carro sem motorista — autonomia é fechar o laço de decisão dentro da própria máquina.
\prov{relações: relaciona robotica; relaciona rede\_autonoma\_broca; usa bai}
\prov{proveniência: wiki \textperiodcentered{} robótica e sistemas autônomos \textperiodcentered{} fato \textperiodcentered{} confiança 1.0 \textperiodcentered{} domínio robotica \textperiodcentered{} história 4 versões no jornal}

%CRISTAL {"arestas": [{"alvo": "modelo_shannon_weaver", "confianca": 1.0, "epistemico": "fato", "origem": "wiki", "rel": "especializa"}, {"alvo": "teoria_da_informacao", "confianca": 1.0, "epistemico": "fato", "origem": "wiki", "rel": "usa"}, {"alvo": "eng_telecomunicacoes", "confianca": 1.0, "epistemico": "fato", "origem": "curriculo", "rel": "parte_de"}], "confianca": 1.0, "contraexemplos": [], "descricao": "A cadeia fonte→transmissor→CANAL→receptor→destino, com ruído injetado no canal. O diagrama de blocos que organiza toda a telecomunicação — codificar, modular, propagar, demodular, decodificar.", "epistemico": "fato", "exemplos": [], "id": "sistema_de_comunicacao", "memoria": "semantica", "meta": {"autor": "Shannon", "dominio": "engenharia", "fonte": ""}, "origem": "wiki", "palavras_chave": [], "sinonimos": [], "tipo": "conceito", "titulo": "Sistema de Comunicação"}
\section{Sistema de Comunicação}
A cadeia fonte\ensuremath{\to}transmissor\ensuremath{\to}CANAL\ensuremath{\to}receptor\ensuremath{\to}destino, com ruído injetado no canal. O diagrama de blocos que organiza toda a telecomunicação — codificar, modular, propagar, demodular, decodificar.
\prov{relações: especializa modelo\_shannon\_weaver; usa teoria\_da\_informacao; parte\_de eng\_telecomunicacoes}
\prov{proveniência: wiki \textperiodcentered{} fato \textperiodcentered{} confiança 1.0 \textperiodcentered{} domínio engenharia \textperiodcentered{} história 4 versões no jornal}

%CRISTAL {"arestas": [{"alvo": "energia_eletrica", "confianca": 1.0, "epistemico": "fato", "origem": "wiki", "rel": "usa"}, {"alvo": "rede_eletrica", "confianca": 1.0, "epistemico": "fato", "origem": "wiki", "rel": "especializa"}, {"alvo": "geracao_de_energia", "confianca": 1.0, "epistemico": "fato", "origem": "wiki", "rel": "usa"}, {"alvo": "transmissao_de_energia", "confianca": 1.0, "epistemico": "fato", "origem": "wiki", "rel": "usa"}, {"alvo": "distribuicao_de_energia", "confianca": 1.0, "epistemico": "fato", "origem": "wiki", "rel": "usa"}, {"alvo": "eng_sistemas_energia", "confianca": 1.0, "epistemico": "fato", "origem": "curriculo", "rel": "parte_de"}], "confianca": 1.0, "contraexemplos": [], "descricao": "A rede inteira como um sistema: geração, transmissão em alta tensão, subestações e distribuição, operados para manter tensão, frequência e balanço em tempo real. O maior objeto de engenharia já construído.", "epistemico": "fato", "exemplos": [], "id": "sistema_eletrico_potencia", "memoria": "semantica", "meta": {"dominio": "engenharia", "fonte": ""}, "origem": "wiki", "palavras_chave": [], "sinonimos": [], "tipo": "conceito", "titulo": "Sistema Elétrico de Potência (SEP)"}
\section{Sistema Elétrico de Potência (SEP)}
A rede inteira como um sistema: geração, transmissão em alta tensão, subestações e distribuição, operados para manter tensão, frequência e balanço em tempo real. O maior objeto de engenharia já construído.
\prov{relações: usa energia\_eletrica; especializa rede\_eletrica; usa geracao\_de\_energia; usa transmissao\_de\_energia; usa distribuicao\_de\_energia; parte\_de eng\_sistemas\_energia}
\prov{proveniência: wiki \textperiodcentered{} fato \textperiodcentered{} confiança 1.0 \textperiodcentered{} domínio engenharia \textperiodcentered{} história 7 versões no jornal}

%CRISTAL {"arestas": [{"alvo": "sistemas_embarcados", "confianca": 1.0, "epistemico": "fato", "origem": "wiki", "rel": "parte_de"}, {"alvo": "sistema_operacional", "confianca": 1.0, "epistemico": "fato", "origem": "wiki", "rel": "especializa"}, {"alvo": "eng_eletronica", "confianca": 1.0, "epistemico": "fato", "origem": "curriculo", "rel": "parte_de"}], "confianca": 1.0, "contraexemplos": [], "descricao": "Quando a resposta certa fora do prazo é resposta errada: escalonamento por prioridade, determinismo, limites garantidos de latência. O SO do embarcado crítico.", "epistemico": "inferencia", "exemplos": [], "id": "sistema_tempo_real", "memoria": "semantica", "meta": {"dominio": "engenharia", "fonte": ""}, "origem": "wiki", "palavras_chave": [], "sinonimos": [], "tipo": "conceito", "titulo": "Sistemas de Tempo Real (RTOS)"}
\section{Sistemas de Tempo Real (RTOS)}
Quando a resposta certa fora do prazo é resposta errada: escalonamento por prioridade, determinismo, limites garantidos de latência. O SO do embarcado crítico.
\prov{relações: parte\_de sistemas\_embarcados; especializa sistema\_operacional; parte\_de eng\_eletronica}
\prov{proveniência: wiki \textperiodcentered{} inferencia \textperiodcentered{} confiança 1.0 \textperiodcentered{} domínio engenharia \textperiodcentered{} história 4 versões no jornal}

%CRISTAL {"arestas": [{"alvo": "sinais_e_sistemas", "confianca": 1.0, "epistemico": "fato", "origem": "wiki", "rel": "usa"}, {"alvo": "transformada_de_laplace", "confianca": 1.0, "epistemico": "fato", "origem": "wiki", "rel": "usa"}, {"alvo": "eng_eletronica", "confianca": 1.0, "epistemico": "fato", "origem": "curriculo", "rel": "parte_de"}, {"alvo": "eng_telecomunicacoes", "confianca": 1.0, "epistemico": "fato", "origem": "curriculo", "rel": "parte_de"}, {"alvo": "eng_sistemas_energia", "confianca": 1.0, "epistemico": "fato", "origem": "curriculo", "rel": "parte_de"}], "confianca": 1.0, "contraexemplos": [], "descricao": "Fechar a malha: medir a saída, comparar com o desejado e corrigir (realimentação). PID, estabilidade, lugar das raízes, resposta em frequência — manter um sistema no ponto apesar de perturbações.", "epistemico": "fato", "exemplos": [], "id": "sistemas_de_controle", "memoria": "semantica", "meta": {"autor": "Nyquist", "dominio": "engenharia", "fonte": ""}, "origem": "wiki", "palavras_chave": [], "sinonimos": [], "tipo": "conceito", "titulo": "Sistemas de Controle"}
\section{Sistemas de Controle}
Fechar a malha: medir a saída, comparar com o desejado e corrigir (realimentação). PID, estabilidade, lugar das raízes, resposta em frequência — manter um sistema no ponto apesar de perturbações.
\prov{relações: usa sinais\_e\_sistemas; usa transformada\_de\_laplace; parte\_de eng\_eletronica; parte\_de eng\_telecomunicacoes; parte\_de eng\_sistemas\_energia}
\prov{proveniência: wiki \textperiodcentered{} fato \textperiodcentered{} confiança 1.0 \textperiodcentered{} domínio engenharia \textperiodcentered{} história 6 versões no jornal}

%CRISTAL {"arestas": [{"alvo": "microcontrolador", "confianca": 1.0, "epistemico": "fato", "origem": "wiki", "rel": "usa"}, {"alvo": "sensor_transdutor", "confianca": 1.0, "epistemico": "fato", "origem": "wiki", "rel": "usa"}, {"alvo": "eng_eletronica", "confianca": 1.0, "epistemico": "fato", "origem": "curriculo", "rel": "parte_de"}], "confianca": 1.0, "contraexemplos": [], "descricao": "Computação dedicada dentro de um produto, com recursos e tempo apertados — do marca-passo ao carro. Hardware e software co-projetados para uma função, não para uso geral.", "epistemico": "inferencia", "exemplos": [], "id": "sistemas_embarcados", "memoria": "semantica", "meta": {"dominio": "engenharia", "fonte": ""}, "origem": "wiki", "palavras_chave": [], "sinonimos": [], "tipo": "conceito", "titulo": "Sistemas Embarcados"}
\section{Sistemas Embarcados}
Computação dedicada dentro de um produto, com recursos e tempo apertados — do marca-passo ao carro. Hardware e software co-projetados para uma função, não para uso geral.
\prov{relações: usa microcontrolador; usa sensor\_transdutor; parte\_de eng\_eletronica}
\prov{proveniência: wiki \textperiodcentered{} inferencia \textperiodcentered{} confiança 1.0 \textperiodcentered{} domínio engenharia \textperiodcentered{} história 4 versões no jornal}

%CRISTAL {"arestas": [{"alvo": "rede_inteligente", "confianca": 1.0, "epistemico": "fato", "origem": "wiki", "rel": "usa"}, {"alvo": "banco_de_dados_grafo", "confianca": 1.0, "epistemico": "fato", "origem": "wiki", "rel": "relaciona"}], "confianca": 1.0, "contraexemplos": [], "descricao": "A cidade instrumentada com sensores, dados e controle que coordena energia, trânsito, água e serviços em tempo real. Estruturalmente é um GRAFO: nós (prédios, carros, subestações, semáforos) e arestas (fios, ruas, redes de dados) — a mesma estrutura da broca-so. A cidade viva opera na borda.", "epistemico": "inferencia", "exemplos": [], "id": "smart_city", "memoria": "semantica", "meta": {"dominio": "engenharia", "fonte": "rede de carregamento e smart cities"}, "origem": "wiki", "palavras_chave": [], "sinonimos": [], "tipo": "conceito", "titulo": "Smart City (a cidade é um grafo)"}
\section{Smart City (a cidade é um grafo)}
A cidade instrumentada com sensores, dados e controle que coordena energia, trânsito, água e serviços em tempo real. Estruturalmente é um GRAFO: nós (prédios, carros, subestações, semáforos) e arestas (fios, ruas, redes de dados) — a mesma estrutura da broca-so. A cidade viva opera na borda.
\prov{relações: usa rede\_inteligente; relaciona banco\_de\_dados\_grafo}
\prov{proveniência: wiki \textperiodcentered{} rede de carregamento e smart cities \textperiodcentered{} inferencia \textperiodcentered{} confiança 1.0 \textperiodcentered{} domínio engenharia \textperiodcentered{} história 5 versões no jornal}

%CRISTAL {"arestas": [{"alvo": "rede_autonoma_broca", "confianca": 1.0, "epistemico": "fato", "origem": "wiki", "rel": "parte_de"}, {"alvo": "semente_booleana_dna_sistema", "confianca": 1.0, "epistemico": "fato", "origem": "wiki", "rel": "usa"}, {"alvo": "estruturas_vivas_sacrificio", "confianca": 1.0, "epistemico": "fato", "origem": "wiki", "rel": "relaciona"}, {"alvo": "automatos_autorreprodutores", "confianca": 1.0, "epistemico": "fato", "origem": "wiki", "rel": "usa"}], "confianca": 1.0, "contraexemplos": [], "descricao": "A semente booleana é o DNA do sistema: um artefato-semente reconstrói o corpo, substrato-independente. Software e arte, aqui, nascem, crescem, reproduzem-se e não morrem — organismos, não produtos. A comunidade orgânica é o ecossistema onde essas sementes circulam e mutam. CONSTRUÍDO (goldchain/semente.py): a semente carrega genoma e MUTA ao ser replantada por um par-ponte (reprodução com variação, régua de ouro determinística, trilha de mutações+linhagem). Verificado ao vivo: arte ◆◆◆→◆◆◇→◆△◇ evoluindo ao cruzar A→B→C. Organismo = circula E evolui, convergindo por dedup+TTL.", "epistemico": "inferencia", "exemplos": [], "id": "software_arte_organismo", "memoria": "semantica", "meta": {"dominio": "arquitetura", "fonte": "semilla_rede_autonoma"}, "origem": "wiki", "palavras_chave": [], "sinonimos": [], "tipo": "conceito", "titulo": "Software e arte como organismos que se reproduzem"}
\section{Software e arte como organismos que se reproduzem}
A semente booleana é o DNA do sistema: um artefato-semente reconstrói o corpo, substrato-independente. Software e arte, aqui, nascem, crescem, reproduzem-se e não morrem — organismos, não produtos. A comunidade orgânica é o ecossistema onde essas sementes circulam e mutam. CONSTRUÍDO (goldchain/semente.py): a semente carrega genoma e MUTA ao ser replantada por um par-ponte (reprodução com variação, régua de ouro determinística, trilha de mutações+linhagem). Verificado ao vivo: arte $\blacklozenge$$\blacklozenge$$\blacklozenge$\ensuremath{\to}$\blacklozenge$$\blacklozenge$$\lozenge$\ensuremath{\to}$\blacklozenge$\ensuremath{\triangle}$\lozenge$ evoluindo ao cruzar A\ensuremath{\to}B\ensuremath{\to}C. Organismo = circula E evolui, convergindo por dedup+TTL.
\prov{relações: parte\_de rede\_autonoma\_broca; usa semente\_booleana\_dna\_sistema; relaciona estruturas\_vivas\_sacrificio; usa automatos\_autorreprodutores}
\prov{proveniência: wiki \textperiodcentered{} semilla\_rede\_autonoma \textperiodcentered{} inferencia \textperiodcentered{} confiança 1.0 \textperiodcentered{} domínio arquitetura \textperiodcentered{} história 6 versões no jornal}

%CRISTAL {"arestas": [{"alvo": "ng_ran_gnb", "confianca": 1.0, "epistemico": "fato", "origem": "wiki", "rel": "parte_de"}, {"alvo": "virtualizacao", "confianca": 1.0, "epistemico": "fato", "origem": "wiki", "rel": "usa"}], "confianca": 1.0, "contraexemplos": [], "descricao": "O gNB se divide funcionalmente: RU (rádio, na torre), DU (camadas baixas, tempo real) e CU (camadas altas, centralizável na nuvem). Permite RAN virtualizada e aberta (O-RAN); a divisão define o fronthaul entre elas.", "epistemico": "inferencia", "exemplos": [], "id": "split_cu_du_ru", "memoria": "semantica", "meta": {"dominio": "engenharia", "fonte": ""}, "origem": "wiki", "palavras_chave": [], "sinonimos": [], "tipo": "conceito", "titulo": "Split CU/DU/RU (Fronthaul)"}
\section{Split CU/DU/RU (Fronthaul)}
O gNB se divide funcionalmente: RU (rádio, na torre), DU (camadas baixas, tempo real) e CU (camadas altas, centralizável na nuvem). Permite RAN virtualizada e aberta (O-RAN); a divisão define o fronthaul entre elas.
\prov{relações: parte\_de ng\_ran\_gnb; usa virtualizacao}
\prov{proveniência: wiki \textperiodcentered{} inferencia \textperiodcentered{} confiança 1.0 \textperiodcentered{} domínio engenharia \textperiodcentered{} história 3 versões no jornal}

%CRISTAL {"arestas": [{"alvo": "camada_fisica_nr", "confianca": 1.0, "epistemico": "fato", "origem": "wiki", "rel": "parte_de"}], "confianca": 1.0, "contraexemplos": [], "descricao": "O Bloco de Sinal de Sincronização que o celular busca ao ligar: sinais de sincronismo (PSS/SSS) + o canal de difusão PBCH com a informação mínima do sistema. A porta de entrada na célula, transmitida em feixes varridos.", "epistemico": "fato", "exemplos": [], "id": "ssb_sincronizacao_nr", "memoria": "semantica", "meta": {"dominio": "engenharia", "fonte": ""}, "origem": "wiki", "palavras_chave": [], "sinonimos": [], "tipo": "conceito", "titulo": "SSB e Sincronização (PBCH)"}
\section{SSB e Sincronização (PBCH)}
O Bloco de Sinal de Sincronização que o celular busca ao ligar: sinais de sincronismo (PSS/SSS) + o canal de difusão PBCH com a informação mínima do sistema. A porta de entrada na célula, transmitida em feixes varridos.
\prov{relações: parte\_de camada\_fisica\_nr}
\prov{proveniência: wiki \textperiodcentered{} fato \textperiodcentered{} confiança 1.0 \textperiodcentered{} domínio engenharia \textperiodcentered{} história 2 versões no jornal}

%CRISTAL {"arestas": [{"alvo": "porta_hodge_correlacao", "confianca": 1.0, "epistemico": "fato", "origem": "wiki", "rel": "usa"}, {"alvo": "sse_push_nao_polling", "confianca": 1.0, "epistemico": "fato", "origem": "wiki", "rel": "usa"}], "confianca": 1.0, "contraexemplos": [], "descricao": "_pulso_vivo é só front/ritmo (não acorda daemon nem altera costura); _pulso mecânico só incrementa quando a assinatura mecânica muda ou a porta abre. Separa deliberadamente o ritmo vivo (barato, alta frequência) do pulso mecânico (caro, só quando a costura muda) — evita virar oscilador. É o batimento da porta.", "epistemico": "inferencia", "exemplos": [], "id": "sse_pulso_vivo_mecanico", "memoria": "semantica", "meta": {"autor": "broca-so", "dominio": "cockpit", "fonte": "goldchain/cockpit_server.py:_disparar_pulso/_notificar_vivo"}, "origem": "wiki", "palavras_chave": [], "sinonimos": [], "tipo": "conceito", "titulo": "Dois Pulsos: Vivo (front) e Mecânico (costura)"}
\section{Dois Pulsos: Vivo (front) e Mecânico (costura)}
\_pulso\_vivo é só front/ritmo (não acorda daemon nem altera costura); \_pulso mecânico só incrementa quando a assinatura mecânica muda ou a porta abre. Separa deliberadamente o ritmo vivo (barato, alta frequência) do pulso mecânico (caro, só quando a costura muda) — evita virar oscilador. É o batimento da porta.
\prov{relações: usa porta\_hodge\_correlacao; usa sse\_push\_nao\_polling}
\prov{proveniência: wiki \textperiodcentered{} goldchain/cockpit\_server.py:\_disparar\_pulso/\_notificar\_vivo \textperiodcentered{} inferencia \textperiodcentered{} confiança 1.0 \textperiodcentered{} domínio cockpit \textperiodcentered{} história 3 versões no jornal}

%CRISTAL {"arestas": [{"alvo": "olho_de_cock_web", "confianca": 1.0, "epistemico": "fato", "origem": "wiki", "rel": "usa"}, {"alvo": "broca_cockpit_engine", "confianca": 1.0, "epistemico": "fato", "origem": "wiki", "rel": "usa"}], "confianca": 1.0, "contraexemplos": [], "descricao": "O cockpit consome o estado por Server-Sent-Events (/sse/broca via EventSource), não por polling: eventos 'vivo' (ritmo front ~0.8s) e 'pulso/hodge/bsd' (só quando o pulso mecânico muda). Bootstrap único por fetch, long-poll dedicado nas portas, fallback re-bootstrap em erro. O browser só rasteriza o que o servidor empurra.", "epistemico": "fato", "exemplos": [], "id": "sse_push_nao_polling", "memoria": "semantica", "meta": {"autor": "broca-so", "dominio": "cockpit", "fonte": "goldchain/cockpit/cockpit.js:connectBroca; cockpit_server.py:_sse_broca"}, "origem": "wiki", "palavras_chave": [], "sinonimos": [], "tipo": "conceito", "titulo": "Consumo por SSE Push (não polling)"}
\section{Consumo por SSE Push (não polling)}
O cockpit consome o estado por Server-Sent-Events (/sse/broca via EventSource), não por polling: eventos 'vivo' (ritmo front \~{}0.8s) e 'pulso/hodge/bsd' (só quando o pulso mecânico muda). Bootstrap único por fetch, long-poll dedicado nas portas, fallback re-bootstrap em erro. O browser só rasteriza o que o servidor empurra.
\prov{relações: usa olho\_de\_cock\_web; usa broca\_cockpit\_engine}
\prov{proveniência: wiki \textperiodcentered{} goldchain/cockpit/cockpit.js:connectBroca; cockpit\_server.py:\_sse\_broca \textperiodcentered{} fato \textperiodcentered{} confiança 1.0 \textperiodcentered{} domínio cockpit \textperiodcentered{} história 3 versões no jornal}

%CRISTAL {"arestas": [{"alvo": "armazenamento_o_cristal", "confianca": 1.0, "epistemico": "fato", "origem": "wiki", "rel": "eh_um"}], "confianca": 1.0, "contraexemplos": [], "descricao": "Guarda energia no campo elétrico entre placas (não em reação química). Muitíssima potência, pouca energia, η_rt~95% e milhões de ciclos — para picos rápidos, não para estoque longo.", "epistemico": "fato", "exemplos": [], "id": "supercapacitor", "memoria": "semantica", "meta": {"dominio": "engenharia", "fonte": "distribuição e armazenamento"}, "origem": "wiki", "palavras_chave": [], "sinonimos": [], "tipo": "conceito", "titulo": "Supercapacitor"}
\section{Supercapacitor}
Guarda energia no campo elétrico entre placas (não em reação química). Muitíssima potência, pouca energia, \ensuremath{\eta}\_rt\~{}95\% e milhões de ciclos — para picos rápidos, não para estoque longo.
\prov{relações: eh\_um armazenamento\_o\_cristal}
\prov{proveniência: wiki \textperiodcentered{} distribuição e armazenamento \textperiodcentered{} fato \textperiodcentered{} confiança 1.0 \textperiodcentered{} domínio engenharia \textperiodcentered{} história 4 versões no jornal}

%CRISTAL {"arestas": [{"alvo": "rede_de_sensores", "confianca": 1.0, "epistemico": "fato", "origem": "wiki", "rel": "usa"}, {"alvo": "internet_das_coisas", "confianca": 1.0, "epistemico": "fato", "origem": "wiki", "rel": "parte_de"}], "confianca": 1.0, "contraexemplos": [], "descricao": "Medir à distância e transmitir: o dado do sensor viaja até onde é usado (como o floco do gex44 sobe pelo canal até o olho). O fio nervoso da IoT — o que liga o colapso na ponta à decisão no centro.", "epistemico": "fato", "exemplos": [], "id": "telemetria", "memoria": "semantica", "meta": {"dominio": "engenharia", "fonte": "IoT e sensores"}, "origem": "wiki", "palavras_chave": [], "sinonimos": [], "tipo": "conceito", "titulo": "Telemetria"}
\section{Telemetria}
Medir à distância e transmitir: o dado do sensor viaja até onde é usado (como o floco do gex44 sobe pelo canal até o olho). O fio nervoso da IoT — o que liga o colapso na ponta à decisão no centro.
\prov{relações: usa rede\_de\_sensores; parte\_de internet\_das\_coisas}
\prov{proveniência: wiki \textperiodcentered{} IoT e sensores \textperiodcentered{} fato \textperiodcentered{} confiança 1.0 \textperiodcentered{} domínio engenharia \textperiodcentered{} história 6 versões no jornal}

%CRISTAL {"arestas": [{"alvo": "protocolo_tcp", "confianca": 1.0, "epistemico": "fato", "origem": "wiki", "rel": "parte_de"}, {"alvo": "eng_telecomunicacoes", "confianca": 1.0, "epistemico": "fato", "origem": "curriculo", "rel": "parte_de"}], "confianca": 1.0, "contraexemplos": [], "descricao": "SYN → SYN-ACK → ACK: como o TCP abre conexão sincronizando números de sequência antes de qualquer dado. O aperto de mão que estabelece o estado dos dois lados.", "epistemico": "fato", "exemplos": [], "id": "three_way_handshake", "memoria": "semantica", "meta": {"dominio": "engenharia", "fonte": ""}, "origem": "wiki", "palavras_chave": [], "sinonimos": [], "tipo": "conceito", "titulo": "Handshake de 3 Vias (TCP)"}
\section{Handshake de 3 Vias (TCP)}
SYN \ensuremath{\to} SYN-ACK \ensuremath{\to} ACK: como o TCP abre conexão sincronizando números de sequência antes de qualquer dado. O aperto de mão que estabelece o estado dos dois lados.
\prov{relações: parte\_de protocolo\_tcp; parte\_de eng\_telecomunicacoes}
\prov{proveniência: wiki \textperiodcentered{} fato \textperiodcentered{} confiança 1.0 \textperiodcentered{} domínio engenharia \textperiodcentered{} história 3 versões no jornal}

%CRISTAL {"arestas": [{"alvo": "imagem_bai", "confianca": 1.0, "epistemico": "fato", "origem": "wiki", "rel": "especializa"}, {"alvo": "rotulagem_e_colapso", "confianca": 1.0, "epistemico": "fato", "origem": "wiki", "rel": "usa"}, {"alvo": "dsl_bai_xadrez", "confianca": 1.0, "epistemico": "fato", "origem": "wiki", "rel": "relaciona"}, {"alvo": "segmentacao_de_imagem", "confianca": 1.0, "epistemico": "fato", "origem": "wiki", "rel": "usa"}, {"alvo": "crescer_regiao", "confianca": 1.0, "epistemico": "fato", "origem": "wiki", "rel": "usa"}, {"alvo": "fusao_e_deny_overrides", "confianca": 1.0, "epistemico": "fato", "origem": "wiki", "rel": "usa"}, {"alvo": "minimax_e_colapso", "confianca": 1.0, "epistemico": "fato", "origem": "wiki", "rel": "relaciona"}], "confianca": 1.0, "contraexemplos": [], "descricao": "A Tiffany segmenta uma imagem JOGANDO: planta sementes (a raiz/autoridade), cresce as regiões (propagação P) por crescimento por semente, e cada FUSÃO é um lance — um colapso Ψ em que a região fraca cede à forte (deny-overrides), até o orçamento δ=k. É o MESMO motor de decisão do xadrez, com outra imagem de fundo: a região que cede é o ramo podado, a fronteira é o efeito horizonte, δ é a profundidade. A narração sai do dicionário imagem↔BAI ao vivo. Motor real em jogo_segmentacao.py.", "epistemico": "fato", "exemplos": [], "id": "tiffany_segmenta", "memoria": "semantica", "meta": {"dominio": "ponte", "fonte": "tiffany segmenta"}, "origem": "wiki", "palavras_chave": [], "sinonimos": [], "tipo": "conceito", "titulo": "A Tiffany Segmenta (o jogo no sandbox)"}
\section{A Tiffany Segmenta (o jogo no sandbox)}
A Tiffany segmenta uma imagem JOGANDO: planta sementes (a raiz/autoridade), cresce as regiões (propagação P) por crescimento por semente, e cada FUSÃO é um lance — um colapso \ensuremath{\Psi} em que a região fraca cede à forte (deny-overrides), até o orçamento \ensuremath{\delta}=k. É o MESMO motor de decisão do xadrez, com outra imagem de fundo: a região que cede é o ramo podado, a fronteira é o efeito horizonte, \ensuremath{\delta} é a profundidade. A narração sai do dicionário imagem\ensuremath{\leftrightarrow}BAI ao vivo. Motor real em jogo\_segmentacao.py.
\prov{relações: especializa imagem\_bai; usa rotulagem\_e\_colapso; relaciona dsl\_bai\_xadrez; usa segmentacao\_de\_imagem; usa crescer\_regiao; usa fusao\_e\_deny\_overrides; relaciona minimax\_e\_colapso}
\prov{proveniência: wiki \textperiodcentered{} tiffany segmenta \textperiodcentered{} fato \textperiodcentered{} confiança 1.0 \textperiodcentered{} domínio ponte \textperiodcentered{} história 8 versões no jornal}

%CRISTAL {"arestas": [{"alvo": "registro_5g", "confianca": 1.0, "epistemico": "fato", "origem": "wiki", "rel": "parte_de"}], "confianca": 1.0, "contraexemplos": [], "descricao": "Inicial (ligou/entrou na cobertura), de atualização por mobilidade (mudou de área de registro), periódico (prova de vida à rede) e de emergência. Cada um reusa o mesmo esqueleto de sinalização NAS.", "epistemico": "inferencia", "exemplos": [], "id": "tipos_de_registro_5g", "memoria": "semantica", "meta": {"dominio": "engenharia", "fonte": ""}, "origem": "wiki", "palavras_chave": [], "sinonimos": [], "tipo": "conceito", "titulo": "Tipos de Registro"}
\section{Tipos de Registro}
Inicial (ligou/entrou na cobertura), de atualização por mobilidade (mudou de área de registro), periódico (prova de vida à rede) e de emergência. Cada um reusa o mesmo esqueleto de sinalização NAS.
\prov{relações: parte\_de registro\_5g}
\prov{proveniência: wiki \textperiodcentered{} inferencia \textperiodcentered{} confiança 1.0 \textperiodcentered{} domínio engenharia \textperiodcentered{} história 2 versões no jornal}

%CRISTAL {"arestas": [{"alvo": "sistemas_distribuidos", "confianca": 1.0, "epistemico": "fato", "origem": "wiki", "rel": "parte_de"}, {"alvo": "modelo_de_atores", "confianca": 1.0, "epistemico": "fato", "origem": "wiki", "rel": "usa"}, {"alvo": "consenso_distribuido", "confianca": 1.0, "epistemico": "fato", "origem": "wiki", "rel": "usa"}, {"alvo": "replicacao_dados", "confianca": 1.0, "epistemico": "fato", "origem": "wiki", "rel": "usa"}, {"alvo": "computacao_em_nuvem", "confianca": 1.0, "epistemico": "fato", "origem": "wiki", "rel": "parte_de"}], "confianca": 1.0, "contraexemplos": [], "descricao": "Projetar o sistema para continuar correto apesar de máquinas que caem: replicação, consenso, failover. Assume que a falha é a regra, não a exceção — a nuvem se mantém na borda estável tratando o caos das falhas como parte do funcionamento normal.", "epistemico": "fato", "exemplos": [], "id": "tolerancia_a_falhas", "memoria": "semantica", "meta": {"dominio": "engenharia", "fonte": "nuvem e data centers"}, "origem": "wiki", "palavras_chave": [], "sinonimos": [], "tipo": "conceito", "titulo": "Tolerância a Falhas"}
\section{Tolerância a Falhas}
Projetar o sistema para continuar correto apesar de máquinas que caem: replicação, consenso, failover. Assume que a falha é a regra, não a exceção — a nuvem se mantém na borda estável tratando o caos das falhas como parte do funcionamento normal.
\prov{relações: parte\_de sistemas\_distribuidos; usa modelo\_de\_atores; usa consenso\_distribuido; usa replicacao\_dados; parte\_de computacao\_em\_nuvem}
\prov{proveniência: wiki \textperiodcentered{} nuvem e data centers \textperiodcentered{} fato \textperiodcentered{} confiança 1.0 \textperiodcentered{} domínio engenharia \textperiodcentered{} história 11 versões no jornal}

%CRISTAL {"arestas": [{"alvo": "manufatura", "confianca": 1.0, "epistemico": "fato", "origem": "wiki", "rel": "explica"}, {"alvo": "corpo_mineral", "confianca": 1.0, "epistemico": "fato", "origem": "wiki", "rel": "usa"}, {"alvo": "numero_de_salem", "confianca": 1.0, "epistemico": "fato", "origem": "wiki", "rel": "relaciona"}], "confianca": 1.0, "contraexemplos": [], "descricao": "A tolerância é o desvio aceitável de uma peça; a capacidade de processo Cpk mede quão bem o processo cabe nela. Cpk<1 CAOS (defeitos escapam), Cpk≈1 BORDA, Cpk≥1.33 CRISTAL (processo capaz). Fabricar bem é manter o processo na borda estável — nem folgado (refugo) nem apertado demais (caro). Verif.: σ=0.02 Cpk=1.67 (cristal), σ=0.05 Cpk=0.67 (caos).", "epistemico": "inferencia", "exemplos": [], "id": "tolerancia_borda", "memoria": "semantica", "meta": {"classe_cpk_no": "caos (defeitos escapam)", "classe_cpk_ok": "cristal (processo capaz)", "cpk_sigma002": 1.67, "cpk_sigma005": 0.67, "dominio": "manufatura", "fonte": "impressão 3D e manufatura"}, "origem": "wiki", "palavras_chave": [], "sinonimos": [], "tipo": "conceito", "titulo": "Tolerância = a Borda (Cpk)"}
\section{Tolerância = a Borda (Cpk)}
A tolerância é o desvio aceitável de uma peça; a capacidade de processo Cpk mede quão bem o processo cabe nela. Cpk<1 CAOS (defeitos escapam), Cpk\ensuremath{\approx}1 BORDA, Cpk\ensuremath{\ge}1.33 CRISTAL (processo capaz). Fabricar bem é manter o processo na borda estável — nem folgado (refugo) nem apertado demais (caro). Verif.: \ensuremath{\sigma}=0.02 Cpk=1.67 (cristal), \ensuremath{\sigma}=0.05 Cpk=0.67 (caos).
\prov{relações: explica manufatura; usa corpo\_mineral; relaciona numero\_de\_salem}
\prov{proveniência: wiki \textperiodcentered{} impressão 3D e manufatura \textperiodcentered{} inferencia \textperiodcentered{} confiança 1.0 \textperiodcentered{} domínio manufatura \textperiodcentered{} história 4 versões no jornal}

%CRISTAL {"arestas": [{"alvo": "poliedros_de_ouro", "confianca": 1.0, "epistemico": "fato", "origem": "wiki", "rel": "usa"}, {"alvo": "o_gato_e_um_mapa_de_anosov", "confianca": 1.0, "epistemico": "fato", "origem": "wiki", "rel": "usa"}, {"alvo": "corpo_estelar", "confianca": 1.0, "epistemico": "fato", "origem": "wiki", "rel": "relaciona"}], "confianca": 1.0, "contraexemplos": [], "descricao": "O motor do vídeo: os poliedros transformando em tempo real, cada transformação lida pelo que PRESERVA (o dicionário 'transformacao'). ROTAÇÃO: ortogonal, det=+1, guarda comprimento e ângulo. ESCALA: det=s³, guarda ângulo. CISALHAMENTO: det=1, guarda o VOLUME mas torce o ângulo. REFLEXÃO: det=−1, guarda comprimento e inverte a orientação. O GATO A=[[1,1],[1,0]] no plano é a transformação HIPERBÓLICA de Anosov (det=−1, preserva área, estica em σ e contrai em −1/σ) — o caos determinístico deformando o sólido, a assinatura da máquina na geometria viva.", "epistemico": "inferencia", "exemplos": [], "id": "transformacoes_em_tempo_real", "memoria": "semantica", "meta": {"dominio": "imagem", "fonte": "poliedros minerais"}, "origem": "wiki", "palavras_chave": [], "sinonimos": [], "tipo": "conceito", "titulo": "As Transformações Geométricas em Tempo Real"}
\section{As Transformações Geométricas em Tempo Real}
O motor do vídeo: os poliedros transformando em tempo real, cada transformação lida pelo que PRESERVA (o dicionário 'transformacao'). ROTAÇÃO: ortogonal, det=+1, guarda comprimento e ângulo. ESCALA: det=s\ensuremath{^{3}}, guarda ângulo. CISALHAMENTO: det=1, guarda o VOLUME mas torce o ângulo. REFLEXÃO: det=\ensuremath{-}1, guarda comprimento e inverte a orientação. O GATO A=[[1,1],[1,0]] no plano é a transformação HIPERBÓLICA de Anosov (det=\ensuremath{-}1, preserva área, estica em \ensuremath{\sigma} e contrai em \ensuremath{-}1/\ensuremath{\sigma}) — o caos determinístico deformando o sólido, a assinatura da máquina na geometria viva.
\prov{relações: usa poliedros\_de\_ouro; usa o\_gato\_e\_um\_mapa\_de\_anosov; relaciona corpo\_estelar}
\prov{proveniência: wiki \textperiodcentered{} poliedros minerais \textperiodcentered{} inferencia \textperiodcentered{} confiança 1.0 \textperiodcentered{} domínio imagem \textperiodcentered{} história 4 versões no jornal}

%CRISTAL {"arestas": [{"alvo": "analise_de_fourier", "confianca": 1.0, "epistemico": "fato", "origem": "wiki", "rel": "especializa"}, {"alvo": "numeros_complexos", "confianca": 1.0, "epistemico": "fato", "origem": "wiki", "rel": "usa"}, {"alvo": "eng_eletronica", "confianca": 1.0, "epistemico": "fato", "origem": "curriculo", "rel": "parte_de"}, {"alvo": "eng_telecomunicacoes", "confianca": 1.0, "epistemico": "fato", "origem": "curriculo", "rel": "parte_de"}, {"alvo": "eng_sistemas_energia", "confianca": 1.0, "epistemico": "fato", "origem": "curriculo", "rel": "parte_de"}], "confianca": 1.0, "contraexemplos": [], "descricao": "Decompor qualquer sinal na soma de senoides — o domínio da FREQUÊNCIA. A ponte tempo↔frequência que torna filtragem, modulação e amostragem tratáveis. Série, transformada e DFT/FFT.", "epistemico": "fato", "exemplos": [], "id": "transformada_de_fourier", "memoria": "semantica", "meta": {"autor": "Fourier", "dominio": "engenharia", "fonte": ""}, "origem": "wiki", "palavras_chave": [], "sinonimos": [], "tipo": "conceito", "titulo": "Análise de Fourier"}
\section{Análise de Fourier}
Decompor qualquer sinal na soma de senoides — o domínio da FREQUÊNCIA. A ponte tempo\ensuremath{\leftrightarrow}frequência que torna filtragem, modulação e amostragem tratáveis. Série, transformada e DFT/FFT.
\prov{relações: especializa analise\_de\_fourier; usa numeros\_complexos; parte\_de eng\_eletronica; parte\_de eng\_telecomunicacoes; parte\_de eng\_sistemas\_energia}
\prov{proveniência: wiki \textperiodcentered{} fato \textperiodcentered{} confiança 1.0 \textperiodcentered{} domínio engenharia \textperiodcentered{} história 6 versões no jornal}

%CRISTAL {"arestas": [{"alvo": "transformada_de_fourier", "confianca": 1.0, "epistemico": "fato", "origem": "wiki", "rel": "generaliza"}, {"alvo": "sistemas_de_controle", "confianca": 1.0, "epistemico": "fato", "origem": "wiki", "rel": "usa"}, {"alvo": "eng_eletronica", "confianca": 1.0, "epistemico": "fato", "origem": "curriculo", "rel": "parte_de"}, {"alvo": "eng_telecomunicacoes", "confianca": 1.0, "epistemico": "fato", "origem": "curriculo", "rel": "parte_de"}, {"alvo": "eng_sistemas_energia", "confianca": 1.0, "epistemico": "fato", "origem": "curriculo", "rel": "parte_de"}], "confianca": 1.0, "contraexemplos": [], "descricao": "Levar EDOs a álgebra no plano complexo s — polos e zeros, função de transferência, estabilidade. A ferramenta padrão para analisar circuitos e sistemas de controle no regime transitório.", "epistemico": "fato", "exemplos": [], "id": "transformada_de_laplace", "memoria": "semantica", "meta": {"autor": "Laplace", "dominio": "engenharia", "fonte": ""}, "origem": "wiki", "palavras_chave": [], "sinonimos": [], "tipo": "conceito", "titulo": "Transformada de Laplace"}
\section{Transformada de Laplace}
Levar EDOs a álgebra no plano complexo s — polos e zeros, função de transferência, estabilidade. A ferramenta padrão para analisar circuitos e sistemas de controle no regime transitório.
\prov{relações: generaliza transformada\_de\_fourier; usa sistemas\_de\_controle; parte\_de eng\_eletronica; parte\_de eng\_telecomunicacoes; parte\_de eng\_sistemas\_energia}
\prov{proveniência: wiki \textperiodcentered{} fato \textperiodcentered{} confiança 1.0 \textperiodcentered{} domínio engenharia \textperiodcentered{} história 6 versões no jornal}

%CRISTAL {"arestas": [{"alvo": "inducao_eletromagnetica", "confianca": 1.0, "epistemico": "fato", "origem": "wiki", "rel": "usa"}, {"alvo": "conversao_eletromecanica", "confianca": 1.0, "epistemico": "fato", "origem": "wiki", "rel": "relaciona"}, {"alvo": "transmissao_de_energia", "confianca": 1.0, "epistemico": "fato", "origem": "wiki", "rel": "usa"}, {"alvo": "corrente_alternada", "confianca": 1.0, "epistemico": "fato", "origem": "wiki", "rel": "usa"}], "confianca": 1.0, "contraexemplos": [], "descricao": "A máquina que sobe ou desce a tensão CA por indução entre dois enrolamentos (razão N₂/N₁). É por causa dele que a rede é CA: sobe a tensão para transmitir com pouca perda e desce para o consumo seguro.", "epistemico": "fato", "exemplos": [], "id": "transformador", "memoria": "semantica", "meta": {"dominio": "engenharia", "fonte": "distribuição e armazenamento"}, "origem": "wiki", "palavras_chave": [], "sinonimos": [], "tipo": "conceito", "titulo": "Transformador"}
\section{Transformador}
A máquina que sobe ou desce a tensão CA por indução entre dois enrolamentos (razão N\ensuremath{_{2}}/N\ensuremath{_{1}}). É por causa dele que a rede é CA: sobe a tensão para transmitir com pouca perda e desce para o consumo seguro.
\prov{relações: usa inducao\_eletromagnetica; relaciona conversao\_eletromecanica; usa transmissao\_de\_energia; usa corrente\_alternada}
\prov{proveniência: wiki \textperiodcentered{} distribuição e armazenamento \textperiodcentered{} fato \textperiodcentered{} confiança 1.0 \textperiodcentered{} domínio engenharia \textperiodcentered{} história 9 versões no jornal}

%CRISTAL {"arestas": [{"alvo": "corpo_mineral", "confianca": 1.0, "epistemico": "fato", "origem": "wiki", "rel": "usa"}, {"alvo": "bai", "confianca": 1.0, "epistemico": "fato", "origem": "wiki", "rel": "usa"}, {"alvo": "newton_acha_o_grupo_discreto", "confianca": 1.0, "epistemico": "fato", "origem": "wiki", "rel": "usa"}, {"alvo": "xadrez_bai", "confianca": 1.0, "epistemico": "fato", "origem": "wiki", "rel": "relaciona"}, {"alvo": "imagem_bai", "confianca": 1.0, "epistemico": "fato", "origem": "wiki", "rel": "relaciona"}], "confianca": 1.0, "contraexemplos": [], "descricao": "A reta mineral, o BAI e o método de Newton são o MESMO processo iterativo classificado pela convergência (a parábola cristal/borda/caos). Três leituras de uma órbita: x→σ·x (mineral), o colapso Ψ com orçamento δ (BAI), z−f/f' (Newton). O dicionário (pivô 'orbita') traduz 7 construtos 1:1; compor() dá as três transições de graça. É o teorema da bússola no coração da máquina.", "epistemico": "fato", "exemplos": [], "id": "transicoes_reta_bai_newton", "memoria": "semantica", "meta": {"dominio": "ponte", "fonte": "transições órbita"}, "origem": "wiki", "palavras_chave": [], "sinonimos": [], "tipo": "conceito", "titulo": "Reta Mineral ↔ BAI ↔ Newton (a mesma órbita)"}
\section{Reta Mineral \ensuremath{\leftrightarrow} BAI \ensuremath{\leftrightarrow} Newton (a mesma órbita)}
A reta mineral, o BAI e o método de Newton são o MESMO processo iterativo classificado pela convergência (a parábola cristal/borda/caos). Três leituras de uma órbita: x\ensuremath{\to}\ensuremath{\sigma}\textperiodcentered x (mineral), o colapso \ensuremath{\Psi} com orçamento \ensuremath{\delta} (BAI), z\ensuremath{-}f/f' (Newton). O dicionário (pivô 'orbita') traduz 7 construtos 1:1; compor() dá as três transições de graça. É o teorema da bússola no coração da máquina.
\prov{relações: usa corpo\_mineral; usa bai; usa newton\_acha\_o\_grupo\_discreto; relaciona xadrez\_bai; relaciona imagem\_bai}
\prov{proveniência: wiki \textperiodcentered{} transições órbita \textperiodcentered{} fato \textperiodcentered{} confiança 1.0 \textperiodcentered{} domínio ponte \textperiodcentered{} história 6 versões no jornal}

%CRISTAL {"arestas": [{"alvo": "juncao_pn", "confianca": 1.0, "epistemico": "fato", "origem": "wiki", "rel": "usa"}, {"alvo": "dispositivo_semicondutor", "confianca": 1.0, "epistemico": "fato", "origem": "wiki", "rel": "especializa"}, {"alvo": "porta_logica", "confianca": 1.0, "epistemico": "fato", "origem": "wiki", "rel": "explica"}], "confianca": 1.0, "contraexemplos": [], "descricao": "Shockley, Bardeen e Brattain (1947): um interruptor/amplificador de estado sólido controlado por tensão — o tijolo de TODA a computação; ele É a porta lógica.", "epistemico": "fato", "exemplos": [], "id": "transistor", "memoria": "semantica", "meta": {"autor": "Shockley/Bardeen/Brattain", "dominio": "engenharia", "fonte": ""}, "origem": "wiki", "palavras_chave": [], "sinonimos": [], "tipo": "conceito", "titulo": "Transistor"}
\section{Transistor}
Shockley, Bardeen e Brattain (1947): um interruptor/amplificador de estado sólido controlado por tensão — o tijolo de TODA a computação; ele É a porta lógica.
\prov{relações: usa juncao\_pn; especializa dispositivo\_semicondutor; explica porta\_logica}
\prov{proveniência: wiki \textperiodcentered{} fato \textperiodcentered{} confiança 1.0 \textperiodcentered{} domínio engenharia \textperiodcentered{} história 4 versões no jornal}

%CRISTAL {"arestas": [{"alvo": "energia_eletrica", "confianca": 1.0, "epistemico": "fato", "origem": "wiki", "rel": "parte_de"}, {"alvo": "transformador", "confianca": 1.0, "epistemico": "fato", "origem": "wiki", "rel": "usa"}, {"alvo": "distribuicao_de_energia", "confianca": 1.0, "epistemico": "fato", "origem": "wiki", "rel": "parte_de"}, {"alvo": "rede_eletrica", "confianca": 1.0, "epistemico": "fato", "origem": "wiki", "rel": "usa"}], "confianca": 1.0, "contraexemplos": [], "descricao": "Levar a eletricidade da usina às cidades em ALTA TENSÃO (centenas de kV). Como a perda por calor é I²R, subir a tensão (e baixar a corrente) na mesma potência corta a perda quadraticamente. As linhas de transmissão são as artérias longas da rede.", "epistemico": "fato", "exemplos": [], "id": "transmissao_de_energia", "memoria": "semantica", "meta": {"dominio": "engenharia", "fonte": "distribuição e armazenamento"}, "origem": "wiki", "palavras_chave": [], "sinonimos": [], "tipo": "conceito", "titulo": "Transmissão de Energia"}
\section{Transmissão de Energia}
Levar a eletricidade da usina às cidades em ALTA TENSÃO (centenas de kV). Como a perda por calor é I\ensuremath{^{2}}R, subir a tensão (e baixar a corrente) na mesma potência corta a perda quadraticamente. As linhas de transmissão são as artérias longas da rede.
\prov{relações: parte\_de energia\_eletrica; usa transformador; parte\_de distribuicao\_de\_energia; usa rede\_eletrica}
\prov{proveniência: wiki \textperiodcentered{} distribuição e armazenamento \textperiodcentered{} fato \textperiodcentered{} confiança 1.0 \textperiodcentered{} domínio engenharia \textperiodcentered{} história 9 versões no jornal}

%CRISTAL {"arestas": [{"alvo": "painel_estelar_obs", "confianca": 1.0, "epistemico": "fato", "origem": "wiki", "rel": "usa"}, {"alvo": "geometria_estelar", "confianca": 1.0, "epistemico": "fato", "origem": "wiki", "rel": "usa"}, {"alvo": "mecanica_estelar", "confianca": 1.0, "epistemico": "fato", "origem": "wiki", "rel": "usa"}], "confianca": 1.0, "contraexemplos": [], "descricao": "Geometria (branca/estática: Lopes3→Hopf→de Sitter, τ,escala), mecânica (negra/dinâmica: o lado c² reversível, fase,θ) e termodinâmica (branca/dinâmica: o lado a=calor, entropia=log φ, classe V/C/F) são o MESMO gato lido em três braços — a álgebra fica no núcleo, o resto é consequência.", "epistemico": "fato", "exemplos": [], "id": "tres_bracos_estelar", "memoria": "semantica", "meta": {"autor": "broca-so", "dominio": "cockpit", "fonte": "papers/painel_estelar.tex §II"}, "origem": "wiki", "palavras_chave": [], "sinonimos": [], "tipo": "conceito", "titulo": "As Três Consequências (os três braços)"}
\section{As Três Consequências (os três braços)}
Geometria (branca/estática: Lopes3\ensuremath{\to}Hopf\ensuremath{\to}de Sitter, \ensuremath{\tau},escala), mecânica (negra/dinâmica: o lado c\ensuremath{^{2}} reversível, fase,\ensuremath{\theta}) e termodinâmica (branca/dinâmica: o lado a=calor, entropia=log \ensuremath{\varphi}, classe V/C/F) são o MESMO gato lido em três braços — a álgebra fica no núcleo, o resto é consequência.
\prov{relações: usa painel\_estelar\_obs; usa geometria\_estelar; usa mecanica\_estelar}
\prov{proveniência: wiki \textperiodcentered{} papers/painel\_estelar.tex §II \textperiodcentered{} fato \textperiodcentered{} confiança 1.0 \textperiodcentered{} domínio cockpit \textperiodcentered{} história 5 versões no jornal}

%CRISTAL {"arestas": [{"alvo": "usina_eletrica", "confianca": 1.0, "epistemico": "fato", "origem": "wiki", "rel": "parte_de"}, {"alvo": "gerador_eletrico", "confianca": 1.0, "epistemico": "fato", "origem": "wiki", "rel": "relaciona"}], "confianca": 1.0, "contraexemplos": [], "descricao": "A máquina que transforma o fluxo de um fluido (vapor, água, gás, ar) em rotação. É o prime mover — o eixo que carrega a energia mecânica até o gerador.", "epistemico": "fato", "exemplos": [], "id": "turbina", "memoria": "semantica", "meta": {"dominio": "engenharia", "fonte": "usinas e geração"}, "origem": "wiki", "palavras_chave": [], "sinonimos": [], "tipo": "conceito", "titulo": "Turbina"}
\section{Turbina}
A máquina que transforma o fluxo de um fluido (vapor, água, gás, ar) em rotação. É o prime mover — o eixo que carrega a energia mecânica até o gerador.
\prov{relações: parte\_de usina\_eletrica; relaciona gerador\_eletrico}
\prov{proveniência: wiki \textperiodcentered{} usinas e geração \textperiodcentered{} fato \textperiodcentered{} confiança 1.0 \textperiodcentered{} domínio engenharia \textperiodcentered{} história 6 versões no jornal}

%CRISTAL {"arestas": [{"alvo": "geracao_de_energia", "confianca": 1.0, "epistemico": "fato", "origem": "wiki", "rel": "usa"}, {"alvo": "energia_eletrica", "confianca": 1.0, "epistemico": "fato", "origem": "wiki", "rel": "explica"}], "confianca": 1.0, "contraexemplos": [], "descricao": "A planta que converte uma fonte primária (calor, queda d'água, vento, sol) em eletricidade. O padrão quase universal: um PRIME MOVER (turbina) gira um GERADOR que induz corrente. Só muda a fonte que move a turbina.", "epistemico": "fato", "exemplos": [], "id": "usina_eletrica", "memoria": "semantica", "meta": {"dominio": "engenharia", "fonte": "usinas e geração"}, "origem": "wiki", "palavras_chave": [], "sinonimos": [], "tipo": "conceito", "titulo": "Usina Elétrica"}
\section{Usina Elétrica}
A planta que converte uma fonte primária (calor, queda d'água, vento, sol) em eletricidade. O padrão quase universal: um PRIME MOVER (turbina) gira um GERADOR que induz corrente. Só muda a fonte que move a turbina.
\prov{relações: usa geracao\_de\_energia; explica energia\_eletrica}
\prov{proveniência: wiki \textperiodcentered{} usinas e geração \textperiodcentered{} fato \textperiodcentered{} confiança 1.0 \textperiodcentered{} domínio engenharia \textperiodcentered{} história 6 versões no jornal}

%CRISTAL {"arestas": [{"alvo": "usina_eletrica", "confianca": 1.0, "epistemico": "fato", "origem": "wiki", "rel": "eh_um"}, {"alvo": "energia_eolica", "confianca": 1.0, "epistemico": "fato", "origem": "wiki", "rel": "relaciona"}], "confianca": 1.0, "contraexemplos": [], "descricao": "O vento gira as pás (o limite de Betz: no máximo 59,3% da energia do vento). Renovável e barata, mas intermitente — exige a rede equilibrar.", "epistemico": "fato", "exemplos": [], "id": "usina_eolica", "memoria": "semantica", "meta": {"dominio": "engenharia", "fonte": "usinas e geração"}, "origem": "wiki", "palavras_chave": [], "sinonimos": [], "tipo": "conceito", "titulo": "Usina Eólica"}
\section{Usina Eólica}
O vento gira as pás (o limite de Betz: no máximo 59,3\% da energia do vento). Renovável e barata, mas intermitente — exige a rede equilibrar.
\prov{relações: eh\_um usina\_eletrica; relaciona energia\_eolica}
\prov{proveniência: wiki \textperiodcentered{} usinas e geração \textperiodcentered{} fato \textperiodcentered{} confiança 1.0 \textperiodcentered{} domínio engenharia \textperiodcentered{} história 6 versões no jornal}

%CRISTAL {"arestas": [{"alvo": "usina_eletrica", "confianca": 1.0, "epistemico": "fato", "origem": "wiki", "rel": "eh_um"}, {"alvo": "energia_hidreletrica", "confianca": 1.0, "epistemico": "fato", "origem": "wiki", "rel": "relaciona"}], "confianca": 1.0, "contraexemplos": [], "descricao": "A queda d'água move a turbina (Pelton/Francis/Kaplan) diretamente — sem ciclo térmico. Alta eficiência, despachável e com armazenamento (reservatório); a maior fonte renovável.", "epistemico": "fato", "exemplos": [], "id": "usina_hidreletrica", "memoria": "semantica", "meta": {"dominio": "engenharia", "fonte": "usinas e geração"}, "origem": "wiki", "palavras_chave": [], "sinonimos": [], "tipo": "conceito", "titulo": "Usina Hidrelétrica"}
\section{Usina Hidrelétrica}
A queda d'água move a turbina (Pelton/Francis/Kaplan) diretamente — sem ciclo térmico. Alta eficiência, despachável e com armazenamento (reservatório); a maior fonte renovável.
\prov{relações: eh\_um usina\_eletrica; relaciona energia\_hidreletrica}
\prov{proveniência: wiki \textperiodcentered{} usinas e geração \textperiodcentered{} fato \textperiodcentered{} confiança 1.0 \textperiodcentered{} domínio engenharia \textperiodcentered{} história 6 versões no jornal}

%CRISTAL {"arestas": [{"alvo": "usina_eletrica", "confianca": 1.0, "epistemico": "fato", "origem": "wiki", "rel": "eh_um"}, {"alvo": "ciclo_rankine", "confianca": 1.0, "epistemico": "fato", "origem": "wiki", "rel": "usa"}, {"alvo": "reator_nuclear", "confianca": 1.0, "epistemico": "fato", "origem": "wiki", "rel": "usa"}], "confianca": 1.0, "contraexemplos": [], "descricao": "Uma termelétrica cujo calor vem da FISSÃO controlada (o reator em k=1), não da queima. Vapor→turbina→gerador, sem emissão de carbono; a fonte é o defeito de massa.", "epistemico": "fato", "exemplos": [], "id": "usina_nuclear", "memoria": "semantica", "meta": {"dominio": "engenharia", "fonte": "usinas e geração"}, "origem": "wiki", "palavras_chave": [], "sinonimos": [], "tipo": "conceito", "titulo": "Usina Nuclear"}
\section{Usina Nuclear}
Uma termelétrica cujo calor vem da FISSÃO controlada (o reator em k=1), não da queima. Vapor\ensuremath{\to}turbina\ensuremath{\to}gerador, sem emissão de carbono; a fonte é o defeito de massa.
\prov{relações: eh\_um usina\_eletrica; usa ciclo\_rankine; usa reator\_nuclear}
\prov{proveniência: wiki \textperiodcentered{} usinas e geração \textperiodcentered{} fato \textperiodcentered{} confiança 1.0 \textperiodcentered{} domínio engenharia \textperiodcentered{} história 8 versões no jornal}

%CRISTAL {"arestas": [{"alvo": "usina_eletrica", "confianca": 1.0, "epistemico": "fato", "origem": "wiki", "rel": "eh_um"}], "confianca": 1.0, "contraexemplos": [], "descricao": "Fotovoltaica (PV): a luz gera corrente direto no semicondutor (sem turbina). CSP: espelhos concentram calor para um ciclo Rankine. A fonte que mais cresce, intermitente e sem partes móveis (PV).", "epistemico": "fato", "exemplos": [], "id": "usina_solar", "memoria": "semantica", "meta": {"dominio": "engenharia", "fonte": "usinas e geração"}, "origem": "wiki", "palavras_chave": [], "sinonimos": [], "tipo": "conceito", "titulo": "Usina Solar (PV e CSP)"}
\section{Usina Solar (PV e CSP)}
Fotovoltaica (PV): a luz gera corrente direto no semicondutor (sem turbina). CSP: espelhos concentram calor para um ciclo Rankine. A fonte que mais cresce, intermitente e sem partes móveis (PV).
\prov{relações: eh\_um usina\_eletrica}
\prov{proveniência: wiki \textperiodcentered{} usinas e geração \textperiodcentered{} fato \textperiodcentered{} confiança 1.0 \textperiodcentered{} domínio engenharia \textperiodcentered{} história 4 versões no jornal}

%CRISTAL {"arestas": [{"alvo": "usina_eletrica", "confianca": 1.0, "epistemico": "fato", "origem": "wiki", "rel": "eh_um"}, {"alvo": "ciclo_rankine", "confianca": 1.0, "epistemico": "fato", "origem": "wiki", "rel": "usa"}], "confianca": 1.0, "contraexemplos": [], "descricao": "Queima combustível (carvão, óleo, gás) para ferver água e mover a turbina a vapor (Rankine). Barata e despachável, mas emite CO₂. A espinha dorsal histórica da rede.", "epistemico": "fato", "exemplos": [], "id": "usina_termeletrica", "memoria": "semantica", "meta": {"dominio": "engenharia", "fonte": "usinas e geração"}, "origem": "wiki", "palavras_chave": [], "sinonimos": [], "tipo": "conceito", "titulo": "Usina Termelétrica"}
\section{Usina Termelétrica}
Queima combustível (carvão, óleo, gás) para ferver água e mover a turbina a vapor (Rankine). Barata e despachável, mas emite CO\ensuremath{_{2}}. A espinha dorsal histórica da rede.
\prov{relações: eh\_um usina\_eletrica; usa ciclo\_rankine}
\prov{proveniência: wiki \textperiodcentered{} usinas e geração \textperiodcentered{} fato \textperiodcentered{} confiança 1.0 \textperiodcentered{} domínio engenharia \textperiodcentered{} história 6 versões no jornal}

%CRISTAL {"arestas": [{"alvo": "fila_de_requisicoes", "confianca": 1.0, "epistemico": "fato", "origem": "wiki", "rel": "explica"}, {"alvo": "lyapunov_da_orbita", "confianca": 1.0, "epistemico": "fato", "origem": "wiki", "rel": "relaciona"}, {"alvo": "numero_de_salem", "confianca": 1.0, "epistemico": "fato", "origem": "wiki", "rel": "relaciona"}, {"alvo": "corpo_mineral", "confianca": 1.0, "epistemico": "fato", "origem": "wiki", "rel": "usa"}], "confianca": 1.0, "contraexemplos": [], "descricao": "A utilização ρ=λ/μ governa a fila: a latência média W=1/(μ−λ) DISPARA ao infinito quando ρ→1. É exatamente o ρ mineral: ρ<1 cristal (fila estável e curta), ρ=1 borda (o polo, latência diverge), ρ>1 caos (fila cresce sem limite). O data center opera na borda — alto ρ por eficiência, sem cruzar 1. Verif.: ρ=0.6→25ms, ρ=0.99→1000ms, ρ=1→∞.", "epistemico": "inferencia", "exemplos": [], "id": "utilizacao_e_borda", "memoria": "semantica", "meta": {"classe_rho1": "caos (fila cresce sem limite)", "dominio": "engenharia", "fonte": "nuvem e data centers", "latencia_ms_rho060": 25.0, "latencia_ms_rho099": 1000.0}, "origem": "wiki", "palavras_chave": [], "sinonimos": [], "tipo": "conceito", "titulo": "A Utilização ρ é a Borda"}
\section{A Utilização \ensuremath{\rho} é a Borda}
A utilização \ensuremath{\rho}=\ensuremath{\lambda}/\ensuremath{\mu} governa a fila: a latência média W=1/(\ensuremath{\mu}\ensuremath{-}\ensuremath{\lambda}) DISPARA ao infinito quando \ensuremath{\rho}\ensuremath{\to}1. É exatamente o \ensuremath{\rho} mineral: \ensuremath{\rho}<1 cristal (fila estável e curta), \ensuremath{\rho}=1 borda (o polo, latência diverge), \ensuremath{\rho}>1 caos (fila cresce sem limite). O data center opera na borda — alto \ensuremath{\rho} por eficiência, sem cruzar 1. Verif.: \ensuremath{\rho}=0.6\ensuremath{\to}25ms, \ensuremath{\rho}=0.99\ensuremath{\to}1000ms, \ensuremath{\rho}=1\ensuremath{\to}\ensuremath{\infty}.
\prov{relações: explica fila\_de\_requisicoes; relaciona lyapunov\_da\_orbita; relaciona numero\_de\_salem; usa corpo\_mineral}
\prov{proveniência: wiki \textperiodcentered{} nuvem e data centers \textperiodcentered{} inferencia \textperiodcentered{} confiança 1.0 \textperiodcentered{} domínio engenharia \textperiodcentered{} história 5 versões no jornal}

%CRISTAL {"arestas": [{"alvo": "minimax", "confianca": 1.0, "epistemico": "fato", "origem": "wiki", "rel": "usa"}, {"alvo": "teoria_dos_jogos", "confianca": 1.0, "epistemico": "fato", "origem": "wiki", "rel": "parte_de"}, {"alvo": "jogo_de_soma_zero", "confianca": 1.0, "epistemico": "fato", "origem": "wiki", "rel": "usa"}, {"alvo": "equilibrio_de_nash", "confianca": 1.0, "epistemico": "fato", "origem": "wiki", "rel": "relaciona"}], "confianca": 1.0, "contraexemplos": [], "descricao": "Num jogo finito de soma-zero e informação perfeita (o xadrez), o teorema minimax de von Neumann garante um VALOR: o resultado com jogo ótimo dos dois lados (vitória, empate ou derrota forçada). O colapso Ψ busca esse valor; é o equilíbrio de Nash do jogo estritamente competitivo.", "epistemico": "fato", "exemplos": [], "id": "valor_do_jogo_von_neumann", "memoria": "semantica", "meta": {"autor": "von Neumann", "dominio": "ponte", "fonte": "xadrez↔BAI"}, "origem": "wiki", "palavras_chave": [], "sinonimos": [], "tipo": "conceito", "titulo": "Valor do Jogo (von Neumann)"}
\section{Valor do Jogo (von Neumann)}
Num jogo finito de soma-zero e informação perfeita (o xadrez), o teorema minimax de von Neumann garante um VALOR: o resultado com jogo ótimo dos dois lados (vitória, empate ou derrota forçada). O colapso \ensuremath{\Psi} busca esse valor; é o equilíbrio de Nash do jogo estritamente competitivo.
\prov{relações: usa minimax; parte\_de teoria\_dos\_jogos; usa jogo\_de\_soma\_zero; relaciona equilibrio\_de\_nash}
\prov{proveniência: wiki \textperiodcentered{} xadrez\ensuremath{\leftrightarrow}BAI \textperiodcentered{} fato \textperiodcentered{} confiança 1.0 \textperiodcentered{} domínio ponte \textperiodcentered{} história 5 versões no jornal}

%CRISTAL {"arestas": [{"alvo": "carro_eletrico", "confianca": 1.0, "epistemico": "fato", "origem": "wiki", "rel": "usa"}, {"alvo": "armazenamento_estabiliza_rede", "confianca": 1.0, "epistemico": "fato", "origem": "wiki", "rel": "eh_um"}, {"alvo": "estabilidade_da_rede", "confianca": 1.0, "epistemico": "fato", "origem": "wiki", "rel": "explica"}, {"alvo": "lyapunov_da_orbita", "confianca": 1.0, "epistemico": "fato", "origem": "wiki", "rel": "usa"}], "confianca": 1.0, "contraexemplos": [], "descricao": "Vehicle-to-grid: um carro plugado pode DEVOLVER energia à rede, não só puxar. Milhões de VEs viram armazenamento distribuído móvel que injeta amortecimento e segura a rede na borda estável contra a intermitência. A frota é um cristal coletivo. Verificado: 1000 carros → λ bem negativo (estável).", "epistemico": "inferencia", "exemplos": [], "id": "vehicle_to_grid", "memoria": "semantica", "meta": {"classe": "cristal (rede estável)", "dominio": "engenharia", "fonte": "mobilidade elétrica", "lambda_1000_carros": -20.05}, "origem": "wiki", "palavras_chave": [], "sinonimos": [], "tipo": "conceito", "titulo": "V2G — a Frota de Cristais Móveis"}
\section{V2G — a Frota de Cristais Móveis}
Vehicle-to-grid: um carro plugado pode DEVOLVER energia à rede, não só puxar. Milhões de VEs viram armazenamento distribuído móvel que injeta amortecimento e segura a rede na borda estável contra a intermitência. A frota é um cristal coletivo. Verificado: 1000 carros \ensuremath{\to} \ensuremath{\lambda} bem negativo (estável).
\prov{relações: usa carro\_eletrico; eh\_um armazenamento\_estabiliza\_rede; explica estabilidade\_da\_rede; usa lyapunov\_da\_orbita}
\prov{proveniência: wiki \textperiodcentered{} mobilidade elétrica \textperiodcentered{} inferencia \textperiodcentered{} confiança 1.0 \textperiodcentered{} domínio engenharia \textperiodcentered{} história 5 versões no jornal}

%CRISTAL {"arestas": [{"alvo": "sistema_autonomo", "confianca": 1.0, "epistemico": "fato", "origem": "wiki", "rel": "eh_um"}, {"alvo": "carro_eletrico", "confianca": 1.0, "epistemico": "fato", "origem": "wiki", "rel": "usa"}, {"alvo": "planejamento_minimax_delta", "confianca": 1.0, "epistemico": "fato", "origem": "wiki", "rel": "usa"}, {"alvo": "malha_de_controle_laco", "confianca": 1.0, "epistemico": "fato", "origem": "wiki", "rel": "usa"}], "confianca": 1.0, "contraexemplos": [], "descricao": "O carro que dirige sozinho: funde sensores (câmera, LIDAR, radar — muitos colapsos), constrói o estado do mundo, planeja a trajetória (minimax contra o trânsito) e atua no volante e nos pedais. É o sistema autônomo completo sobre o carro elétrico — o laço sensório-motor a 100 km/h.", "epistemico": "fato", "exemplos": [], "id": "veiculo_autonomo", "memoria": "semantica", "meta": {"dominio": "robotica", "fonte": "robótica e sistemas autônomos"}, "origem": "wiki", "palavras_chave": [], "sinonimos": [], "tipo": "conceito", "titulo": "Veículo Autônomo"}
\section{Veículo Autônomo}
O carro que dirige sozinho: funde sensores (câmera, LIDAR, radar — muitos colapsos), constrói o estado do mundo, planeja a trajetória (minimax contra o trânsito) e atua no volante e nos pedais. É o sistema autônomo completo sobre o carro elétrico — o laço sensório-motor a 100 km/h.
\prov{relações: eh\_um sistema\_autonomo; usa carro\_eletrico; usa planejamento\_minimax\_delta; usa malha\_de\_controle\_laco}
\prov{proveniência: wiki \textperiodcentered{} robótica e sistemas autônomos \textperiodcentered{} fato \textperiodcentered{} confiança 1.0 \textperiodcentered{} domínio robotica \textperiodcentered{} história 5 versões no jornal}

%CRISTAL {"arestas": [{"alvo": "syscall", "confianca": 1.0, "epistemico": "fato", "origem": "curriculo", "rel": "define"}, {"alvo": "kernel", "confianca": 0.703, "epistemico": "fato", "origem": "manual", "rel": "referencia"}, {"alvo": "word", "confianca": 0.5, "epistemico": "fato", "origem": "wiki", "rel": "relaciona"}, {"alvo": "vontade_de_poder", "confianca": 0.5, "epistemico": "fato", "origem": "wiki", "rel": "relaciona"}, {"alvo": "while", "confianca": 0.5, "epistemico": "fato", "origem": "wiki", "rel": "relaciona"}, {"alvo": "vita_contemplativa", "confianca": 0.5, "epistemico": "fato", "origem": "wiki", "rel": "relaciona"}, {"alvo": "sistema_operacional", "confianca": 1.0, "epistemico": "fato", "origem": "wiki", "rel": "usa"}, {"alvo": "processo", "confianca": 1.0, "epistemico": "fato", "origem": "wiki", "rel": "usa"}, {"alvo": "data_center", "confianca": 1.0, "epistemico": "fato", "origem": "wiki", "rel": "usa"}, {"alvo": "computador_mineral", "confianca": 1.0, "epistemico": "fato", "origem": "wiki", "rel": "relaciona"}], "confianca": 1.0, "contraexemplos": [], "descricao": "Fatiar um servidor físico em muitas máquinas virtuais / contêineres isolados. É o que torna a nuvem elástica e densa — e o data center é, em escala, o computador mineral materializado num galpão: ULAs, memória e canais alugados por fatia.", "epistemico": "inferencia", "exemplos": [], "id": "virtualizacao", "memoria": "semantica", "meta": {"dominio": "engenharia", "fonte": "nuvem e data centers"}, "origem": "wiki", "palavras_chave": [], "sinonimos": [], "tipo": "conceito", "titulo": "Virtualização e Contêineres"}
\section{Virtualização e Contêineres}
Fatiar um servidor físico em muitas máquinas virtuais / contêineres isolados. É o que torna a nuvem elástica e densa — e o data center é, em escala, o computador mineral materializado num galpão: ULAs, memória e canais alugados por fatia.
\prov{relações: define syscall; referencia kernel; relaciona word; relaciona vontade\_de\_poder; relaciona while; relaciona vita\_contemplativa; usa sistema\_operacional; usa processo; usa data\_center; relaciona computador\_mineral}
\prov{proveniência: wiki \textperiodcentered{} nuvem e data centers \textperiodcentered{} inferencia \textperiodcentered{} confiança 1.0 \textperiodcentered{} domínio engenharia \textperiodcentered{} história 17 versões no jornal}

%CRISTAL {"arestas": [{"alvo": "armazenamento_o_cristal", "confianca": 1.0, "epistemico": "fato", "origem": "wiki", "rel": "eh_um"}, {"alvo": "numero_de_salem", "confianca": 1.0, "epistemico": "fato", "origem": "wiki", "rel": "relaciona"}], "confianca": 1.0, "contraexemplos": [], "descricao": "Guarda energia na rotação de um disco (E=½Iω²). Muita potência por pouco tempo, η_rt~88% — segura a frequência da rede em segundos. Literalmente uma rotação armazenada (o Salem que gira guardando).", "epistemico": "fato", "exemplos": [], "id": "volante_de_inercia", "memoria": "semantica", "meta": {"dominio": "engenharia", "fonte": "distribuição e armazenamento"}, "origem": "wiki", "palavras_chave": [], "sinonimos": [], "tipo": "conceito", "titulo": "Volante de Inércia (flywheel)"}
\section{Volante de Inércia (flywheel)}
Guarda energia na rotação de um disco (E=\ensuremath{\tfrac12}I\ensuremath{\omega}\ensuremath{^{2}}). Muita potência por pouco tempo, \ensuremath{\eta}\_rt\~{}88\% — segura a frequência da rede em segundos. Literalmente uma rotação armazenada (o Salem que gira guardando).
\prov{relações: eh\_um armazenamento\_o\_cristal; relaciona numero\_de\_salem}
\prov{proveniência: wiki \textperiodcentered{} distribuição e armazenamento \textperiodcentered{} fato \textperiodcentered{} confiança 1.0 \textperiodcentered{} domínio engenharia \textperiodcentered{} história 6 versões no jornal}

%CRISTAL {"arestas": [{"alvo": "http", "confianca": 1.0, "epistemico": "fato", "origem": "wiki", "rel": "usa"}, {"alvo": "eng_telecomunicacoes", "confianca": 1.0, "epistemico": "fato", "origem": "curriculo", "rel": "parte_de"}], "confianca": 1.0, "contraexemplos": [], "descricao": "Um canal full-duplex persistente sobre uma conexão HTTP promovida — para tempo real (chat, jogo, cotação) sem o custo de reabrir requisições. A web que empurra, não só responde.", "epistemico": "fato", "exemplos": [], "id": "websocket", "memoria": "semantica", "meta": {"dominio": "engenharia", "fonte": ""}, "origem": "wiki", "palavras_chave": [], "sinonimos": [], "tipo": "conceito", "titulo": "WebSocket"}
\section{WebSocket}
Um canal full-duplex persistente sobre uma conexão HTTP promovida — para tempo real (chat, jogo, cotação) sem o custo de reabrir requisições. A web que empurra, não só responde.
\prov{relações: usa http; parte\_de eng\_telecomunicacoes}
\prov{proveniência: wiki \textperiodcentered{} fato \textperiodcentered{} confiança 1.0 \textperiodcentered{} domínio engenharia \textperiodcentered{} história 3 versões no jornal}

%CRISTAL {"arestas": [{"alvo": "camada_de_enlace", "confianca": 1.0, "epistemico": "fato", "origem": "wiki", "rel": "especializa"}, {"alvo": "acesso_ao_meio", "confianca": 1.0, "epistemico": "fato", "origem": "wiki", "rel": "usa"}, {"alvo": "ofdm", "confianca": 1.0, "epistemico": "fato", "origem": "wiki", "rel": "usa"}, {"alvo": "mimo", "confianca": 1.0, "epistemico": "fato", "origem": "wiki", "rel": "usa"}, {"alvo": "eng_telecomunicacoes", "confianca": 1.0, "epistemico": "fato", "origem": "curriculo", "rel": "parte_de"}], "confianca": 1.0, "contraexemplos": [], "descricao": "A LAN sem fio: enlace por rádio com CSMA/CA, reconhecimento de quadros e associação a um ponto de acesso. Usa a camada física OFDM/MIMO da trilha de princípios — o enlace mais usado do mundo.", "epistemico": "fato", "exemplos": [], "id": "wifi_802_11", "memoria": "semantica", "meta": {"dominio": "engenharia", "fonte": ""}, "origem": "wiki", "palavras_chave": [], "sinonimos": [], "tipo": "conceito", "titulo": "Wi-Fi (IEEE 802.11)"}
\section{Wi-Fi (IEEE 802.11)}
A LAN sem fio: enlace por rádio com CSMA/CA, reconhecimento de quadros e associação a um ponto de acesso. Usa a camada física OFDM/MIMO da trilha de princípios — o enlace mais usado do mundo.
\prov{relações: especializa camada\_de\_enlace; usa acesso\_ao\_meio; usa ofdm; usa mimo; parte\_de eng\_telecomunicacoes}
\prov{proveniência: wiki \textperiodcentered{} fato \textperiodcentered{} confiança 1.0 \textperiodcentered{} domínio engenharia \textperiodcentered{} história 6 versões no jornal}

%CRISTAL {"arestas": [{"alvo": "algebra_posicional_unificada", "confianca": 1.0, "epistemico": "fato", "origem": "wiki", "rel": "usa"}, {"alvo": "tiffany", "confianca": 1.0, "epistemico": "fato", "origem": "wiki", "rel": "parte_de"}], "confianca": 1.0, "contraexemplos": [], "descricao": "Veredito geral (agente): o ganho do M×M é CONCEITUAL/extensibilidade (uma reta nova = uma lista de coefs, zero código novo), não de runtime (~zero). A recomendação de projeto — 'álgebra M×M geral + caminho rápido 2×2' — é EXATAMENTE o estado atual: metal_pos.py carrega a álgebra geral (passo_companheira, graus 2..8, 128 bits) e ula/metais.h/Word carregam o 2×2 rápido. Os dois níveis já coexistem corretamente. Não há colapso a fazer; o 2×2 é o motor, o geral é a moldura. [I] o desenho já está no ponto — a varredura confirma, não muda.", "epistemico": "inferencia", "exemplos": [], "id": "word_m_geral_fastpath_2x2", "memoria": "semantica", "meta": {"autor": "Lopes/broca-so", "dominio": "arquitetura", "fonte": "agente veredito: metal_pos.py (geral) + metais.h (2×2 rápido) = a recomendação, já implementado"}, "origem": "wiki", "palavras_chave": [], "sinonimos": [], "tipo": "conceito", "titulo": "A arquitetura JÁ é 'Word_M geral + fast-path 2×2' — o desenho certo, não colapsar"}
\section{A arquitetura JÁ é 'Word\_M geral + fast-path 2$\times$2' — o desenho certo, não colapsar}
Veredito geral (agente): o ganho do M$\times$M é CONCEITUAL/extensibilidade (uma reta nova = uma lista de coefs, zero código novo), não de runtime (\~{}zero). A recomendação de projeto — 'álgebra M$\times$M geral + caminho rápido 2$\times$2' — é EXATAMENTE o estado atual: metal\_pos.py carrega a álgebra geral (passo\_companheira, graus 2..8, 128 bits) e ula/metais.h/Word carregam o 2$\times$2 rápido. Os dois níveis já coexistem corretamente. Não há colapso a fazer; o 2$\times$2 é o motor, o geral é a moldura. [I] o desenho já está no ponto — a varredura confirma, não muda.
\prov{relações: usa algebra\_posicional\_unificada; parte\_de tiffany}
\prov{proveniência: wiki \textperiodcentered{} agente veredito: metal\_pos.py (geral) + metais.h (2$\times$2 rápido) = a recomendação, já implementado \textperiodcentered{} inferencia \textperiodcentered{} confiança 1.0 \textperiodcentered{} domínio arquitetura \textperiodcentered{} história 29 versões no jornal}

%CRISTAL {"arestas": [{"alvo": "xadrez", "confianca": 1.0, "epistemico": "fato", "origem": "wiki", "rel": "usa"}, {"alvo": "teorema_geral_traducao", "confianca": 1.0, "epistemico": "fato", "origem": "wiki", "rel": "usa"}, {"alvo": "bai_como_bussola", "confianca": 1.0, "epistemico": "fato", "origem": "wiki", "rel": "usa"}, {"alvo": "laco_traducao_bai_peirce", "confianca": 1.0, "epistemico": "fato", "origem": "wiki", "rel": "relaciona"}, {"alvo": "teoria_da_decisao", "confianca": 1.0, "epistemico": "fato", "origem": "wiki", "rel": "relaciona"}], "confianca": 1.0, "contraexemplos": [], "descricao": "Xadrez e a Biblioteca de Autorização Isomórfica descrevem o MESMO objeto: uma busca em árvore, limitada por um orçamento de profundidade, que colapsa numa decisão e fecha (conservador) quando o orçamento esgota. Duas simbologias de um significado — o teorema da bússola. O dicionário traduz 11 construtos 1:1; a diferença (roque, tenant…) é só bump.", "epistemico": "inferencia", "exemplos": [], "id": "xadrez_bai", "memoria": "semantica", "meta": {"dominio": "ponte", "fonte": "xadrez↔BAI"}, "origem": "wiki", "palavras_chave": [], "sinonimos": [], "tipo": "conceito", "titulo": "Xadrez ↔ BAI (a mesma árvore de decisão)"}
\section{Xadrez \ensuremath{\leftrightarrow} BAI (a mesma árvore de decisão)}
Xadrez e a Biblioteca de Autorização Isomórfica descrevem o MESMO objeto: uma busca em árvore, limitada por um orçamento de profundidade, que colapsa numa decisão e fecha (conservador) quando o orçamento esgota. Duas simbologias de um significado — o teorema da bússola. O dicionário traduz 11 construtos 1:1; a diferença (roque, tenant…) é só bump.
\prov{relações: usa xadrez; usa teorema\_geral\_traducao; usa bai\_como\_bussola; relaciona laco\_traducao\_bai\_peirce; relaciona teoria\_da\_decisao}
\prov{proveniência: wiki \textperiodcentered{} xadrez\ensuremath{\leftrightarrow}BAI \textperiodcentered{} inferencia \textperiodcentered{} confiança 1.0 \textperiodcentered{} domínio ponte \textperiodcentered{} história 6 versões no jornal}

\end{document}
